% !TEX program = lualatex
% Continuous Chinese translation without facing-page English source plates.
\documentclass[UTF8,11pt,a4paper,openany,fontset=fandol]{ctexbook}
\usepackage[margin=2.35cm,headheight=25pt]{geometry}
\usepackage{amsmath,amssymb,bm,mathtools}
\usepackage{fancyhdr,microtype,xcolor,hyperref,xurl,needspace}
\setlength{\parindent}{2em}
\setlength{\parskip}{0.55em}
\setlength{\emergencystretch}{3em}
\linespread{1.22}\selectfont
\clubpenalty=10000
\widowpenalty=10000
\displaywidowpenalty=10000
\raggedbottom
\pagestyle{fancy}
\fancyhf{}
\fancyhead[C]{\small\nouppercase{\leftmark}}
\fancyfoot[C]{\thepage}
\renewcommand{\headrulewidth}{0.3pt}
\hypersetup{hypertexnames=false,colorlinks=true,linkcolor=black,urlcolor=blue!60!black,pdftitle={计算流体动力学前沿（1998）},pdfauthor={D. A. Caughey、M. M. Hafez 主编}}
\title{计算流体动力学前沿（1998）\\[0.45em]\Large 中文译本}
\author{D. A. Caughey、M. M. Hafez 主编\\\small 中文初译经 CFD 专业术语表统一}
\date{}
\begin{document}
\frontmatter
\maketitle
\chapter*{译本说明}
本译本按原书章节顺序翻译并连续排版，正文不附英文原文，供中文阅读和检索使用。作者姓名、文献题名、专业缩写以及数学符号按学术写作惯例保留必要的拉丁字符。

中文初稿由离线神经机器翻译批量生成，并通过 CFD 专业术语表统一 Navier-Stokes 方程、有限体积法、伴随方法、多重网格、预处理、湍流、磁流体动力学等词汇。由于全书篇幅大且原 PDF 文本层存在断行与公式混排，本文件属于可继续校订的学习译本，不替代正式授权译本。
\tableofcontents
\mainmatter
\chapter*{前页}
1998年计算流体动力学的前沿

此页面是故意留空的

Frontiers of 计算流体动力学 1998 编辑 D A Caughey 康奈尔大学 M M Hafez 加州大学，戴维斯世界科学新加坡 • New JerseM • 伦敦 • 香港

由世界科学出版公司出版。 Pte。 有限公司 P O Box 128, Farrer Road, Singapore 912805 USA 办公室：Suite IB, 1060 Main Street, River Edge, NJ 07661 英国办公室：57 Shelton Street, Covent Garden, London WC2H 9HE 大英图书馆编目出版数据 本书的目录记录可从大英图书馆获得。 FRONTIERS OF 计算流体动力学 1998 版权所有 © 1998 世界科学出版公司。 Pte。 有限公司 保留所有权利。 未经出版商书面许可，不得以任何形式或任何方式（电子或机械）复制本书或其部分，包括复印、录音或任何已知或即将发明的信息存储和检索系统。 如需复印本卷材料，请通过版权清关中心支付复印费，地址为222 Rosewood Drive, Danvers, MA 01923, USA。 在这种情况下，不需要出版商的复印许可。 ISBN 981-02-3707-3 由Uto-Print在新加坡印刷

奉献 本卷包括在研讨会上发表的论文，以纪念Earll Murman，并表彰他过去三十年来对跨声速空气动力学和计算流体动力学（CFD）的开创性贡献。 题为“CFD和跨声速流动的三十年”的研讨会于1997年6月24日至26日在华盛顿州埃弗雷特举行。 作者是从从事空气动力学和CFD领域的国际知名研究人员中选出的，Murman的贡献对其中的影响非常重要。 作者和编辑很高兴把这本书献给Earll，以表彰他在我们的技术和生活中发挥的重要作用。 Earll Murman出生于1942年5月12日。 他在旧金山长大，但去了普林斯顿大学，分别于1963年、1965年和1967年获得了学士学位、硕士学位和博士学位。 他的博士研究，在S教授的指导下。 M。 Bogdonoff，导致了题为“层面高超音速锥形唤醒的实验研究”的论文。 他于1967年加入波音科学研究实验室（BSRL），并一直待到1971年。 在此期间，他与朱利安·科尔教授一起工作，朱利安·科尔教授正在加州大学洛杉矶分校在波音实验室休休假。 他们的合作产生了被称为Murman-Cole方案的突破，它允许对稳定的跨声速流动场进行首次实际计算，该场包含嵌入亚音速区域的超声速流动区域。 1971年，Earll加入了美国宇航局艾姆斯研究中心，1973年，在加利福尼亚州棕榈泉举行的第一届AIAA关于CFD的会议上，他提出了他完全保守的Murman-Cole计划版本。 他在风洞墙干扰以及设计和优化方面的开创性工作也是当前开发多学科优化技术的尝试的显著先驱。 描述Murman-Cole计划的论文发表在1971年1月的《AIAA杂志》上，并被《Current Contents, Vol.》确定为引用经典。 27，不。 45，1987年11月9日。 Frontiers of 计算流体动力学 - 1998 编辑：David A. Caughey k，Mohamed M. 哈菲兹 ©1998 世界科学

VI DEDICATION Murman继续与Cole的合作，并制作了另一篇经典作品，即他们关于Transonic Speeds的Inviscid Drag的AIAA论文，于1974年在帕洛阿尔托发表。 那一年，Earll加入了华盛顿州肯特的Flow Research公司，并于1977年成为副总裁兼总经理。 1980年，他转到麻省理工学院担任航空航天学教授，并于1990年至1996年担任系主任。 Earll曾担任Calspan、United Technologies Corporation、Pratt \& Whitney、General Electric Corporation、Stellar Computer Company、Kendall Square Research和Microcraft Technology的顾问。 在麻省理工学院，Earll于1980年至1990年担任计算流体动力学实验室主任。 他主持了雅典娜项目资源委员会，最近，他是Todor背后的动力，Todor是一个旨在使用计算机增强流体力学课程的软件包。 Earll曾教授过超音速空气动力学、计算流体动力学、黏性流体、热量和传质、飞行流体动力学和再入飞行器等课程，并监督了本科生的数值和实验项目。 他为26名本科生、20名硕士生和8名博士生提供过建议。 他在麻省理工学院的研究旨在解决Euler 方程，包括模拟涡流。 他还研究了边界层及其与欧拉计算、化学反应流、Navier-Stokes 方程和黏性高超音速流以及流可视化技术的耦合。 他对工程教育真正感兴趣，从他的出版物中可以看出。 他积极参与了美国航空航天研究所、美国国家航空航天局、国防部和航空航天业的许多委员会，自1995年以来一直担任精益飞机倡议的主任。 他是AIAA的院士和美国国家工程学院的成员。 在本书的第一章中，将更详细地讨论Murman的技术贡献，特别是它们对跨声速空气动力学和CFD的影响。 Earll的贡献不仅限于他的技术理念、他的领导力、他教授的课程或他对麻省理工学院许多有才华的学生的监督。 我们社区有幸拥有像Earll这样的人，他不仅影响了多年来与他直接合作的人，还影响了许多其他从未有幸亲自见到他的人。 据说，Murman为他最初的贡献之后的大量活动打开了大门。 由于他的远见卓识、个性和慷慨的天性，他受到整个航空航天界的高度尊重。 我们希望未来将举行第二次会议，以表彰他未来30年的持续贡献。

伯爵M。 默尔曼

此页面是故意留空的

奉献................................ X 对Earll Murman对跨声速流动和计算流体动力学 Hafez d Caughey的贡献的回顾................................. 1.1 介绍................................ 1.2 对跨声速流动的贡献...................... 1.3 对CFD的贡献....................... 1.4 穆尔曼早期贡献的影响............. 1.5 Murman最近的贡献................. 1.6 结束说明.............................. 1.7 由Earl1 Murman监督的论文.................. 1.8 Earl1 Murman的出版物...................... 2个最佳高超音速锥形翅膀三叉戟。 Schwendeman €4 Cole....................... 2.1 介绍................................ 2.2 高超音速小扰动理论；锥形翼...... 2.3升力和阻力系数；功图：锥形翼...... 2.4 数值和优化方法................. 2.5 扁平翼和矟翼的计算结果；最佳翼...... 2.6 确认和免责声明.................... 参考资料................................. 3 理论几何学。 应用。和教育流体动力学................................ Sobieczky................................ 3.1 导言 3.2 跨声速知识库................................................. 3.3 几何生成器...... 3.4 通用运输飞机............................................................ 3.5 结论参考文献................................

目录4 轴对称喷管流动的计算.......... 厨师。 纽曼。 Rimbey d Schleiniger...... 57 4.1 介绍................................ 57 4.2 问题表述........................... 57 4.3 数值方法............................. 62 4.4 结果...............:................... 62个参考资料................................ 64 5 超聚焦声爆问题的分析和数值模拟 Cheng €4 Hafez................................ 67 5.1 介绍................................ 67 5.2 声爆在伽利略框架中分析的非等温大气中...................................... 68 5.3 超聚焦声爆问题...................... 74 5.4 闭幕词............................ 102 参考资料................................ 103 6 跨声速流动 Chen d Garabedian的复杂分析............................... 107 6.1 介绍................................ 107 6.2 霍多图变换..................... 108 6.3 奇点................................ 109 6.4 非线性边界值问题................... 111 6.5 算法的结构...................... 114 6.6 整合路径............................... 116................... 6.7 坐标的傅里叶分析 118 6.8 设计与分析的比较................. 120................................ 参考资料123 7跨声速小横向扰动方程及其计算罗。 Shen. d Liu................................ 125 7.1 介绍................................ 125 7.2 超音小横向扰动方程.......... 126.................. 7.3 轴对称入口127的计算................ 7.4 翼型和Wing 132的多重解决方案 7.5 结论................................ 137个参考资料................................ 139 8 边界层流中绝对不稳定扰动的激发 Ryzhov d Terent'ev................................ 141 8.1 介绍................................ 141 8.2 扩展三层甲板模型...................... 143 8.3 线性分析.............................. 144 8.4 计算结果................................ 146 8.5 理论论点................................ 148

内容............................... 8.6 结论... 致谢............................... 参考资料................................. 9 在T1622的伴随方程上，CFD GQes中的最优网格自适应............................................ 9.1 介绍................................ 9.2 有限体积分析................................ 9.3 有限元分析................................ 9.4 一些总结性评论...................... 致谢................................ 参考资料................................. 10 流量计算中增加的消散 MacComack.................................... 10.1 介绍................................ 10.2 黏性应力张量........................ 10.3 添加数值黏性的选择.................................... 10.4 计算结果 10.5 结论................................ 参考资料................................. 11 A 四运算符 守恒格式 用于 Euler 方程...................................... Chattot 11.1 介绍................................ 11.2 Burgers 方程的混合方案................. 11.3 1-D Euler 方程............................... 11.4 一个特征性的盒子方案...................... 11.5 使用离散特征值和特征向量............. 11.6 一些结果................................ 11.7 结论................................ 参考资料................................. 12 带分叉和翼襟翼的自动阻挡 Eberhardt €4 Wibowo.............................. 摘要.................................. 12.1 介绍................................................................ 12.2 方法 12.3 翻盖翅膀的结果........................ 12.4 结论................................ 参考资料................................. 13 局部预处理：操纵大自然母亲来愚弄她的时间 Damofal d Van Leer............................ 13.1 介绍................................

目录 13.2 设计注意事项................................ 212 13.3 当前状态.............................. 230 13.4 未来发展........................... 236 致谢................................ 237... 参考资料.................................' 237 14 放松重温-重新审视稳定流动的Multtgrid Roberts，Sidilkover d Swanson........................ 241 14.1 介绍................................ 241 14.2 多网的方法...................... 243 14.3 数学表述...................... 244 14.4 解决方案程序................................ 246 14.5 结果................................ 249 14.6 扩展到可压缩流动.................... 257 14.7 结论................................ 258 致谢................................ 258 参考资料................................ 258 15 平行计算机上的航空航天工程模拟Morgan。 Weatherill。 哈桑。 Broolces。 Manzari d Said........... 261 15.1 介绍................................ 261 15.2 求解算法................................ 262 15.3 内存要求................................ 266 15.4 结论................................ 273 致谢................................ 274 参考文献................................ 274 16 优化CFD代码和算法，用于Cray Computers Wigton...................................... 277 16.1 介绍................................ 277 16.2 Cray计算机，包括T9O................... 278 16.3 TRANAIR................................ 284 16.4 TLNS3DMB............................... 289 致谢............................... 291个参考资料................................ 291 17最近在空气动力学中的应用与NSMB结构化多块求解器韦伯，Gacherieu。 Rizzi et a1................................ 293 17.1 介绍................................ 293 17.2 物理模型.............................. 295 17.3 数字模型................................ 302 17.4 程序结构和多块实施......... 308 17.5 高性能计算策略................. 308 17.6 计算结果................................ 311

目录xiii 17.7 总结和结论...................... 325 致谢................................ 327 参考资料................................. 328 18 航空航天应用及其他领域的不可压缩Navier-Stokes计算................ 夸克。 Kiris。 达克勒斯。马里亚尼。 Rogers B Yoon 333................................ 18.1 介绍 333................................. 18.2 解决方法 334............................ 18.3 计算结果 343.............................. 18.4 结束 349 参考文献................................ 350 19 翼型优化 Drela的利弊...................................... 363 19.1 介绍................................ 363 19.2 方法总结................................ 364 19.3 低 雷诺数 翼型 应用............... 365 19.4 Transonic 翼型 应用...................... 372................................ 19.5 结论 377 参考文献................................. 380 20 走向工业实力 Navier-Stokes 代码。 重访Jou...................................... 383 20.1 介绍................................ 383 20.2 邮轮设计............................... 384 20.3 高升机翼分析和设计.................. 389 20.4 预测处理质量和控制效果...... 391 20.5 结论................................ 392 致谢.............................. 394 参考资料................................ 394 21 我们从计算流体动力学火车空气动力学研究中学到了什么？ 藤井和小川................................ 395 21.1 介绍................................ 395 21.2 数值法................................ 396 21.3 结果和讨论......................... 400 21.4 一种新的准一维预测方法........... 409 21.5 结论................................ 414 参考资料................................ 415 22 关于“与CFD Rubbert一起追求价值................................ 417 22.1 介绍................................ 417 22.2 空气动力学设计过程必须努力...... 419 22.3 导致价值的压倒性指标............. 420

目录 22.4 管理设计的等式................... 421 22.5 1990年之前CFD的影响或价值................. 422 22.6 20世纪90年代的时代................................ 423 22.7 关键推动因素............................ 424 22.8 未来方向................................ 425 22.9 结束说明........................... 427个参考资料................................ 427 23 CFD在十字路口：行业视角 Raj...................................... 429 23.1 介绍................................ 429 23.2 CFD在飞机设计中的作用..................... 430 23.3 CFD 飞机设计的有效性................. 431.................. 23.4 CFD 进步：过去三十年 432................ 23.5 CFD在哪里？ 未来二十年 439 23.6 结束之词........................... 444个参考资料................................ 445 24 航空航天工程 2000：综合。 实践课程 Seebass b Peterson................................ 449................................... 24.1“有史以来最大的礼物”449 24.2“最好的努力是不行的！”........................ 450............... 24.3 我们的客户和其他人怎么说？ 451.............. 24.4 综合教学实验室 452........... 24.5 航空航天工程科学课程 2000 455 24.6 两个警告.............................. 459 24.7 结论................................ 459 确认............................... 460个参考资料................................. 460 25 基于计算机的流体力学教科书 CaugheybLiggett................................ 465 25.1 介绍................................ 465 25.2 导航和集成...................... 467 25.3 插图和数据...................... 468 25.4 计算和实用工具...................... 472 25.5 教学经验.............................. 480 致谢............................... 481个参考资料................................ 481 作者字母列表....................... 483

\clearpage
\chapter{回顾Earll Murman对跨声速流动和计算流体动力学的贡献}
回顾Earll Murman对跨声速流动和计算流体动力学 Mohamed M的贡献。 哈菲兹1和大卫A。 Caughey 2 1.1 导言 本章将讨论Earll Murman过去三十年的工作，并将详细研究他的一些主要想法。 还包括他截至1996年的出版物的完整列表，以及他学生论文的标题。 （值得注意的是，Earll的博士论文是关于层压、高超音速唤醒的实验研究。） 默尔曼从普林斯顿大学毕业后，可以将他的论文分为两组；第一组属于1967年至1980年期间，当时他在波音、美国宇航局艾姆斯研究中心和Flow Research, Inc. 几乎所有这些论文都基于势流的跨声速、小干扰方程。 第二组的论文是从1980年他加入麻省理工学院后撰写的，它们涉及许多主题，包括欧拉和Navier-Stokes 方程的解决方案。 例如，如果将1962年跨声速研讨会的会议记录与1975年后续会议的内容和加州大学戴维斯分校机械和航空工程系的内容进行比较，那么Murman对跨声速流动的贡献的重要性就很明显了，加利福尼亚州戴维斯，95616。 2 西比利机械与航空航天工程学院，康奈尔大学，伊萨卡，纽约 14853-7501。 Frontiers of 计算流体动力学 - 1998 编辑：David A. Caughey h Mohamed M. 哈菲兹 ©1998 世界科学

2 HAFEZ \& CAUGHEY 1988。 Murman工作的影响体现在后两次会议上涉及跨声速流动模拟的论文的质量和数量上。 他在第三届研讨会（1988年）的会议记录中关于三角翼上的涡流的论文反映了他对相关现象的物理学和数值的强调。 另一方面，Murman对CFD的贡献更令人印象深刻。 介绍类型依赖性差异化的概念、上风方案和人工黏度之间的关系、适当的线性化，并将松弛视为人工时间依赖过程，然后是守恒的重要性--不仅在差分水平，而且在离散水平上--以及正确捕获激波的适当开关的构建；这些是一些新诞生的概念，首次使用20世纪70年代初的计算机，在空气动力学配置上为跨声速流动产生二维和三维解决方案。 合理的预测，在跨声速制度中，不仅是升力，也是非黏性阻力。 可以说，他的工作提供了第一个令人信服的证据，导致航空航天业认真考虑CFD作为补充风洞的分析工具。 1.2 对跨声速流动的贡献 关于跨声速流动的理论工作可以追溯到二十世纪初，Chaplygin关于气体喷射的工作。 到20世纪40年代末，Oswatitsch、Guderley、Busemann和von Karman独立推导出平面流（K*x - \#/2）x + <t>vy = 0的非线性、超音速、小扰动方程。 在这个方程中，<j>是缩放扰动速度势，x和y分别是平行和垂直于自由流方向的笛卡尔坐标，K是跨声速相似性参数。 3 1947年，冯·卡曼推导出了上述方程形式，确定了跨声速相似性参数并建立了跨声速相似性，同时发现了马赫数冻结原理。 Oswatitsch \& Keune基于细长体理论发现了跨音域规则，并由Whitcomb进行了实验验证。 Guderley在1957年写了第一本关于跨声速流动理论的书，Ferrari \& Tricomi的Transonic Aerodynamics于1962年出版。 Landahl's 3 对于平面流，缩放是K oc (1 - M 2 )/T 2 / 3；对于跨音子小扰动理论的轴对称版本，相似性参数是K «■ (l-M 2 )/r 2

EARLL MURMAN 3 Unsteady 跨声速流动的贡献于1960年出版，Ashley \& Landahl的《翅膀和身体的空气动力学》于1965年出版。 冯·米塞斯的可压缩流体流动数学理论和贝尔斯的亚音速和跨声速气体动力学数学方面都于1958年出版。 Bitzadze的混合类型方程和Manwell的Hodograph方程分别出现在1964年和1971年。 在这方面，还必须提到Germain在1964年的评论文章。 1986年，Cole \& Cook在《跨声速航空动力学》上发表了迄今为止对跨声速、小干扰理论的最全面的处理。 早期关于跨声速流动的大部分工作都基于hodograph方法，其中因变量和自变量互换，结果是平面流动方程变得线性，允许通过叠加构建解决方案。 这种方法仅限于等熵、非旋转、二维稳定流动，冲击的处理很复杂。 Guderley \& Yoshihara和Cole在20世纪50年代初研究了有限楔形的简单几何学。 Vincenti \& Waggoner在统一的自由流马赫数上计算了双楔形剖面和倾斜平板的流量。 一般翼型形状的固体表面的边界条件在hodograph平面上极其难以处理。 Nieuland \& Boerstoel在20世纪60年代通过扩展Lighthill \& Cherry的早期工作，设计了无冲击机翼。 在20世纪70年代初，Garabedian \& Korn成功地应用了复杂特征的方法来处理混合型方程，并能够设计超临界（无冲击）机翼部分，包括被广泛研究的Korn 翼型。 （见本卷中Garabedian和Chen的相关作品。） 鉴于对Morawetz和其他人早期工作的普遍误解，暗示不存在连续的跨声速流动过去配置文件，这些贡献尤为重要。 4 事实上，这项工作并没有排除没有相邻无冲击解的无冲击奇异解的存在。 皮尔西和后来的惠特科姆成功地通过实验证明了基本上无冲击的机翼可以实现和使用。 从工程角度来看，允许存在非常弱的冲击，因为它们不会破坏阻力上升马赫数以下设计的效率。 另一种公式，仅限于稳定的二维（包括轴对称）流动，是基于流函数的使用。 在20世纪40年代初，埃蒙斯使用这种公式计算了通过机翼的跨声速流量。 他的有限差分解是用激波拟合的松弛程序手工计算的。 同样在20世纪40年代，Macoll解决了最近，计算机模拟被用来研究存在和独特性问题。 SteinhofT\&Jameson，以及后来的Salas等人，展示了势流解决方案的非独特性。 1991年，詹姆森对Euler 方程的稳态解提出了类似的问题。 正如Hafez \&; Guo所证明的那样，即使对于高雷诺数、Navier-Stokes 方程的附着流动解决方案，这个问题仍然存在。

4 HAFEZ \& CAUGHEY分离电击问题。 莫雷蒂和他的学生后来研究了冲击装束：例如潘多尔菲、萨拉斯、格罗斯曼等人。 在他1969年的博士论文中，Steger试图使Emmons的计算自动化，并报告了收敛的困难。 由于同样的困难，Sells的分析仅限于纯粹的亚音速流动。 5 此外，20世纪50年代引入了基于流线型曲率方法的相关公式，特别是用于内部流，但直到20世纪80年代Drela和Giles的工作，将其扩展到跨声速流才成功。 即使使用势流公式，其中密度是根据流体速度（使用伯努利方程）唯一定义的，许多模拟混合流动的尝试也惨遭失败。 为了从角度看待Murman的贡献，我们首先讨论其中的一些尝试。 为了解释可压缩性效应，1939年引入了Karman-Tsien（切线气体）近似法。 正如Lighthill在他的《普林斯顿高速空气动力学系列》的《高近似》部分中所讨论的，即使是纯粹的亚音速流（有升力），这种近似也注定要失败，因为它与基本物理不一致。 Latoine在1951年引入的另一种临时方法是基于Prandtl-Glauert变换，将可压缩流动的压力系数与相关的不可压缩流动联系起来。 如果使用本地而不是自由流马赫数，则可以找到更好的近似值。 Spreiter \& Alksne后来也采用了类似的局部线性化方法。 与此同时，超音速流得到了更系统的处理，从Ackeret对Prandtl-Meyer流的线性化近似，以及Busemann对斜冲击关系的二阶和三阶展开开始。 Rayleigh \& Janzen和Hentzsche \& Wendt分别进行了自由流马赫数 M\textasciicircum{}和厚度比T的膨胀，Imai试图扩大这些膨胀的有效性范围，以处理更高马赫数的亚音速流动。 Oswatitsch使用带有椭圆内核的积分方法进行亚音速流动。 非线性项滞后，每次迭代时都会求解泊松方程。 对于超音速流，在每次迭代时用从上一个迭代中评估的非线性项（强迫函数）求解非均匀波方程（具有恒定系数）。 同样，通过Oswatitsch的《流函数方法》获得的序列主要困难源于这样一个事实，即流体特性（如密度和压力）与质量通量密度pq之间的关系不是单值的。 在声波点有一个平方根奇点，其中质量通量密度为最大值（对应于窒息现象）。更大的通量密度值会产生密度的复杂值。 这个问题可以被规避，正如哈菲兹在20世纪80年代所展示的那样，可以计算出带有冲击的非旋转和旋转混合流动。

一般来说，对于混合（亚音速和超音速）流动，EARLL MURMAN 5方法的贡献不会收敛。 Oswatitsch引入了一种以声速旋转的抛物向方法。 在这种方法中，加速度项（流动方向上势函数的二次导数）是滞后的，并且在每次迭代时求解抛物线方程。 （Ryzhov、Cole等人后来开发了声波流动问题的优雅渐近理论和马赫数冻结原理。） Zierep \& Gullstrand将积分方法应用于跨声速流，后来在20世纪70年代，Hancock \& Nixon和Spreiter \& Ogana将集成方法应用于跨声速。 这些贡献的一个关键点是在音速条件下正确处理鞍座点。 使用了具有正常冲击的准一维近似，并获得了合理的结果。 有趣的是，对于一个小型的嵌入式超声速流动区域，如果将人工黏度项作为强迫函数的一部分添加到方程中，基于泊松方程的迭代确实会收敛。 然而，在这种情况下，收敛取决于添加的人工黏度。 森野在20世纪80年代引入的一个更通用的方法是基于时间依赖（非稳定）方程，其中积分公式是根据局部条件具有可变系数的波方程构建的。 原则上，莫里诺的公式对跨声速（混合）流是有效的，但由于评估与时间相关的内核需要密集的计算，因此成本很高。 也许跨声速小扰动理论最重要的发展是科尔的渐近推导，其中自由流马赫数 M\textasciicircum{} ->■ 1和厚度比r -► 0，而跨声速相似性参数K保持不变。 适当的排序，并与边界条件匹配，产生了著名的非线性、跨声速、小扰动方程作为一阶近似。 跨声速流动本质上是非线性的，所有基于线性化的尝试都失去了基本现象的这一重要特征。 在Cole \& Messiter和Cole's Twenty Years of 跨声速流动中，有两个最终结果相同的推导；一个从全势方程开始，另一个以Euler 方程作为基本守恒定律。 这个练习不仅证明了一致性，还阐明了阻力的评估，稍后将讨论。 1965年，Hayes对跨声速流动的更高近似的完整研究表明，在第一个近似中，流动仍然是潜在的。 还表明，由守恒形式推导方程的弱解所接受的跳跃条件与兰金-胡戈尼奥特关系的扰动形式一致。 另一方面，特征形式揭示了稳定的跨声速流动的混合型性质。 科尔系统地开发了在亚音速、声波和超音速流中提升机翼和机翼的解决方案的远场行为。

6 HAFEZ \& CAUGHEY Cole \& Sichel推导出黏性跨声速方程，用于描述冲击结构和冲击交叉问题。 在消失黏度的极限中，超音速小扰动方程被恢复。 研究了与瞬态现象相关的非稳态方程。 非线性仍然是低频制度的重要组成部分。 由于其中一个特征与流动方向对齐，科尔将该方程命名为“顺交”。 根据指定的约束，时间依赖方程可以描述一个具有可能稳态解的良好初始值问题。 需要注意的是，<f>t = (1 - M 2) 4> XI + <j>yy或4>tt = (l - M 2) 4> xx + <j)yy形式的方程，其中1 - M 2 = 1 - M£，- (7 + V)M\textasciicircum{} X和M是局部马赫数，一般来说，不能描述良好的非线性初始值问题，而由a<t>xt = (l - M 2) cj) xx + <j)yy给出的非稳态、超音素、小扰动方程可以。 事实上，这一观察解释了一些标准（Jacobi和Liebman）迭代方法在应用于跨声速流动问题时的失败。 继加拉贝迪安根据它们近似的时间依赖性偏微分方程（所谓的时间依赖类比）对连续线过度放松（SLOR）迭代程序的解释和分析之后，詹姆森解释了穆尔曼选择随流行进的线松弛的成功，因为它模仿了非稳定的跨音小扰动方程，特别是与其时间依赖的模拟相对应的修改方程中与4>xt成正比的项的出现。 因此，将稳定的跨音子方程嵌入到适当类型的高维（时间依赖）问题中对迭代过程的收敛至关重要。 以守恒形式添加黏性术语对于捕捉具有正确强度和位置的冲击至关重要。 组合方程由a4>xt = (K4> x - 4> 2xl2)x + <t> yy + (v(f> xx)x或更通用的形式6，包括高频项<t>u + a<t> It = \{K4> x - 4> 2x/2)x + 4> yy + \{v<j> xx)x 此形式允许基于时间显式的三级方案对跨音流进行向量平行计算（见Keller、Jameson、South和Hafez）。

EARLL MURMAN的贡献 7 人们可以将人工黏度视为一种正则化工具，它排除了膨胀冲击，并通过声学条件强制实现平稳加速。 冲击将被涂抹，在每次迭代时都可以使用牛顿方法和基于高斯消除的直接求解器获得数值解，因为即使不存在非稳定项，包括离散黏性项的贡献在内的矩阵也是可逆的。 直接解决方案是可行的，至少对于适度的二维问题。 然而，对于大型问题，迭代方法是必要的（在除最强大的计算机外的所有计算机上），并且必须隐含（如线放松与流进行）或显式地包括与时间相关的术语。 非线性的重要性和控制方程的局部特征、黏度的作用以及理解与非稳定方程相关的瞬态过程是成功模拟稳定、跨声速流动的必要因素。 拥有这种知识体系，Murman于1968年开始与Cole合作。 Murman在正确的时间是正确的人。 他负责使用波音公司可用的计算机获得稳定的跨声速流动问题的数值解。 上述许多概念在20世纪60年代末并不明确，实施问题可能会为更多的障碍打开大门。 穆尔曼对物理学和数值的深刻理解导致了第一个稳定的实际计算机模拟，跨声速流动与激波。 数字的细节将在下一节中讨论，特别是Murman-Cole方案的数值稳定性、从一个区域到另一个区域的开关设计以及Murman的完全保守运算符。 在本节的其余部分，讨论了Murman关于超音速空气动力学的论文。 一旦跨声速流动模拟工具可用，显然需要准确测定空气动力学载荷。 升力可以从表面压力分布或平衡与飞行方向垂直的动量计算，在翼型周围的控制体上。 由于使用了线性化边界条件和解在前缘的奇异行为，控制体方法可能更适合基于小扰动方程的解。 如果控制面是在远域中选择的，它可以显示，假设只有正常的激波，L = - />oo<7ooor where /? <»和?„，分别是自由流密度和速度，V是环流，可以计算为尾缘上表面和下表面的电位值之间的差值。 因此，Kutta-Joukowsky定理对超音子流有效。 然而，非黏性阻力的测定更加微妙。 无冲击的

8 HAFEZ \& CAUGHEY机翼的阻力为零，这与D'Alembert的悖论一致。 假设品质、动量和能量是守恒的，Oswatitsch和Liepmann独立地在熵上升方面，按一阶计算了因冲击而引起的阻力。 7 然而，对于势流，质量、熵、涡量和能量是守恒的--垂直于激波的动量不是。 阻力可以通过流中出现的所有激波的净动量不平衡来计算。 或者，人们可以想到，如果正常动量得到守恒，就会在整个冲击中产生的熵上升，从而导致相同的结果，到一阶。 对于正常冲击来说，这很容易显示出来，因为冲击时电位公式中的净动量是7+1 3 [pu(u + l)+p] ss jj-H，其中u和p分别是（非维）扰动速度和压力，7是比热的比率。 即使对该领域的一些专家来说，这种发展也不清楚，他们声称势流计算的阻力应该因为等熵条件的强制执行而消失。 Germain和Cole给出了正确的表述和解释。 Steger \& Baldwin分析了势流中的阻力，但他们的差异方案并不保守；因此冲击强度和位置不一定正确。 对于质量不守恒的解决方案，阻力评估必须按照Garabedian的描述进行校正。 在Murman \& Cole（1974）的一篇杰出的AIAA论文中，这些问题得到了澄清，Murman的守恒格式用于计算表面压力分布和冲击中“熵”上升的积分的阻力。 包括风洞边界的流动，其中必须考虑通过风洞壁的动量通量的贡献，以及自由空气中的机翼，并将结果与实验数据进行了比较。 如果将跨音小扰动方程乘以流向速度u，然后从不旋转条件乘以跨流速度v中减去，则以发散形式得到动量方程。 然而，这并不符合物理保护原则。 默尔曼和科尔整合了这个方程，解释了跳跃7，这里给出了一个简单的证明：从克罗科关系中，由于熵变化导致的扰动速度u由-u y;s £=■=\_ SAs Qr - \textasciicircum{} \_Xsa给出。 因此，阻力是-pqudy=-\textasciicircum{}- / pqA ■I q°°./休克冲击

EARLL MURMAN 9在激波上的贡献，以获得阻力积分\textasciicircum{}-/\{('T-T- 2 \textasciicircum{})*-'-"\}- 3 \textasciicircum{}/\textasciicircum{}-!” *在这方面，人们应该提到Engquist和Osher在构建1980年计划时使用的“熵”不等式。 他们将跨音小扰动方程视为Burgers 方程的非标准二维扩展。 在任何局部控制体 Q上，不包括固体表面，在\{ \textbackslash\{\}KT + T + \textasciicircum{}T\textasciicircum{}"3) dy + \{uv) dx\textbackslash\{\} \textasciitilde{} °中，平滑流的等式成立。 Engquist \& Osher的计划旨在保证不会发生膨胀冲击。 还设计了二阶上风方案，并证明了其非线性稳定性。 8 还表明，不可能得到高于二阶的这种稳定的近似值。 Murman用他的方案计算了通风风洞壁对提升跨声速翼的干扰效应，以取代线性亚音速理论中不充分的经典公式。 后来，Murman引入了一种跨声速风洞壁干扰的评估/校正方法。 该方法基于寻找自由空气马赫数和攻角，在测量和校正的压力分布之间提供最佳最小二乘匹配。 该程序扩展用于评估三维模型的墙体干涉校正。 对于接近单位的马赫数的细长配置，采用了等价规则公式。 相反，对于高长宽比的机翼，可以使用跨声速提升线理论。 Murman率先将数值优化技术应用于翼型设计。 他与美国宇航局Ames的Vanderplaats和Hicks联手，研究具有现实客观函数和几何约束的自动翼型设计的实用程序。 Murman与Seebass合作，计算了不连续信号在声学附近的行为。 他还使用谐波扰动方法解决了非稳态超音子方程。 Murman为跨声速激波-边界层相互作用的位移厚度效应引入了一个现象学模型。 9 1976年，Warming \& Beam引入了Murman类型的二阶迎风方案。 矛盾的是，即使在引入正确的保守形式来区分势流方程之后，许多工业工程师更喜欢使用非保守的数值方案解决跨声速势流问题，因为它们

10 HAFEZ \& CAUGHEY 最后，我们想提一下，Murman教授的第一门课程是1973年在斯坦福大学教授的“超音速空气动力学”。 他有效指导的第一个博士生是詹姆斯·克鲁普，他在波音公司工作，于1971年毕业于华盛顿大学。 克鲁普的论文题为“平面的数值计算，稳定，跨声速流动通过薄，提升机翼。” 1.3 对CFD的贡献 在本节中，与前一节一样，我们将首先简要介绍CFD的历史（因为它与跨声速流动模拟有关），特别是其在20世纪60年代末的状态，然后将讨论Murman的贡献及其对该领域的影响。 我们将从双曲和抛物线偏微分方程组开始，然后是椭圆方程和混合方程。 1.3.1 双曲方程 也许在这方面的第一个重要贡献是Courant, Friedrichs \& Lewy在1928年的论文，其中推导出了CFL稳定性条件。 10 在20世纪40年代末，特征方法被开发出来，并应用于二维（平面和轴对称）问题，包括稳定、旋转流。 在20世纪50年代，有限差分方法开始吸引科学家和数学家的注意，正如Richtmyer \& Morton的第一版所指出的那样。 线性初始值问题的Lax等价定理是这一时期的重要贡献之一。 正如Richtmyer \& Morton第二版所讨论的那样，俄罗斯的Oleinik和美国的Lax \& Wendroff对非线性守恒定律及其弱解的理论研究，以及Lax-Wendroff方案的出现，是20世纪60年代的重要里程碑。 （Lax-Wendroff方案取代了Lax-Friedrichs方案，后者具有高度耗散性，只有一阶精度，人工黏度项的系数与通常与风洞数据更好的光谱半径成正比。 这既是由于非保守方案在定性上模仿激波-边界层相互作用效应的倾向，也是由于潜在解决方案中的冲击比相应的兰金-胡戈尼奥特跳跃更强的倾向。 Melnik ei al.表明，包括边界层效应在内的保守计算产生了正确的结果，阻力极性与实验数据一致。 1 后来，冯·诺伊曼·傅里叶分析出现了。 Godunov和Ryabanikii在关于有限差分的书中，以及Kreiss Gustaffson、Osher、Beam \&; Warming等人深入研究了稳定性问题。

BARLL MURMAN 11 Jacobians的贡献。） Lax（以及Hopf和Kruzkov）在20世纪70年代初研究了熵不等式的作用。 Friedrichs和Friedrichs \& Lax研究了对称双曲系统，Mock等人给出了其“数学”熵函数的系统构造。 在20世纪80年代，Harten引入了熵变量来对称Euler 方程，后来Hughes表明，物理熵是Navier-Stokes 方程具有传热的唯一可接受的函数。 11 MacCormack开发了一种高效的Lax-Wendroff方案形式，该方案避免了Jacobians的计算，他的方案在20世纪70年代被使用有限差分和有限体积方法的研究人员广泛用于可压缩流动计算。 后来，Briley \& McDonald和Beam \& Warming推出了他们的交替方向隐式（ADI）方案。 在这一类别中，Lerat最近的隐性因子方案变体被认为是最先进的。 与此同时，詹姆森在20世纪80年代引入的显式方案的改进，特别是他使用多重网格法成功地将收敛加速到Euler 方程的稳态，至少在跨声速流动计算方面取得了显著的进步。 （詹姆森使用了直线法，其中空间导数，包括二阶和四阶耗散项，首先被离散化，然后使用多阶段Runge-Kutta方法求解由此产生的常微分方程（时间）系统。 二阶和四阶黏度术语是由MacCormack首次引入的；事实上，使用二阶黏度来捕获激波可以追溯到von Neumann \& Richtmyer，Kreiss倡导使用四阶数值黏性进行平滑流动；Jameson引入了一种新的混合（或开关）公式，对锐化激波特别有效。） 另一方面，戈敦诺夫于1959年在俄罗斯引入了基于特征的方案。 Godunov方案直到20世纪60年代末才在西方流行（见Holt关于CFD的书）。van Leer和其他人试图构建高阶扩展（因为一阶Godunov方案“不够好！”）。 在20世纪80年代初，Osher引入了基于近似黎曼问题的精确闭合形式解的方案，而Roe引入了基于一种新的线性化导致黎曼问题的近似解的通量差分分裂。 在这方面，人们应该提到Steger \& Warming的开创性工作，以及van Leer在构建通量向量分割方案方面的修改。 （拆分是基于雅各布矩阵在每个方向上的特征值的符号。） 最近，Liou引入了替代拆分方案--称为“AUSOM”--具有非常理想的属性。 人们注意到，在20世纪60年代末，上风11 Magnus U Yoshihara是第一个将跨声速小扰动方程写成一个具有人工时间依赖项的对称、双曲方程组的人。 哈菲兹直到1998年才将整个三维问题扩展到。

12 HAFEZ和CAUGHEY方案自然导致对角线主导系统，因此放松算法能够为Euler 方程提供有效的解决方案。 Hirsch的第二版书详细讨论了其他贡献，包括限制器、总变异减小（TVD）和基本非振荡（ENO）方案。 最近，重点放在构建完全避免黎曼求解器的真正多维方案上。 这项工作正在进行中，陪审团仍未确定这种方法是否会占上风。 显然，Murman不知道Godunov的工作，但他确实在1971年参加了加州大学洛杉矶分校关于双曲系统和守恒定律系统的短期课程。 这些讲座由Lax（应Cole的邀请）发表，随后由SIAM出版。 当时，试图将Lax-Wendroff方案（包括MacCormack的变体）应用于Euler 方程以解决跨声速流动问题，但并没有产生令人印象深刻的效果。 计算似乎完全不切实际，因为当时演算法和计算机的发展状态。 然而，Murman成功地将数值守恒的概念应用于跨声速小扰动、势计算，并引入了Murman-Cole方案的完全保守版本，这对于以数值方式捕捉正确的弱解（具有正确的激波位置和强度）至关重要。 后来，Murman为模拟涡流的Euler 方程解决方案做出了贡献，如下文将讨论。 1.3.2 抛物价方程 Euler 方程是Navier-Stokes 方程的隐性极限。 即使在解决完整的Navier-Stokes 方程时，也必须为非黏性方程有一个好的方案，否则人工黏度项可以主导真正的黏性效应（特别是对于高雷诺数流量）。 当然，湍流建模仍然是高雷诺数实际应用的主要绊脚石。 有趣的是，求解可压缩Navier-Stokes方程的尝试是基于抛物线样方程组（包括连续性方程中的人工黏度），而大多数求解不可压缩Navier-Stokes 方程的方法（至少直到最近）都基于压力校正程序，每次迭代都会求解泊松方程。 即使是Chorin、Temam等的人工可压缩性方法，将与时间相关的压力项添加到连续性方程中--为了有一个良好的初始值问题--通常被解释为将松弛（Jacobi）方法应用于压力的泊松方程。 对于不可压缩流动，连续性方程被认为是一个约束--即压力被确定为满足质量守恒。 人们可以把压力想象成拉格朗日

变分框架中EARLL MURMAN 13乘数的贡献。 有关更多详细信息，请参阅本卷中Temam关于Navier-Stokes 方程的书和Kwak等人的文章。 最近，人们对研究可压缩流动计算的不可压缩极限非常感兴趣。 欧拉和Navier-Stokes 方程的现有可压缩流动代码可以使用预处理技术（基本上）产生不可压缩流动解决方案（见van Leer、Turkel和Muradoglu \& Caughey）。 另一方面，根据合适的变量重新表述问题，包括适当的非维度化，人们可以有一个统一的代码，能够处理所有速度制度下的流动（见Merkle，Hafez）。 此外，一些基于声学模式熵对流和总焓对流分解的公式允许不可压缩的解（见Roe、Deconnick、Tassan和Hafez）。 对于高雷诺数流量，Navier-Stokes 方程解的替代方法是黏性/非黏性相互作用过程，其中边界层方程的解与无黏流动模型的解相匹配。 在这方面，Murman \& Bussing做出了开创性的贡献，他们展示了边界层解决方案如何与Euler 方程的解决方案相结合。 此外，Murman \& Lloyd为Navier-Stokes计算引入了隐式Runge-Kutta方法，其中方程隐式在边界层和尾迹区域的垂直方向上求解。 1.3.3 椭圆方程 对于稳定的二维亚音速流动，流函数已被用于求解非黏性、非旋转和旋转流动。 流函数涡流公式被解为两个椭圆方程组或单个双谐波方程。 对于二维流，可以使用特殊技术（例如，复杂变量）。 对于轴对称和三维流动，必须使用更通用的方法，如格林函数或积分方程方法。 过去，面板方法在空气动力学界被广泛用于不可压缩的势流计算。 其中一些想法最近在其他学科中以边界元素方法的名义被重新发现和发展。 对于椭圆问题，还有经典的变分方法，包括瑞利-丽兹、加勒金和最小二乘公式。 对未知的局部表示的有限元素进行了严格的研究，并出现了一个令人印象深刻的理论。 许多将有限元方法应用于流体流动的尝试也取得了成功。 然而，有限元法应用于双曲方程并不具有相同的严谨性。 Donnea介绍了Taylor-

14 HAFEZ \& CAUGHEY Galerkin方案，这是Lax-Wendroff方案的有限元模拟，而Hughes引入了最小二乘Galerkin过程。 对于对称双曲系统，这两种方案可以是相同的。 流线型上行也引入了有限元计算，基于特征的方案也引入了。 （见本卷中Morgan等人的文章。） 无论如何，有限元在使用非结构化网格处理复杂几何时的优势是显而易见的。 最近，网格生成（结构化和非结构化）以及网格的调整，以最佳方式解决解决方案的特征，本身就已成为一个重要主题。 目前，存在基于分析（复杂变量）、代数、变分和微分方程公式的方法；后者包括基于椭圆（泊松和双谐波）、抛物线和双曲方程的方法。 各种几何方法可用于非结构化网格。 此外，经典优化和伴随方程方法被用于适应。 变分公式可用于获得椭圆问题的离散化方案。 同样，有限差分方案可以基于泰勒级数展开或积分公式。 后一种公式对于构建有限体积方案很有用，特别是在非结构化网格上。 正如Jameson \& Caughey的全势方程计划所证明的那样，这些方案不一定局限于椭圆问题。 当然，人工黏度、适当的开关和仔细开发迭代算法对于跨声速流动计算是必要的，就像小扰动方程一样。 还有一个完整的迭代算法理论，用于求解由椭圆问题产生的一致离散方程。 变量方案（梯度和共轭梯度）以及松弛方法（Jacobi、Gauss-Seidel、连续超过松弛以及用于并行计算的利布曼和棋盘方法）都得到了很好的研究。 基于块隐式方法、外推技术和ADI（对其参数进行周期性评估）的收敛加速随后出现。 然而，最成功的加速技术是基于共轭梯度类型方法（例如，Saad的GMRES）和多网格（见Fedorenko、Bakhvalov和Brandt）。 这两种策略都已成功扩展到非椭圆问题。 1.3.4 混合型方程 描述稳定跨声速流动的方程是混合型的。 最简单的例子来自跨音小扰动近似，与Euler 方程不同，不清楚如何将其嵌入描述良好初始/边界值问题的人工时间依赖系统中。

EARLL MURMAN 15的贡献 如前所述，最早的尝试是基于hodograph变换，它产生了一个线性方程（尽管仍然是混合类型）。 Tricomi研究了后者的模型方程。 Waggoner \& Vincenti和Filipov报告了使用亚声速流动的有限差异和超音速带的特征方法对Tricomi方程进行数值求解的尝试。 正如Oswatitsch的《气体动力学》一书中所讨论的那样，特征方案可以被解释为有限差分方案。 Katsanis使用Lesaut开发的特殊有限元法来求解Tricomi方程。 Murman和Cole是第一个在物理空间中求解非线性方程并自动捕获激波的人。 由于该方程是混合类型的，因此在亚音速和超音速区使用不同的近似值是很自然的（例如，正如Oswatitsch所建议的）。 显然，中心差异应该用于亚音速（椭圆）区。 然而，在超音速（超斜）区域使用明确的上风方案在声线上是不稳定的（因为当本地马赫数接近统一时，流动方向的网格大小必须消失）。 Murman和Cole引入了隐式后向差分方案，尽管它只有一阶准确（或者用希尔德布兰特的术语来说过度阻尼），但基于任意网格大小和（超音速）马赫数的线性分析，它是无条件稳定的。 这是Murman和Cole必须克服的第一个障碍，因此他们的工作与Emmons或Oswatitsch的提议不同。 另一个问题与从一个区域切换到另一个区域有关，即使是没有冲击的流动。 必须使用测试函数（基于局部马赫数）来检查流量是亚音速还是超音速。 Murman和Cole决定使用中心差分近似值（在均匀网格上是二阶精度）来计算马赫数。 根据现有信息，该点被标记为椭圆或双曲，并使用适当的模版来离散方程。 考虑超音速带的第一个点。 测试函数（T = 1 - M2）将基于速度的中心差近似值为负数。 然而，如果使用后向差分方案来评估测试功能，它将是积极的。 切换导致的这种差异将导致局部近似不一致，并最终导致数值不稳定。 为了避免这种困难，Murman和Cole假设马赫数此时接近于单位，并忽略了系数（1 - M 2），因此方程退化为<f> yy = 0的近似值。 他们称这个点为抛物点。 事实证明，正如Murman后来在他1974年的论文中展示的那样，这种近似值是数值保守的。 Murman \& Cole还决定使用线松弛，沿着流动的方向行进，垂直线上的所有未知数同时被解决。 大于1的过度松弛参数仅用于椭圆点。 对于跨声速流动病例，收敛有可能在几百次扫描中稳定状态，从均匀流动开始

16 HAFEZ和CAUGHEY的初步猜测。 对于超音速自由流，Murman发现了收敛的困难，他没有得到正确的解决方案。 这导致他研究了方案的数值守恒，后来他引入了他的冲击点算子（这是双曲和椭圆算子的总和--不是平均值）。 冲击点算子与局部正常冲击相匹配，因此质量是守恒的；如果没有这个特殊算子，净质量源会以数值方式生成，由此产生的冲击通常比它应该的要弱。 通过这种修改，Murman能够为亚音速和超音速自由流产生正确的解决方案。 他还能够在（无冲击）设计和非设计条件下计算出Korn 翼型的流量。 尽管Murman-Cole和Murman的完全保守计划已经存在了近三十年，但描述它们总是耐人寻味的。 也许是因为这两个方案都基于不止一个概念，而且实施方式很好地反映了每个概念的作用。 也许也是因为这一贡献影响了许多研究人员的职业生涯，包括本卷的作者和编辑。 1.4 穆尔曼早期贡献的影响 穆尔曼工作的直接影响是巨大的。 他在跨声速小扰动方程方面的开创性工作提供了高效模拟机翼和机翼上的跨声速流动的能力，不仅在学术界和政府机构，而且在工业界（特别是在波音公司），以及许多航空航天发动机公司。 一波活动也在欧洲蔓延，例如，英国的皇家飞机机构、荷兰的NLR、法国的ONERA和德国的多尼尔。 由Harv Lomax领导的加利福尼亚州美国宇航局艾姆斯研究中心的CFD分部，以及弗吉尼亚州美国宇航局兰利研究中心的Jerry South领导的类似小组，开始密集地进行跨声速流动模拟。 Ballhaus \& Goorjian致力于将Murman方案应用于不稳定流动。 Jameson \& Caughey为全势方程开发了扩展。 Hafez、South和Murman以及Hoist根据Eberle的有限元计算引入了延迟密度公式。 在法国，Chattot引入了他的盒子计划，Engquist \&; Osher开发了他们的熵不平等满足计划。 人们注意到，Murman的方案允许膨胀冲击（（j> yv = 0的解可以是x的不连续函数），12和Engquist和Osher引入了2，一个简单的修复方法是将原始方程以抛物线方程的形式放在x和

EARLL MURMAN 17的贡献 具有连续导数的数值通量函数（包括耗散）的近似值，并证明了他们的方案不能发生膨胀冲击。 Osher、Hafez和Whitlow将这项工作扩展到二维势流方程，Osher将这个方案推广到Euler 方程。 与此同时，Roe介绍了他的计划，这可以被解释为将Murman计划推广到Euler 方程。 （Chattot在本卷中描述了这个方向的另一个尝试。） 即使在今天，许多人过去和现在都在使用Murman的计划。 在Murman的工作激励下，Hafez和Cheng使用类似艾特肯的外推技术和激波拟合方案进行跨声速流动计算，研究了收敛加速。 Cheng \& Hafez在跨声速等价规则和跨声速提升线理论的工作中也使用了Murman的方案。 哈菲兹基于跨声速小扰动方程的变分公式，研究了穆尔曼类型依赖方案的有限元模拟。 最近，Cheng \& Hafez在研究超聚焦声爆问题时使用了Murman的计划，正如本卷文章中所讨论的那样。 Cole \& Cook在他们的工作中也使用了Murman的方案，包括在hodograph平面上设计无冲击翼型，以及解锥形坐标中的非稳定方程。 Rusak还用它来计算真正的气体效应。 在文献中可以找到许多其他应用程序，但没有试图在这里提供完整的列表。 Murman工作的间接影响也很重要。 它迫使增强基于Euler 方程的稳定和非稳态流动的计算，例如，在Caughey关于应用多网格加速Euler 方程解决方案的收敛的工作中。 随后，Caughey将他的工作扩展到Navier-Stokes的计算。 其他人试图纠正经典的电位公式，以包括熵和涡量效应。 哈菲兹最近展示了一种有前途的方法，至少在外部空气动力学方面。 它可以作为异构区域分解实现，基于流动物理学的自适应建模，包括黏性效应。 事实上，很难解释Murman工作的所有影响，我们将这一部分限制在我们自己的个人观点上。 1.5 Murman在加入M的教师后，Murman的最新贡献。 我。 T.，他把注意力转向了Euler 方程的解决方案。 特别是，他对三角洲翅膀上的流动以及捕捉尾迹及其卷起感兴趣。 y的隐形极限；见本卷中Cheng \&; Hafez的文章。

18 HAFEZ \& CAUGHEY 提升翼的尾部是接触不连续性--在许多方面，这是一个比激波更困难的问题，因为压力是连续的，而速度的切向分量可能会跳跃。 Murman \& Powell研究了具有连续速度曲线（由于人工黏度）的有限厚度后的总焓损失。 他们的工作基于以下观察。 如果整个尾部的压力和密度没有变化，总焓不能保持恒定，因为f-lp 2因此，尾部的结构是人工的，类似于无黏流动计算中的冲击结构问题。 Murman \& Modiano和Murman \& Darmofal研究了三角翼和轴对称流的涡流分解。 后者提供了一个更简单的模型来理解故障机制。 Murman和Roberts研究了一架悬停直升机上的流动。 Murman \& Bussing引入了点隐式方法来计算化学反应的流量。 Murman \& Shapiro使用自适应有限元算法进行欧拉计算。 Murman \& Usab使用多重网格法进行尤拉计算的收敛加速。 13 Murman是1987年AIAA应用空气动力学年度杰出论文和1992年美国宇航局兰利研究中心航空局年度杰出论文的合著者。 1989年，他获得了美国宇航局兰利颁发的团体成就奖，并于1989年获得了高等教育软件奖。 Murman还研究了三角洲机翼上黏性高超声速流动的解决方案，并获得了美国宇航局太空法案技术公司关于这项工作的论文。 1991年短奖。 Murman的几位学生已经成为教授，他们取得了优异的成绩，包括David Darmofal（德克萨斯A\&M）、Mark Drela（M. 我。 T.）、Scott Eberhardt（华盛顿大学）、Mike Giles（M. 我。 T.和牛津大学，U. K.）和Kenneth Powell（密歇根大学）；在本卷中查看他们的贡献。 其他人正在工业和政府机构追求他们的职业生涯--例如，托马斯·罗伯茨，他的贡献也出现在本卷中。 简而言之，Murman的贡献在影响的广度上是独一无二的，也是我们领域影响最深远的贡献之一。 Usab在1983年的论文中，根据Sells的工作，为提升翼型开发了一个远场边界条件，作为均匀流和可压缩点涡的叠加，其强度由翼型上的计算升力决定。 1986年，Thomas \& Salas使用了类似的想法。 有了这个边界条件，远场边界可以更接近翼型。 这些想法最近由Ryabanikii h Tsynkov扩展到亚音速自由流的三维Euler 方程，基于其差分电位法。

EARLL MURMAN的贡献 19 1.6 结束语 Murman对跨声速流动和CFD研究的原始贡献永远改变了这些领域。 Murman的工作引发了航空航天业的密集努力，用非线性场计算取代面板方法，包括对波阻的准确评估。 穆尔曼的成功也为全潜力和Euler 方程开发新方法提供了动力，特别是现代上风方案。 据说，回想起来，Murman-Cole方案是显而易见的：局部亚声速流动使用中心差异，局部超声速流动使用后向差异。 的确，Murman的想法很简单--事后！ 然而，多年来，许多人没有看到Murman首先看到的东西，14，一旦他让它变得如此明显，其他人就向许多方向提供了重要的扩展和应用。 一个人可能会幸运一次（或两次），但不会那么多次。 我们希望Earll在未来许多年继续为新的突破打开大门。 1.7 由Earll Murman监督的论文 1.7.1 S.B. 论文 以下学生在Earll Murman的监督下完成了他们的高级论文。 除非另有说明，否则学位是航空和航天学。 1. 丽莎·理查森，前扫翼与后扫翼的装载中心，1981年5月。 2. 胡安·R。 克鲁兹，帆翼机身的阻力减少，1981年12月。 3. 迈克尔·P。 Garrity，近地升降体设计中的计算机方法，1982年6月（机械工程）。 4. 约翰·吉布森，1982年12月，1982年12月，《降水》。 5. George Hoehn，1982年12月，来自降水的翅膀上的结冰。 6. William Giuffre，直升机叶片尖端旋旋，1983年5月。 7. 奥格登·琼斯，直升机叶片尖端旋向，1983年5月。 8. Sherri Hess，来自尖锐边缘三角洲翼的前沿电流的流动可视化，1985年12月。 9. 卡洛斯·费雷拉，飞机尾翼的前缘涡旋诱导装载，1987年5月。 10. James Weygandt，飞机尾鳍上的前缘涡旋诱导加载，1987年5月。 也许Murman看到了很远，因为他“站在巨人的肩膀上”。

20 HAFEZ \& CAUGHEY 11。 威廉·L。 约翰逊，涡流爆裂：铉叶前锋的预防，1987年5月。 12. 安德鲁·基思，涡流爆裂：铉叶前锋的预防，1987年5月。 13. 亚历山德拉·兰斯伯格，使用前缘襟翼控制涡流破裂，1987年5月。 14. 斯科特·米勒，使用前缘襟翼控制涡流破裂，1987年5月。 15. 泰德·A。 曼宁，涡流发生器在风洞扩散器中的应用，1988年5月。 16. Ali Zandieh，涡流发生器在风洞扩散器中的应用，1988年5月。 17. 罗伊·L。 酿酒师，电梯与 1988年5月，低速飞行前扫翼在各种角度的卡纳德事件的阻力性能。 18. 起亚·弗里曼，Lift vs. 1988年5月，低速飞行前扫翼在各种角度的卡纳德事件的阻力性能。 19. Melanie Glenn，翼尖涡总压力损失，1988年5月。 20. 詹姆斯·维奥莱特，翼尖涡总压力损失，1988年5月。 21. Marc Lie，使用Triple-翼型，可变几何机翼增加升力阻力比，1988年5月。 22. 蒂莫西·斯通，使用三重翼型，可变几何机翼来增加升力与阻力比，1988年5月。 23. J。 N。 Rees，用于数值分析和演示的X工具包小部件，1989年8月（计算机科学与工程）。 24. Jeanne Wei，使用数据流架构实现二维欧拉求解器的多网格技术，1990年6月，（机械工程）。 25. Shervin Limbert，F-22快速投球期间的非稳态空气动力学，1992年12月。 26. Robert Wickham，F-22快速投球期间的不稳定空气动力学，1992年12月。 1.7.2 S.M. 论文 以下学生在Earll Murman的指导下进行了硕士论文研究。 1. 保罗·M。 Stremel，应用于常规和旋转翼的非平面涡旋流场的计算方法，1982年2月。 2. 托马斯·理查德·阿瑟·巴辛，一维、非稳定的双模Scramjet操作模型，1982年6月。

EARLL MURMAN 21 3的贡献。 约翰·D。 Blascovich，分离流翼型分析方法的特点，1983年1月。 4. Martin Francis Peeters，使用数值优化的空气动力学设计的计算方法，1983年2月。 5. 格雷戈里·劳伦斯·拉森，通过空气动力学和优化迭代的交错加速空气动力学设计，1983年5月。 6. 托马斯·韦斯利·罗伯茨，用嵌入式涡旋环计算电位流动，并应用于直升机转子唤醒，1983年9月。 7. 埃里克·佩雷斯，前缘流锥流的计算，1984年8月。 8. 伯纳德·洛伊德，可压缩边界层方程的有限体积解，1985年5月。 9. 乔纳森·罗伯特·韦伯，使用MacCormack的隐式方法的一维Euler 方程的数值解，1985年6月。 10. 菲利普·保罗·博尚，欧拉对波浪墙问题的解决方案，1986年1月。 11. Robert Edward Malecki，1986年8月，使用Cyber 205的裁剪三角翼的欧拉方程计算。 12. Stephen Merrick Ruffin，高超音速黏性流动在Delta Wings上的解决方案，1987年6月。 13. Aga Myung Goodsell，三角洲翼上涡流的3D尤拉计算，1987年7月。 14. David Modiano，三维CFD结果的可视化，1989年5月。 15. 豪尔赫·佩雷斯-桑切斯，轴对称涡旋减速的数值模拟，1989年5月。 16. 托德·M。 贝克尔，NTF 三角翼上跨声速涡流的混合3D欧拉/Navier-Stokes计算，1989年8月。 17. Kuok-Ming Lee，高超声速流动在钝前缘三角翼上的数值模拟，1989年10月。 18. 亚历山德拉·兰斯伯格，在三角洲翼上对涡流黏性流动的适应，1990年。 19. David Darmofal，三维电流计算的分层可视化，1991年3月。 20. 萨宾·A。 Vermeersch，F-117A的自流特性调查，1993年5月。 1.7.3 博士论文 以下学生在Earll Murman的指导下进行了博士论文研究。

22 HAFEZ \& CAUGHEY 1. 威廉·J。 Usab，使用多网格方法对Euler 方程的嵌入式网格解决方案，1983年。 2. 托马斯·A。 Bussing，用于有限率化学的Navier-Stokes 方程的有限体积法，19853。 3. 托马斯·W。 罗伯茨，悬停直升机转子上流动的尤拉方程计算，1986年。 4. 肯尼斯·鲍威尔，锥形Euler 方程的螺旋解，1987年。 5. 理查德·A。 Shapiro，Euler 方程的自适应有限元解算法，1988年。 6. 伯纳德·洛伊德，半隐性Navier-Stokes求解器及其在钝三角翼分离流动研究中的应用，1989年。 7. 大卫·L。 Modiano，三角翼流中涡流分解的自适应网格尤拉方程计算，1993年2月。 8. 大卫·L。 Darmofal，轴对称涡流分解机制研究，1993年11月。 1.8 Earll Murman的出版物 以下是Earll Murman的出版物，直到1996年。 1.8.1 评审期刊中的论文1。 默尔曼，E。 M.和Hurlbut，R. L.，对M = 9的二维圆柱的唤醒的探索性研究，AIAA学生杂志，卷。 2，不。 1，1964年5月17日至22日。 2. Vas，我。 E.，穆尔曼，E. M.和Bogdonoff，S. M.，无支撑球体的觉醒研究，AIAA杂志，卷。 3，不。 1，1237-1244，1965年7月。 3. 默尔曼，E。 M.，层锥形尾波的实验研究，AIAA杂志，卷。 7，不。 9，1724-1730，1969年9月。 4. 默尔曼，E。 M.和Cole，J. D.，平面计算，稳定，跨声速流动，AIAA杂志，卷。 9，不。 1，114-121，1971年1月。 15 5. 克鲁普，J。 A.和Murman，E. M.，通过升降机翼和纤细体的跨声速流动的计算，AIAA杂志，卷。 10，不。 7，880-886，1972年7月。 6. 默尔曼，E。 M.，通过松弛方法计算的嵌入式激波分析，AIAA杂志，卷。 12，不。 5，626-632，1974年5月。 7. 默尔曼，E。 M.和Steinle，F. W.，评论“在1马赫时评估风隧道壁干扰的标准”，AIAA飞机杂志，卷。 11，不。 12，782，1974年12月。 15 本文被1987年11月9日第27卷第45期《当前内容》选为引用经典。

EARLL MURMAN的贡献 23 8. 哈菲兹，M。 M.，穆尔曼，E. M.和South，J. C，人工可压缩性超音速全势方程的数值解法，AIAA杂志，卷。 17，不。 8，838-844，1979年8月。 9. 阿巴伯内尔，S。 S.和Murman，E. M.，显式和隐式方案的二维双曲初始边界值问题的稳定性，计算物理学杂志，卷。 48，不。 2，160-167，1982年11月。 10. 里兹克，M。 H.和Murman，E. M.，跨声速制度中飞机模型的风洞壁干涉校正，AIAA飞机杂志，卷。 21，不。 1，54-61，1984年1月。 11. 米勒，R. H.，罗伯茨，T. W.和Murman，E. M.，悬停和前方飞行中唤醒建模和叶片气载测定的计算方法，计算机和数学应用，卷。 12A，不。 1，29-52，1986年。 12. 理发师，T。 J.，穆勒，G. L.，Ramsay，S. M.和Murman，E. M.，混合器中的三维无黏流动，第一部分：使用笛卡尔网格的混合器分析，AIAA推进和功率杂志，卷。 2，不。 3，275-281，1986年。 13. 理发师，T。 J.，穆勒，G. L.，Ramsay，S. M.和Murman，E. M.，混音器中的三维无黏流动，第二部分：涡轮风扇强制混音器分析，AIAA推进和动力杂志，卷。 2，不。 4，339-344，1986年。 14. 博尚。 P.和Murman，E. M.，Euler 方程的波浪墙解决方案，AIAA杂志，卷。 24，不。 1986年12月，2042-2045年。 15. 鲍威尔，K.，默尔曼，E. M.，Perez，E.和Baron，J.，锥形Euler 方程的螺旋解中的总压力损失，AIAA杂志，卷。 25，不。 3，360-368，1987年3月。 16. 洛伊德，B.和穆尔曼，E. M.，可压缩边界层方程的有限体积解，AIAA杂志，卷。 25，不。 4，525-526，1987年4月。 17. Bussing，T.和Murman，E. M.，2维高空气火焰在坡道和向后方台阶上的数值调查，AIAA推进和功率杂志，卷。 3，不。 5，448-454，1987年9月至10月。 18. 鲍威尔，K。 G.，默尔曼，E. M.，伍德，R. M.和Miller，D. S.，带有涡流襟翼的三角洲翼的实验和数值结果的比较，AIAA飞机杂志，卷。 25，不。 5，405-412，1988年5月。 19. Bussing，T.和Murman，E. M.，用于计算化学反应流的有限体积法，AIAA杂志，卷。 26，不。 9，1070-1078，1988年9月。 20. 麦克米林，S。 N.，托马斯，J. L.和Murman，E. M.，欧拉和Navier-Stokes·利赛德流在超音速三角洲之翼上，AIAA杂志

24 HAFEZ \& CAUGHEY飞机，卷。 26，不。 5，452-459，1989年5月。 21. 默尔曼，E。 M.和Powell，K. G.，自流中的轨迹积分，AIAA杂志，卷。 27，不。 7，982-984，1989年7月。 22. 默尔曼，E。 M.，贝克尔，T. M.和Darmofal，D.，前沿涡流的计算和可视化，工程计算系统，卷。 1，不。 2-4，341-348，1990年。 23. 莫迪亚诺，D。 L.和Murman，E. M.，涡流分解三角翼周围的流动的自适应计算，AIAA杂志，卷。 32，不。 7，1545-1547，1994年。 24. 梅里特，T。 R.，默尔曼，E. M.和Friedman，D. L.，通过顾问研讨会吸引新生，工程教育杂志，接受出版，1997年1月。 1.8.2 评审会议记录1。 默尔曼，E。 M.，彼得森，C. W.和Bogdonoff，S. M.，高超音速锥形唤醒的诊断研究，AGARD CP No. 1967年5月19日。 2. 默尔曼，E。 M.，通风跨声波风洞中壁效应的计算，AIAA论文72-1007，加利福尼亚州帕洛阿尔托，1972年9月。 3. 默尔曼，E。 M.和Cole，J. D.，跨声速下的隐性阻力，AIAA Paper 74-540，加利福尼亚州帕洛阿尔托，1974年6月。 4. 哈菲兹，M。 M.，里兹克，M. H.和Murman，E. M.，非稳态超音子小扰动方程的数值解，AGARD Conf。 程序。 226，论文18，葡萄牙里斯本，1977年4月。 5. 默尔曼，E。 M.，跨声速风洞壁干扰的校正方法，AIAA论文79-1553，弗吉尼亚州威廉斯堡，1979年7月。 6. 默瑟，J。 E.和Murman，E. M.，跨声速势流计算在飞机和风洞配置中的应用，AGARD Conf。 程序。 285，论文1，德国慕尼黑，1980年5月。 7. 朱，W。 H.和Murman，E. M.，跨声速激波 边界层相互作用的位移厚度影响的现象学模型，AGARD Conf。 程序。 291，论文215，科罗拉多州科罗拉多斯普林斯，1980年9月。 8. 里兹克，M。 H.，哈菲兹，M. M.，穆尔曼，E. M.和Lovell，D.，三维模型的跨声速风洞壁干涉校正，AIAA CP-82-0588，弗吉尼亚州威廉斯堡，1982年3月。 9. 默尔曼，E。 M.和Stremel，P. M.，用于势流计算的涡流唤醒捕获方法，AIAA论文82-0947，AIAA/ASME联合会议，St. 密苏里州路易斯，1982年6月7日至11日。 10. 公交车，T。 R。 A.和Murman，E. M.，双模式Scramjet操作的一维模型，AIAA论文83-0442，内华达州里诺，1983年1月。

EARLL MURMAN的贡献 2.5 11. 罗伯茨，T。 W.和Murman，E. M.，直升机涡流唤醒的计算方法，AIAA论文84-11554，科罗拉多州斯诺马斯，1984年6月。 12. 罗伯茨，T。 W.和Murman，E. M.，使用Euler 方程的悬停直升机转子的解决方案方法，AIAA论文85-0436，内华达州里诺，1985年1月。 13. 默尔曼，E。 M.，Powell，K.，Wood，T.和Miller，D.，前缘涡流的计算和实验数据的比较-偏航和涡旋瓣的影响，AIAA论文86-0439，内华达州里诺，1986年1月。 14. 默尔曼，E。 M.和Rizzi，A.，Euler 方程对带前缘螺旋的尖刃三角翼的应用，AGARD CP 4112，论文15，法国普罗旺斯地区艾克斯，1986年4月7日至10日。 15. 罗伯茨，T。 W.和Murman，E. M.，围绕悬停直升机转子流动的欧拉解决方案，AIAA论文86-1784 CP，加利福尼亚州圣地亚哥，1986年6月。 16. 默尔曼，E。 M.，Powell，K. G.，Goodsell，A. M.和Landahl，M.，具有大总压力损失的前沿涡流解决方案，AIAA论文87-0039，内华达州里诺，1987年1月。 17. 沙皮罗，R。 A.和Murman，E. M.，Euler 方程的笛卡尔网格有限元解决方案，AIAA文件87-0559，内华达州里诺，1987年1月。 18. 麦克米林，S。 N.，托马斯，J. L.和Murman，E. M.，欧拉和Navier-Stokes以超音速在三角洲翼上流动的利赛德解决方案，AIAA论文87-2270 CP，1987年8月。 16 19. 默尔曼，E。 M.，LaVin，A. R.和Ellis，S. C，使用基于工作站的软件加强流体力学教育，AIAA Paper 88-0001，内华达州里诺，1988年1月。 20. 沙皮罗，R。 A.和Murman，E. M.，Euler 方程的自适应有限元法，AIAA论文88-0034，内华达州里诺，1988年1月。 21. 拉芬，S。 M.和Murman，E. M.，高超音速黏性流动在三角洲翼上的解决方案，AIAA Paper 88-0126，内华达州里诺，1988年1月。 22. 鲍威尔，K。 G.和Murman，E. M.，细黏性涡流核心模型，AIAA论文88-0503，内华达州里诺，1988年1月。 23. Modiano，D.，Giles，M.和Murman，E. M.，三维CFD解决方案的可视化，AIAA论文89-0138，内华达州里诺，1989年1月。 24. 沙皮罗，R。 A.和Murman，E. M.，高阶和S-D有限元方法，AIAA论文89-0655，内华达州里诺，1989年1月。 本文被AIAA应用空气动力学技术委员会评选为年度杰出论文。

26 HAFEZ \& CAUGHEY 25. 兰茨伯格，A。 M.和Murman，E. M.，使用Id World实现二维尤拉求解器，一个平行编程环境，AIAA论文89-1941，1989年6月。 26. 洛伊德，B.，李，K. M.和Murman，E. M.，半隐性Navier-Stokes求解器（SINSS）钝三角翼周围分离流动的计算，AIAA论文90-0590，内华达州里诺，1990年1月。 27. Rizzi，A.，Murman，E. M.，Eliasson，P.和Lee，K. M.，钝三角翼上高超音速背位涡的计算，AGARD CP 494，涡流空气动力学，论文8，1990年10月。 28. McCune，J.，Drela，M.，Ellis，S. C.和Murman，E. M.，新的教学模块“Polhausen”、“Kelvin”和“Vortex Lattice”，为Todor Fluids Software开发，Proc。 ASEE会议，1992年6月。 29. 莫迪亚诺，D.和穆尔曼，E. M.，涡流破裂的三角翼周围流动的自适应计算，AIAA论文93-3400，加利福尼亚州蒙特雷，1993年8月。 30. Vermeersch, S., Modiano, D., Harizopoulos, I., Murman, E. M.和Peraire，J.，F-117A周围流场的实验和计算调查，AIAA论文93-3508，加利福尼亚州蒙特雷，1993年8月。 31. 达莫法尔，D。 L.和Murman，E. M.，关于轴对称涡旋分解的陷阱波性质，AIAA论文94-2314，科罗拉多州科罗拉多斯普林斯，1994年6月。 1.8.3 其他主要出版物1。 默尔曼，E。 M.和Krupp，J. A.，使用混合有限差分方案的跨声速势流方程的解，Proc。 第2国际。 流体力学数值方法会议，加利福尼亚州伯克利，物理学讲义，8:199-206，1970年9月。 2. Seebass，A。 R.，默尔曼，E. M.和Krupp，J. A.，在苛性附近不连续讯号行为的有限差分计算，Proc. 第三届声爆研究会议，美国宇航局SP-255，1971年。 3. 默尔曼，E。 M.，嵌入式激波的隐形跨声速流动的计算方法，流体力学中的数值方法，AGARD LS-48，1972年。 4. 默尔曼，E。 M.，用独立船头冲击计算跨声速流动的松弛方法，Proc. 第三届流体力学数值方法会议，物理讲义，19-11，201-205，1973年。 5. 默尔曼，E。 M.，超临界机翼截面回顾，计算数学，卷。 28，671-672，1974年4月。 6. 默尔曼，E。 M.，贝利，F. R.和Johnson，M. C，TSFOIL - 二维跨声速计算的计算机代码，包括风洞壁效应和波阻评估，在需要高级计算机的空气动力学分析中，美国宇航局SP 347，1974年3月。

EARLL MURMAN的贡献 27 7. 范德普拉茨，G。 N.，希克斯，R. M.和Murman，E. M.，数值优化技术在翼型设计中的应用，在需要高级计算机的空气动力学分析中，美国宇航局SP 347，1974年3月。 8. 默尔曼，E。 M.，跨声速小扰动方程的一些数值解回顾，研讨会Trarissonicum II，K。 Oswatitsch和D。 Rues，编辑，Springer-Verlag，415-422，1976年。 9. 默尔曼，E。 M.，涡轮机械中超音速电位流的计算，跨声速流动涡轮机械中的问题，T. C. 亚当森和M。 F。 Platzer，Eds.，Hemisphere，115-138，1977年。 10. 哈菲兹，M。 M.，Wellford，L. C.和Murman，E. M.，跨声速流动计算的有限元和有限差，流体中的有限元，卷。 3，R。 H。 Gallagher等人，Eds.，John Wiley and Sons，219-254，1978年。 11. 默尔曼，E。 M.和Bussing，T.，关于边界层和欧拉方程解的耦合，Proc。 空气动力学流动的数值和物理方面第二研讨会，T. Cebeci，Ed.，1983年1月。 12. Usab，W。 J。 Jr.和Murman，E. M.，使用多网格方法的Euler 方程嵌入式网格解决方案，流体中数值方法的最新进展，卷。 Ill，W。 G。 Habashi，Ed.，1985年。 13. 默尔曼，E。 M.和Abarbanel，S.，超级计算机对计算流体动力学未来十年的影响，计算流体动力学的进步和超级计算，Birkhauser-Boston, Inc.，1-10，1985年。 14. 默尔曼，E。 M.，Rizzi，A.和Powell，K.，计算流体动力学中涡流、进度和超级计算的Euler 方程的高分辨率解决方案，Birkhauser-Boston, Inc.，93-113，1985年。 15. Drela，M.和Murman，E. M.，欧拉CFD直升机涡流分析的前景，美国直升机协会会议记录，德克萨斯州阿灵顿，1987年2月。 16. 默尔曼，E。 M.和Powell，K. G.，三角翼前缘涡流中测量和计算的皮托压力的比较，涡流主导流动研究，M. Y。 Hussaini和M。 D。 Salas，Eds.，Springer-Verlag，270-284，1987年。 17. 洛伊德，B.，默尔曼，E. M.和Abarbanel，S. S.，纳维尔·斯托克斯方程的半隐式方案，Proc。 1987年9月，第7届GAMM流体力学数值方法会议，比利时Louvain-la-Neuve。 18. 鲍威尔，K。 G.和Murman，E. M.，前缘涡流的嵌入式网格程序，美国宇航局兰利跨声速研讨会论文集，美国宇航局CP-3020，231-259，1988年4月19日至21日。

28 HAFEZ \& CAUGHEY 19. 默尔曼，E。 M.，Modiano，D.，Haimes，R.和Giles，M.，并行处理器上多个CFD回路的性能，Proc。 第三届超级计算国际会议，卷。 II，372-379，1988年5月。 20. 默尔曼，E。 M.，Goodsell，A. M.和Malecki，R. E.，跨声速流，跨声速III研讨会，哥廷根，1988年5月。 21. Rizzi，A.和Murman，E. M.，涉及边缘流的流动计算，计算流体动力学和反应气体流动，IMA数学卷及其应用第12卷，Springer-Verlag，307-332，1988年。 22. 默尔曼，E。 M.，雅典娜项目-目标、理念、现状和经验，第二届北欧大学和高等教育计算机辅助学习会议，挪威特隆赫姆，1989年6月。 23. 默尔曼，E。 M.，Orcutt，R. L.和VanHaren，C，雅典娜项目的观点，科学和技术计算机辅助培训国际会议记录，西班牙巴塞罗那，1990年7月。 24. Landsberg，A.和Murman，E. M.，涡流建模中的数值扩散控制，航空航天工程计算非线性力学第6章，S. Atluri，Ed.，AIAA航空和航天学进展，1992年。 25. 莫迪亚诺，D。 L.和Murman，E. M.，时间精确计算的加速技术，Proc。 第11届AIAA CFD会议，佛罗里达州奥兰多，1993年7月。 26. 魏斯，S。 I.，Murman，E. M.和Roos，D.，空军和工业思维精益，美国航空航天，32-38，1996年5月。 1.8.4 内部备忘录和进度报告1。 希克斯，R.，默尔曼，E. M.和Vanderplaats，G.，通过数值优化评估翼型设计，美国宇航局TMX 3092，1974年7月。 2. 哈菲兹，M。 M.，Wellford，L. C，默克尔，C。 M.和Murman，E. M.，用有限元和有限差分法对跨声速流进行数值计算，美国宇航局CR 3070，1978年12月。 3. 哈菲兹，M。 M.和Murman，E. M.，瞬音速翼型有限元分析的最新发展，美国宇航局CP 2045，281-296，1978年3月。 4. 洛伊德，B.和穆尔曼，E. M.，可压缩边界层方程的有限体积解，美国宇航局CR 4013，1986年10月。 5. 李，K。 M.和Murman，E. M.，近耦合鳅翼的实验和初步数值结果的比较，麻省理工学院CFDL TR-88-8，1988年5月。

\clearpage
\chapter{最佳高超音速锥翼}
最佳高超音速锥形翅膀Susan A. Triantafillou，1 Donald W. 施温德曼1和朱利安·D。 Cole 1 2.1 导言 本文的目的是为高超音速范围设计最佳机翼。 为了简化问题，提出了各种假设，然而，这仍然允许我们获得具有实际价值的结果。 这个问题在经典空气动力学、无黏性完美气体的框架内被认为。 这限制了马赫数字如此之高，以至于真正的气体效应很重要，尽管就升力和波阻而言，它们并不发挥重要作用。 考虑锥形翅膀有助于减少变量的数量，这些变量是实用的形状。 此外，重点是提升翼下表面的压力分布，该提升翼提供了提升的主要组成部分。 进一步的高超音速小扰动理论用于描述流动。 该理论在适当的范围内非常准确，并透过提供相似性定律、与非定常流动的类比和减少因变数的数量来帮助。 有了所有这些假设，尽管仍然需要大量的计算，但寻找最优值的问题变得非常容易处理。 在以下章节中，给出了高超音速小扰动方程的公式。 给出了边界条件和冲击关系。 We 1 数学科学系，Rensselaer理工学院，1Yoy，纽约12180-3590 Frontiers of 计算流体动力学 - 1998 编辑：David A. Caughey和Mohamed M. 哈菲兹 ©1998 世界科学

30 TRIANTAFILLOU、SCHWENDEMAN和COLE希望将给定电梯的波阻最小化。 这样做的方式有些任意性。 定义了一个独立于机翼的平分攻角的优点。 概述了数值和优化程序，得出了最后一节中描述的结果。 平坦的三角翼和caret翼波骑手都被认为是比较案例。 有关此处显示的结果的更多详细信息，请参阅[4]。 2.2 高超音速小扰动理论；圆锥翼高超音速小扰动方程可以作为小流量偏转（5 -» 0）和高马赫数（Moo -> °°）的渐近极限过程扩展的第一个项获得。 这个扩展的想法是由范戴克[6]给出的。 另见[2]。 身体和冲击之间的区域是O(S)。 这个和弦是统一的。 极限过程有0，（x = x，y= y\textasciitilde{}；H=\textasciitilde{}\textasciicircum{}\&xedy x，y，z是物理坐标。 展开的形式为q(i,y,z;M 00,J), = i(l + 5 u\{x, y, z; H) + • • •) + Jv T (z, y, z\textbackslash\{\} H) + • • • • p(x,y,g\textbackslash\{\}M„,S) \_ 2 Poof 2 p\{x,y,z\textbackslash\{\}M00,5) = 5 2(P(x,y,z;E) + --- (2.1) = <r(x,y,z;E)--- Pooo where v T = v(x,y, z;H)j + w(x,y,z;E)k = 横平面上的流动扰动。 将表达式（2.1）代入Euler 方程产生的基本方程是连续性、横向动量和熵。 这些方程与二维非稳态Euler 方程完全相同，x对应于时间，因为th a:-momentum方程退出。 边界条件和冲击关系也对应如下图所示。 物体表面切流的边界条件 B H'£)=!- F H) =O (2 -2) 可以用 q-V5 = 0 表示，它变成 dF dF »(*, F\{x, z), z) = - + w(x, F(x, a), z)- ( 2,3)

最佳高超音速翼 31 图1 典型的锥形翼（a）在（x，y，z）坐标下，以平面形式和（b）在（Y，Z）平面上显示。 当机翼具有圆锥对称性时，方程和边界条件允许具有圆锥对称的解。 因此，身体形状可以表示为y = F\{x，z）= xW（Z），或Y = W（Z）（2.4），其中Y = y/x，Z = z/x是圆锥坐标。 切流（2.4）的边界条件变成圆锥形变量v(W(z), z) - W(z) = \{w\{W\{z), z) - z)W z。 （2.5）相应地，弓形冲击表面由y = g（x，z）= xG（Z）或Y s = G（Z）给出（2.6）

32 TRIANTAFILLOU, SCHWENDEMAN \& COLE 弓形冲击关系变成 2 G-ZGz 2 H s \textasciitilde{} 7 = 1 l + G| 7 + 1 G - ZGz ws = -v sGz P3=(G- ZG z)vs (2.7) 1 = 7\textasciitilde{}1 ■ 2 1 + G| cr3 7+1 7+1 (G-ZGz) 其中i> s = v(G(Z),Z) 等。 锥形投影的一般设置如图所示。 1. (Y, Z)中的运动方程的形式为,, rt da „. fdV dw\textbackslash\{\} (v-Y)W + \{w -Z)dZ+a\{oY + dz) =° (((.-v)£+ (-*>£)(P + F/7 O-T = 0 (2.8) 最后一个（2.8)表明熵在流线的圆锥投影上是恒定的：P + Hh dY v-Y,„, = const on- = (2.9) <TT dZ w - Z v; 另一对特征曲线可以用标准方式找到'dY\textbackslash\{\} \_ \{v-Y) 2 -\textasciicircum{}-,dZjch\&T- \{v-Y)(w-\textasciicircum{}-\textasciicircum{}Mf\textasciicircum{}i (2 - 10) 其中i fr i / v- y 对于Z中的光滑翼对称，Wz = 0在Z = 0和切线流动（2.5）的边界条件，v = Y = W（0）。 因此，Mc = 0，该方程是椭圆的。 另一方面，M c > 1对于足够大的Y，Z，因此系统是双曲的。 因此，总的来说，我们正在处理一个混合类型的问题。

最佳高超音子翼33 可以从冲击关系中得到一些有趣的结果。 局部研究可以回答锥形边缘附着电击的可能性的问题。 在边缘Z = A附近，冲击可以表示为直Y，= G\{Z）=W（A）+Q（A-Z）+ ---（2.11）在冲击后面的点和表面，正则解必须同时满足冲击关系（2.7）和切流边界条件（2.6）。 这导致d " «W(A» (*gg± \textasciitilde{} W\textasciicircum{}) - \textasciicircum{} (-(A) - -'<A» - 0 (2.12) 这是一个6的三次方程，参数为HA, W(A), W'(A)。 一个解决方案降低了熵，并被排除在外。 从另外两个中，我们假设较弱的冲击（熵增加较少）是发生的。 在这种情况下，局部流动在冲击的后方是双曲的。 当当地流变成椭圆时，冲击被分离。 发生这种冲击的特殊机翼是插入式机翼，由在山脊线相交的平面组成。 从（2.12）开始，可以绘制边缘的附着波和分离波的区域，如图所示。 2. 类似的考虑也适用于任意形状的锥形机翼。 我们主要对计算升降翼高压侧的流量感兴趣，并优化形状，以产生给定升降机的最小压力阻力。 下一节给出了升力和阻力系数的公式以及最佳值的适当度量。 2.3 升力和阻力系数；功德图：锥形翼升力和阻力系数是常规定义的，但考虑的力只是下表面的力 \textasciicircum{} = -4*1' \textasciicircum{} = T4I ( 2-13 > Poo 2 A "° 2 其中 L = 升力 = f f\{pi- Poo)dxdz, D = 阻力 = / / (pi - p x)5-dxdz (2.14) A = 翼平面面积，pi = 下表面的压力。

34个TRIANTAFILLOU、SCHWENDEMAN和COLE卡套翅膀，包含H = 2.0的（Y，Z）对（1,0）和（W（A），±A）。 附着的冲击区域包含恒定的0等高线。 在高超音速坐标中，这给出了JL CL o 2 4 / / Pidx,dz, C D = 5 3j I f P,\textasciicircum{}dx,dz (2.15) 其中 A = A/5 = 缩放平面区域。 一般来说，机翼将由有限数量的设计变量iV或设计变量向量D（d1，d2，-- ■，\textasciicircum{}N）来描述。 CD = 5 3CD\{B; H), C L = 5 2CL(D; H) (2.16) 我们通过使用拉格朗日乘法A，寻求将给定升力系数 CL的CD最小化的最佳值。 最小 K = J 3 C D (D; H) - XS 2CL(D; H) K 静止的条件是 \textasciicircum{}-=352CD-2\textbackslash\{\}5CL = 0 oS (2.17) (2.18) 图 2 分离冲击区域（实体曲线下方）和附着冲击

最佳高超音功率翼35 dK \_ 3dCD，c2dCL消除A表明d（C D ddj \textbackslash\{\}c 3L/2 = 0，j = l，2，---，N。 因此，通过选择dj、di、••、d\textasciicircum{}来找到最优值，以便将si fit = figure of merit = -\textasciicircum{} (2.20)最小化。 当使用锥形对称性时，功德形象的公式变成A /0A Pi\{W\{Z) - ZW '\{Z))dZ fM = \textbackslash\{\}l2 / A - x3/2 ( 2 -21 ) 插管翼波骑手或设计上的插管翼是那些在下侧流动均匀和双曲的。 弓形减震器是直的，并连接到两个前缘。 因此，这种翅膀的形状是Y = W\{Z) = l + -*j\textasciicircum{}-\textbackslash\{\}Z\textbackslash\{\}, \textbackslash\{\}Z\textbackslash\{\}<k (2.22) 其中A > v\textasciicircum{}\textasciicircum{}V* -1) 这种翅膀的优点数字很容易计算，可以写成IM =■ -\textasciicircum{}= = i j - (2.23) \textasciicircum{}? +\textbackslash\{\}7FP 对于给定的H，设计上的通用机翼的优点与A无关。 对于扁平机翼（W = 1），类似的计算给出了图中附着和分离的冲击区域。 3和图中的一般流动结构。 4. 2.4 数值和优化方法 本节简要概述了用于获得以下结果的数值方法和优化程序

36 TRIANTAFILLOU, SCHWENDEMAN \& COLE 图3 分离冲击区域（实心曲线上）和扁平机翼的附着冲击区域B = Y - 1，具有各种A和h值。 附着的冲击区域包含恒定的0等高线。部分。 有关更多详细信息，请参阅[4]。 锥形方程（2.8）以u< + fy + gz = c（2.24）的形式进行，其中u = \{a，a（v - Y），a\{w - Z），oE -（P + Hh））f =（cr（v - Y），a\{v - Y）2 + P + H/f，a\{v - Y），aE（v - Y））g =（a（w - Z），a（v - Y）（w - Y）（w - g），a\{w - Z）2 + P + H/-y，aE\{w - Z）w = \{a，a\{v - Y），a（w - Z），aE -（p + Hh）））c =（-2<7，-3<T（V - Y，-3<r（w - Z），-a（2E +（v - Y） 在这种守恒形式中，方程类似于具有广义源项的非稳定二维Euler 方程。 数值方法使用二阶精确的Godunov方案和Roe的近似Riemann 求解器求解该系统的映射版本。 由于伪时间的引入，系统总是双曲的。 冲击

最佳高超音索克翼流线型。-... 图4平翼以下流动的溶液特征，B = Y - 1121 < A。 Y可以用守恒形式的数值方法捕获。 物理平面中的一个域在翼尖外和弓形冲击下方延伸一段距离，用于计算目的，映射到一个矩形。 时间精确的步进用于将解决方案及时推进到稳定状态。 使用预测-校正方案。 必要时引入上卷。 对于初始条件，使用未受干扰的流动或之前的解决方案。 切流的边界条件是通过使用幽灵点来适应的。 当解接近稳态时，它由牛顿方法提出。 方程（2.24）用ut = 0求解。 一旦找到数值近似值，就从（2.21）计算出优点fM。 沿着半跨度对不同数量的细胞进行计算，例如，20、30、40和超过翼尖的足够数量的细胞。 在分离的情况下，在翼尖之外的域部分使用了诺伊曼条件。 事实证明，与同时计算上区域和下区域的解决方案相比，这令人满意。 为了最佳化，在Idwer表面坡度和翼的平均偏转为6的限制下修改了初始翼形。 优化程序通过迭代初始设计来减少优点。 考虑的每个设计都需要进行流量计算。 设计的改进取决于一组设计变量的优点图的导数。梯度计算遵循想法--....熵区域双曲和I I椭圆区域之间的边界I -A 0 A z

38 TRIANTAFILLOU, SCHWENDEMAN \& COLE of Baysal and Eleshaky [1]。 域内所有细胞的稳定流动的离散解V =（Vi，V2，■ ■ ■，VQ）取决于一组设计变量Np =（d\textbackslash\{\}，\$2，• • ■，dND）-优点值取决于（V，D），因此'3F\textbackslash\{\}dV dfu（2.25）在没有约束的情况下，设计变化的方向是最陡峭的下降，向量e 0 = I|VD/M|（2.26）用于移动到改进的设计。 如果违反了约束，则使用可行方向[5]方法的适应。 最陡峭的下降方向被改变，以满足约束。 在此处完成的计算中，翅膀的锥形投影由（NJJ + l）-节点立方样条定义。 W\{Z) = W,(Z) 0 < Z < A (2.27) 节点(Yj,Zj) j = 0,---,ND- Np = 3的翅膀示例如图所示。 5. 如上所述，迭代了设计，以考虑到限制，以产生下一节中讨论的最佳值。 图5 Exa，[Nd = 3的ple wi ng design。

最佳超音带翼 39 H =.9199，A = 3.205，N w = 40。 图（a）显示了恒定的P轮廓，图（b）显示了P（虚线）以及前缘附近恒定压力区域的翅膀和P（实心）的确切值。 2.5 平面和毛毛翼的计算结果；最佳翼 平面三角翼以下的流量计算结果，高超音速相似性参数H =.9991和A = 1.447如图所示。 6. 这对应于Hillier [3]计算的三角翼在15°攻角下扫荡50°的情况，M\textasciicircum{} = 4。 对于扁平机翼，图中高超音速相似性参数H的各种值的扫荡角度参数A绘制了优点图/M。 7. 当下部表面的冲击附着时，JM几乎是恒定的，类似于线性化理论的结果。 随着A的减少，冲击分解，/M减少到最小，然后增加较小的A。 因此，有一个最佳的扫荡。 还对插花机翼进行了一些计算。 图。 8a显示了通过最佳卡啫翼的流量，该机翼与设计上的卡巷不同（H = 2.0，YTIP = W（A）= 1.75）。 由于侧面压力较低，f M的值为0.4809，比设计的侧翼高0.6\%。 图6 平坦三角翼以下流量的结果

40 TRIANTAFILLOU, SCHWENDEMAN \& COLE 图 7 具有各种半长、A和各种H值的扁平翅膀的优点图。 这些计算是使用牛顿方法完成的，用于a上翼下方的域，网格为N m = 40。

最佳高超音速机翼41 图8（a）H = 2.0，A = 1.9，W（h）= 1.75的切形机翼的流量结果；（b）H = 2.0，A = 2.7的平滑优化机翼；（c）H = 2.0，A = 2.4的平滑优化机翼。 上图显示压力轮廓，下图显示下表面的压力。

42 TRIANTAFILLOU, SCHWENDEMAN \& COLE 按照优化程序，进行了一系列计算，以优化圆锥形机翼，并具有平滑坡度。 优化中使用的约束是||>0对于下表面的所有点。 这成为圆锥形坡度W(Z)-ZW'(Z) > 0（2.28）的条件。此外，还应用了平均偏转约束，使物理坐标下表面的攻击角等于8。 在圆锥形变量中，这读作i /-A ZW')dZ = 1 (2.29) kl> 从H= 1.0、1.5和2.0的扁平翼开始的平滑最优的摘要，图中显示了几个A。 9. 设计空间需要几千步，才能将最初平坦的机翼变形为最佳机翼。 对于H = 2.0，A = 2.7，优化的机翼的优点比原始机翼好1.8\%。 初始和最终机翼都有附着的减震器。 压力P接近设计上的柱塞翼P = 2.14，H = 2.0（见图。 8b）。 另一个案例从H = 2.0开始，A = 2.4，这是一个最初最优的平翼。 最终结果如图所示。 8c。 带有分离冲击的机翼变形了，从而产生了附着的冲击。 2.6 承认和免责声明 由美国空军空军材料司令部空军科学研究办公室赞助，赠款编号为F49620-97-1-0141。 美国政府有权出于政府目的复制和分发重印本，尽管上面有任何版权说明。 此处包含的观点和结论是作者的观点和结论，不应被解释为必然代表空军科学研究办公室或美国政府的官方政策或认可，无论是明示的还是暗示的。 参考资料1。 Baysal，O.和Eleshaky，M.E.，可压缩尤拉方程的空气动力学灵敏度分析方法，J. 流体工程。 113，1991年，第681-688页。 2. Hayes，W.和Probstein，R.，高超声速流动理论，学术出版社，1959年。

最佳高超音波翼 43 3. Hillier，R.，使用二阶精确的Godunov方法计算流过锥形超音速翼，AGARD CPP 428，1987年。 4. Triantafillou，Susan A.，高超音速翅膀的优化，博士。 论文，Rensselaer理工学院，纽约特洛伊，1997年6月。 5. 范德普拉茨，G。 N.和Moses，F.，通过可行方向、计算机和结构方法进行结构优化，3，1973年，pp. 739-755。 6. 范戴克，M。 D.，对高超音速小扰动理论的研究，NACA Tech。 代表。 1194，1954年，第885-905页。 图9 平滑最优和相关扁平翼的总结。

\clearpage
\chapter{理论、应用和教育流体动力学的几何}
理论、应用和教育流体力学几何学 Helmut Sobieczky1 3.1 导言 受本书主题及其出版事件的激励，在本文中，试图说明我们在可压缩流体力学领域积累的一些知识如何保持活力，并对航空航天工程现代工具的开发有用。 在计算流体动力学（CFD）解决复杂的空气动力学问题，而对潜在现象的详细了解似乎变得没有必要的时候，这可能需要。 跨声速空气动力学中的数值建模对Ear 11 M的工作产生了决定性的影响。 穆尔曼对超音速的关键现象进行了建模，在数学上通过从椭圆到双曲模型方程的类型变化来描述，在数值差分方案的简单变化中[5]。 在最近的一次活动中，Murman指导了一个项目，通过在计算机上解决一系列流体力学和空气动力学问题，帮助新一代学生了解流体力学和气体动力学。 我们试图从此类活动中学习，利用过去的工作成果，通过信息、通信和软件技术工具在未来保持熟悉。 这里说明了一些空气动力学设计工具的背景，这也是我们知识库的一部分，可以实施1 DLR德国航空航天研究机构，Bunsenstr。 10，D-37073 哥廷根。 Frontiers of 计算流体动力学 1998 编辑：David A. Caughey和Mohamed M. 哈菲兹 ©1998 世界科学

46 SOBIECZKY在机身和涡轮机械设计的现代高效优化工具中。 流动现象的细节导致了数学函数的正确描述。 由于需要达到逼真的三维配置，理想的3D流动元素的结构应与其边界条件的几何描述保持一致。 因此，几何预处理器工具可能会将流体力学知识库转移到支持跨学科工作的CFD分析、实验模型制作和组件数据定义的CAD系统。 案例研究说明了我们目前朝着这个方向开展的活动。 它可以用于航空航天研究、应用和教育的实际用途。 3.2 跨声速知识库 在过去的50年里，直到30年前，跨声速的先驱者（Guderley、Oswatitsch、Frankl、Tricomi、Germain......）开发了流动运动的基本方程模型，以便有可能找到孤立的或整个特定解决方案系统。 这些可以解释在接近跨声速飞行马赫数的更高速度制度中，这些影响是奇怪的，有时是危险的。 三十年前，势能理论被接受为建模翼型流动基本空气动力学成分的好工具。 此时，数字计算机的可用性使第一次（相对）“大规模”数值模拟成为可能。 使用混合亚音速/超音速域模拟跨声速流动的问题，通过Murman和Cole [5]的方法解决了，一旦超声速流动似乎嵌入亚声速流动模型中，将离散化的方案从中心差异更改为正向差异。 从那里开始，跨声速流动的分析是CFD开发的主要目标。 上述基本微分方程的解还包括高亚音速马赫数的2D 翼型流量，该流量没有冲击，因此承诺了更高的升力/拖力比。 花了一段时间，在实验和分析中，改进的技术才能证实这种无冲击流动的理论概念，因为基本的数学解决方案在数学上是孤立的，需要非常准确的分析和精确的风洞技术，因此，在一段时间内，它们对塑造真正的机翼没有实际价值。 但是，一旦被证明对应用空气动力学有用，系统设计方法就需要创建跨声速-“超临界”-翼型，并在它们的帮助下，在跨声速飞行系统中获得更好的机翼，具有高空气动力学效率。 三十年来，作者有幸参与了跨声速流动设计方法的开发，见证了更快的计算机和分析可压缩流动的精炼软件的出现。

流体动力学几何学 47 使用上述先驱模型的教育使作者使用一个简单的模拟静电设置来解决映射工作平面（hodograph平面）中的潜在边界值问题。 例如，这种方法允许相对容易理解Garabedian [1]在数学上优雅但困难的方法，即在复杂特征（4D）工作空间中设计无冲击翼型，并用2D静电网络[6]中的边界条件识别解的2D实部分。 用计算网格取代后者，用数值放松来代替后者，这在七十年代当然是及时的，但这种早期“玩”超音速的真正好处是它的教育价值，让我们对正在发生的流动现象有“感觉”，并更深入地理解其他人的相关工作。 此外，在模拟设置中解决映射边界值问题的技术在当时基于潜在理论操作的许多数值分析方法中得到了实施，并以这种方式将这些方法转化为无冲击2D机翼和3D机翼的设计工具[7]，[8]。 一旦流场内的静压低于临界值，这种方法就需要暂时改变理想气体特性，这是一个简单的软件修改。 因此，这种半逆技术被称为“Fictitious Gas”（FG）方法，在过去几年中也在现代欧拉[4]和Navier-Stokes[9]分析代码中实施。 随着这些工具的可用以及飞机机翼和涡轮机械叶片重新设计的许多案例研究，很明显，在给定的初始配置中观察到的减少冲击损耗，导致系统的形状修改，这可能允许直接的数学建模仅应用于表面几何的局部域。 由于几何预处理是参与航空航天产品设计过程的各个学科中任何精细优化的起点，因此，关于提高空气动力学效率的表面变化质量和数量的知识可能会大大减少优化过程中需要变化的参数数量。 3.3 几何发生器作为一项对CFD 边界条件和电网的经济生产越来越重要的活动，15年前开始开发具有特殊输入控制的灵活几何发生器代码，用于超音速设计改进，同时已发展为上述建议的预处理器软件，补充和使用设计知识库不仅用于跨声速设计和分析，还支持超音速应用和低速空气动力学和流体力学的形状定义。 在最近的一份汇编中，作者试图说明这个链接

48 SOBIECZKY在气体动力学和几何学之间更详细[10]，[11]。 CFD的最新发展显示了使用块状结构（六面体）和非结构化（terahedral）网格的进展。 对完整飞机等复杂配置的CFD分析在非结构化网格[2]中似乎更成功，使周围大部分空间离散化，并将结构化网格的使用限制在靠近表面的黏性流动层。 这种混合网格目前可能已经由市售软件[3]生成，我们的几何生成器作为预处理器：从目录中选择的分析功能能够对大多数复杂的细节进行建模。 对给定曲面指标的显式和非迭代评估允许非常快速地产生密集的曲面网格，这些网格用于确定所有组件和远场边界曲面之间空间三角测量的三维边界。 有一些重要的应用激励了引入一个类似时间的单一“超级参数”，以代表3D配置建模中的第四维度：配置可以被视为随着时间的推移而改变其形状，例如在自然界中发生的周期性运动中与动物飞行，或在具有流体-结构相互作用的空气弹性和转子空气动力学中。 或者形状在进化优化过程定义的一系列受控变化中以非周期模式变化。 这与实验模型或真实配置组件的机械适应过程平行。 在我们的软件中，定义形状细节的选定输入参数在间隔内给出，超参数控制其变化。 因此，这种移动形状的合适数据定义对所有非稳态过程都非常有吸引力，用于优化适当控制的参数数量，以及对控制执行器的微型计算机进行编程，以使机翼几何形状适应不同的操作条件。 飞机设计中的许多创新项目需要像这种几何生成器这样的工具。 在下一章中，报告了新定义的测试配置的工作状态，作为使用具有高度灵活性的几何建模的示例，以观察跨声速知识库的许多细节。 建议将其作为各种概念、计算和实验研究的测试台，旨在完善创新飞机设计的工具。 3.4 示例：通用运输飞机配置 流体动力学和空气动力学中的测试案例，以验证理论和实验一直受到欢迎，其中一些已经使用了多年。 他们作为各种研究人员和开发人员的共同基础的成功曾经取决于两者之间的平衡

流体动力学几何学图1通用高翼运输飞机的数学简单性是创建边界条件所需的，几何复杂性足够高，对应用具有吸引力。 经过多年将机翼作为测试案例处理，我们觉得是时候提供更复杂的配置了，包括所有组件的整架飞机。 利用现代软件和通信技术提供的可能性，这个目标似乎很合理。 我们设想通过互联网提供的工具来创建案例研究并与感兴趣的合作伙伴沟通数据。 基于几何生成器的合适软件工具正在开发中[13]，以下示例是完善它们的练习。 受当前一些设计活动以及飞机行业和我们研究小组之间的相关合作的激励，几何生成器被用于对通用运输机进行建模（图。 1）。 由于这些工具更早地应用于中翼或低翼安装配置的组件连接，对翼根区域的形状细节进行建模的高翼几何建模是一项雄心勃勃的任务，特别是如果在有大车轮齿轮箱的情况下，在跨声速流动中对组合的机翼-车身形状进行空气动力学优化。 我们的跨声速知识库建议应用面积规则来实现截面面积的平滑轴向分布。 机身横截面在轴向上变化，以促进翼根前缘附近的圆角形状，并提供在机身屋顶上延伸的平坦上翼表面（图。 2）。 用高升力系统扩展基本的机翼体几何数据库需要提供板条和襟翼运动的运动几何，后者沿着机身的平坦部分移动，以提高其空气动力学效率。 一个尾翼和四个螺旋桨发动机完成了这里显示的飞机--几乎，因为仍然需要添加襟翼轨道整流罩。 计划对相对较小的哥廷根跨声速风洞（TWG）进行实验，仅使用这架通用飞机的机身和清洁机翼的数据。 使用IGES格式，几何预处理器为CATIA CAD系统提供了表面数据，这允许定义

50 SOBIECZKY 图2 结构的机翼体连接细节，并用表面补丁替换输入数据，用于模型组件的数控铣削。 很明显，只有为形状细节所需的准确性提供非常密集的支撑网格才能确保可接受的表面质量。 然而，由于用于定义曲面的分析显式函数，无需任何迭代或插值，在用户友好的交互式软件中实现，在工作站上几秒钟内就会产生密集的数据。 在第一个实验中，将仅使用不完整模型进行局部流量质量分析的概念变为现实：应用关于提升机翼理论的经典空气动力学知识，在大长宽比机翼上适当形状的端板，其尖端被夹住，应该允许在剩余机翼上保持机翼的循环，其跨度缩短。 这将确保在对称平面上保持不变的流动细节，即在机翼体交界处的区域。 在本例中，跨度减少到原始机翼的60\%，长宽比为9，“循环控制分离器叶片”（CCSB）像小翼一样安装，图。 3. 参考中描述了他们使用2D扫翼理论和3D CFD模拟的设计。 [14]。 减少80厘米的跨度使模型（DLR-F9）适合lxl m超音速风洞，用于调查机翼根部区域的流动细节。 这个风洞有自适应的上壁和下壁，但封闭和固定的侧壁，CCSB不仅应该补偿升力产生跨度的减少，还应该补偿侧壁干扰。 图。 4显示了完整的安排，包括模型支持：为了用这个测试案例创建一个未来的“虚拟风洞数据库”，还对TWG中给出的刺和剑的附加组件进行了建模，其组件由创建为“翅膀”或“身体”类型的形状组成，这是迄今为止使用几何生成器定义的仅有的两种表面拓扑。 风洞侧壁和上方和下方的远场边界完成了模拟自适应风洞的流动空间，

流体动力学的几何学51图3配置DLR-F9之后，来自壁适应的数据也将建模为流动边界。 在尤拉模式下使用非结构化网格代码DLR-t [1]进行了几次运行，以估计CCSB所需的调整范围。 图。 5显示了这个数值模拟的图形后处理。 改进CFD工具，用于应用混合网格，以模拟N/S版本代码的流程是下一步。 实验很容易允许在隧道运行时改变攻角，但到目前为止，CCSB的适应需要运行中断。 因此，在实验期间开发并使用了一种在线数据处理技术[12]，以最佳攻击角选择流动条件，模拟整个机翼上的跨度循环分布，从靠近CCSB的夹翼尖端的根部区域。 四个带有静压测量的跨度站被用来控制跨度负载分布（图。 6），选定的案例证实了相对不受机翼夹持干扰的压力分布，至少占被夹翼跨度的75\%或近全翼跨度的近一半。 这项实验表明，带有剪翼的模型非常适合在翼根区域进行详细调查。 我们现在就可以开始了

52 SOBIECZKY 图4 DLR-F9包括跨声速风洞哥廷根的模型支持，表面参数变化。 未来的测试将致力于使用替代连接模块获取数据。 后来，由于模型机身内空间很大，可以进行必要的伺服设备，因此设计具有弹性或气动表面元件的自适应模块似乎是可能的。 该模型的这种目的可能有助于发展表面适应控制系统的知识，通过基于选定的表面压力测量的预编程来优化流动质量。 前面提到的通过FG方法的系统表面重新设计对无冲击流动进行建模的理论工作已经提出了这种方法[8]。 数值模拟将允许研究比实验所能承受的更多变化，参数改变将为优化策略提供输入条件，建模机翼变形为流体-结构相互作用创建界面数据。 这样，这个的4D扩展和

流体动力学的几何学53图5 CFD模拟DLR-r代码：中心平面中的网格和等边形颜色等边，由我们的几何工具创建的类似测试配置可能对空气动力学技术以外的开发有用。 3.5 结论“几何”一词用航空航天工程的语言描述了一门科学学科以及特定配置的形状。 因此，本贡献的标题模棱两可，但这是一个值得欢迎的效果：试图将学科几何放在设计空气动力学知识库的重要元素和实际应用（包括向年轻一代传授知识）之间进行调解的位置，是这里的一个信息。 尝试具体也是通过一个例子来适当地实现的。 这里提出了一个通用的飞机配置，作为数值和实验空气动力学中各种技术的测试案例。 信息和通信技术的新方法开辟了令人兴奋的全球合作工作选择，将使我们能够以真正的多学科方法使用大型数据集实现分析和设计的雄心勃勃的目标。

54 SOBIECZKY 图6 四个翼截面的测量静压（Re = 2.3 Mill.） 参考资料。 鲍尔，F.，加拉贝迪安，P. R.k. 科恩，D。 G.，超临界机翼部分。 经济学和数学系统讲义，卷。 66，Springer Verlag，1972年。 Friedrich，0.，Hempel，D.，Meister，A. \&c Sonar，Th.，使用DLR-t-代码对非定常流动场进行自适应计算。 第77届AGARD FDP会议Conf。 程序。 CP 578，1995年。 ICEMTM，CFD/CAE手册集，版本3.1.2.1，参考。 控制数据的手册，。 李，P。 \& Sobieczky，H.，用Euler 方程计算虚构气体流。 Acta Mechanica（Suppl.）4，1994年，pp. 251-257。 默尔曼，E。 M。 \& 科尔，J. D.，平面稳定跨声速流动的计算，AIAA J. 卷。 1971年9月，第114-121页。 Sobieczky，H.，跨声速翼型设计中的相关分析、模拟和数值方法。 AIAA-79-1556，1980年。 Sobieczky，H.，Fung，K-Y.，Seebass A. R。 \& Yu，N。 J.，设计无冲击超音速配置的新方法。 AIAA J. 卷。 17，不。 7，1979年，pp. 722-729。 Sobieczky，H。 \& Seebass，A。 R.，超临界翼型和机翼设计。 安。 牧师。 流体机械。 1984年16月，pp. 337-63。 索比茨基，H.，盖斯勒，W. \& Hannemann，M.，不稳定的数值工具

流体动力学几何学55 黏性流动控制。 程序。 第15届Int。 关于Num的Conf. 流体动力学中的冰毒。 物理学讲义，ed。 K。 W。 Morton，Springer Verlag，1997年10。 Sobieczky，H.，高速流动建模的气体动力学知识库，在：“高速航空运输的新设计概念”，CISM系列讲座，卷。 366，Springer Verlag，维也纳，1997年，第105-121页11。 Sobieczky，H.，空气动力学设计的几何生成器，在：“高速航空运输的新设计概念”，CISM系列讲座，卷。 366，Springer Verlag，维也纳，1997年，第137-158页12。 Trapp，J.，Pagendarm，H.-G. \& Zores，R.，控制DLR-F9运输飞机模型风洞实验的高级数据处理。 AIAA-97-2221，13。 Trapp，J.，Zores，R.，Gerhold，Th。 \& Sobieczky，H.，Geometrische Werkzeuge zur Auslegung Adaptiver Aerodynamischer Komponenten。 程序。 DGLR JT 1996 14. 佐雷斯，R。 \& Sobieczky，H.，在小型风洞中使用流量控制设备。 AIAA-97-1920，1997年

\clearpage
\chapter{轴对称喷管流动的计算}
轴对称喷管流动 L的计算。 Pamela Cook，1 Elsa Newman，2 Scott Rimbey 1 \& Gilberto Schleiniger 1 4.1 导言 在本文中，我们计算了来自喷嘴的轴对称、非黏性、可压缩、稳定流动。 Legendre，或接触，潜在配方用于hodograph平面。 我们强调了轴对称情况，因为hodograph方程是非线性的，而对于更常研究的二维流动来说，它是线性的。 结果和公式与Alder[1]的早期工作进行了比较和对比。 最终目标是获得洞察力，并为处理表现出跨声速行为的轴对称流提供一个可行的框架。 4.2 问题公式化 管理非黏性、可压缩、稳定、等熵、非旋转流的方程是非旋转性：u r -t> x = 0，（4.1）1 特拉华大学数学科学系，纽瓦克，DE 19716 2 Marymount大学数学系，阿灵顿，VA 22207；作者的工作得到了国家科学基金会在DMS赠款\#9796052 计算流体动力学-1998年前沿下的支持，作者：David A. Caughey U Mohamed M. 哈菲兹 ©1998 世界科学

58 COOK, NEWMAN, RIMBEY \& SCHLEINIGER 图1 喷气流连续性：(a - u )u x - uv(u r + v x) + (a - v )v r + X- v = 0» (4-2) r和伯努利关系：u2 + v 2 + 1+ 1 «-。 7-1 2（7 -l（4.3）在这里，u和v分别是速度的a：和r分量，a是局部音速，a*是声波声速，7是比热比值的比率（空气为7 = 1.4）。 对于轴对称流动，x = 1和d（z，r）是圆柱坐标；对于二维流动，x = 0和（x，r）表示矩形坐标。 如图1所示，这项工作的重点是喷射流。 气体占据了一个由喷嘴表面OB、自由流线BC、均匀下游状态CD、x轴DO和上游无限O所界定的区域。 显示了典型的流线型。 喷嘴壁OB与水平形成一个角度S，由r = H - tan<5(x - xo)规定，其中H是喷射开口的半径，选择x 0，以便H + tan 5 x 0 = 0。 边界条件在r = 0上为v = 0，在OB上为v/u = - tanJ，q = y/u 2 + v 2 -> 0为x -> -00。 (4.4) (4.5) (4.6)

轴对称喷管流动 59在自由流线BC q = v u 2 +v 2=q0和- = -. ax u上，假设go的值最多为a*。 在声波射流的这种限制性情况下，Ovsiannikov [9]和Guderley [5]分别表明，二维流动在喷嘴出口下游的有限距离上达到均匀的声波状态。 Cole [3]在小扰动理论（6 -¥ 0）的背景下对轴向流动显示了相同的结果。 请注意，对于亚音速射流，均匀状态仅作为x -¥ oo出现。 如果我们在达到这种均匀状态时用h表示喷射的半径，那么可以定义收缩系数C c = h/H。 这将证明对以后的比较很有用。 上述公式的一个主要缺点是自由流线在物理平面上的位置未知。 然而，在hodograph或速度平面中克服了这一困难。 在这个框架中，我们有x(u,v)和r(u,v)。 将方程（4.1）-（4.3）直接变换到hodograph平面给出xv-ru= 0，（4.7）（a2 - u 2）rv + uv\{x v + r u）+（a 2 - v 2）xu + X - J = 0，（4.8），其中J = x urv - x vru是雅各布。 请注意，对于二维流动（x = 0），系统（4.7）-（4.8）是线性的，但对于轴对称流动（x = 1），它是非线性的。 图1中描述的物理平面中的域映射到图2所示的hodograph平面中的域。 所有流线都从原点（驻点）发出，并流入CD（统一下游状态）。 霍多图平面中BC的边界条件确定如下。 由于u 2 + v 2 = q\% =常数，udu = -vdv。 此外，dr/dx = v/u需要v(xudu + x vdv) = u(r udu + r vdv) 结合这些产量v2xu - uv\{r u +x v) + u2rv = 0 on q = qo- (4.9) 流量可以有几种方式计算：（1）用上述边界条件求解系统（4.7）-（4.8）；（2）使用流函数公式；或（3）使用Legendre或contact，

60 库克、纽曼、RIMBEY和施莱尼格的潜在配方。 流函数设置由Rimbey [11]在小扰动理论的上下文中使用，Alder [1]用于假定8不是小的更一般情况。 在前一种情况下，方程是显式的，在变数变换后，是准线性的。 二维流动中的全流函数方程是众所周知的Chaplygin方程，它是显式和线性的，但全轴对称方程是非线性的，而且，其系数只能隐式确定。 然而，接触电位公式是明确的，也是我们选择追求的方法。 我们现在将开发合适的方程。 方程（4.7）意味着存在一个势<j>，这样x = (/>„, r = 4> v。 这个Legendre势4>与物理势\$相关联，其中u = \$ x和v = \$ r，关系<f> = xu + rv - \$。 就<j>而言，连续性方程（4.8）成为一个广义的Monge-Ampere方程：（a 2 - u 2）<j）vv + 2uv4> uv + (a 2 - v 2)4>uu + X-r-\{<t>uu<t>w - 4>l v) = 0 (4.10) 图2 Hodograph Plane

轴对称喷管流动 61 这个方程是完全非线性的，因此人们可能会质疑使用接触电位公式获得了什么优势。 然而，它具有Courant和Hilbert[4]中指出的几个特殊属性。 特别是，它基本上是一个准线性方程，其类型由线性项决定：L[4>] = (a 2 - u 2)<t>vv + 2uv4> uv + (a 2 - v 2)4>uu。 （4.11）因此，该方程分别为椭圆或双曲，如马赫数 M = q/a < 1或M > 1。 此外，双曲区域中的柯西问题和椭圆区域中的狄利克雷问题的独特性得到了保证[8] [10]。 尽管有这些优势，但至少从计算的角度来看，该方程似乎很少受到关注[6][2]。 对于手头的问题，喷气流，用（q，9）替换（u，v）坐标很方便，其中tan 9 = v/u和q = \textbackslash\{\}/u 2 + v 2。 这产生L[4] + XN[<I>] = 0 (4.12)，其中L[4>] = q 2<j>qq + f\{q)\{q4> q + 4>ee), n l q 2 \_ (1 + l)(g* 2-q2) a2 ( 7 + l)a* 2 -( 7 -l) 9 2 and q sin 6 I. / 1 -jT-T- < 4> qq\{q(pq + <P9») - \textbackslash\{\}4> qe - -N\textbackslash\{\}i>] = hnr \textbackslash\{\} 4>M, + <M - [ 4> v« - z+') f = °> A[4>] = sin 0\textasciicircum{}, H <f>0。 请注意，smt x = 4> u - cos0cf>q <f>e, (4-13) r = 4>„ = smdcl> q + \textasciicircum{}-<f>»。 (4.14) q 从方程(4.4)-(4.6)和(4.9)中可以看出<j>„ = 0 on 9 = 0,-6 (4.15) <j>e = q 0h when 9 = 0,q = q 0 (4-16) <l>9e + qo<i>q = 0 on q = qo (4.17) cf> = 0 on q = 0 (4.18) 这些边界条件如图3所示。 9 = -8条件的同质性是通过允许x 0的值随变化来实现的

62 COOK、NEWMAN、RIMBEY和SCHLEINIGER的墙角6。 Q = 0时的狄利克雷条件源于该点流动的源性质。 由于x -► -oo (q -* 0)，该流是轴对称的不可压缩源流：§(x,r) \textasciitilde{} \{x 2 + r 2)'1! 2. 对于接触电位，这需要<f> \textasciitilde{} y/q作为g -> 0。 右侧不寻常的二阶边界条件值得特别提及。 虽然有一个二阶导数，但它是一个切向导数，因此实际上提供的信息比（j>替换它（罗宾条件）少（不是更多）。 4.3 数值方法 有限差分用于解决图3中的边界值问题。 建立了（q，6）坐标的均匀网格，并建立了网格节点电位值cj>ij的差分方程组。 请注意，只有右上角的非同质条件才能阻止该系统有一个微不足道的解决方案。 这种情况涉及未知数量h，该数量被分配为1的值。 此选择设置了喷气口的尺寸H。 对于二维问题，我们求解L[4>ij] = 0。 右侧的边界条件导致系数矩阵是非正定的，因此实施了直接而不是迭代的方法。 我们使用相当粗糙的24 x 24网格取得了良好的效果。 对于非线性方向对称问题，我们使用线性计算的结果<p\textbackslash\{\}-，作为初始猜测，然后迭代地进行求解。我们期望这种简单的策略有效，因为三维性的效果相对较小[7]。 事实上，除非我们用线性解平均，否则第一次校正太大，会导致发散。 然后，该计划按计划进行，通常在十次迭代内收敛。 4.4 结果 在获得网格点的接触电位后，物理坐标x和r的值从（4.13）和（4.14）中找到。 这使我们能够绘制计算的自由流线和其他恒定速度线。 壁角为S = n/2的轴对称声波射流的典型结果如图4所示。 坐标已根据喷气式飞机开口的H大小进行缩放。 测量不同喷射的方便参数是放电系数Cd-它被定义为实际质量通量与

轴对称喷管流动 63理想一维射流（8 = 0）。 对于二维流d，等于先前定义的收缩系数C c。 对于轴对称射流，它等于C\textbackslash\{\}（即面积比）。 表1显示了不同壁角度的二维和轴对称声波射流的C<*值。 前四个价值也由Alder [1]运行，并且是一致的。 最后一个值（8 = 5°）在小干扰理论中，与Rimbey非常吻合[11]。 请注意，对于三维情况，系数略小，并且随著墙角的减少而增加。 前者行为是由利普曼[7]预测的，后者是预期的，因为流动越来越接近一维。 8 二维轴对称 90° 0.7455 0.7292 60° 0.8153 0.7944 30° 0.9063 0.8891 15° 0.9574 0.9461 5° 0.9893 0.9853 表1 声波喷射的放电系数值 这些情况中的自由流线在形状上与图1所示相似。 然而，二维声态总是可以实现的

64 COOK、NEWMAN、RIMBEY和SCHLEINIGER图4 5 = 7r/2的恒定速度线，距离比相应的轴对称情况稍长。 此外，随着壁角的减少，这些均匀状态发生在靠近喷嘴出口的地方。 我们还可以改变自由流线的速度Qo。 表2显示了其影响。 同样，这些结果与Alder的结果一致。 不出所料，放电系数随着qo的减少而减少。 对于这些情况，理论上只有x -¥ oo才能实现统一状态。 由于离散化，我们计算了一个有限值，但值确实随着qo的减少而增加。 此外，这两种情况都接近适当的不可压缩（qo -»■ 0）极限：二维流动0.611，可以使用共形映射得出，轴对称[1]为0.591，参考资料Alder，G。 M.，通过孔和收敛锥形喷嘴的窒息和超临界理想气体流动的数值解，机械杂志。 英语。 科学。 21（3）（1979），pp. 197-203。 樱桃，T。 M.，通过Hodograph方法发现的一些喷嘴流，J.

轴对称喷管流动 65 10二维轴对称a* 0.7455 0.7292 0.8a* 0.6936 0.6748 0.6a* 0.6603 0.6389 0.4a* 0.6395 0.6163 0.2a* 0.6280 0.6039 表2 8 = 7r/2 澳大利亚数学的放电系数值。 社会。 1（1959）pp. 80-94。 3. 科尔，朱利安·D.，关于轴对称声波射流的说明，SIAM应用数学杂志，43（4）（1983），pp. 944-948。 4. Courant，R.和Hilbert，D.，《数学物理方法》卷。 II，偏微分方程，威利，1962年，纽约。 5. 古德利，K。 G.，跨声速流动理论，佩加蒙出版社，1962年。 6. V。 Krzywoblocki，M。 Z.，可压缩流体流理论中的伯格曼线性积分算子方法，Springer-Verlag，1960年，奥地利维也纳。 7. 利普曼，H。 W.，孔流的气动力学和气动力学，流体力学杂志，10（1961），pp. 65-79. 8. 纽曼，E.，私人交流。 9. Ovsiannikov，L。 V.，具有直线过渡线的气体流动，NACA TM 1295（1951）。 10. 波戈洛雷夫，A。 V.，椭圆型蒙格-安培方程，诺德霍夫有限公司，格罗宁根，1964年。 11. 林比，S。 E.，超音速喷射：霍多图平面的数值方法，R.P.I. 数学报告编号 145，1/21/85和博士论文，数学系，加州大学洛杉矶分校1984年。

\clearpage
\chapter{超聚焦声爆问题的分析和数值模拟}
超聚焦声爆问题H的分析和数值模拟。 K。 程1kM。 M。 Hafez 2 5.1 介绍 很高兴参加这次研讨会，以表彰Earll Murman，我们的工作涉及Earll和同事们早些时候为研究的一个方面做出了贡献[31]。 “超聚焦声爆”早期在海耶斯作品[20,21]中用于指代飞机马赫数接近并通过阈值时在地面上发生的异常强烈的声爆，这取决于靠近地面和巡航高度的声速之比。 阈值马赫数对应于靠近地面的波场的声音（Machone）条件，温度和声速通常高于巡航高度。 从声学的角度来看，声爆强度增加的现象是波合或光线聚焦在“caustic”的结果，在这种情况下，是双曲-椭圆过渡边界。 如参考[20,21]，我们将在过渡边界及其附近解决和处理标准声波爆炸理论中的故障，以区别于可能在其他地方可能发生的故障（由波融合和射线聚焦引起的）。 1 南加州大学航空航天工程系，洛杉矶，加利福尼亚州 90089-1191。 2 加州大学戴维斯分校机械与航空工程系，加利福尼亚州戴维斯，95616。 Frontiers of 计算流体动力学 - 1998 编辑：David A. Caughey和Mohamed M. 哈菲兹 ©1998 世界科学

68 CHENG \& HAFEZ 从这个意义上说，超聚焦声爆很可能被称为“过渡吊杆”。 Murman的跨声速流动计算开创性工作[24，25]对CFD开发和空气动力学设计应用的重大影响是世界范围内无可争议的。 Earll Murman比任何人都更有助于提高CFD主要进入航空的闸门，这让人想起Keith Stewartson在20世纪70年代初曾经说过的话。 事实证明，他的作品对我们自己的研究的影响也出奇地深刻。 不仅我们早期关于收敛加速和激波拟合 [18, 26]主题的工作是与穆尔曼的跨声速计算程序的增长，而且我们关于跨声速等价规则与升力[3]和跨声速升力线理论[10]的两项工作所依据的计算程序基本上都是穆尔曼的计算程序。 在本文中，我们将再次发现Earll的影响力是无处不在的。 演讲将侧重于超聚焦声爆域方程的推导，以及三维非线性混合型Tricomi方程的计算研究。 作为一个新的理论发展，分析显示了车辆加速/减速对中场故障行为的关键影响，因此对超级/过渡吊杆特性的关键影响。 这里讨论的工作代表了博士做出的重大贡献。 C.J. 李，W.H.先生 Guo和C.Y. Tang作为合作者，仅以会议文件和报告的形式记录下来[4-9]。 这些和最近研究的更完整的论文版本正在准备中。 5.2 在伽利略框架中分析的非等温大气中的声爆 5.2.1 初步说明 按照惯例，我们使用近场和中场的概念分别表示车辆附近的波场（长度尺度为L）和距离足够远（r >> L）的波场，累积的非线性效应可能导致声学特征失真，而大气分层可能导致波衰减和折射。 然而，远场的概念，被一些早期的工作者用来表示远距离波场的渐近/极限状态，例如发生在等温大气中，在目前的调查中不需要考虑。 声爆理论和预测方法的核心是对中场结构的分析。 简化了接近飞机时中场故障的分析，与近场结构的适当匹配

超聚焦声爆问题的分析和模拟69图1 超聚焦声爆问题的坐标对于确定中场（波传播）计算的初始数据或F函数是必要的。 近场匹配引起的一个问题是，近场的非线性（Euler）计算是否实际上可以提高中场分析的准确性水平[4]。 然而，这不是本文的主要关注点。 第二个更重要的问题是上述超级/过渡热潮的中场故障。 本文的主要重点是这个问题的分析和数值模拟。 波传播问题将在固定在车辆（移动声源）上的伽利略坐标系中进行分析，这与大多数现有声爆程序的射声学论文不同[2、14、22、28、29、32]。 该公式使用与时间相关的特征表面（Mach锥形）作为坐标参考，并将揭示射线-声学方法中不明显的波场属性。 它还大大促进了中场两端的匹配分析。 笛卡尔系统（x、y、z、t）和圆柱极坐标系（x、r、w、t）将交替使用，其中原点固定在飞机的天头（机头），选择正z轴指向地面（见图、la、b，z和w的方向和原点）。 5.2.2 二阶小扰动方程5.2.2.1 涡量校正虽然从射线声学方法中不明显，但对中场结构的一致描述必须解释在大气熵不均匀的情况下气压产生的涡量。 根据开尔文-亥姆霍兹方程确定的一阶涡量，控制中场的二阶小扰动方程可以修正为a'VV-m D D DtDt • + H0 Dt |W|2 + 7-1 2a 2 D (5.1)

70 CHENG和HAFEZ，其中<j>和D/Dt是与扰动速度u'相关的电位，分别修正为涡量和对流导数*=v*-w\textasciicircum{} \textasciicircum{} = l+(c7+tr)（5.2）。 在未受干扰的速度向量上方，U = Ui + Vj + Wk，Hp和H分别表示基于基本大气密度和熵梯度H;1 - \textasciicircum{}fjnp, n:\textasciicircum{}\textasciicircum{} = \textasciicircum{}lnp-\textasciicircum{}lnp的尺度高度。 请注意，Eq.中u的涡量组件 （5.2），p' = -pD<f>/dt，很容易将Kelvin-Helmholtz 涡量方程满足到前序，并且方程左侧的最后一个项。（ 5.1）结合了压力梯度和熵梯度的贡献的结果。 对于中场分析，H p,Ha和H p = (dlnp/dz)* 1将与参考刻度高度A进行比较。 p上的下标“A”和其他热力学变量表示与位置z处静止分层大气有关的值。 5.2.2.2 时间依赖性、水平、直线性飞行（V = W = 0）在下面，我们分析了伽利略系统中时间依赖性声爆问题，这是我们以前的研究中没有尝试过的。 分析仅限于仅涉及水平、直线性车辆运动（V = W = 0）的情况，但足以证明加速/减速对超/过渡吊杆波场影响的关键性质。 在这种情况下，二阶PDE变成o？ V2S-U2 2U\textasciitilde{}4>xt - 4>u - ( -77" bx + HD - U - dt dx (5.3) 线性理论的马赫锥形为了捕捉缓慢演变的非线性中场，分析将使用线性理论x-x lc(y, z,t) = 0中的马赫锥体（特征曲面）作为参考表面，该曲面分别对应于U(t)和a(z)中的时间和高度依赖性扭曲的马赫锥体。 功能

超聚焦声爆问题71 xlc\{y,z,t)的分析和模拟由非线性一阶PDE控制，该PDE可以写成1 + |V xzlc|2 - Wl - \textasciicircum{}T) =0（5.4），其中|V ± x lc | 2 = \{x\textbackslash\{\} cf + (x\textasciicircum{}) 2，M = U(t)/a(z)。 该方程可以通过沿着其自身的双特征（根据Courant和Hilbert的术语）积分ODE集来求解，p 2 = x\textbackslash\{\}c>Pz = x\textbackslash\{\} c，p 4 = x] c: dy dz dt dp2 dp3 iz)(U-Pi)2 (5.5) P2 P3 \{U-p\textasciicircum{}a- 2 0 -\{l/2)a- i\{daldz)\{U -p 4 ) 2 dpi \{dU/dt)\{U - Pi)a-2' 为了方便起见，双特征将被称为“bi-char”。 因此，沿bi-char x* c和[xj c - U(t)]是不变的，并将保留各自的初始值，在t = £;，即x\textbackslash\{\}c = p.，x] c - U(t) = (x] c)t - U(U) s -U„(5.6) y \textasciicircum{}ic - -而，x\textbackslash\{\} c = p-j可以是 从Eq中评估。 （5.4）作为？ Lc = ±VR = ±V(U,/a( 2 )) 2 -l-p 2 (5.7) xr 从现在开始，所有初始值都将用下标星号表示。 显然，R的消失确定了过渡边界，越过该边界，Eq的双曲波场。 （5.3）不可能存在。 因此，函数x lc(y,z,t)是从Eq.获得的。 （5.5）以x lc（（z，p«，<«））的形式，其中参数p*和t*扮演两个特征变量的角色。 顺便，我们注意到p 2 = [(U,/a«) 2 - l]cos 2o;,,在每个高度，与横向截止相对应的最大p 2是[(U*/a(z)) 2 - 1]。 令人欣慰的是，（5.6）的U*中的（xj c）»完全为零，因此U，= U（t t）= U t，如果t = t的初始数据在原点x = y = z = 0 3规定。 它也源于Eq。 （5.6）U，是所有在t = i处发出的信号将相对于介质在（负）x方向移动的速度。 出波的中场方程 为了得出一个适合出波锋旁的中场更强烈的部分的控制PDE，我们引入了四个变量（= x-x lc（y，z，t），y' = y，z' = z，t' = t（5.8）3 C.J.首先注意到这一观察将大大促进后续的后续计算工作。 李。

72 CHENG \& HAFEZ，将假设£范围与L相当，而y'和z'与L相比更大。 因此，与y'、z'和适当比例的t'相比，£中的偏导数将被视为大。 变换等式。 （5.3）在变量£、y'、z'和£'中一，并利用参考曲面的属性，Eq。 （5.4），等式左侧（LHS）中£的最高（第二）偏导数。 （5.3）完全消失，而其右侧（RHS）变成RHS = dt 7 + (U „lc >l \{[ i + iv±xicr + 1(U \textasciicircum{}c\textbackslash\{\}2 +2(v L*le.v'to + \textasciicircum{}J 7 (Er *Y)4>t + |v±\textasciicircum{}|2 + 7-1 2a2a2 ' n'v (5.9) 中场中场出波的非线性演变的时间尺度预计可与A/a*相当，£、y'和z'的长度尺度可以分别为L、A和A。 我们规定，波形V的实际s长度不能比飞机L的长度大得多；这是在后验证的。 在上面，Vi和Vj\_分别是（y,z）和（y',z'）的横向场梯度。 遵循方程第二行的条款。 （5.9），以及第一行由d/dt'操作的那些，都可以因为相对较小而被删除。 在分别用刻度集[T'U 0L,L, A, A, A/a„; A, A/a*]规范化[<f>,(,y',z'; x lc,y,z,t]后，最好在变量集[v-idiii 2,i; x lc,y,z,i]中研究得出的中场方程。 然后，时间依赖性中场的方程（5.3）假设归一化形式T(z,U)u-=u+xlc =lc dy U + X f -KZU + OZ d。 U.a* d „ 1 zU + dt 2 LHPX'Z + \textbackslash\{\}dy* + dz* +jjLQ(z,t*,p*) ■u = dy -9 ' c\textbackslash\{\}-r> A d 2 Ql 其中\{/* = dU(t t)/dt*，系数T和Q是（5.10）r = A\textbackslash\{\} 7 + 1 fU tU0 Q(z,U,P*) = U»a» dt r和U 0是恒定参考速度。 （5.10）的RHS显示了订单L/A的所有条款高于LHS上的条款，并将在随后的分析中省略。 “...”表示未写出的条款，包括在不消失的T'L/A假设下与（L/A）2相当的条款。 在上面，变量（y,z,t）和（y,z,i）可以互换。 省略RHS项的方程（5.10），概括了Cheng等人的中场方程对于稳态情况[4，6-9]。

分析和模拟超聚焦声爆问题73沿双特征的ODE解u的一阶准线性PDE可以再次通过其双字符£ - £ c(y,z,i) = £„的ODE系统解决，相应的x - x c(y,z,t) = x„, d£ \_ dy \_ di du N \textasciitilde{} if \textasciitilde{} if \_ u.a. a-» - \_i j\textasciicircum{}pn + v2 ±f ic + u,Q] a (5-n) 其解可以正式表示为沿双字符的线积分，这在随后分析过渡边界附近的行为时将变得至关重要 i-L= [fuR-'\textasciicircum{}dz, y= [p.R-'\textasciicircum{}dz, t-U = \textbackslash\{\}Jtat I a- 2R-l'2dz Jo Jo u = U(£„U,p.) J\textasciicircum{}exp \textbackslash\{\}\textasciitilde{}\textbackslash\{\} JQ\textasciitilde{}Kd\textasciitilde{}z \textasciitilde{}\textbackslash\{\}jj J Xdz\textbackslash\{\} (5.12) 其中2,1。 ,I, - \_ I U*a*\textbackslash\{\} fdt t *,,. 使用了K.Vi\textasciicircum{}/\textasciicircum{}、\textasciicircum{}[\textasciicircum{})|f)/J和双字符属性x| c = pt，以及x\textbackslash\{\} c = VR。 就变量z、参数£»、£*和pt而言，它们实际上是特征变量，解u与线性理论的解相同，在时间依赖的情况下也有效。 在通过特征变量在x、y、z和t中重新表达u时，线性理论的唯一变化是通过\{* = £- I" TuR-l,2dz找到双特征的向前位移。 Jo Shock跳跃中场PDE的弱解承认冲击不连续曲面£ - £D\{y,z,i) =0满足跳跃关系dyD p-x/2 rfy lc di D U tat 1/2 \_dP° dz d« d£ a 2 dz，其中上标“C”，不要与上标“lc”混淆，指的是非线性系统的双色，符号<>表示跳跃的算术平均值。

74 来自C.J.基于上述ODE系统的计算机程序的CHENG \& HAFEZ解决方案示例。 正如Cheng等人[4，5，8]所报告的那样，Lee与通过准确的特征方法[13]和基于射线声学方法[28，32]的PC Boom程序变体获得的结果进行了很好的比较。 5.3 超聚焦声爆问题 对相关现象的彻底调查必须分析故障边界处的中场解决方案行为，并得出重新缩放的控制方程。 在下一部分，我们将研究困难来源及其故障性质的中场解决方案（第3.1节），由此可以得出与渐近匹配原则一致的控制方程和边界条件（第3.2节）。 本节（第3.3节）的主要部分将致力于我们对简化方程系统的计算研究，以改进该问题的数值模拟，Earll和同事在27年前就进行了[31]。 Onyeonwu [27]基于射线声学方法，对飞机机动产生的波聚焦问题进行了半分析处理。 然而，从参考中，下面的3.1.1和3.1.2节中提出的各种明确的中场行为并不明显。 [27]。 5.3.1 中场故障：分析性研究 虽然有关现象无疑是波合或光线聚焦，但从第2.2节的中场解决方案中确定困难的来源和故障行为的类型是有指导意义和有用的。 5.3.1.1 过渡边界和其他地方的中场奇点 中场击穿的两个来源 对于稳定、水平、直线性运动的车辆，中场解u预计仅在双曲椭圆过渡边界R = (U*/a) 2 - 1 - pi =0处分解。 然而，在非均匀的车辆运动（U \textasciicircum{} 0）下，u中的另一种奇点可能发生在中场的其他地方，这与“极限线”的出现相一致[11,33]。 这些非均匀性的来源可以追溯到解形式方程的指数因子中两个（行）积分的积分。 （5.12）：B = VI*-/4 C和x S (\textasciicircum{}) (§) /*J« (5.13)

分析和模拟超聚焦声爆问题75必须注意，V\textbackslash\{\}x lc和dU/di都不能轻易从Eq的ODE解中确定。 （5.5）与x\textbackslash\{\} c,x\textbackslash\{\}c和x\textbackslash\{\}c一起沿着单个双字符轨迹，沿着该轨迹，参数对（t*，p»）是固定的。 [在我们的计算机解决方案中，这个ODE系统的稳定案例[6]，C.J. Lee通过在相邻的双字符处使用和区分x\textbackslash\{\}c和x\textbackslash\{\}c值，成功地获得了Vj\_x lc的准确数值，因此K。] V\textbackslash\{\}x lc和di t/di的评估为了实现目标，必须用沿双字符的已知值来表达V\textasciicircum{}\_5 lc和di t/\&t，从中分解行为可能会变得明显。 Dtt/dt的横向拉普拉斯，需要沿着相邻双字符的数据，可以通过组合方程两侧的yand 2-部分导数来评估。 （5.4），即（4C）2 + (\textasciicircum{}c）2 - (§ " l ) = °' (5-14)，并考虑到U 2对<»的依赖性，以及a 2对z的依赖性。 结果是VP3 P*J y P3 a a 2 a 2 U„ [p 3 \textbackslash\{\}dzj p t\textbackslash\{\}dyj (5.15)，其中a = a(z),U* = U(i„)和U* = dU(t*)/dt*。 偏导数XgCy = dp t/dz、dit/dz和dit/dy都可以由ODE Eq产生的t = i(z,it,pt)和y = y(z,t t,p*)的逆关系来确定。 （5.11），即，其中从(z,t,p t)到(y,z,t)的逆变换为三对fdit dp A (du dpA (du dpA UF'aFJ' \textbackslash\{\}Wdg)' \textbackslash\{\}dr as)'产生三组同时代数方程，其中一个得到

76 CHENG \& HAFEZ U AU, HK N >+<-\textasciicircum{}£(I+¥ -L * V-tt K J St* U, I + p*K (5.16) (5.17) dt a* A 其中A是变换的雅各布行列式（不要与刻度高度A混淆）A= - (I + p,K)-\textasciicircum{}[(N-J)(I+p,K)-p 2 L 2 ] (5.18)，对应于可压缩流体流的经典理论中“极限线”的出现[11,33]。 上面的I、J等是z、i t和p* 1 = f'Z dzt 2 f dh f* dzj Jo v\textasciicircum{}7' *Jo al\textasciicircum{}R,' K -io liJil 3 /2 '的已知函数的积分，它们在过渡边界R = 0附近的行为，其中z = z t*对感兴趣的u行为至关重要。 让m = -(dR/dz) tt,z" = z tt - z，我们有，对于足够小的z" R \textasciitilde{} (dR/dz)„\{z - z„) = mz"。 在极限z" ->■ 0中，I和J仍然是有限的，但K、L和N变得无限：K - m- 3 /VT 1/2 + CK，L \textasciitilde{} 2U 2 a- 2 m- 3 /V|" 1/2 + C L，N \textasciitilde{} 2Uta- 4 m- 3 / 2 |z"|- 1/2 + C N (5.20)因此像|z"| -1 / 2是无限的。 它源于Eqs。 （5.13）以及V\textbackslash\{\}x lc和di t/di对R和A的依赖性，中场解u只能在R或A消失的位置不受限制。 可以预期三类u分解：（a）A = 0，R远非零，（b）R = 0，A远非零，（c）A = 0，R小。 它源于Eqs。 （5.18）和（5.19）A可以随着加速度\{U* > 0消失），U„/U*在控制这三种发生方面的重要作用是显而易见的。 要发生类型（a），需要大加速度U»/U，= 0（1）。 然而，在以下分析中，我们将更加关注小|U,/U„|对超级/过渡繁荣的关键影响问题，这将被视为属于上述（b）和（c）类型。 顺便说一下，我们可以指出，对于导致类型（a）的单位序U»/U*加速度，预计会出现代数类型分解（变量z“' = z ttt - z，其中z = z tt*在A = 0）。

超聚焦声爆问题的分析和模拟77 5.3.1.2 过渡边界附近的中场行为：对车辆加速度的关键依赖 分析将研究车辆加速对过渡边界附近（或不远）波场影响的几个案例，其中0 < z" = «„ -z «1 while，对于发生在z = z* + *** i 0<z'" = z ttt-z<0(z")« 1 注意z,,, < 2»*和z" = z" tt z'" + (z ttt - •? ***)■从Eqs中可以看出。 （5.16,5.17）和（5.18），|U*/U»|对于正在研究的问题来说必然很小。 然而，除了加速度（U»>0）外，波场马赫数 U,/a不需要关闭为1。 这些情况，即U，= 0，U«/U* < 0和U„/U» > 0，将分开考虑。 （i）稳态（U，=0）此态包括横向截止，不限于超级/过渡臂的情况。 在这种情况下，x lc和A在过渡边界附近的跨场拉普拉斯是，R \textasciitilde{} mz", 1*-h\textasciicircum{}H('4)' \textasciicircum{}>*>在这里，A不会消失，但鉴于边界附近的K行为，Eq。 （5.20），它可能变得无界，实际上有助于降低V 2j\_xlc的奇点强度。 方程函数k的结果。 （5.12），不限制p2的相对大小顺序，和\textbackslash\{\}/R，I ndz \textasciitilde{} -Zn|z"| - In z" + 2p 2 /m 3 / 2 I„ + 0(1)，将稳定情况下过渡边界附近的u减少到- frit « \textbackslash\{\} [P* Co(z„,pt) V P \textbackslash\{\}y/z" + 2pi/m i/zItt\textbackslash\{\}1/'，其中Co是方程中积分有限部分的指数。 （5.12）。 结果恢复了早期在Hayes的[20, 21] 超聚焦声爆研究中推断的，其中p\textbackslash\{\} = [(U,/a*) 2 - 1]COSUJ*与\textbackslash\{\}fz T\textbackslash\{\}相比必然小，至少在对称平面上。 结果将其适用性域扩展到不仅限于小p»，因此不仅限于小[(U t /a,) 2 - 1]。 有趣的是，对于非消失的pi，u在过渡边界处有界，但鉴于该极限中x\textbackslash\{\}c的奇点，u和中场方程的导数仍然必须进行修改。

78 CHENG \& HAFEZ（ii）车辆减速的影响：U,/U* < 0 对于较早时间的车辆减速，即U,/U* < 0，在t = i»，A不能在过渡边界附近完全消失。 考虑到K、L和N的行为，R和\textbackslash\{\}的线积分可以评估为小R，因为I adz m 3/2 A 'u。 I«+p.-.u, \textasciicircum{}2 U, U a. ln\textbackslash\{\}z"\textbackslash\{\} Xdz I -- Jo \textbackslash\{\} a** \textbackslash\{\} a* / A** \textbackslash\{\}V m P2t. Inz 这导致向 R - 0 指数衰减，这是由于 t = i, u \textasciitilde{} U\{£ t,tt,pt] P* C3(Z ttt,tt,pt) \_x/2v /p7 9... b"|-+l/4 (5.21)，其中C3是方程指数因子中积分的有限部分。 （5.12），和A = 1 m3 /2A U，l.，--（- u，/u。 U. la，1+？ Ait 这个减速的结果不限于小\textbackslash\{\}U*/U*\textbackslash\{\}，只要U。 4 1 + U2和A不会在其他地方消失（由于较早的加速）。 而在这种情况下，指数衰减表明中场描述可能不会在过渡边界附近崩溃，但可能仍然存在R = 0的邻域，其中指数小u的控制方程必须从Eq的控制方程中修改。 （5.10）。 对于小p 2，u行为对弱减速的敏感性，从方程中取决于z"的因素的替代表示中更清楚地表现出来。 （5.21），即z"\textbackslash\{\}-1/i™p(-filn\textbackslash\{\}Z"\textbackslash\{\}-\textbackslash\{\}\textbackslash\{\}\textbackslash\{\}z»\textbackslash\{\}-\textasciicircum{}

分析和模拟超聚焦声爆问题79，该问题简化为Hayes [20]结果，对于微不足道的U/U和p 2，但在z范围内变得指数小，足够接近z tt，其中>> Vz（Hi）早期车辆加速的影响U«/U* > 0，A可以在到达过渡边界（R = 0）之前消失。 考虑到方程中K、L和N的奇异行为。 （5.20），行列式A在z = z„*附近，其中R = 0 I - )1 a* U, I,-\textasciicircum{}-\textasciicircum{}.. m3/2Vz" [\textbackslash\{\}j t 其中S*» = J»* - CN - CKU\textbackslash\{\}/«*»>，可以在4 / a* II \_ II \_ \textasciitilde{}* / "* \textbackslash\{\} f-2 I \textasciitilde{}r ut ut\_2 -p*（5.23）表达式Eq. （5.22）A可以用变量z'" = z tt* - z交替书写，该变量与A一起消失。 也就是说，A 1U *T Z a* 1 V\textasciicircum{}ZT 函数«和x采取形式，对于小A和小R = mz" = m U。 1 / 1 + pi 2R a, A \textbackslash\{\}VS \textbackslash\{\}R\textbackslash\{\} 3'2 U« 1 J /U I, 2 /U, u, A [ V<w l-R| 3/2 m v a*» 1 + u. i? «M/a。 V 2 1 2p* v\textasciicircum{} 77 - v 7 \textasciicircum{} ™ 3/2 \textasciicircum{}(v/i 77 - xAST) 以与以前类似的方式进行积分，在使用z"的定义后，导致I Rdz \textasciitilde{} -\textbackslash\{\}ln\textbackslash\{\}z''\textbackslash\{\}+-\textasciicircum{}-(Vz\textasciicircum{}+Jz\textasciicircum{}ln\textbackslash\{\}Vz\textasciicircum{}-./zJZ\textbackslash\{\})-ln 4 V mI«» a / v>i

T ¥T A** m U* a* CHENG \& HAFEZ 这些在加速度影响下产生了ti的分解行为/p7 C 4(z*»,£*,p*)，其中，要确定z" = z"\# \# + 2'" > z'" tt，以及(5.24) e» = y/rn ul(qj\{a*J \textbackslash\{\}/*t" + m*/>)，在t = i时加速度足够小，«"" << 1 [cf. 等式，（5.23）]。 因此，e t很小，就像（[/*/£/*）的平方一样。 在极限U -> 0中，Eq。 （5.24）在稳定的情况下恢复。 然而，在z 1'的范围内，加速效应接管了u行为，并赋予了u一个完全不同的代数行为，其特征是i4U(i/4)+e -iv\textasciicircum{}-v/\textasciicircum{}r 1+e*>>i-注意y/z\textasciitilde{}" - Jz 1\textasciicircum{} \textasciitilde{} \{l/2)JzJ\textasciicircum{}- z'" for \textbackslash\{\}z'"\textbackslash\{\} « <'„。 5.3.2 超聚焦声爆问题的重新缩放：简化方程对兴趣的表述至关重要的是将现象视为中场描述的不均匀性。 我们将首先考虑在稳定、水平、直列飞行中发生的超级/过渡吊杆问题。 5.3.2.1 稳定、水平、直线性飞行：超聚焦声爆接近阈值马赫数我们使用双星号下标来表示B 2 = 0的局部条件。 假设位置z„和z*»的所有有限流动/波属性都是从中场解中知道的。 重新调整超聚焦声爆区域，我们引入局部笛卡尔变量x lc >r，z y =y

超聚焦声爆问题的分析和模拟81可以争辩说，在感兴趣的区域，x"c必须与x" = x - x\textasciitilde{}ll相当；后者不应该比LjA大多少，除非（T'A/L）变得无界。 然后，马赫圆锥体在稳态中的行为表明，z"\_被重新缩放为（L/A）2/3，|V\_L£ 1C |和\textbackslash\{\}V]\_x lc\textbackslash\{\}的相应大小分别为（L/A）1/3和（L/A）"1 / 3。 达到此比例集的另一种方法，允许更通用的x lc行为，在后面的第3.2.2节中显示。 因此，有关过渡区的新变量集可以引入为£ = £, (y',S')=(\textasciicircum{}j (»',-*"), * = A"V (5.25)，以及x" c = x lc - x\textbackslash\{\}ct及其导数的重新归一化为7\textbackslash\{\}-' /r\textbackslash\{\}-l/3 *'L (*i°,*i c)=(f) (4 C.4C)，因此，所有带有插入帽的数量都将假定属于该域中的单位顺序。 3D非线性Tricomi方程在这个过渡区的新变量中，在取A = (r')- 1/3 (r'j) ''(7 + 1)- 1 (5.26)并将两边乘以A \_1 (L/A) 1 / 3, Eq. （5.10）及其在RHS上的余数，在这种情况下变为l\textasciicircum{} 2 + 2 (x\textbackslash\{\}cX- + x\textbackslash\{\} c \textasciitilde{}\textbackslash\{\} *, + Vix lc • *i = V 2 ± \$ + • • • (5.27) V v dy' dz'J f i，其中中场方程的拉普拉斯\$从余数中出现，成为一个一阶重要性的项。 等式中省略的其余项。 （5.27）最多与顺序（L/A）2 /3相同，在约束r'（A/L）下。 将Vj\_\$添加到其他中场方程中，允许控制方程更改为混合（超椭圆）类型，这在跨声速流动理论中很熟悉。 必须强调的是，等式中选择的尺度。 （5.25）主要需要从剩余组中取出\$的拉普拉斯，无论

82 CHENG和HAFEZ在LHS上存在非线性项。 等式中的因子A。 （5.26）已被选择以允许在该域中完全的非线性效应，尽管非线性项的相对重要性在现阶段尚不清楚，并将取决于中场的入射波强度，稍后将对此进行更仔细的检查。 重新表达等式更有用。 （5.27）在笛卡尔变量中，通过逆变换£=£ + xlc(y',z')，y = y'，z = z'。 （5.28）在这些笛卡尔变量和马赫圆锥形属性中，PDE（5.35）从超音速空气动力学中转变为一种熟悉的形式：dx（£X*+H a2 d 2，其中方括号内的\$£系数可以用超音速参数K 2 /A\textbackslash\{\} 2 / 3 M 2-l识别，这个PDE是椭圆还是双曲的，取决于（dB 2 /dz）t*z + \$£是负的还是正的，这取决于（dB 2 /dz）t*z + \$£是负的还是正的。 方程的推导。 （5.29）来自Eq. （5.27）假设B2 > 0，因此受到限制。 它允许椭圆域的扩展可以通过在同一组假设下使用变量x、y、z和\$直接从原始二阶PDE推导来验证。 Eq.的2D版本。 （5.29）是非线性Tricomi方程，已被用作超聚焦声爆问题的模型，并由Hayes [20，21]和Seebass [30]分析研究，并由Seebass、Murman、Krupp [31]进行计算。 上面的方程（5.29）代表了解释波折射三维方面的版本。 我们可以从Eq中推断。 （5.25）与中场尺度A相比，域的横向长度尺度小，但比近场尺度L大得多。 等式的弱解。 （5.27）允许像跨声速小干扰理论一样进行冲击跳跃。 在双曲区域，\$j的中场的一阶PDE从Eq中再现。 （5.27）：！ <•«>■ (+2(ii'|+1"a)*«-\textasciicircum{}i"-*«=°提供了变换等式。 （5.28）与相同的x lc（反向）一起使用。 因此，与中场的匹配得到了保证，但在渐近理论[11，34]的接受意义上，更合适的匹配必须是或大z，如下所示。

分析和模拟超聚焦声爆问题83的远边界条件，x lc的行为接近z" = 0，K或K的行为被视为-p 1/2；然后u中的指数函数与|.z"| -1 '4-成正比，因此，中场接近过渡边界为z" -» 0为«\textasciitilde{}(-) U((,p t)Co(pt,z„)\textbackslash\{\}z"\textbackslash\{\}-\textasciicircum{}\textbackslash\{\}\textbackslash\{\}rlP/ **，其中U可以与F函数直接相关[4-6]。 因此，u的大小接近水平（A/L）1 /6，适合描述性术语“超聚焦声爆”。 这种行为已被数值证实[5，6]。 就适当的超级/过渡繁荣变量而言，这个表达式变成u = \textasciicircum{} \textasciitilde{} (|) (r'jj (7+l)tf(f.,p»)C7.(p„*o) (fy (5)" 1/4 (5.30)，并提供了PDE (5.29)所需的远场行为，确保与中场的良好匹配。 要匹配的两个解决方案的共同域必须是1 « * « \{L/A)- 2'3，或（L/A）2 /3 « \textbackslash\{\}z"\textbackslash\{\} « 1。 在地面平面上，要施加不渗透性条件。 在计算中，只要在上述z范围内完成，直接与（L/A）1 ' 3（r'A/L）（7+ l）u匹配也是合理的。 与Hayes和Guiraud的相似方程（5.29）与远场条件（5.30）、地面平面条件以及上游和下游的消失扰动要求的关系，构成了一个包含两个相似性参数的系统，即z gr标识了地面平面与B = 0的距离。 这个系统等同于Hayes的[20]，其中具有因子\{L/R)1的类似参数！ 6出现在还原的PDE中，而不是在远边界条件下，其中R是因果几何中的曲率半径，先验未确定。 继Gill和Seebass[15]和Plotkin[28，29]的Hayes之后的作品没有考虑（或正确解释）地面影响，而地面影响本应是声波-繁荣影响分析的一个重要方面。 是否需要重新定位坐标系中的地面平面

Eq的84 CHENG和HAFEZ。 （5.29）对于涉及垂直运动的更一般案例，也有待质疑。 在Guiraud的[16]研究中，因子（L/A）1'6也出现在u中，对于正在加速的车辆来说，这不能严格正确（见下文）。 等式中的条件。 （5.30）表明，对于此处假设的有界（T'A/L），\$£可以被视为渐近小，如（L/A）1'6，其中前导近似值为等式。 （5.27）将成为真正线性三科米方程的三维版本[33]。 Guiraud[16]考虑了消失（L/A）1/6的极限，其中由此产生的线性系统的形式主义允许（L/A）1'6完全扩展。 然而，这种形式主义忽略了非线性在过渡边界z = 0的几个不同区域中的作用，这反过来又产生了简化的相似性中不包含的反射波结构，稍后将进一步阐明。 然而，这种基于渐近小（L/A）1/6的处理对于计算目的来说可能不是绝对必要的，因为顺序（L/A）1'6项仍然可以与单位顺序项一起包含在分析中，因为方程的其余部分属于高阶（L/A）2/3。 关于波场结构的进一步讨论将在稍后的Sec末尾以更广泛的语境进行。 3.2.2。 5.3.2.2 加速 超聚焦声爆：水平、直线性运动 第 3.1 节前面显示的更复杂的代数行为，用于与时间相关的案例需要偏离匹配的 渐近展开 的标准程序。 以下分析仅适用于仅涉及水平、直线车辆运动的问题，第3.1节对中场故障行为进行了相当彻底的分析。 引入规范函数，从前面关于稳态的子节中可以看出，关键是确定横向空间尺度，这将把拉普拉斯V j\_y>从中场PDE的剩余组提升到一阶重要性的等级。 分析将关注与z = z ttt的奇异边界的小距离，如前面第3.1节所述，该奇异边界被识别为R = 0或A = 0。 我们将使用第3.1节中前面介绍的变量x"' = x -x„», z'" = zttt - z, y'" = y-y t„ 请注意，如果崩溃发生在z = 2»*» < z ttf，我们有z" = |„ - z = z'" + (z* t - z»*,)。 与第3.1节一样，我们允许A在R之前消失在双曲椭圆过渡边界附近，其中R和z"很小，而不限制z",z'"和z",„ = z tt - z t„的相对大小。 这将允许明确地描述从稳态状态的过渡

分析和模拟超聚焦声爆问题85 超聚焦声爆到以加速度为主。 为了完成与第3.2.1节类似的任务，需要一个特定的尺度，然而，由于缺乏中场行为的简单代数表达式，这个尺度并不明显。 为了克服这一困难，我们引入了两个量表函数：一个是“J”，它测量了该区域的横向尺寸，并且需要渐近地消失（L/A），另一个是“g”，它测量了与时间相关的参考马赫锥体的dx lc/dz或ft（并且可能在极限（L/A）-4 0中消失、有界或无界）。 如第3.1节前面所示，dx lc/dz或ft在焦点附近（\textbackslash\{\}z'"\textbackslash\{\} = \textbackslash\{\}z - z***| << 1）的行为是（假设）通过分析或准确的数值手段已知/建立的。 这意味着q作为5的函数将采取与dx lc/dz作为fi\{\textbackslash\{\}z'"\textbackslash\{\})的函数相同的形式。 因此，我们可以把/i作为q，写ft = \textasciicircum{}！ \textasciitilde{}<K|z"'|)，更重要的是假设，在大多数情况下，x lc行为显然满足了这个规则，这在\textbackslash\{\}z'"\textbackslash\{\}-中是代数的，我们注意到q(z'")不需要在z"' = 0时完全消失。 根据第3.1节，它的分解行为可以通过代数函数u \textasciitilde{} g*„v(z )来表征，其中g ttt独立于z'"，因此在超/过渡域中，u必须重新缩放为u = 0\textbackslash\{\}g.„u(S)]，（5.33）允许数值因子g*„与1大不相同。 例如，v（5）用于加速过渡的行为已被建立为z"和z" tlf I/ = |z"|" 1/4的代数函数，并且可以转换为z'" = z" - (z»* - z t„) = z" - z"*«- 到目前为止，x lc的大小尚未得到保证。 有趣的是，如果人们推测x lc的大小是物理上可接受的顺序，即L/A，Eq。 （5.33）然后可以确定z'",5的尺度的大小，从5q(S) = | (5.35) q\{\textbackslash\{\}z'"\textbackslash\{\}) = 0\{q(S)\} (5.31 = 0[Sq(S)] (5.32) J-T\textasciicircum{}c, (5.34)

86 CHENG \& HAFEZ Since q \textasciitilde{} y/x m + z”，规则Eq。 （5.31）适用于Eq。 （5.35）在这种情况下，产生8的三次代数方程，可以方便地表示为（£），+（£）'--«••'-©■ <-» 它给出k = z'l\textasciicircum{}/S作为y的代数函数，= (z" t „)(L/A) 2 ' 3，即k-3+k-2 = m- 2\{iJ,)-3。 （5.37）对于小z" ttt与（L/A）2 /3相比，从稳定状态J =（L/A）2 /3，q（S）=（L/A）1/3，\textasciicircum{}（L/A）”1 / 6的标度以上恢复；对于4'，，与（L/A）2 /3相比，等式。 （5.36）或（5.37）产生S = (4**r 1/2 (£/A), q(5) = \{z'Uf'\textbackslash\{\} v = (\textasciicircum{}J\textasciicircum{}OL/A)” 1 必须验证这个尺度5和q(5)的假设是否会导致PDE的适当重新缩放兴趣问题；然而，通过随后的审查，这一猜想是正确的。 新的变量和S的确定：带有等式的拉普拉斯的出现。 （5.31,5.32）、（5.33）和（5.15），我们将x'\{ c, V ± x lc, V]\_x lc和u重新归一化，因为要进一步检查的原因，并接受后续验证，我们将变量归一化为t = Z，（yJ）\textasciicircum{}S- 1（y'”，z“）和£=f；0 = [G...！ /\textasciicircum{})]" 1 • y> 其中£和i被假定不受重新放的影响，G„» = f0g***d£。 时间依赖的中场PDE Eq。 （5.10）水平直线性运动现在可以减少到A\textasciicircum{}J 5l (\textasciicircum{} + \textasciicircum{}J\textasciicircum{}°fe'rA\textasciicircum{}'J (5'38) 其中\textasciicircum{}B) 1 \textasciicircum{}-

分析和模拟超聚焦声爆问题87和三重星号表示过渡位置的条件。 现在很明显，在\textbackslash\{\}z'"\textbackslash\{\} = 0(8)的范围内，S由Eq.确定。 （3.3）S-q（5）=\textasciitilde{}等式RHS上的拉普拉斯。 （5.38）将从中场剩余组中出现，与左侧的项排名相同，而（时间依赖）中场方程中的分层和非稳定校正下降到Eq的RHS上的余数水平。 （5.38）。 方程中余数的参数形式。 （5.38）表明等式的有效性。 （5.38）将需要r L，5 5 = -- << 1和- << 1，qA q 上述第二个要求是由于忽略了原始中场方程中的时间依赖项。 关于比例选择的更多评论，在方程的推导中对y'"和t的重新缩放做出假设的原因。 （5.38）需要更多的审查。 测试ODE系统Eq。 （5.11）以及PDE等式。 （5.14）表明y应与z共享相同的尺度因子，即5；事实上，一旦扰动进入跨音波场，就无法防止横向传播，除非施加了特殊类型的边界条件来防止此类事件的发生。 如果不允许y和z的相同比例，超聚焦声爆域的横向范围永远无法正确恢复。 接下来，对ODE系统的因果检查可能表明，dt必须缩放并减少一个系数S/q。 但这个建议是错误的，因为沿着双车的dt只是指信号旅行时间（事实上，即使在稳定极限中也会保持），因此无关紧要。 相关时间尺度由加速度比U t/Ut决定，该加速度出现在声源强度的变化率中，比系数5/q增加的时间尺度要长得多。 因此，时间尺度与中场的时间尺度没有变化。 简化为Tricomi型边界值问题 变量£、y和z中的PpE（5.38），省略了RHS上的剩余组，只是变相的跨声速小扰动方程的扩展版本。 透过逆变换x =( + x lc(y,z)重新转换为笛卡尔形式

88 CHENG和HAFEZ，y和z保持不变，如Eq. （5.28）在稳态情况下，超聚焦声爆域的二阶PDE变为1，M-l „（|r + |r）\textasciicircum{}与a\{z）q，而参数A和q由Eq常数固定。 （5.31,5.32）、（5.33）和（5.35），一旦U和U在更早的时间被知道，这个PDE中的M，与跨声速小扰动方程[11]不同，必须以与z的单位阶范围内的尺度和马赫锥形属性一致的方式变化。 在这里，商（M 2 - l）/q2可识别为跨声速相似性参数K的变体，即M2 - 1 \textasciicircum{} mz" \_ z + k = q 2 \textasciitilde{}! 根据Eq.，f\textasciitilde{} \textasciitilde{} T+k，其中k = z"„/8已被证明是fi = £»*»(■£/A) -3 /3的代数函数。 （5.37）。 在极限k -> 0中，-Jf变为5，Tricomi方程[20,21,30]的非线性三维版本被恢复。 在相反的（加速度主导的）极限中，-K接近单位和等式。 （5.39）被发现是标准跨声速小扰动方程，其特定值为K = - 1（在超音速范围内）。 Tricomi方程的双曲椭圆过渡边界（通常在z = 0）被加速度和非线性移位移到z = -k - A(\textasciicircum{}i(l + k)在z小于这个值时，方程。 （5.39）将分析延续到椭圆域；通过分析更完整的小扰动方程来验证其有效性，使用变量x、y、z和i。 从Eq的弱解中推断出的激波属性。 （5.39）遵循跨声速小扰动理论中的那些[11]。 边界条件 混合椭圆双曲问题将作为每个t的准稳题来解决，双曲边y \textasciitilde{} u(5z)/u(8) (5.40)的远边界条件在K<f<<5 -1或8 «\textbackslash\{\}z'"\textbackslash\{\}«l (5.41)上得到满足

分析和模拟超聚焦声爆问题89，以便与Eq中所示的小z"'处的中场双曲解相匹配。 （5.24）或（5.33）。 事实上，PDE（5.39）可以承认Eq.描述的远场行为。 （5.40）由Eq保证。 （5.38），可以从Eq中恢复。 （5.39）并恢复了中场方程的形式，一旦其中RHS上的tp的跨场拉普拉斯变得不那么占主导地位。 v(8z)和v(5)的更明确的形式可以从Eq中推断出。 （5.34）其中z'"分别被Sz和5取代。 在双曲域中，不允许远在风中扰动，扰动必须在椭圆域中远在上下风中消失。 假设z = z gr处的（水平）地面平面上的不渗透边界条件，即y = 0。 （5.42）在地面平面远远低于过渡边界的情况下，需要在</>上进行易逝行为。 应更仔细地检查高k或高/u下的边界/匹配条件以及场结构（见下文）。 顺便，人们观察到，如果地面不像条件Eq所暗示的那样平坦和光滑。 （5.42），它与入射波的相互作用无法在伽利略框架中保持准静止波场。 相似性参数；对Hayes和Guiraud理论的限制。在水平加速度下，上述研究揭示了，除了A和z gr之外，超聚焦声爆域的相似性还有两个参数。 即，\textbackslash\{\} - F vS -a A = 1 »*»，[I - Z ttt q，其中q是\textbackslash\{\}A(1 + k)，S = z'l tt/k，可以从Eq中确定。 （5.37），v可以从Eq中推断出来。 （5.34）其中z"被5 + z"„ t取代。 从Eq中回忆是有帮助的。 （5.23）z'U = o，其中A、a、U和t以物理尺寸给出；其大小从HSCT的10-3[35]到高速冲刺中的战斗机的10\_1不等[36]。 使用这四个参数，波场相似度以重新缩放变量x、y、z、t;ip和（pi. 后者可以与压力系数相关。-\textasciicircum{}/ o -gr和t t 0 a * dt

90 CHENG \& HAFEZ £« = i* = a»t., t /A被列为上述参数，以表示U t = Z7"(i*)和£/»/[/*进入A,/i, S和q的确定中的简化问题的准稳定性质。 尽管i表示比t更早的时间瞬间（半分钟左右），但人们可以从z = z ttt（中场分解位置）的马赫锥体x\textbackslash\{\}c的瞬时x位移率中正确推断/推断出U t和U*，因为U表示马赫锥体相对于介质的水平推进率（在Eq.给出的段落中提前。 （5.7）和Eq。 （5.8））。 与目前的做法相反，上述分析表明，加速超音速飞行的聚焦声爆不能仅由A和z gT表征。 添加/i和t，以及相关的缩放，对于描述相似性至关重要。 尽管PDE（5.39）中K的比率（M 2 - 1）/<J 2的线性i依赖性将允许其转换为更接近Refs理论中的Tricomi形式。[20，21，30]，但更一般的匹配条件（5.40，5.41）不能像稳态那样简单地要求u RS（Z）-1 ' 4来满足，除非fi小到可以忽略不计，即，这为Hayes的相似性的适用性提供了一个标准/限制。 当不满足这一要求时，即使通过某些临时程序绕过远场匹配问题，也无法正确确定相似性参数A。 一个未回答的问题是参数A是否控制了方程中的非线性效应。 （5.39）可能仍然像稳态一样保持渐近小。 检查8、A和q对fi的依赖性表明，无论\textbackslash\{\}x的大小顺序如何（衡量加速度影响的相对重要性），它确实如此。 即使在无界JU的极限中，参数A仍然很小，如（z”，，）1 /4或（iX/U，）1 / 2因此，波场再次由线性、类似Tricomi的域主导，其中完整的非线性效应将主要表现在入射波形的前缘和后缘与过渡边界的交叉点附近，该边界现在位于z = - k。 其中，预计在每个交叉点的小邻域中都会有一个具有自身非线性跨声速相似性的波场，这反过来又决定了反射波前部和后部附近的场结构。 很明显，严格基于消失A极限的相似性无法充分描述跨声速相互作用区的结构和反射波细节，这可能对理解/解释测量的波场记录很重要。 在加速度控制状态下（/j >> 1），超聚焦声爆波场

超聚焦声爆问题91的分析和模拟的特点是本质上是线性的（超曲面的）Prandtl-Glauert区（前面提到的），其中解及其次阶校正与一侧的奇异中场行为相匹配，另一侧的线性Tricomi方程匹配。 后一个域在维度上比Prandtl-Glauert区大得多。 有趣的是，行为等式。 （5.40）位于Prandtl-Glauert区的两侧。 在这里，非线性主要表现在过渡边界沿线的一些微小邻域中。 关于数值模拟和一些感兴趣的案例的评论 除了大气温度反转外，车辆加速度预计将大大影响超聚焦声爆强度及其波场结构，除非（U»/U*）2与（L/A）2 /3小得可以忽略不计。 准稳态非线性3D PDE方程。 尽管如此，混合类型的（5.39）可以通过适用于稳态示例的相同数字程序解决每个英镑（将在Sec中讨论。 3.3如下），具有适当的边界条件和参数。 事实上，PDE可以被一个更完整的方程的跨音小扰动取代模拟研究，但如果没有与上述类似的分析的指导，就无法正确推断远边界的条件和位置。 除了需要完成上述发展的几个部分的分析外，实际感兴趣的相关问题的未处理方面是1。 超聚焦声爆涉及车辆垂直运动和转弯操作的问题。 2. 聚焦声爆在过渡边界附近以外的位置导致中场故障。 3. 横向截止边界附近的中场崩溃和过渡问题。 5.3.3 计算研究 下面描述的计算研究将主要关注解决稳态非线性三科米类型方程的混合型三维边界值问题。 尽管简化方程中的缩放和参数可能有很大不同，但相同的数值程序和算法应该同样适用于水平加速度的情况。 然而，如前所述，在更一般的情况下，当加速涉及垂直和横向分量时，预计会偏离PDE（5.29）的形式。 对于下面要讨论的参考[5]中的数值示例，SR-71的中场计算修改了入射N波的案例

92 CHENG \& HAFEZ 图2 考虑了计算域稳定飞行[17]。 5.3.3.1 程序和算法计算域；边界条件 计算工作将提供感兴趣的非线性三维波场，在适当缩放的无量笛卡尔变量的矩形域中进行，如图所示。 2. 与前面讨论的研究的理论部分不同，正z轴将朝着与引力加速度相反的方向走。 上边界的边界条件由分布在区域的中场入射波数据提供。 后者的几何形状非常像具有弯曲中心线的高长宽比机翼。 双曲域的所有边界也需要适当的辐射条件，这将允许outcfotnci WELVGS的传输，并让传入事件通过入事件，而不是规定此类辐射条件集，更权宜的条件d> x = 0 0as sbee nsed on nhe euper rnd dide eoundaries sx除了边界上的“冲击区”，需要6

超聚焦声爆问题93的分析和模拟与N波形相对应的弦曲线。 由于双炭在椭圆-双波过渡边界附近的3-D折射特性，只有相对较小的入射波在上界能够到达感兴趣的非线性混合型区域的内部。 因此，在这种情况下，不正确的边界条件只能影响撞击区下游的波场结果。 在整个地面平面上，施加了不渗透性条件（d<p/dz = 0）。 为了方便起见，附加到变量x、y、z和it、Eq上的caret符号。 （5.27），将被省略。 使用了基于Murman的[24，26]（类型依赖）一阶方案及其高阶版本的加权平均值的有限差分方案和计算程序线松弛程序。 在高阶版本中，额外的上风点被添加到超音速区域的差分操作模版中，而亚音速、声波和冲击点的模版保持不变。 本方案通过与一阶方案（高阶格式）•（1 - e）+（1阶方案）• e与e = O（Ax）相结合，补偿了高阶方案的分散和不足的耗散误差。 这种组合产生了一个二阶准确的整体方案。 发现二阶方案需要一个大的t来补偿其色散余数。 具有耗散余数的三阶方案在大多数计算中都是首选和使用的。 与原始的Murman程序一样，从一种操作员类型切换到另一种操作员类型的标准至关重要。 对于后者，在詹姆森[23]之后对紧凑型切换算法进行了修改。 算法和计算程序的主要特征如下所示。 穆尔曼方案采用于非线性三维三角方程（i）基本程序 计算程序是在默尔曼[25] fx + 4>yv + 4>zz = 0与/ = K(z)4> x - \$1/2，K(z) = K= K'(0)z，对应于PDE中的有限差分形式，在默曼[25]fx + 4>yv + 4>zz后，应用于PDE的PDE的有限差分差分法形式中的线松弛行进形式，K为K(z) = \$1/2，K为K(z) = K'(0)z，对应，对应于PDE中的(dB 2/dz)ttz5.29。 一阶和二阶部分y和z导数由中心差商表示，而第一部分x导数由中心差商计算，差运算模版中/和<f> x的x导数由

94 CHENG \& HAFEZ 四种不同的运算符，取决于问题点，例如i，可能是亚音速、音速、超音速或（a的下游）冲击。 中心差应用于亚音速点的f x，这使得方案二阶准确；f x在声点设置为零，这是数值稳定性所必需的；后向差分算子（相当于中心差商向上移一个网格）应用于超音速点的f x，这使得方案耗散和局部一阶准确。 在冲击点，f x被计算为冲击前部和后部的中心差商的总和，这给出了局部正常冲击的正确跳跃条件。 不足为奇的是，该解决方案无法准确捕捉斜冲击，因此，基本程序在双曲区域的准确性较差。 （ii）高阶后向差分为了提高双曲域中的差解精度，可以使用二阶方案（fx）i - 2（fx）i-l - (f x)i-2，但其余部分是色散的。 提供耗散余数的三阶方案是（/*）* = 3(/,)<\_i - 3\{f x)i-2 + (f x)i-3。 这两种方案都用于研究非线性三角问题，结合前面描述的基本一阶方案。 我们早期研究中的大多数计算都是用二阶和一阶组合方案进行的。 （Iii）操作员切换方法 在Murman的原始程序中，切换标准基于本地马赫数的标志。 Jameson [23]引入了一个开关函数e= J 1，如果（***££*-K）>1；（否则为0。 这将差分方程与适当开关以紧凑的形式放置。 当前程序采用詹姆森的变体，允许更渐进的过渡，并增强数值稳定性，如下所示：fi+i/2 \textasciitilde{} fi-i/2 4>j+\textbackslash\{\} - 2<t>j + 4>j-i。 4>k+i - 2<f> k + \textasciicircum{}fc\_i \_ Ax Ay 2 + Az 2 fi+1/2 - /i -1/2 fi-1/2 - fi-3/2 Ax €' Ax

分析和模拟超聚焦声爆问题95分钟。 (£;,£;\_!, £i\_ 2) fi-l/2 - fi-3/2 \_ /; -3/2 \textasciitilde{} /i-5/2 Ax Ax -min。 (E,-,e;\_i,ej\_2,ej\_ 3) fi-3/2 \textasciitilde{} fi-S/2 fi-5/2 \textasciitilde{} fi-7/2 Ax Ax 左侧是中心差分近似；右侧的第一个项导致一阶上风方案，二阶和第三项是高阶校正。 请注意，二阶和三阶校正在数值上不是保守的。 （四）对穆尔曼方案的修改 穆尔曼方案承认了扩张冲击。 声波点运算符的简单修改在这方面是有用的。 我们提出以下/rt n -1，n t\_n Wxx • <Px - Vyy，其中（p = 4> - Kx和n是迭代级别。 中心差分近似用于y导数，而\textasciicircum{}导数则由上风差分近似。 在声线中，<p xx对于加速流动总是正的，上述方程的抛物线性质不允许在声线附近的速度场中发展不连续性。 数值通量不再是保守的，但是守恒的误差是局部且小的，因为<p x在这个区域几乎为零。 为了实现高阶方案，需要在声线附近逐步过渡方案的顺序，以保持数值稳定性，图。 3说明了拟议的综合计划。 此外，在通过声线平稳减速期间，Murman的冲击点算子可能会被中心差分方案取代（如Murman和Cole的原始非保守方案）。 同样，数值通量守恒的误差是局部且小的，但近似值至少与微分方程一致。 严格要求在整个激波进行保护。 另一种方法是Engquist和Osher方案。 他们的通量计算是统一保守的，不允许膨胀冲击。 更多详情可以在本卷Engquist的论文中找到。 计算结果和讨论 本工作中的计算非常像高长宽比、高扫翼的跨声速空气动力学逆问题。 上界的输入数据是从中场计算中获得的[5]，其中F函数从SR-71飞行实验[17]的记录中采用并略有修改，重新缩放（保持r'固定）到

96 CHENG \& HAFEZ 图3 通过声波边界加速流动的拟议方案。 符号表示应用了六个difFernet运算符的网格点，o椭圆；•声波点；-f一阶上风；X二阶上风；*三阶上风；0冲击点。飞行马赫数非常接近标准大气中的阈值M th = 1.1533。 3D计算；横向截止区 3D计算使用了两个（x、y、z）域。 较大的一个延伸到-0.1 < x < 32；0 < y < 40；0 < z < 5.926，较小的范围超过-0.1 < x < 25；0 < y < 30; 0 < z < 5.926。 见图中的草图。 3.1 较大的跨度宽度比上边界冲击区的跨度长；跨度较短的跨度宽度没有完全覆盖上部冲击区。 通过对对称平面连续z高度下过压波形中较大和较小域的计算的计算，验证了具有较短跨度的修改冲击区的充分性，即y = 0（图 4a），在离飞行轨道更远的另外两个跨度站（对称平面），y = 10（图。 4b）和y - 20（图。 4c）。 两个数据集之间的差异太小，无法从数字中看出，这表明排除上部撞击区尖端附近的部分入射波数据是不必要的。 两个计算集都使用了（800 x 81 x 201）的网格。 相同的结果通过一个稍细的网格（1000 x 121 x 201）密切地再现，这表明计算域假设的上边界的高度

分析和模拟超聚焦声爆问题97图4压力分布的比较是充分的。 然而，需要一个更细的网格来解决过压峰值及其相邻的描述（见下文）。 然而，地面超压模式的横向范围已经以合理的准确性确定，可用于定义阈值马赫数处/接近的“横向截止区”。 后者因光线声学而丢失，在当前大多数程序中使用“焦点区域模型”[28]。 图中显示了阈值马赫数处地面过压轮廓的示例。 5. 在飞行轨道附近证实的二维假设 二维描述的假设隐含在估计/预测超/焦点吊杆强度的现有程序中；通过在参考[28，29]中与几个飞越实验的比较，发现了其明显的有用性。 通过比较图中前面给出的基于对称平面（y = 0）附近的非线性超聚焦声爆理论的三维计算，详细检查了这个假设及其有效性域的充分性（在对称平面的稳态情况下）。 4a和使用相同的x和z范围内的（800 x 201）网格进行相应的二维计算，因此具有相同的机械间距。 这里，在图中。 6，我们观察到，在对称平面上，2-D和3-D结果之间的差异似乎非常小，并且只在地面附近变得明显（但很小）。 因此，除了峰值压力和波形的不确定性外，二维假说可以被认为足够用于大多数实际目的（至少沿着飞行轨道）

98 CHENG \& HAFEZ 图5 峰值附近的阈值马赫数处地面的过压轮廓。 波形的详细比较表明，在这种情况下，假设可以一直工作到y>10，前提是波形中的下风偏移也用于解释冲击区的局部扫荡，而现有的二维模型无法合理确定。 在y = 10之外，非线性和三维效应在“横向截止区”中建立过压行为具有决定性意义。 网格加密的重要性；独特的折射/相互作用特征，而图中的比较。 4a和图中。 固定网格和固定计算（x，z）域的6可能足以表明修改后的上边界条件（利用较短的跨度）和对称平面中的二维假设的充分性，它们未能正确指示极限PDE解，因为网格大小（网格点间距）消失。 一系列二维计算可用于指示附近的波形如何

超聚焦声爆问题的分析和模拟99图6 2D和3D压力分布（线性化）的比较在z = 0和地面（z = min z）的过渡边界可能会发生变化，因为网格越来越细；这个过程确实带来了与超聚焦声爆现象相关的特征。 在无花果这里。 7a，b，我们重复图的二维结果。 6用于三组网格，网格大小Ax通过减少z范围连续减少，而不改变x和z范围内的网格点数量。 后者固定在（800 x 201）。 图7a将粗网格结果（-0.1 < x < 25）与细网格结果（-0.1 < x < 16）进行了比较，图。 7b将精细网格结果与更精细网格结果（-0.1 < x < 9）进行比较。 几个特征是超聚焦声爆波形的熟悉特征，特别是前部的过压峰值和真实冲击后的超压，

CHENG \& HAFEZ Cp Distribuflons的比较 0.0 5.0 IOD 图7 网格分辨率对二维压力分布的影响随着网格尺寸的减少而逐渐锐化。 现在可以看到，超压峰值是上界峰值的四倍多。 就物理变量而言，本例中的繁荣强度在最大值下达到每平方英尺4.4-5.0磅（psf）。 还获得了与低于阈值的飞行马赫数相对应的亚临界情况的解决方案，并揭示了更明显的U形过压曲线（\#波形的倒U形）的趋势，但在地面上强度较小，分布更平滑。 压力过高，类似于兔子的耳朵从反射波形的两端升起，被认为是与双曲-椭圆过渡边界的存在以及（相关）线性双曲波场与跨越过渡边界（线性理论）的非线性混合层区域的相互作用明显相关。

超聚焦声爆问题的分析和模拟101由线性Tricomi域主导的不同波场结构本研究提出的关于超聚焦声爆波形结构的独特概念是存在两个明显不同的相邻波场：一个由线性Tricomi方程支配的大空间区域和一个不太广泛的区域，其中非线性对从双曲到椭圆行为的过渡至关重要。 这个概念可以从第3.2.1节末尾的讨论和等式之后的备注中推断出来。 （5.38）对于相对较小的Su/q。 计算研究详细证实了这一点。 为此，我们首先研究了非线性项4> x(f>Xx与控制方程中Tricomi项K'(0)z ■ cf> XI的相对重要性；它们的比率是<f> x/K'(Q)z，如图所示。 8a作为z的函数，对于三个x。 非线性在Tricomi域的广泛z范围（1 < z < 5.925）中不重要，从跨度站y = 4.5的3D计算中显示的比率值非常小可以看出。 在其他跨度站和相应的二维解决方案中也发现了类似的行为。 必须注意的是，i w 5.925的上边界位置是根据理论对中场分析分解位置的估计来选择的，该理论与非线性无关。 这个概念的意义在图中更直接地被提出。 8b对于二维情况，线性和非线性解之间的明显差异只能在较小的范围内找到，z < 1。 结果表明，开发更有效的计算程序以及基于渐近方法的场结构研究的可行性。 图中显示的非常精细的计算。 6a，b和无花果。 另一方面，7a、b表明，线性和非线性之间的明显差异导致反射/折射斜冲击附近的区域，这些区域无法由Tricomi方程的纯线性版本解释。 冲击拟合问题 尽管事实证明网格加密有助于带出超聚焦声爆域中的这些独特特征，但它在从过渡边界中去除的减震的更斜的部分效果不好。 这是由于所使用的冲击点运算符的固有不足。 通过采用Hafez和Cheng[18]的冲击拟合算法，基本问题已经克服了，然而，事实证明，无论发生多重冲击并变得过于接近，如从墙壁反射的冲击，都很难实现。 如果冲击太弱，无法与压缩波区分开来，它也会变得无效。 正在开发一个程序，当冲击不连续性变得太弱或当墙时，程序将从冲击拟合切换到冲击点运算方案

到达附近。 Mllo.vs的情节。 Z at y = 4.5 Comparfson of Cp DlsbibuUots 5.4 结束语 我们讨论了最近关于声爆传播工作的主要特征，这是我们对超级/过渡繁荣现象研究的支撑。 该分析采用伽利略框架与飞机一起移动，并使用线性理论的马赫锥形作为坐标参考（允许时间依赖）。 该理论明确界定了过渡边界附近的奇异中场行为，以及水平车辆加速和减速对波场行为的关键影响，之前没有具体考虑过。 已经设计了一种匹配策略，允许在嵌入式非线性3D域中适当缩放和还原/恢复修改后的Tricomi类型方程。 通过为超音速小扰方程开发的松弛方法，简化的混合、椭圆-超音问题适合于数值求解。 除了现在可以识别地面撞击区域的3D波场描述外，该解决方案允许（LWX2DlZDl.raordadn.r(r)--02\textasciitilde{}）8.0 -0 5.0 10.0 0.0 4.0 2.0 0.0 ww-O.z\textasciitilde{}W X 图8 非线性的重要性lo00 X 81 X101. rcmdadn.ltI--OZaI) a - 8 - 6。 WW 1-8.812284 -.-- - I r 7.819774 - --. -4- -5-c-30 -1.0 1.0 5.0

分析和模拟超聚焦声爆问题103在阈值条件附近/处的重要地面和声波边界相互作用，并正确满足冲击跳跃（通过冲击捕获/拟合算法），在现有（焦点区）模型中没有考虑。 接近阈值的飞行马赫数的计算分析细节揭示了几个预期和意想不到的特征，包括2D假设的充分性，对其在对称平面上的准确性的评估，以及存在明显的线性、双曲域，以及一些孤立的、相对较小的、非线性的混合型区域。 这项研究的实际结果的主要结论是，车辆加速度对超聚焦声爆结构和相似性的影响是显著的。 然而，只要根据理论选择缩放和远边界条件，为稳态情况开发的相同的数值算法和程序预计将同样适用。 对于空间发射应用来说，对涉及爬升和垂直加速度的问题进行分析是显而易见的，这可能会改变控制方程和修改后的Tricomi方程的边界条件。 所呈现的大部分结果都是与C.J.合作的成果。 李，C.Y. 唐和W.H. 郭，我们要向他表达我们最感激的感谢。 我们要感谢M。 唐宁，J。 爱德华兹，G。 哈格伦德，H。 林，K。 Plotkin和R。 Seebass提供许多刺激性的讨论和有用的建议。 感谢AF阿姆斯特朗实验室OEBN噪音效应部门在1996年3月19日至9月18日期间监测的国防部SBIR计划的支持。 参考资料1。 阿什利，H.和兰达尔，M.，（1965）翅膀和身体的空气动力学，艾迪生-韦斯利。 2. Carlson，H.W.和Maglieri，D.J.，（1972）“声波-Boom生成理论和预测方法的回顾”，J. 声学社会。 美国，卷。 51，不。 2（pt. 3），pp. 675-685。 3. Cheng，H.K.和Hafez，M.M.，（1975），“跨声速等价规则：涉及升力的非线性问题”，J. 流体机械，卷。 72，pt。 1，pp. 161-187。 4. Cheng，H.K.，（1996）“支持HSCT声波冲击研究的超音波场分析”，Proc。 1995年美国宇航局高速分辨率。 计划：声爆研讨会，在美国宇航局兰利研究。 中心，9月 12-13，1995年，卷。 1，pp. 136-150。 5. 程，H.K. 哈菲兹，M.M. Lee，C.J.和Tang，C.Y.，（1996）“声波轰预测和潜艇撞击研究的进一步程序开发的改进理论评估”，最终报告提交给AF Armstrong实验室，OEBN噪音效应分部，HKC Research 11月。 28 1996。 6. Cheng，H.K.和Lee，C.J.，（1996）“声爆理论中的问题被研究为奇异扰动问题”，AIAA论文96-2165。

104 CHENG \& HAFEZ 7. Cheng, H.K., Lee, C.J., Hafez, M.M. and Guo, W.H., (1995) “分层大气中声波波传播的CFD问题”，第一届亚洲CFD会议，香港科技大学，香港，1月。 16-19。 8. Cheng, H.K., Lee, C.J., Hafez, M.M., and Guo, H., (1996) “声爆 传播及其潜艇影响：理论和计算问题研究”，AIAA论文96-0755（1996）。 9. Cheng, H.K., Lee, C.J., Hafez, M.M. and Tang, C.Y., (1996) “非等温大气中的聚焦声爆作为混合双曲椭圆、非线性、三维问题：”Proc. 第二届亚洲CFD会议，东京，12月。 1996年，卷。 2，pp. 337-342。 10. Cheng，H.K.和Meng，S.Y.，（1980）“斜翼作为跨声速流动中的提升线问题”，J. 流体机械，卷。 97，pt。 3，pp. 531-556；另见Cheng，H.K.，（1982），“跨声速流动中高长宽比机翼渐近流动结构的计算研究”，跨声速冲击和多维流动，（ed. R.E. Meyer），Acad。 新闻，pp. 107-146。 11. Cole，J.D.，（1968）应用数学中的扰动方法，Blaisdell Pub. 公司。 12. Courant，R.和Hilbert，D.，（1965）《数学物理学方法》，卷。 II，偏微分方程（由R. Courant），Interscience Pub. pp. 75-88，173-174。 13. Darden，CM.，（1984）“冲击凝聚分析，包括应用于声爆外推的三维效应”，美国宇航局TP2214。 14. Darden，CM.，（1988）（编辑） “声爆方法和理解的现状”，美国宇航局会议。 酒吧。 3027. 15. Gill，P.M.和Seebass，A.R.，（1975）「苛性的非线性声学行为：近似解决方案」，AIAA Astro.和Aero.的进步，ed. H.T. 长松，麻省理工学院出版社。 16. Guiraud，J.P.，（1965）“Acoustic geometrique, bruit ballistique des avions supersoniques et focalisation”，J. de Mecanique，卷。 4，pp. 215-267，17。 Haering，E.A.，Ehernbergen，L.J.和Whitmore，S.A.，（1995）“SR-71 声爆传播实验的初步机载测量”，美国宇航局技术。 备忘录。 104307。 18. Hafez，M.M.和Cheng，H.K.，（1977）“休克拟合应用于跨声速小扰动方程的松弛解”，AIAA J.，卷。 15，不。 6，pp. 786-793。 19. Hayes，W.D.，（1968）“几何声学和波理论”，Proc。 关于声爆研究的第二次会议（ed. 我。 R。 施瓦茨）NASA SP-180，pp. 158-。 20. Hayes，W.D.，（1969）“通过苛性非线性声学传播的相似性规则”，第2届会议。 声爆研究，美国宇航局SP-180，pp. 165-171。 21. Hayes，W.D.，（1971）“声爆”，年度修订版。 流体机械，卷。 3，pp. 269-290。 22. 海耶斯，W。 D.，Haefeli，R. C，以及Kulsrud，H。 E.，（1969）“声爆计算机程序在分层大气中的传播”，美国宇航局CR-1299。 23. 詹姆森，A.，（1975）“跨声速势流计算，使用守恒形式”，Proc。 第二届AIAA CFD会议，哈特福德，1975年6月，pp. 148-161。

分析和模拟超聚焦声爆问题105 24。 Murman，E.M.和Cole，J.D.，（1971）“跨声速流动飞机的计算”，AIAA J.，卷。 9，pp. 114-121。 25. Murman，E.M.，（1974）“通过松弛法计算的嵌入式激波分析”，AIAA J.，卷。 12，pp. 626-632。 26. Hafez，M.M.和Cheng，H.K.，（1977）“收敛跨声速计算的放松解决方案加速”，AIAA J.，卷。 15，不。 3，pp. 329-336。 27. Onyeonwu，R.C.，（1973）“声爆聚焦飞机机动影响的数值研究”，多伦多大学航空航天研究所，UTIAS报告编号 192. 28. Plotkin，K.J.，（1985）“声爆焦点区预测模型的评估”，Wyle研究实验室。 报告WR 84-43，2月。 1985年。 29. Plotkin，K.J.，Downing，M.和Page，J.A.，（1994）“美国空军单事件声爆预测模型：PC Boom 3”，高速研究美国宇航局声爆研讨会，美国宇航局会议。 酒吧。 3279，pp. 171-184。 30. Seebass，R.，（1970）“苛性声学的非线性声学行为”，Proc。 第3个。 会议。 声爆研究，美国宇航局SP-255，pp. 87-120。 31. Seebass，R.，Murman，E.M.和Krupp，J.A.，（1970）“不连续信号附近行为的有限差分计算”，Proc。 第3次会议。 声爆研究，10月pp. 19-30，1970年（ed. 磅。 施瓦茨），美国宇航局SP-255，pp. 361-371。 32. Thomas，C.L.，（1972）“用波形参数法推断声爆压力特征”，NASA TND-6832。 33. Tricomi，F-G.，（1957）“Equazioni a Derivate Parziali”，Cremonese，罗马；另见Ferrari，C.和Tricomi，F-G.，Transonic空气动力学，（Transl. R.H. Cramer）学术出版社，pp. 105-117（1968）。 34. Van Dyke，医学博士，（1975）流体力学中的扰动方法，Acad。 出版社（1964年）；抛物石出版社，注释版，pp. 106-120。 35. Haglund，George T.，私人通信（1996）。 36. Downing，M.和Zamot，（1995）“美国空军聚焦声波爆炸的飞行测试调查”，J. 声学社会。 上午，卷。 97，不。 5，pt。 2，Proc。 第129届会议声学社会。 上午，1995年5月。

\clearpage
\chapter{跨声速流动的复杂分析}
跨声速流动 Connie K的复杂分析。 陈 1 k 保罗 R. Garabedian 1 6.1 导言 自复杂特性的方法导致设计看起来像实用机翼部分的无冲击机翼以来，已经过去了近三十年[1,10]。 大约在同一时间，Murman和Cole发明了他们简单、高效的上风差分方案来计算跨声速流量[9]。 很快，计算流体动力学成为超临界机翼技术最喜欢的工具[2]。 然而，真正有效的是Murman的分析概念，而只有少数专家能够从设计方法示例中提取像Korn 翼型一样成功。 本文的目标是描述新FLOW代码的基本特征，该代码为复杂特征方法提供了更强大的实现。 首先，我们将揭开两个复杂变量域的一些奥秘，并展示数值计算是如何组织的。 然后将举例说明最近取得的进展。 也许要讨论的最重要的贡献是利用流动方程解的分析性来定义更准确的剖面坐标的过滤器。 1 Courant数学科学研究所，纽约大学，纽约，NY 10012 Frontiers of 计算流体动力学 1998 编辑：David A. Caughey和Mohamed M. 哈菲兹 ©1998 世界科学

108 CHEN \& GARABEDIAN 6.2 霍多图变换 由于我们关注的是无冲击流，而且流体通过冲击前部时熵的变化是冲击强度的三阶，我们假设熵是守恒的，这意味着流动是等熵的。 因此，存在速度势 <p(x,y)和流函数ip(x,y)。 让q是q 2 = u2 + v2的流速，让c是c 2 = dp/dp的局部音速，让M = q/c是局部马赫数。 对于tp，我们得到二阶准线性偏微分方程（c2 - u 2）ipIX - 2uvip xy +（c2 - v 2）ipyy = 0。 （6.1）由于我们正在处理跨声速流动，这个方程是混合型的，当M2 < 1时是椭圆的，当M2 > 1时是双曲的。 我们的目标是获得部分规定了速度分布的无冲击翼形。 为了实现这一点，我们需要解决翼型的跨声速流动场的设计问题，因此我们使用hodograph方法将（6.1）转换为线性系统。 为了方便计算，我们将流动区域的共形映射引入到复特征平面的椭圆上。 然后，我们在指定区域上计算了一组完整的特殊解，以在椭圆上用切比雪夫多项式表示<p和tp，其中可以有效地研究具有复杂几何的流。 配置文件是流线型的，其中\%j）=0。 一旦获得ip和ip，翼型的坐标将从公式x + iy= I -- (dip + i\textasciicircum{}-\textbackslash\{\} = I - (dip + z\textasciicircum{}-j. （6.2）我们可以应用测距变换来获得线性系统9，M 2 - 1，，„ s ¥e = -V q，<p q = ipe，（6.3）p qp，其中9是流动的角度。 从现在开始，我们将把（M，v）平面或（q，0）平面称为hodograph平面。 定义特征的常微分方程是（l-M 2 W<j 2,„, n, „ (i-M 2 wy 2, s \textasciicircum{} TJ-L- + d92 = 0, dip 2 + i \textasciicircum{}- = 0. 6.4）q2 pi在亚声速流动区域，其中M 2 < 1，系统（6.3）是椭圆的，其系数是分析的。 因此，ip和t/>在所有参数中也是分析的[5]。 因此，我们可以分析地将所有因变量和自变量扩展到复杂域中。 在这种扩展之后，方程的类型不再被区分。

跨声速流动 109的复数分析 我们可以在复域中引入特征坐标，但它们不是唯一的。 假设hi和h 2是任何非平凡的分析函数，并且£和r\textbackslash\{\}ax是特征。 那么a = hi(()和j3 = h,2(r])仍将是一组特征坐标。 测度图方法的一个困难是测度图平面中的流动区域未知，可能具有复杂的几何形状。 我们通过引入共形映射a = J y- dq + i\$ = f(0, (6.5) Vl-M2 f VI - M 2 P = I q dq-i0 = f\{fj)来克服这一困难，至少在亚音速情况下，它将平面£ = fj中的椭圆转换为平面a = ft中的流域。 这里/是一个分析函数，a和（3成为新的特征坐标。 有一个由6、q、tp和ip组成的四个规范方程组，具有新的自变量£和rj，因此每个方程中只有一个特征方向的微分。 当«/>(£, f?)的特征初始数据时，这种公式允许通过连续近似得到系统的解 I)和4>((i, rf)沿着两条复杂的马赫线rj = r]i和£ = £i给出-当初始数据被切比雪夫多项式[4]定义为(和rj的函数时，可以应用高效的有限差分方案，在£平面的椭圆及其在r\textbackslash\{\}-平面中具有£ = fj的反射图像的笛卡尔积中构建一个完整的解系统。 适当的积分常数（6.5）意味着当a +（5 = 0时，M 2 = 1。 这定义了四维复域C2中的二维声波位点。 它与实物理平面的交点对应于声线。 由于向特征坐标的变换（6.5）中的奇点，复杂特征的方法在声学位点上崩溃。 我们用来获得数值解的差分方案在M 2 = 1的点上变得奇异，在附近点条件不好。 因此，必须构建所有计算网格，以避免声波位点。 6.3 奇点 我们一般介绍二阶偏微分方程的基本解，以分析跨声速流动的奇点的具体行为，这些奇点代表霍多图平面中的源、汇和偶极子。

110 CHEN \& GARABEDIAN 让我们考虑线性椭圆微分算子的形式！[ *] = jA* - A*，（6.6）其中A是x和y的实解析函数。 L[F] - 0的基本解被定义为F = F(x,y;x,y) = R\{x,y;x,y)\textbackslash\{\}og - + Q(x,y;x,y), (6.7) r的解，其中r = yj(x - x)2 + (y - y)2，其中x，y是参数，其中R和Q在参数点（x,y）的某个附近是规则的。 系数R将是R(x, y; x, y) = 1的黎曼函数。 让我们对变量x、y、x和y的复值执行F的分析延续，并引入复数替换£ = x + iy，T] = x -iy，| = x + iy，fj = x -iy。 （6.8）复变量£和77实际上是（6.6）的特征坐标。 很容易验证！[ *] = 9 (n - A* = 0 (6.9) 就这些特征坐标而言。 我们现在使用乘法代入\textbackslash\{\}1/ = cr\{q)ip来简化hodograph平面中流函数tf>的流方程，以便它获得规范形式（6.9）。 如果我们插入基本解决方案，并注意r2 = (£ - £)(? 7 -fj），我们得到恒等式w\textasciicircum{}-\textasciicircum{}-w\textasciicircum{}+w-0-（6-10）（6.10）中的前三个项是奇异的。 它们不容易相互取消，因为对数函数是多值的，但极点不是。 为了满足（6.10），我们必须同时拥有L[R] = 0和R<\{t,f)>iv) = a, Rviln,Lv) = 0- (6.11，但这些正是黎曼函数的属性。 我们可以通过有限差分或连续近似求解R。 当我们这样做时，（6.10）中三个奇点之和变为正则，剩余项Q也可以计算。 在上一节中，我们介绍了映射功能/在平面上和£平面中的椭圆之间。 我们需要在椭圆内选择点来表示机翼的无限速度，或者代表叶片级联的入口和出口流量。 这些点是指定翼型的偶极子或级联刀片的源和沉的奇点。 很明显，这种流动可以从基本解决方案的角度进行分析计算。

跨声速流动 111 6.4 非线性边界值问题的复杂分析 我们使用逆法设计无冲击翼型。 这涉及一个自由边界问题，其中翼型的轮廓将被确定为流线\$ = 0，但速度分布q = q(s)将沿着边界规定，其中s是弧长。 当速度足够高时，流动将变成超音速，偏微分方程将是混合型的。 我们根据规则<p(s)= I q(s)ds, T= (bq(s)ds)确定沿剖面的速度势 tp和循环T。 （6.12）通过这种方式，提出了一个可以用复杂特征方法处理的非线性边界问题。 特殊解决方案的线性组合用于解决hodograph平面中无冲击流动的逆问题。 实数亚声速流动区域对应于平面a = j3，或平面£ = fj。 为了使ip和ip在这个区域是实的，施瓦茨反射原理被用来构造作为实函数的规范方程的解。 我们现在引入切比雪夫多项式，以在焦点为±1的椭圆上表示解析函数，其中它们构成了一个完整的正交系统。 在变换£=4(C+Vi, (6.13) 2 V C)下，椭圆可以在£平面上用一个圆环来识别，焦点之间的直线对应于单位圆|(| = 1。 在这种情况下，切比雪夫多项式表示为（6.14）T.«） = 5 [C + \textasciicircum{}) ■ 让我们用洛朗级数/(« los (frf) + £amr = log (frf) + £anTn(0 (6'15)与m = a-m)表示分析图函数/流动问题，其中c\textbackslash\{\}是与驻点对应的椭圆上的点，其中c 2是常数，\textbackslash\{\}c 2\textbackslash\{\} > | Ci | 选择来改进收敛。 观察，如果5是环外圆的半径，w是一个合适的角度，我们有H(\textasciicircum{} +£ f) (6-l6)

112 CHEN \& GARABEDIAN在椭圆上。 为了确定映射函数/在（6.15）中定义的值，我们需要一个解ip，在椭圆边界上ip = 0的流动方程的ip。 另一方面，为了计算解决方案tp，ip，我们还需要一个映射/来满足要求（6.12）。 这导致了一个非线性边界值问题，其中Re[/]和Re[ip]同时在|（\textbackslash\{\} = 5. 使用光谱方法，我们将问题分解为映射函数/和流函数ip的一对耦合方程，这些方程将通过迭代方案解决，我们将更详细地描述。 首先，我们在椭圆中找到一个不可压缩的解ip，它可以在没有地图函数/的情况下确定。 £平面中的椭圆已被识别为£平面中的圆形环。 <p的计算是在环上使用快速傅里叶变换进行的。 通常，128个术语足以提供足够的准确性。 下一步是应用规定的函数q = q(s)来确定关系q = q(tp)，然后将其与不可压缩解ip一起使用，沿着椭圆确定q。 根据（6.5），q的值在映射函数上提供了一个边界条件。 同样，快速傅里叶变换给出了/的系数。 确定/后，我们返回并通过受边界条件ip = 0约束的规范方程的特殊解的线性叠加来计算椭圆上的tp和ip。 从解ip中获得一个新的映射函数/，并且可以重复迭代。 该过程收敛得足够好，因此在大多数情况下，一次迭代会产生合理的翼型。 在跨音域的情况下，我们需要修改边界数据和映射函数/以使问题摆好。 如果M>1，则hodograph平面中的超音速点不再映射到复域中的平面£ = fj。 地图函数/将把亚声速流动区域转换为由声波位点M 2 = 1界定的£ = fj的椭圆的一部分。 在这个椭圆的其余部分，q和0变得复杂，实超声速流动区域对应于在音速线上穿过平面£ = fj的复域C 2中的表面。 在这种情况下，我们的边界值问题变得更加复杂。 为了求解跨声速流动中的流函数ip，我们在音位点M 2 = 1之外的椭圆部分规定了一个人工边界条件。 沿着与实数亚声速流动相对应的椭圆弧线，我们仍然强加条件ip = 0，但在边界的其余部分，我们设置Re </>(£, £) + CIm »/>(£, 0 = 0, (6.17)，其中常数C是经验确定的，以提高分辨率。 事实证明，此类问题提出得很好，计算证实，通过将边界条件强加在圆\textbackslash\{\}(\textbackslash\{\} = 8等距点的特殊解的线性组合上获得的线性方程组合是非常有条件的。 至于地图功能/，我们

跨声速流动 113的复杂分析 图1 采用复杂特性方法设计的对称翼型周围的无冲击流动中的马赫线。 新的FLOW代码的强大性能使得有可能在相对厚的轮廓上获得一个非常大的超音速带。

114 CHEN \&GARABEDIAN需要修改我们在亚音速情况下制定的边界条件，以确保它适用于混合类型。 在椭圆的亚音速弧上，我们像以前一样将RAf\textbackslash\{\} = J y/l-M2 dq - log e-ci i-c2 (6.18)，但在其他地方，我们分配了值Re[f] log e-cj t C2 (6.19)，其中K\textbackslash\{\}和Ki再次是经验常数，q = q(Re[ip])。 我们在这里描述的复杂特征方法在FLOW代码中得到了有效的新实现。 图。 1显示了由该代码设计的对称、10\%厚的无震翼型，通过绘制马赫线来显示异常大的超音速区域。 6.5 算法结构 FLOW代码是从早期工作[3,4]中开发的，为无冲击翼形的计算复杂特性方法提供了强大的实现。 它以规定的速度分布q = q(s)读取，沿轮廓的弧长s作为参数。 由于输入数据在某些点可能是密集的，在其他点可能是稀疏的，因此使用样条例程来获得重新缩放的、更均匀分布的值集。 速度作为速度势 ip的函数q = q(<fi)迭代计算，这是在上一个周期中计算的。 映射函数的实部分的边界值/是从规定的速度分布中找到的。 然后，通过计算与椭圆相关的切比雪夫多项式的级数展开系数来解决Re[/]的狄利克雷问题。 为了构建流函数和速度势，我们解决了特征初始值问题，既要找到一个完整的特殊解系统，又确定指定无限流的奇异项。 双曲系统采用的稳定有限差分方案((n+1 - (n)(rjn+l \textasciitilde{} Vn) (6.20)沿特征求解微分方程，其中(mVn是复平面路径上的网格点。 从一整套特殊解决方案中，一个线性的离散系统

跨声速流动的复杂分析图2椭圆图，其中流动方程的解在复杂特征E平面上计算，用于提供图中所示的对称无冲击翼型的FLOW代码的运行。 1. 积分的亚音速路径由这个椭圆内的实线表示。 虚线代表一对超音速路径，椭圆本身的阴影标记是边界值问题的插值点。 驻点和后缘的图像由最左端和最右端的十字表示，星号指定偶极子<a = [b和初始点5的位置。流函数的系数获得方程。 解的实部分a.在边界条件中用于椭圆弧，其中流动是亚音速，但剩余的弧需要复杂的值。 定义解系数的线性方程组是有条件的，因为我们在计算中使用切比雪夫多项式。 对于亚声速流动，我们的方法与通过特殊函数的叠加来解决线性边界值问题的更传统的程序没有什么不同。 然而，在跨音域的情况下，该方法具有一个新的特征，因为边界条件沿着椭圆弧施加，对应于复杂域C2中的非平凡曲线。 对于规定的速度或压力分布的任何合理选择，当自由流马赫数不是太大时，我们描述的计算结果是无压力的轮廓。 通常，为了更好地适应规定的速度分布，迭代几次地图功能是多余的，因为

116 CHEN和GARABEDIAN沿着轮廓的超音速弧的完美一致是无法期待的。 然而，有必要引入一个过滤器来消除与复方根A = Vl - M 2的不准确值相关的错误点。 算法中最困难的步骤是确定正确的分支。 数值实验表明，一个有效的全局分支位于半平面Re[A] < -Im[A]中。 构建复杂特征方法的积分路径非常重要，以至于我们将下一节专门讨论这个问题。 那里将展示如何避免交叉声波位点M 2 = 1，我们的方程变得奇异。 为了确定超音速区域轮廓的坐标，我们奠定了一组超音速积分路径，并沿着实特征计算解。 在跨声速情况下，这些特征绘制在超声速流动的真实区域中。 我们沿着它们进行整合，通过寻找流函数tp符号的变化来确定轮廓上的点。 6.6 集成路径 让我们讨论计算级联压缩机叶片所需的集成路径。 单个翼型的路径与级联的路径遵循相同的原则，只有轻微的变化。 表示级联的源和汇的奇点位于特征£平面的点£a和£j点。 在单个翼型的限制情况下，£„和£l之间的距离消失，源和汇重合形成偶极子。 路径从£a和£j绘制到初始点£c，从中追踪进一步的路径，以解决流动方程的特征初始值问题。 需要从£tt和£j到£c的初步路径来处理黎曼函数的特征初始值问题，黎曼函数是定义源和汇的奇异项的系数。 在亚音速情况下，相对于实对角线£ = fj，路径在复域上对称，这只是意味着77平面轴中的路径反映了£平面中路径的图像。 因此，我们只需要计算网格点三角形上的有限差异，因为较大的方形网格上的值可以反射到对角线上。 这对于跨音子路径来说是不可能的，因为它们必须绕过复域中的音位点M2 = 1。 因此，£-平面和77平面的跨声速路径不是彼此的反射图像，而是以相反方向穿过的椭圆上的弧线结束（参见 图。 2）。 当两条路径在垂直于对角线的段上相遇时，计算就会停止。 这些是椭圆上的点，稍后用于解流函数的边界值问题。 流函数的值在弧上不再是实值

跨声速流动 117的复杂分析 图3（a）特征a平面中的超音速路径和（b）实霍多图平面中的相应特征的示意图。 虚线是声域，为了避免M 2 = 1，路径不得相互交叉。在声域之外的椭圆。 值得注意的是，无论如何，这些点导致了复杂域中产生有用的无冲击跨声速流的良好边界值问题。 边界值问题解决后，必须单独找到真正的超音速解。 超音速带的真实特征不能通过直接越过音速线来确定，因为当M 2 = 1时，规范方程会失败。 为了避免我们选择由特征a平面中水平线段和垂直线段联合组成的多边形超音速路径，其中声波轨迹位于虚轴线上（参见 图。 3）。 每个多边形都旋转，以避免M 2 = 1的点，然后以声波位点的弧线结束。 由于声波位点只是一个二维曲面，而计算是在四维复域C 2中完成的，因此总是可以移动

118 CHEN kGARABEDIAN的整合路径，以便它们绕过声波位点。 图中的示意图。 3a显示了a平面和/3平面中的两条路径是如何构建来实现的。 粗线表示平面中的路径，而细线表示/5平面中的路径，虚线表示声波位置。 图。 3b显示了真实测速图平面的相应特征。 来自两个超音速路径的声波位点对对应于Hodograph平面中超声速流动的实点。 当这些点沿着声波位点移动时，特征a平面中的相应点在虚轴上向上或向下移动。 因此，a的虚部分的行为就像实特征坐标。 成功取决于选择路径的早期部分，以便最终在复杂域中找到解的右分支。 亚音速和超音速积分路径都绘制在特征£平面上，并由代码显示，如图所示。 2用于孤立的翼型。 椭圆上的阴影标记指定了边界条件用于插值的网格点。 椭圆内显示的多边形曲线是用于通过复特征方法解决特征初始值问题的积分路径。 它们在£ a = £j处加入偶极子到初始点£c，从那里进一步的段向椭圆周围的几个弧线前进。 超音速路径沿着声波位点终止，声波位点是一条从椭圆边缘切割出一条碎片的曲线。 翼型的前缘和后缘的图像由图中左侧和右侧的大十字表示。 6.7 坐标的傅里叶分析 对于跨声速流动，复杂特征方法的一个缺点是它在声波位点M 2 = 1处引入奇点。 跨声速和超音速路径的构造是为了避开这些点，这些点是对角线平面£ = fj中相应的解析曲线M 2 = 1相互反射。 如果声波点出现在必须计算解决方案的点上，则在计算轮廓的可靠坐标时会遇到问题。 我们引入了一个特殊的过滤器，通过利用边界处解决方案的复杂分析性，消除了x、y坐标中的此类错误数据。 首先，让我们重新缩放轮廓周围的弧长5，使其像环中的角度1/5 < | (\textbackslash\{\} < 5，因为x和y在那里是规则的。 我们选择这个类似角度的变量ui，因此s = s(w)由整数a = Ki (1+e- UJ) 2\{L> + efdw + K 2, (6.21)定义

跨声速流动 119的复杂分析 图4 在筛选剖面坐标之前，在粗网上设计的无冲击翼型的图。 错误点的影响是显而易见的。其中e是小于0.2的输入参数，K\textbackslash\{\}和Ki是调整的常数，将w放在1和-1之间。 该公式已被安排，使x和y的傅里叶级数快速收敛。 回顾6是流动向量的角度，我们有dx.dy ds. ds., ds,.„. - +i-r = -cos9 + i-sm8 = - exp(i0)。 (6.22) dur du> dui dui 因为速度分量u和v是由FLOW代码可靠地计算的，我们能够获得函数dx的傅里叶级数。dy v\textasciitilde{}V i i ■ \textbackslash\{\} - + i- = 2\_\textasciicircum{}(a n + ib n) exp [mnuj)

跨声速流动 121的复杂分析图5图5图所示的翼型的垂直重新缩放坐标。 4. 小圆圈表示未经过滤的坐标，实线表示最终的、更准确的轮廓。 下面的点是没有垂直重新缩放的过滤坐标，这显示了无冲击翼型的真实形状。引入使我们能够从相对原始的FLOW代码运行中获得无冲击翼型的更好坐标，其设计属性通过BGKM分析得到确认。 FLOW代码产生了图中所示的无冲击翼型。 4 在过滤过程之前计算的坐标被认为相当不令人满意。 输出的上表面在超音速带的后面附近并不光滑。 这是因为其中一个超音速路径末端的网格点落在声波位点上，变得错误。 图。 5显示了翼型的垂直拉伸图，其中我们不仅看到了坏点，还看到了更远的点，振荡不均匀。 放大图片中的小圆圈代表过滤前的数据输出，但实体曲线的点和下方图片中的点代表更多

122 CHEN \& GARABEDIAN 图6 图中无冲击翼型上的压力分布。 4和5由BGKM代码计算。 霍多图法和欧拉计算非常一致，有助于确定过滤溶液是真正无冲击的。由过滤过程计算出的精确坐标。 一旦过滤器被包括在内，我们就会通过运行BGKM代码获得一个平滑的设计。 BGKM结果如图所示。 6，人们可以看到翼型是无冲击的，由hodograph法和Euler 方程计算的压力分布一致。 McGrattan [8]发现了一种对称的无冲击翼型，它产生了Euler 方程的非唯一解决方案，但它相当薄。 我们获得了另一个对称的无冲击翼型，其厚度与弦比更大，为10\%，它是由利用复杂特性方法的新的、更强大的FLOW代码设计的。 对于零攻角的对称翼型，我们知道有一个零升力的解决方案。 因为

跨声速流动 123 BGKM代码的复杂分析也产生了一个非零升力的解决方案，这必须是第二个解决方案。 通过对称性，具有负升力的反射流是第三个解决方案。 最初的无震翼型及其hodograph在图中早些时候已经展出了。 1和2。 在零攻角下可以用非平凡的升力找到解决方案的自由流马赫数的范围非常狭窄。 在实际应用中，人们可能会担心非唯一解决方案，因为具有这种设计的翼型可能会表现为不可预测。 然而，这些解决方案发生在形状、马赫数和攻击角度的有限范围内，没有观察到不利的物理后果。 超临界机翼部分现在在许多商用飞机上成功运行。 我们希望基于复杂特征方法的改进的FLOW代码将成为这项被普遍接受的技术的一部分。 我们感谢Frances Bauer、Margaret Bledsoe和Mark McConnell的贡献。 这项工作得到了国家科学基金会DMS-9420499赠款的支持。 参考资料1。 F。 鲍尔，P。 加拉贝迪安和D。 Korn，Supercritical Wing Sections，Spring-Verlag，纽约，1972年。 2. F。 鲍尔，P。 加拉贝迪安和D。 Korn，Supercritical Wing Sections II，Spring-Verlag，纽约，1975年。 3. F。 鲍尔，P。 加拉贝迪安和D。 Korn，Supercritical Wing Sections III，Spring-Verlag，纽约，1977年。 4. M。 Bledsoe和P。 Garabedian，跨声速翼型设计的复杂特性方法，应用于压缩机，计算跨声速的进步4，由W编辑。 G。 哈巴希，松树岭出版社，斯旺西，1985年，pp. 111-129。 5. P。 Garabedian，偏微分方程，切尔西，纽约，1986年。 6. P。 加拉贝迪安和G。 McFadden，超临界扫翼的设计，AIAA杂志30，1982年，pp. 289-291。 7. A. 詹姆森，W 施密特和E。 Turkel，有限体积法使用Runge-Kutta时间步进方案对Euler 方程的数值解，AIAA论文81-1259，1981年6月。 8. K。 McGrattan，跨声速流动模型的比较，AIAA杂志30，1992年，pp. 2340-2342。

124 CHEN \& GARABEDIAN 9. E。 穆尔曼和J。 科尔，平面稳态跨声速流的计算，AIAA杂志9，1971年，pp. 114-121。 10. E。 Swenson，跨声速流动中复杂特征的几何，Comm。 纯Appl。 数学。 21，1968年，pp. 175-185。

\clearpage
\chapter{跨声速小横向扰动方程及其计算}
超音速小横向扰动方程及其计算 罗世军 1，沈慧丽 2 和刘平 3 7.1 导言 小扰动近似在亚音速和超音速空气动力学中被广泛接受，用于计算复杂配置，因为它在计算效率方面具有优越性。 在超音速空气动力学中，小扰动近似也是可行的。 经典跨声速小扰动（TSP）方程是由冯·卡曼制定的[6]。 Murman和Krupp[9]进行了修改，以提高临界压力表达式的准确性。 另一种允许大纵向扰动速度分量的修改已被用于[7]中计算机翼-机身组合。 它消除了驻点处TSP方程的奇点。 本文由三部分组成。 首先，描述了跨声速小横向扰动（TSTP）方程的一般属性。 在第二种中，Murman-Cole混合差分方案[8]用于计算入口周围的跨声速流。 Steinhoff和Jameson[12]为中国西安西北理工大学飞机工程系的全部潜力（FP）1发现了跨声速流动在翼型周围的多个解决方案。 2 中国西安西北理工大学航空发动机工程系。 3 工程发展中心，成都飞机工业公司，中国成都。 Frontiers of 计算流体动力学 - 1998 编辑：David A. Caughey和Mohamed M. 哈菲兹 ©1998 世界科学

126 LUO、SHEN和LIU方程和Chen[2]的TSP方程。 在第三部分，研究了跨声速流动在翼型和机翼上的TSTP方程的多重解。 7.2 跨音小横向扰动方程 通过小横向扰动的假设，全势方程简化为\{l-M2)<f>xx + 4>yy + 4>zz = Q (7.1) 其中1 \_ M2 - da + DM 2 **- - \textasciicircum{}M2 (\textasciicircum{}) 2 1 - M 2 = °° l7+J "'-: 2 f<\textasciitilde{}> (7.2) x、y和z是笛卡尔坐标，x平行于身体纵轴或翼的弦，y与翼跨平行。 M x和<？ oo是自由流马赫数和速度。cf>是扰动速度势。 下标x、y、z分别表示x、y、z的偏导数。 7是比热的比例。 M具有本地马赫数的含义。 TSTP流的压力系数是Cp \textasciitilde{} iMl As M = 1，来自Eq.（ 7.2) \textasciicircum{} + A = i-*C (74) 9oo + 2<\& ( 7 + l)Afi [ ' 替换为Eq.( 7.3），在小横向扰动和弱激波的假设下获得精确的临界压力系数，兰金-胡戈尼奥特休克关系变成l-M\textasciicircum{}-(7 + l)Mi4±\textasciicircum{}-\textasciicircum{}Mj(4±\textasciicircum{} \_ 2 l-(7-l)\textasciicircum{}\textasciicircum{}-Y\textasciicircum{}(\textasciicircum{}) 3 \textasciicircum{} \textasciicircum{} +(v[ - v 2)2 + (u>i - w' 2)2 = 0 (7.6)

小横向扰动方程127图1冲击极点在M x = 1.2，其中（u'\textasciicircum{}v\textasciicircum{}wj和（u'2,v'2,w'2）是激波上游和下游侧的扰动速度分量。 事实上，Eq.（ 7.6）也是TSTP微分方程（7.1）的差分方程。 对于自由流中的激波，u[ = v[ = w[ = 0，Eq.( 7.6）减少到[l-Ml-\{1+l)Ml\textasciicircum{}-\textasciicircum{}M2(\textasciicircum{}-)2\textbackslash\{\}ui L zqoo 2 \textbackslash\{\}2q 00/ J +[i- (\textasciicircum{}- 1)M -i- 7TiM-(£) 2]\textasciicircum{} 2+\textasciicircum{}= 0 \textasciicircum{} 这是M„处的TSTP冲击极点。 TSP流和FP流在M x的相应冲击极性分别是[1 - Mi - (7 + 1)\textasciicircum{}4, \textasciicircum{}] "2 2 + v? + w' 22 = 0（7.8）和=0（7.9）等式。（ 7.9）是Af„的确切冲击极点。 等式。（ 7.7）-（7.9）如图所示。 Moo = 1.2的1，其中a*是临界音速。 7.3 轴对称入口的计算 入口周围的外部和内部流动由TSTP和FP方程计算。 只考虑轴对称流。

128 LUO、SHEN和LIU 7.3.1 控制方程和边界条件 在圆柱坐标x、r和6中，轴对称流动的TSTP和FP方程是2 / (a2 - u2)4>xx + (a2 - fj,v2)4>rr - 2fiuv<jjxr -\textbackslash\{\} = 0 (7-10)，其中a是局部音速，a2 = al + \textasciicircum{}l\{ql-u2-\textasciicircum{}) (7.11) aoo是自由流中的声速，u和v分别是x和r方向的速度分量，/J, = 0和1分别给出了TSTP和FP方程。 相同的边界条件用于TSTP和FP方程。 罩的外表面和内表面以及中心体表面的边界条件，r = f(x)是精确的。<£r(*, /(*)) = [«« + *«(*, /(*))]/' (*) (7.12) 如果罩唇是钝的，上述边界条件在钝唇点被4x = -9oo ( 7-13)取代。 入口轴上的控制方程，r = 0，中心体外，被（a 2 - u 2）cj>xx + 2a 24>rT = 0（7.14）取代，在那里（j>r = 0（7.15）入口出口的平均速度应用于入口圆柱形延伸的远下游端，作为下游内部边界条件。 自由流条件<j> = 0应用于外部流的外部边界，该边界距离入口足够远。 在本文中，对于M\textasciicircum{} < 1，上游和下游边界距离入口5D（D = 入口直径），侧向边界为r = 25Z；对于M\textasciicircum{} > 1，侧向和下游边界的距离延伸到500D，这对于数值稳定性是必要的，上游边界距离入口保持5D。

129 小横向扰动方程图2 在钝唇周围生成非均匀网格 7.3.2 差分方程 Murman-Cole非保守混合差分方案[8]用于TSTP和FP方程。 首先求解TSTP方程，其解被用作求解FP方程的起始域。 通过这种技术，非旋转混合差分方案对FP方程是稳定的。 为了方便计算，在罩子的钝唇周围生成的网格如下。 对于给定的Ax序列，绘制了一系列垂直于x轴的网格线。 相应的Ar序列由网格线和罩子午线的交点决定（图。 2）。 生成的网格是不均匀的。 在罩唇附近，网格拉伸不应过度，例如1/2 < Ax i-1J/AxlJ < 2，其中i和j分别是x和r方向的网格点数。 如果网格拉伸过度，松弛计算可能无法收敛。 网格在罩唇附近很好，在外边界附近很粗糙。 7.3.3 松弛法 连续松弛线（SLOR）方法用于求解差分方程。 不受干扰的流动，<j> = =是外部场的松弛计算，但对内部场不好。 对于内部场，一维流动解决方案用作起始场。 对于带有收敛扩散器的皮托式入口，发现唯一的r方向松弛可能不稳定，x和r交替方向松弛稳定了计算并加速了收敛。 R方向的松弛使入口的侧面的扰动迅速传播到流场中，而x方向的松弛使入口出口的扰动在x方向上迅速传播。

130 LUO, SHEN, \& LIU 对于中心体的入口，扩散器有收敛和发散部分。 一些通过出口平面的x方向松弛线被身体的侧面切割，因此rc方向松弛无效。 相反，使用区域多重松弛技术，其中扩散器内部的r方向扫描比外部多，扩散器的发散部分的r方向松弛扫描比收敛部分更多。 为了确保收敛的放松，在放松的早期使用下放松。 cf>的收敛标准是，在所有网格点的连续迭代中，4>的最大相对增量不大于10-3。 收敛还需要入口所有横截面之间最大1\%的质量流速相对差异。 多网格技术用于加速收敛并评估数值解的准确性。 7.3.4 数值示例 计算了皮托类型入口、IE和IC[10]。 图。 3显示了IE罩子钝唇周围的流动，Moo = 0.7197，出口马赫数，M e = Moo和流量系数，CA = 1。 FP解决方案有一个等熟超音速口袋，然后是一个更大的超音速口袋，以激波结尾。 TSTP方程很好地预测了下游超音速口袋，但错过了上游超音速口袋。 图。 4显示了Moo = 1.14，M e = 0.7197和CA = 0.91的入口IE的弓激波和音速线。 与FP解决方案相比，TSTP弓激波和声线在对称轴附近很好。 图。 5显示了压力系数在入口IC罩的外部和内部表面的分布，在M x = 1.14，M e = 0.8215和C A = 0.98。 TSTP溶液与FP溶液一致，但外表面钝唇附近除外。 与现有实验数据[10]和尤拉解[ll]的比较如图所示。 4和5。 对于上述计算，使用了41 x 41和70 x 80两个网格。 发现粗网格提供了足够的准确性。 计算了中心体60-40（13）的入口，M\textasciicircum{} = 1.27，M e = 0.5413和CA = 0.655。 图。 6给出了中心体前部和罩子外表面的压力分布，其中大便是自由流中的压力。 TSTP解决方案与FP解决方案和实验数据一致，但罩子边缘附近除外。 中心体鼻子上的潜在解与圆锥流动理论非常吻合[3]。 使用的两个网格是42 x 51和80 x 100，粗网格提供了足够的准确性。 粗网格和细网格上的收敛迭代数分别为300和400。 西门子7760上的相应计算机时间，以10 6个浮点运算的速度工作

小横向扰动方程131图3 绕入口唇的流动IE，M\textasciicircum{} = M e = 0.7197 图4 绕入口的流动IE，Moo = 1-14，M e = 0.7197

132 LUO、SHEN和LIU 图5 入口IC罩上的压力分布，Moo = 1.14，M，= 0.8215每秒，分别为500秒和1000秒。 7.4 翼型和Wing的多重解 通过改变松弛计算的起始场、松弛方法、松弛参数和差分方案，研究了TSTP方程的跨声速流动在翼型和wing的多重解。 7.4.1 边界条件 使用了笛卡尔坐标x、y和z。 控制方程是Eq.（ 7.1）。 将坐标平面xy靠近机翼表面，边界条件用于以小攻角围绕薄机翼的流动，a如下所示。 在机翼的上表面，z = F(x,y) 4>z(x, y, 0+) = [.joo + (f> x(x, y, 0+)]F x(x, y) - q xa (7.16) 类似的条件在下表面。 在钝前缘，使用精确的边界条件。

小横向扰动方程133图6 入口60-40的罩子和中心体的压力分布，M\textasciicircum{} = 1.27，M e =0.5413 机翼尾的库塔条件是<l>z(x,y,0+) =<t> z(x,y,0-)和<j>(x,y,0+) - 4>(x,y,0-) = <j>(x t,y,0+) - <j>(x t,y, 0- (7.17) (7.18)，其中x t是机翼后缘的x坐标。 在远场边界上，<f> = 0，下游边界除外，其中4>* = o。 7.4.2 松弛方法 差分方程与Murman-Cole非守恒格式[8]形成。 对于翼型的流量，差分方程由两种松弛法求解。 它们是SLOR方法和隐式近似因数分解方法AF-2[1]。 对于机翼上的流动，使用平面AF-2松弛法。 松弛平面垂直于翼展，在平面上使用AF-2方法。 在SLOR方法中，引入了人工时间阻尼项[5]，因此过松弛可用于亚音速和超音速点来加速收敛。 在这项调查中，收敛标准是，在所有网格点的连续迭代中，<j>的最大相对增量不大于10 -5。

LUO、SHEN和LIU图7 翼型 NACA 0012的多重解，M，= 0.85，a = 0 7.4.3 翼型 NACA 0012 对于翼型 NACA 0012在M，= 0.85和a = 0的流动，得到了三个不同的数值解，A、B和C。 压力系数、C在翼型表面的分布主要在激波的位置和强度上有所不同，如图所示。 7，其中c是翼型的和弦。 如表1所示，这些多重解是通过具有固定松弛参数的松弛计算的不同起始场获得的。 计算网格是45 x 39。 在表1中，起始字段的下标A和B分别表示弱和强激波解，对于AF-2，松弛参数Q = 0.8，加速度收敛参数的最低和最高值，UL = 0.5和UH = 65，以及圆形o序列中的数字，K = 5。 上述多重解也可以通过AF-2方法的不同松弛参数获得，起始域相同，如表2所示。 在这里，起始字段为C\$ = 0，对于AF-2，UL = 0.5，UH = 65和h' = 5，fl各不相同。

135小横向扰动方程表1不同起始场启动场松弛方法的多重解解决方案\textasciicircum{}0 AF-2或SLOR A A </>|M 00=.9，Q=O，A AF-2 B 0|M ro =.875，g=O，g AF-2或SLOR C表2不同松弛参数的多重解松弛参数，Q，解.795，.8，.805 A.6，.7，.75，.79，.81，.825，.85 B 数值实验表明，上述多重解也可以通过改变AF-2方法的其他参数获得。 广泛的数值实验，包括使用精炼网格，89 x 39和111 x 39，证实了这些解决方案的持久性。 本调查发现，对于a = 0的翼型 NACA 0012，多个解出现在区域中，0.825 < M x < 0.900。 表3列出了多个波阻系数，CD w0。 可以看出，CDWO的差异是显著的。 Murman-Cole保守差分方案也为TSTP方程提供了多个解。 表3零升的多重波阻系数\textasciicircum{}foo C Dw0由TSTP 0.840 0.0274，0.0292，0.0320 0.850 0.0325，0.0362，0.0407 0.875 0.0480，0.0540 0.900 0.0673，0.0690，0.0710

LUO、SHEN和LIU图8 Engquist-Osher方案的多重解，NACA 0012，Mm = 0.85，cu = 0 7.4.4 Engquist-Osher方案与Murman-Cole差分方案相比，Engquist-Osher差分方案[4]满足熵，而且是保守的。 研究Engquist-Osher方案是否产生了多种解决方案会很有趣。 在这项调查中，TSTP方程被建立Engquist-Osher方案的TSP方程所取代。 对于M, = 0.85和CY = 0的翼型 NACA 0012，Engquist-Osher方案通过改变放松计算的起始场来进行yieId多个解，如表4所示。 一个解决方案是对称的，另一个解决方案是不对称的。 这里的网格是45 x 39，放松方法是AF-2，52 = 0.8，\textasciitilde{}\textasciitilde{} = 0.5，\textasciitilde{}\textasciitilde{} = 50和K = 5。 相应的压力分布如图所示。 8. 它们主要在激波的位置和强度上有所不同。

小横向扰动方程137表4 Engquist-Osher方案的多重解启动场解0 = 0对称解<t> = </,|M00=.85,a=i»,A不对称解表5机翼起始场解的多重解4> = 0对称解A 4> = 0|Moo=.87,a=o对称解B 4> = </,|Moo=.85,a=20,/i不对称解7.4.5三维机翼为了调查跨声速流动的TSTP方程在三维机翼上是否存在，考虑了长宽高比的未扭曲矩形翼，AR=4和翼型 NACA 0012在a = 0处。 采用Murman-Cole非保守方案和平面AF-2松弛法。 数值实验表明，0.84 < M\textasciicircum{} < 0.90区域确实存在多个解。 例如，表5和图。 9给出了Moo = 0.85和a = 0的矩形翼的结果。 在图中。 9，SP是机翼的半跨度。 有四种不同的解：两个对称解，A和B，以及两个不对称解。 两个不对称的解决方案是彼此的镜像。 发现不对称溶液的上下表面压力系数与两个对称溶液的上下表面压力系数一致，多个溶液主要在激波的位置和强度上有所不同。 使用的网格是45 x 39 x 16。 对于AF-2，ft = 0.8，a h = 0.5，a H = 65和K = 5。 7.5 结论跨声速小横向扰动方程适用于纵向扰动速度分量不小的流动。 它

138 LUO、SHEN和LIU图9未扭曲矩形翼的多重解，长宽比4，翼型 NACA 0012，M\textasciicircum{} = 0.85，a = 0得出精确的临界压力系数和近似的激波关系。 在入口周围内流和外流的连续线过松弛计算中，纵向和横向交替方向松弛和区域多重松弛技术有助于稳定计算并加速收敛。 非均匀网格的数值经验表明，在大梯度区域中，超过2的网格拉伸比可能会引起振荡。 离散小横向扰动势流方程的多个解不仅适用于翼型上的稳定二维流动，也适用于机翼上的稳定三维流动。 广泛的数值实验，包括不同的松弛起始场、不同的松弛方法和参数以及网格加密，证实了这些解决方案的持久性。 多重解主要在激波的位置和强度上有所不同。Murman-Cole保守和非保守方案以及Engquist-Osher方案的三种差异方案都产生了多重解。 因此，似乎多重解是势流微分方程所内在的，而不是由于数值方法。

小横向扰动方程139参考1。 Ballhaus，W.F.，Jameson，A.和Albert，J.，稳定跨声速流动问题的隐式近似因子化方案，AIAA杂志16，1978年6月，pp. 573-579。 2. Chen，T.M.，二维稳态跨声速流动有限差分计算中的起始场问题，Ada Aerodynamica Sinica 1，No. 3，1982年，pp. 67-73。 3. Dailey，C.L.和Wood，F.C.，可压缩流体问题的计算曲线，John Wiley \& Sons，纽约，1949年。 4. Engquist，B.和Osher，S.，跨声速流动计算的稳定和熵满足近似，数学计算34，1980年，pp. 45-75。 5. 詹姆森，A.，跨声速流动在机翼和机翼上的迭代解，包括马赫1的流动，关于纯数学和应用数学的通信27，1974年，pp. 283-309。 6. von Karman，Th.，跨声速流动的相似性定律，数学物理学杂志26，1947年，pp. 182-190。 7. Luo，S.J.，Zheng，Y.W.，Qian，H.和Wang，D.Q.，跨声速稳势流的有限差分计算，应用力学和工程中的计算机方法21，1981年7月，pp. 129-138。 8. Murman，E.M.和Cole，J.D.，平面稳变超音速流的计算，AIAA杂志9，1月。 1971年，pp. 114-121。 9. Murman，E.M.和Krupp，J.A.，使用混合有限差分系统解决跨声速势流方程，Proc。 第二届。 关于Num的Conf. 流体动力学中的冰毒，物理学讲义8，1971年。 10. Olstad，W.B.，跨声速风隧道研究唇钝和形状对旋转体中正常冲击鼻入口阻力和压力恢复的影响，N AC A RM L56C28，1956年。 11. Rizzi，A.W.和Schmidt，W.，使用有限体积方法研究皮托型超音速入口流场，AIAA论文78-1115，1978年。 12. Steinhoff，J.和Jameson，A.，Transonic 势流方程的多重解，AIAA杂志20日，11月。 1982年，pp. 1521-1533。 13. WooUett，R.R.，Meleason，E.T.和Choby，D.A.，跨声速非设计阻力和三个混合压缩轴对称入口的性能，NASA TM X-3215，1975年。

\clearpage
\chapter{边界层流动中绝对不稳定扰动的激发}
边界层流动中绝对不稳定的扰动激发Oleg S. Ryzhov1和E。 D。 Terent'ev 2 摘要 提供了回顾，以显示当前不稳定概念与边界层在扫翼上过渡开始的实验结果之间的不一致。 事实证明，三层甲板理论的扩展版本是解决该问题的适当手段，假设雷诺数具有足够大的值。 上游前进的波包会导致绝对不稳定，任何具有交叉流的三维边界层都提前分解。 理论论证将计算结果置于坚实的基础上。 8.1 导言 如果我们仅限于在经过仔细控制的条件下进行的风洞测试，并具有相当弱的人为干扰，那么不同类型的边界层过渡到湍流的基本属性可以总结如下。 平板上二维边界层中TS波放大的线性阶段是温和的，在周期性振动带的下游延伸了几百个波长。 1数学科学系，Rensselaer理工学院，特洛伊，纽约12180-3590。 2 计算中心，俄罗斯科学院，瓦维洛夫街40号，117333莫斯科，俄罗斯联邦。 Frontiers of 计算流体动力学 1998 编辑：David A. Caughey和Mohamed M. 哈菲兹 ©1998 世界科学

142 RYZHOV \& TERENT'EV TS级缩短到几十分之一的波长，外部源在脉冲模式下运行。 一般来说，在布拉修斯流中引入人工干扰会极大地改变在很短的距离内突然发生的过渡位置。 在过渡开始时，强烈的非线性效应几乎与速度场中的三维性进化同时发挥作用。 整个过程是由对流不稳定的扰动驱动的，这些扰动在它们被激发的区域下游席卷。 相反，在扫过的平板（顶部有一个位移体）上的三维边界层的过渡以强烈的非线性相互作用为主，几乎从扰动增长到实验可检测的大小的弦形位置。 根据测试环境的不同，它们可以以静止的顺流涡旋和共存的行波的形式产生。 这些模式的前者源于三维边界层中固有的横流。 后一种模式被认为受二维边界层典型的TS类型机制的支配。 此外，横扫翼的流动容易受到前缘不稳定和污染以及离心不稳定的影响，这超出了本讨论的范围。 大多数实验者都认为，流体动力学稳定性理论可以很好地预测移动波的频率和相位速度以及静止横流涡的波长。 然而，他们无法就组速度和振幅放大率达成良好的共识（例如，见Bippes、Mtiller和Wagner[1]；Kachanov[3]；Reed、Saric和Arnal[9]）。 类似的开放式问题的特点是旋转盘上密切相关的边界层的过渡工作（Wilkinson \& Malik [13]；Corke \& Knasiak [2]）。 然而，大多数可用的实验数据都有一个共同点。 无论使用扰动源来产生脉动，临界雷诺数在过渡开始时始终只显示513的平均值周围的小散射。 换句话说，发现流动的最终状态对初始自然或人为干扰的确切形式不敏感。 这一观察导致Lingwood [5]认为，向湍流的过渡是由于旋转盘边界层内在的绝对不稳定性，而不是对流性。 在随后的论文中，Lingwood [6]通过直接测量证实了她的理论发现。 下面列出的分析旨在提供一个严格的证明，证明在任何具有交叉流动的三维边界层中出现绝对不稳定的振荡的可行性，前提是雷诺数足够大。 这一假设使得在扩展三层方案的框架中利用渐近方法成为可能。 计算证明了扰动源的下游和上游的过压分布以高度的形式

绝对不稳定扰动的激发143个调制信号。 交叉流涡流被证明会产生额外的特征模式，在整个过渡过程中带来新的物理学。 能够逆流移动的波包是新特征模式触发的绝对不稳定的核心。 8.2 扩展三层模型 基础运动应该是一般稳定的三维边界层，上部的马赫数 M\textasciicircum{} < 1。 三层甲板用于描述自激发TS波和横流涡流形式的下分支不稳定性，其耦合导致上游前进扰动的新特征模式。 我们选择了这种渐近方法中采用的特殊单位系统来缩放和归一自变量和所需函数（Stewartson [12]；Mesister [7]）。 让i是时间，x、y、z指定笛卡尔坐标，其中x轴与局部外部流的方向对齐，y代表法线到墙的距离，z指向局部跨度方向。 相应的速度用u、v、w表示；过压用p表示。 在薄黏性近壁子层中，最初的Navier-Stokes方程简化为一组更简单的Prandtl方程du dv dw dx ay dz du du du dp d 2u dt dx dy dz dx dy 2 dw dw dw dp d 2w dt dx dy dz dy 1 由于正常压力梯度dp/dy = 0，我们有p = p(t, x, z)。 将过压与瞬时位移厚度相关的相互作用定律-A（t、x、z）以ir*r的形式呈现？ M°l - d( \_ £* (8.2)，以便在线性近似中保持柯西问题（Ryzhov \& Terent'ev 1997）。 （8.2）中的小参数e与成正比。 R -1 /8，其中R表示参考雷诺数，跨度动量厚度roo £>(«) = / My2)U 2z0(y2)dy2 Jo

144 RYZHOV和TERENT'EV通过密度分布Ro和交叉流速度U zo表示，在大多数初始边界层中。 在扩展方法的框架内，整个三层计划保持不变（Ryzhov \& Terent'ev [10]）。 因此，在近壁子层的上边缘，我们有« - T xy -► r xA，w - r zy -► T ZA as y -► oo。 （8.3）在这里，T X和T Z表示归一化的皮肤摩擦成分，因此振动带通常用于风洞测试，将人工扰动引入边界层（Saric \& Yeates [11]；Radeztsky，Reibert，Saric \& Takagi [8]）。 为了为这些实验设置提供一个相关的数学模型，我们从同质的初始数据u = w =p = A = Q开始，在t = 0（8.4），并通过f Jsin(tJoi）/(i)cos(moz), <>0 y = Vw = \{ (8.5) \{ o, t < 0，函数/在流坐标的有限区间内有效非零）指定与时间相关的障碍物的形状。 移动表面上的边界条件在y = y w (8.6)时变成u = w = 0，v = -g-，其中y w在（8.5）中定义。 接受性问题（8.4）、（8.5）和（8.6）允许我们通过在横流的局部跨度方向上拉伸的丝带的开关和随后的单色振动来追踪在三维边界层中激发的各种类型的扰动的诞生和发展。 然而，我们分析了在丝带初始脉冲运动中辐射的以下波浪系统。 8.3 线性分析 使用S -► 0，我们得出了一个典型的接受性过程的公式，其中扰动机构的振幅被假定为小。 因此，我们把（u - r xy,v,w- r zy,p,A) = SRe [(T xuc,vc,Tzwc,pc,Ac)eimoz] (8.7)，并简化了Prandtl方程（8.1）以及边界条件（8.3）、（8.6）和初始数据（8.4）。 所需的复杂功能

绝对不稳定扰动的兴奋145 uc、v c、w c、p c、A c通过[fie(w、k、y)、v c(u>、k、y)、W C(UJ、k、j/)、p c(w, k)、Ac(u, k)] = CO /'CO dx / e-< wt+i * I)(u c(i,x,j/)、D c(i,a;,y)、w c(t,x,t/)、p c(i,x)、A c(i,x)]df转化为t中的拉普拉斯积分和x中的拉普拉斯积分。 （8.8）将（8.8）替换为线性化普兰特尔方程组，结果为一组齐次的常微分方程。 为了进一步简化分析，让我们定义一个减少的波数K = kr x + m QTz，一个新的自变量Y = 0，+ i 1\textasciicircum{}K1liy,9, = i' 2l3wK-2l3和一个函数F = kT xuc + moT twc（Ryzhov \& Terent'ev [10]）。 它满足第一导数dF/dY的Airy方程d3F „dF, s dYs--YdY=° < 8"9>。 边界条件 F = -tf-J\textasciicircum{}/, f£j = iV*K-W (fc« + m 2) p c at F = 0 (8.10) for (8.9)可以从（8.5)的线性化版本（8.6）和确定F的不均匀二阶方程中推导出来，为了简洁而省略。 这里f(k)是振动器形状f\{x)的傅立叶变换。 限制条件F\textasciicircum{}KAC asF\textasciicircum{}-oo（8.11）来自（8.3）。 相互作用定律（8.2）变成了一个简单的关系+ imlA c（8.12）[(l-M\textasciicircum{})P+mg] 1/2，通过§\textasciitilde{}R 1 /8根据雷诺数以A表示p c。 受（8.10）、（8.11）、（8.12）的解读作j \textasciicircum{}/»(fc»+mg)fc»[(l-Mi)fc' + mg]- 1/a +gm§ f Y AimdY \textbackslash\{\} P + K*l*Hp) *(n)-Q\{fc,mo;JI/oo,e;r„r.) J Q \{ ' f 其结果是，（8.7）中定义的复杂压力p c由以下逆拉普拉斯-傅里叶变换Pc = \textasciicircum{}f dke\textasciicircum{}Kk）\{k 2 [(1 - Ml) *" + ml] " 1/2 + eml) x 将（8.8）替换为线性化Prandtl方程结果系统

146 RYZHOV \& TERENT'EV i-e+ico x <koe ut\$\{Q.)\{u2 +tjZ)- 1\{\$(Q)-Q\{k,mo;M00,i;Tx,Tz)\}-1。 （8.13）Je-ioo Here Ai（Y）指定Airy函数，\$（ft）和/（ft）通过其一导数和不当积分*（fi）= \textasciicircum{}jP-m\}-1，W）= J\textasciitilde{}MY）dY（8.14）表示，数量Q（k，mo;M 00,i;Tx，Tz）代表Q = \{\textasciicircum{}（k 2 + ml）K-“3 [k 2 [（1 - Ml）k 2 + ml] ” 1/2 + emg\}（8.15）只要w-频谱中缺少连续部分，我们可能会将（8.13）中的逆拉普拉斯变换扩展为其极点整数的余数级数级数。 如果我们把自己局限于时间振荡p co和快速增长的不稳定扰动p c\textbackslash\{\}，逆拉普拉斯变换的系列只能减少为三个项，导致一个显式表示p c c = p co -\textbackslash\{\}--p c\textbackslash\{\}，其中Pco -- -£ /\textasciicircum{} dke\textasciicircum{}f(k) \{k 2 [(1 - Ml) k 2 + ml] \textasciitilde{} 1/2 + em 2\} > l \$(-Oo) iw nt figo) *(-n0)-Q(ifc,mo) *(no)-Q(*,»»o) Pel =-\textasciicircum{}/\_° °O \textasciicircum{}e\textasciicircum{}/(fc)\{A; 2 [(l-Mi)\textasciicircum{}+mg]- 1/2 + mg\} x [\textasciicircum{}g\textasciicircum{}O] f ° r Suffidently \textasciicircum{} *• (8-17) Herefto = i 1/3u0K-2/3, §(Tt) = Q(k,m 0;Moo,e;Tx,Tz)。 （8.18）在表示Q的参数（8.16,17）时，为了简洁起见，省略了四个参数Moo、£；TX、TZ。 8.4 计算结果 让我们把MQO = 0.2来处理典型的低马赫数制度，并选择S = 0.01，r x = 2v2/3，r 2 = 1/3作为一般三维边界层的特征值。 （8.5）中的丝带由u> 0 =3和/ = 7r \_1 /2 exp(-a: 2 )固定，其结果是/ = exp(-/fc 2/4)。 我们

绝对不稳定扰动的激发147图1带状诱导的压力分布与mo = 4的下游扫荡波包中的横流坐标的压力分布。下面专注于在初始脉冲模式打开外部源期间激发的指数级增强扰动，并抛开，出于这个原因，从（8.16）开始的时间期间的振荡。 为了将蜡包与丝带单色运动的影响隔离，它们以不会太大的瞬时t = 10显示。 图1展示了跨度波数值mo = 4的复杂压力分布的实值和虚部分。 发现一个高度调制的波包以群速度Vg* = 4.2向下游扫荡，位于距离x = 40 t 44的两个中心周期中振荡的摆动原来是带状振幅的0.5倍4倍。 这种类型的扰动放大是对流不稳定性的核心，在二维剪切流中产生了传统的过渡路径。 然而，波包的快速增长使它们与发展速度慢得多的周期性波列车不同，这是风洞测试中利用的规则。 与流行的看法相反，依靠对二维边界层的长期观察经验，计算揭示了一个完全不同的波包，它正在丝带前向迎面而来的水流前进。 群速度V；可以大致估计在-0.22+-0.28范围内。 调制信号的压力变化如图所示。 2. 在复杂压力的实值和虚值部分，中央循环中的脉动摆动波动都高达0.5。 与下游扫荡波包中的振荡大小相比（图。 1）这个幅度几乎要小4倍。 因此，需要放大比例，以使上游移动信号完全可识别。 尽管其幅度要小得多，但波包传播着

148 RYZHOV \& TERENT'EV 图2 带状诱导压力分布相对于mo = 4的上游前进波包中的溪流坐标。迎面而来的溪流在确定初始层流分解的位置方面可以发挥极其重要的作用。 根据这一猜想，在风洞测试中发现扫翼上三维边界层的非线性相互作用几乎从扰动增长到实验可检测大小的弦位置变得占主导地位（Bippes等人。 [I]；Radeztsky等人。 [S]；Reed等人。 [9]）。 图3描绘了跨度波数的更大值mo = 9的振荡模式。 在这里，扰动以单波包的形式出现，该包以相同的组速度向下游移动，就像之前的情况一样，mo = 4。 然而，高调制信号尾端的长波长周期穿透了带状的上游区域。 图中的上游前进波小集。 2和图中讨论的长波长周期。 3在其发展过程中呈指数级增长，触发了任意类型的三维边界层的绝对不稳定性，并在分解过程中带来了新的物理学。 横流在推动这条以前未知的过渡路径方面起着决定性作用，这显然归因于横流，横流是横扫机翼上黏性流体运动的组成部分。 8.5 理论论证 上面计算的结果可以被视为导言中所述实验结果的附加数据。 在这方面，出现了一个问题，即是否有可能使用相同的渐近方案将绝对不稳定的概念放在坚实的基础上。

绝对不稳定扰动的兴奋Z I图3大值mo = 9的全带状诱导压力分布与溪流坐标。 为了解决这个问题，让我们用函数@（a）和右侧Q（k，mo；M，E”；T，T，T，分别定义为（8.14）中，分析色散关系（8.18）的属性。 我们把自己限制在第一根Ql = Rl (k, mo; M,, E"; T,, 7;) 引起不稳定的扰动。 图中显示了与该根相关的复杂频率平面中第一个色散曲线的行为。 4为mo=4。 与彻底研究的mo = 0的二维振荡的情况不同，曲线分崩离析成位于上半部和下半叶的单独分支。 反过来，每个分支由两个裂片组成。 两个左叶在无限远近处接近虚轴，两个右叶的渐近线距离该轴（4/2）\textasciitilde{},3'\textasciitilde{}(1- M;)'/\textasciitilde{}。 请注意，在下分支上形成一个向实轴方向延伸的细闭环，以顺利地连接其左右叶。 Re(wl)的全局和局部最大值，中间有一个局部最小值，是下分支右叶的突出特征。 一个更小的Re（w1）正丘来自远离虚轴的薄环的尖端，这个丘属于左叶。 从理论上讲，图中所示的波小集的激发。 1和2与复数频率的实数部分所达到的最大值直接相关（例如，见Landau \& Lifshitz [4]）。 我们还强调，dIm(wl)/dk在右叶上< 0，而两个分支的左叶上相同的导数dIm(wl)/dk > 0构成了第一个分配曲线。 下分支右叶上Re(wl)的全球最大值位于丝带下游的波包生成的中心。 事实证明，来自相邻局部最大值的贡献不足以诱导一个单独的清晰切割子包（在使用一个时，它变得可辨别

绝对不稳定的扰动151的激发，导致能够向迎面而来的水流前进的弱波包的激发。 图中波包理论上预测的群速度V d* = -dlm(ui)/dk。 2大约是-2.5；正如我们所看到的，计算给出了一个以相同值为中心的更广泛的估计值。 此外，波包内振荡周期的理论参考长度也与计算结果没有明显偏差。 当跨度波数上升到值mo = 9时，复杂频率平面中第一个色散曲线的行为会发生如图所示的显著变化。 5. 下分支右叶上的Re(u>i)的全局最大值和局部最大值合并在一起，形成单个最大值。 连接右叶和左叶的薄闭环从下分支的形状上消失。 分散曲线的这种修改为图中不可分割的扰动模式提供了解释。 3. 由于唯一最大值的Re(uj\textbackslash\{\})发生在右叶上，而左叶上没有正的Re(u>i)，图中的波包。 3以组速度V* = -dlm(u>i)/dk整体向下游移动，这与4.2的相应计算值非常一致。 扰动的大小减少了一千倍以上，因为跨度波数从mo = 4增加到mo = 9。 波包振幅的急剧下降可以归因于这样一个事实，即由mo固定的第一个色散曲线右叶上Re(u>i)的全局最大值的大小几乎是mo = 9的第一个色散曲线的单个最大值的两倍。 理论参考长度单调地从波系统的前部到其尾部。 预测与图中的计算密切相关。 3，在其中，出现调制信号尾部的振荡周期穿透了色带上游的区域。 它们是由iZe（a>i）沿第一色散曲线下分支和上分支左叶的单调变化的贡献驱动的。 因此，一种替代机制控制着外部源上方和前方的脉动发展。 8.6 结论 上游带发出的扰动类型在很大程度上取决于定义特定黏性交叉流特征模式的跨度波数。 在中等值（mo = 4）下，它们以单独的波包的形式出现，逆向迎面而来的流前进。 大值（mo = 9）在整个下游传播的不可分割信号的尾端产生振荡周期。 在这两种情况下，上游扰动都会随着时间的推移呈指数级增长。 它们的放大导致了三维边界层的绝对不稳定，这带来了新的物理学

152 RYZHOV \& TERENT'EV 图5 大值mo = 9的复频平面中的第一个色散曲线。 薄而封闭的环从下分支消失。过渡到湍流的过程。 致谢 两位作者都非常荣幸地将这份作品作为他们对教授的高度敬意。 Earl1 Murman的科学活动。 作者们还想向教授表示感谢。 朱利安·D。 科尔为许多讨论和有益的评论提供。 由美国空军空军材料司令部空军科学研究办公室赞助的努力，赠款编号为F49620- 97-1-0141。 美国 政府有权出于政府目的复制和分发重印本，尽管上面有任何版权说明。 此处所载的观点和结论是作者的观点和结论，不应被解释为一定代表科学空军的官方政策或认可，无论是明示的还是非正式的

绝对不稳定的干扰的兴奋 153 美国政府的研究。 参考资料1。 Bippes，H.，Mtiller，B. \& Wagner，M.，不稳定三维边界层中扰动生长的测量和稳定性计算，物理。 流体A，流体Dyn。 3，1991年，pp. 2371-2377。 2. 科克，T。 C，Knasiak，K。 F.，具有周期性分布粗糙度的交叉流动不稳定。 程序。 IUTAM糖浆，关于三维边界层的非线性不稳定性和过渡（ed. P。 W。 鸭子和P。 大厅），Kluwer Academic，1996年，pp. 267-282。 3. 卡查诺夫，于。 S.，边界层三维不稳定性的实验研究，AIAA论文96-1978，1996年。 4. 兰道，L。 D。 \& 利夫希茨，E。 M.，流体力学，佩加蒙，1959年。 5. 灵伍德，R。 J.，边界层在旋转盘上的绝对不稳定，J. 流体机械。 299，1995年，pp. 17-33。 6. 灵伍德，R。 J.，旋转盘边界层的绝对不稳定性的实验研究，J. 流体机械。 314，1996年，pp. 373-405。 7. 梅西特，A。 F.，平板后缘附近的边界层流动，SIAM J。 应用程序。 数学。 18，1970年，pp. 241-257。 8. Radeztsky，R. H.，Jr.，Reibert，M. S.，Saric，W. C. \& Takagi，S.，微米大小的粗糙度对扫翼流过渡的影响，AIAA论文93-0076，1993年。 9. 里德，H。 L.，Saric，W. C. \& Arnal，D.，线性稳定性理论应用于边界层，Ann。 牧师。 流体机械。 28，1996年，pp. 389-428。 10. 里佐夫，O。 S。 \& 特伦特耶夫，E. D.，稳定三维亚音速边界层中波运动的复合渐近模型，J。 流体机械。 337，1997年，pp. 103-128。 11. 萨里克，W。 C. \& 叶芝，L。 G.，横流涡旋在扫翼流中稳定性的实验，AIAA论文85-0493，1985年。 12. Stewartson，K.，在平板后缘附近的流动中。 II，Mathematika 16，1969年，pp. 106-121。 13. 威尔金森，S。 P。 \& Malik，M. R.，在旋转盘上流动的稳定性实验，.4L4.4 J。 23，1985年，588-595。

\clearpage
\chapter{在误差分析的伴随方程和CFD中的最优网格自适应上}
在误差分析的伴随方程上，在CFD Michael B中的最优网格自适应上。 Giles 1 9.1 导言 CFD面临的一个挑战是能够给出严格的误差界限，以便工程师知道计算结果的准确性。 抛开评估湍流和过渡预测导致的建模误差大小的难题，这需要准确估计p.d.e.系统离散化导致的误差。 有了这些知识，人们可以希望开发一种严格的方法来优化网格自适应，为给定的计算成本提供最准确的解决方案，或者在实现给定的准确性水平时将计算成本降至最低。 目前，数学理论和工程实践之间仍然存在相当大的差距。 当使用平滑结构网格进行无奇点的平滑流场时，可以从分析截断误差中推断出准确性的顺序。 特定网格大小的误差的绝对大小可以根据过去对测试问题的网格加密研究的经验来估计。 然而，当人们开始使用网格重新分配（运动网格点）来提高流动特征的分辨率时，网格不再光滑，如果使用局部网格加密（添加额外的网格点），网格就会变得非结构化，至少从理论分析的角度来看，如果不是从编程的角度来看。 1 劳斯莱斯阅读器在CFD，牛津大学计算实验室，牛津，英国，电子邮件：giles@comlab.ox.ac.uk 计算流体动力学的前沿 - 1998 编辑：David A. Caughey和Mohamed M. 哈菲兹 ©1998 世界科学

156 GILES 对于非结构化网格，网格加密的当前做法仍然是使用启发式方法。 其中一些是基于有根据的物理推理，即人们需要对冲击、边界层、唤醒和自由剪切层等特征有良好的分辨率。 然而，完全有可能将过多的计算精力投入到这些特征的分辨率上，而牺牲了流场其他部分的分辨率不足，例如翼型前缘的平稳但快速膨胀。 其他方法是基于减少截断误差幅度的想法，但这没有考虑到截断误差引起的解误差的幅度。 直到最近，非结构化网格误差分析的严格数学方法涉及使用Aubin-Nitsche技术来推导模型问题的有限元近似的误差边界，如对流扩散方程[1]。 然而，当应用于双曲p.d.e.的近似时，这通常会导致使用负索博列夫规范（例如[2]）的误差边界，这些规范几乎没有工程意义。 因此，这通常对工程师帮助不大，尽管它导致了实用的网格加密指标[3，4]。 然而，最近，Becker \& Rannacher、Still \& Giles、Paraschivoiu、Patera和Peraire引入了一种对误差分析和最佳网格加密的有前途的新方法。 这始于观察，即在许多情况下，工程兴趣的关键量是解决方案的功能，例如翼型的升力和阻力。 因此，解误差最相关的度量是这些推导量中的绝对误差。 这导致了涉及伴随p.d.e.与不均匀项和适合特定函数的边界条件的数学分析。 由此产生的伴随解定义了函数误差和有限元残差误差之间的关系，即有限元解不是原始分析问题的解的程度。 因此，伴随解的估计值以及局部有限元残差误差可用于定义最佳网格加密策略，以获得给定计算成本的升力和阻力的最准确预测。 这一研究仍处于起步阶段。 根据结构分析中为椭圆p.d.e.开发的理论[5, 6, 7]，Becker和Rannacher为不可压缩的Navier-Stokes 方程[8]开发了事后误差估计。 除了类似的事后误差估计外，Siili、Giles等人还为不可压缩的Navier-Stokes 方程开发了先验误差估计，证明了一个有趣的超收敛属性[9]。 Paraschivoiu、Patera和Peraire还开发了一种略有不同的分析，使用伴随解来获得椭圆p.d.e.的函数的上下限[10]，目的是进入Navier-Stokes 方程。 本文的目的是向误差分析解释这种方法，并展示

伴随方程 FOR 误差分析 157如何用于最优网格自适应。 第一节概述了Euler 方程在结构化和非结构化网格上有限体积离散化的先验误差分析。 误差估计值可用于改进函数的计算值，也可以通过重新分配或细化作为网格自适应的基础。 此外，研究表明，对于非结构化网格，使用保守的离散化确保了函数的精度顺序比有限体积离散化的截断误差的序大一。 第二部分介绍了简单椭圆模型问题的有限元离散化理论，讨论了先验误差分析产生的超收敛属性，以及网格自适应使用后验误差分析。 这为有限元误差分析的文献提供了介绍；为将分析扩展到对流/扩散和不可压缩Navier-Stokes 方程提供了参考。 9.2 有限体积分析 分析从原始流体动力学方程的有限体积近似产生的离散方程开始，Rh(Uh) = 0。 其中U h是离散流解，方程直接来自通量天平，而不是由计算单元的面积或体积归一化。 解误差e h由Uh = U + e h定义，其中U是分析流解。 离散方程的线性化给出了其中R h(U)是将解析解替换为离散算子而获得的截断误差的向量。 这个方程描述了相对容易估计的截断误差和解误差之间的关系，解误差是更大的兴趣量。 如果I（U）是基于分析解感兴趣的标量函数（例如升力或阻力），那么相应的离散近似值Ih（Uh）中的误差可以分为两个部分，Ih（Uh）-I（U）=（I h（Uh）-Ih（U））+（I h（U）-I（U））。

158 GILES 第二个项是近似算子I的截断误差。 第一个项是由于离散解U h的误差，可以近似如下，dlh fdR h dlh dUh Rh\{U) dUh \textbackslash\{\}dU h, = V TRh\{U），其中向量V是伴随流方程的解，T = 0。 （dRky（di h\textbackslash\{\} \textbackslash\{\}dUh）+\textbackslash\{\}dUh）因此，伴随流解将升力和阻力等数量误差与有限体积单元残差评估中的潜在截断误差联系起来。 伴随流动解决方案在最佳设计中的作用现已确立[11、12、13、14、15]。 同样的伴随解决方案在误差分析中发挥着至关重要的作用，这一事实应该不足为奇。 在设计中，人们关注的是由于几何形状变化而导致的函数扰动；在误差分析中，人们关注的是截断误差导致的函数扰动。 9.2.1 结构化网格 对于希望通过网格重新分配来提高准确性的结构化网格，移动网格点以更好地解决大流量变化区域，该分析可用于定义最佳适应策略。 考虑Euler 方程在平滑结构网格上的2D或3D离散化。 泰勒级数展开可用于分析离散化的截断误差。 考虑到残差不是由计算单元的面积或体积归一化的，截断误差为h p+d阶，其中h是单元格的长度（假定具有固定的长宽比以简化分析），p是准确性的顺序，d是问题的维度。 在计算并伴随解并使用流动解估计截断误差中的高阶导数后，函数中的总体误差可以表示为以下形式的总和：£ h*+dTj。 在这里，索引j表示流动变量的单个位置（通常在节点或细胞中心），hj是相关细胞的长度。 Tj是一个scalar

伴随方程 FOR 误差分析 159，涉及来自截断误差估计的伴随解的乘积和流解的高阶导数。 通过使用高阶多项式的局部最小二乘近似，可以从计算流解中获得这些高阶导数值的数值估计。 伴随流计算提供了伴随解，因此可以在每个节点上获得数量T的准确估计。 函数的误差估计值可以用两种方式使用。 第一个是将其用作函数计算值的校正，以伴随流计算为代价提供更准确的值。 如果一个人对多个函数感兴趣，那么每个函数都需要进行伴随流计算，但这在计算上可能仍然比通过使用更精细的网格进行原始流计算来提高准确性更便宜。 第二个选项是通过重新分配使用最优网格自适应的错误估计，运动网格点，以更好地解决流特征。 目标是最小化hj是单元格的面积\{d = 2）或体积（d = 3），这个求和可以用整数Jh"\textbackslash\{\}T\textbackslash\{\}dV近似，其中\textbackslash\{\}T\textbackslash\{\}是\textbackslash\{\}Tj\textbackslash\{\}的连续近似值。 同样，细胞总数大约以f h\textasciitilde{}d dV给出。 使用拉格朗日乘数将第一个积分最小化，同时保持后者固定，导致要求h p+d \textbackslash\{\}T\textbackslash\{\}应该是均匀的。 在实践中，人们通常会关注不止一个功能，例如提升和拖动。 为了处理这个问题，可以修改适应标准，以便策略是确保hP+d Y\textasciicircum{}\textbackslash\{\}T im)\textbackslash\{\}大致均匀，其中T（m）是关注函数的相应误差分量。 原则上，每个T（m）的构造需要伴随方程的解。 然而，在实践中，它可能

160 GILES是，基于对其起源和定性性质的详细理解，这些伴随解可以通过简单的分析函数很好地近似于适应目的[16]。 这种策略似乎与当前旨在使整个网格的截断误差均匀的适应实践没有太大区别。 然而，关键的区别是包含并伴解决方案，这反映了一个事实，即并非所有截断误差对升力和阻力等工程兴趣量的影响都是相等的。 一个很好的例子是翼型后面的截断错误。 由于网格生成难以预测唤醒的轨迹，在翼型的和弦或更下游时，尾音解析不良的情况并不罕见。 然而，尽管由此产生的截断误差可能相对较大，但升力和阻力函数的伴随流动解决方案相对较小，这反映了这些误差不会显著影响翼型附近的流动。 因此，仅基于截断误差估计的适应程序可能会过度解决尾迹区域，而那些包括相邻解决方案影响的适应程序将正确地强调减少靠近翼型的那些对升力和阻力影响最大的细胞中的误差。 9.2.2 非结构化网格 对于非结构化网格，需要进行一些进一步的分析才能获得良好的误差估计。 考虑使用基于节点的变量、基于边缘的数据结构、非线性通量项的标准有限体积离散化（也可以解释为Galerkin有限元离散化）加上特征平滑的添加，对2D或3D Euler 方程的典型离散化。 假设所有单元格都具有有界的长宽比，内部节点j的截断误差可以按以下方式表示为从节点j出来的边缘集Ej的总和。k€E\}在这里hk是边缘的长度，并且有一个隐含的假设，即所有单元格都是有界的长宽比，因此与边缘关联的面具有长度/面积0（/i t \_1）。 Sjk由流动解的导数组成；这个结果来自面上通量的线性表示的标准积分误差。 请注意，由于单元格的面积/体积为0（hj），其中hj是单元格的代表长度尺度，因此由单元格区域归一化的局部截断误差为0（h）。 截断误差的一阶精度导致一些人认为解精度也可能是一阶。 的确，

伴随方程 FOR 误差分析 161 如果使用上述获得的误差表示，VTR\{u) = Y JV\textasciicircum{}Rj，那么由于节点总数为0(h\textasciitilde{}d)，其中h是整个网格中hj的平均值，因此似乎误差为0(h)。 然而，数字证据表明，这种方法是二阶准确的[17]。 如果考虑覆盖面积/体积为0（1）的区域的相邻单元的结合，那么将这些单元相加，内部通量的截断误差因守恒而抵消。 边界有0（h\textasciitilde{}+1）面，每个面的截断误差为0（h d+1），因此该单元格聚合的总体截断误差为0（h 2）。 Giles试图使用截断误差和由此产生的解误差的傅里叶分解来完善这一论点，但分析并不严谨[18]。 要使用伴随方法恢复误差分析中的二阶精度，需要对误差求和进行简单的重新排列。 它至关重要地取决于守恒，因此对于连接节点i和j的边缘k，与沿边缘通量相关的截断误差对于两个单元格是相等且相反的，即 Sik = -Sjk- 因此，所有单元格的误差求和可以重新排列为所有边缘E和所有边界节点B的求和，形式为Sk是沿边缘k的通量的截断误差，AVJt是边缘连接的节点相邻解的差值。 假设分析伴随解是可微的，AVJt应该是0（hk），边缘总数为0（h\textasciitilde{}d），这意味着第一个和是0（h 2）。 边界节点的数量是0（h\textasciitilde{}d+1），因此第二个和也是0（h 2）。 结论是，尽管局部截断误差是一阶，但整数的总体误差，如升力和阻力，是二阶误差。 专注于边缘引起的误差贡献，这现在是k€E的形式，其中Tk是沿边缘的伴随解梯度和来自用于评估通量截断误差的泰勒级数展开的流动解的高阶导数的乘积。

162 GILES 与结构化网格误差分析一样，此误差估计可用于提高函数的计算值。 或者，它可以通过添加额外的网格节点来用于网格自适应，从而减少单元大小hk-通过将额外的节点引入/i fc+2|Tjt|平均大小最大的区域，可以最大限度地减少误差。 随着反复的提炼，这个数量最终应该在网格上变得大致均匀。 改进可以继续进行，直到误差估计小于用户定义的公差，从而实现将给定精度水平的计算成本降至最低的目标。 9.3 有限元分析 为了简化分析的细节，我们将考虑限制在域ft上的2D或3D泊松方程-V\textbackslash\{\} = /，该域是单位平方或立方体，取决于尺寸，并受边界dQ上的狄利克雷条件u = g的约束。 9.3.1 符号和定义 我们定义了两个内积，一个在域上，(u,v) = / uvdV, Jn，另一个在边界上，(f,i>)an - / uvdA。 Jan 定义以下双线性泛函a(u,v) = (Vu, Vt>)也很方便。 Z/2规范、H 1半规范和H 1 Sobolev规范由ll"llL -("'")' Mffi =a(u,u), \textbackslash\{\}\textbackslash\{\}uf Hi =(u,u) + a(u,u)给出。 高度的半规范和索博列夫规范的定义相似。 9.3.2 标准f.e.分析 这个问题的标准误差分析是既定的；有关详细信息，请参阅Strang \& Fix的教科书[1]。

伴随方程 FOR 误差分析 163 函数空间ff 1(f2)由那些函数u组成，其中||«||\#i < oo。 此外，子空间HQ（Q）包含那些在边界dfl上为零的函数，而子空间-ffj（ft）具有满足Dirichlet b.c. u = g的函数。 问题的弱解由函数u € HUO给出。）这样a（u，w）=（/，w），Vui£ Hi（ft）。 现在让S h是if 1（ft）的有限元子空间，由连续函数组成，这些函数在单位正方形（或立方体）的三角测量Th的每个三角形（或四面体）上是线性的，h是任何单个单元的最大直径。 子空间SQ由边界上为零的函式组成。 此外，我们将假设边界数据g是分段线性的，因此也存在一个子空间Si 1，其函数满足Dirichlet 边界条件。 有限元问题的解是函数u h S Si 1，使a(uh, w h) = (f,w h), \textbackslash\{\}/w heS£。 对于任何w h € S\# C H\&(Sl)，a(u,wh) = (f,w h)，因此我们得到正交属性a\{u-uh,wh) = 0, \textbackslash\{\}/w h€S\textasciicircum{}。 现在，存在独立于h的正常数C\textbackslash\{\}，C2，因此a(w,w) > CIHIHI, VtoeiT\textasciicircum{}n), a(v,w) < C 2 H| H1 |M| Hl Vv,weHl(Cl)。 因此，使用正交属性，||w-w A ||\textasciicircum{}, < C\textasciitilde{} la\{u-uh,u-uh) = Ci 1a\{u-uh,u-wh), < C\textasciicircum{}C 2\textbackslash\{\}\textbackslash\{\}u-uh\textbackslash\{\}\textbackslash\{\}H1\textbackslash\{\}\textbackslash\{\}u-wh\textbackslash\{\}\textbackslash\{\}H1，对于任何w h € 5*。 因此，此时w h被选为u的插值。 然后，有关插值精度的标准结果可用于推断||it - if || H1 < C± C2C3 h \textbackslash\{\}u\textbackslash\{\} Hi，

164 GILES，其中C3是独立于h的另一个常数。 这证明了H 1规范中的一阶精度。 要证明L 2规范中的二阶精度，需要使用伴随问题的Aubin-Nitsche技术。 函数v £ HQ（ft）由弱问题a\{v,w）=\{u-u h,w）、Vw £\#£（«）定义，根据标准椭圆规律性结果，存在第四个常数C4，使|t>| H2 <C 4 ||u-V|| l2。 然后，再次使用正交属性，II W \textasciitilde{} U 1L 2 = <*(«-«'» = a(u-u h,v-v-vk)，Vv h £ S h，因此，选择v作为v的插值，h\textasciitilde{}uh\textbackslash\{\}\textbackslash\{\}2L2 < c 2 \textbackslash\{\}\textbackslash\{\}u-u h\textbackslash\{\}\textbackslash\{\}H1\textbackslash\{\}\textbackslash\{\}v-vh\textbackslash\{\}\textbackslash\{\}H1 < C\textasciicircum{}C\textbackslash\{\}C\textbackslash\{\}h 2 \textbackslash\{\}u\textbackslash\{\}H2\textbackslash\{\}v\textbackslash\{\}H> < Cr 2C*CjCt h 2 \textbackslash\{\}u\textbackslash\{\} m\textbackslash\{\}\textbackslash\{\}u\textasciitilde{}uh\textbackslash\{\}\textbackslash\{\}h2 < C\textasciicircum{}Clld h 2 \textbackslash\{\}u\textbackslash\{\} H2.,h\textbackslash\{\} 9.3.3 线性泛函的原/对偶公式 假设现在我们有兴趣获得以下函数的值，该函数的值，该函数是域和边界上的内积的组合，I = (u,d) + (-Tj-,e) 对于任意函数w £ if*(fi)，按部分积分产生(f,w) = a\{u,w) + (-\textasciitilde{},c) 8 n，其中u仍然是原始（原始）问题的弱解。 因此，确定线性泛函值的问题可以用弱形式表示为：解析原始V：给定函式d,e,f,g，找到I(d,e,f,g)和u £ H](U)，使/ + a\{u, w)-\{f,w)-(u,d) = 0, Vu; £ Hi (Q)。 相应的对偶问题的弱公式是：

伴随方程 FOR 误差分析 165 解析对偶V：给定函数d, e,f,g, find I(d,e,f,g)和v € H\textbackslash\{\} (ft)，使7 + a\{w,v) - (/,«) - (w,d) = 0, \textbackslash\{\}/w € ff*(ft)。 这两个问题的等价性，即它们产生相同的线性泛函I(d,e, /,<?)，直接源于在原始问题中考虑w = v，在对偶问题中考虑w = u。 通过用5 efc和S*替换77*（ft）和7f*（ft）来获得原始和对偶问题的有限元近似。 离散原始V h：给定函数d,e,f,g，找到I h(d,e,f,g)和u' 1 € Sj，使7fc + a(u\textbackslash\{\}w ft ) - \{f,w h) - (u\textbackslash\{\}d) = 0, Vw h e Si离散双V h：给定函数d,e,f,g，找到I h(d,e,f,g)和u' 1 g S*，使7h + a(u; fe,z;ft) - (/,v' 1) - (w ft ft,d) = 0，Ww h € 5 ffft。 从这两个有限元问题中获得的线性泛函的等价性再次直接源于在原始问题中考虑w h = v h，在对偶问题中考虑w h = u h。 9.3.4 函数误差表示 让u,v和u h分别是问题V,V和V h的解，让w h是S\textasciicircum{}中的任意函数。 然后，从V、V和V h的定义来，a\{u-uh,v-wh) = a(u,v) - a(u h,v) -a(u,w h) + a(u h,wh) = -(/-(/,«) -(u,g)) + (/-(/,«) -(« A,3)) + (7 -(/,«") -(u,g)) - \{I h-\{f,wh)-\{u\textbackslash\{\}g))) = I-I h。 同样，如果u、v和v h分别是V、V和V h的解，并且u'eSj，那么a(u-w) A,v-v /! ）= 7-7 fc。 因此，我们在线性泛函的有限元近似中有两个不同的误差7-7表示。 9.3.5 先验误差分析 先验误差分析从7-7 fe = a\{u-u h,v-vh)开始。

166 GILES 使用之前给出的标准误差分析的误差边界，使用分段线性函数的有限元空间时的函数误差由|,\_/H| < \textbackslash\{\}\textbackslash\{\}u-u»\textbackslash\{\}\textbackslash\{\} H1\textbackslash\{\}\textbackslash\{\}v-vh\textbackslash\{\}\textbackslash\{\}H1 < Cr 2C22Cih2\textbackslash\{\}u\textbackslash\{\}H2\textbackslash\{\}v\textbackslash\{\}H*限制。 如果使用具有更高阶多项式的有限元空间，使\textbackslash\{\}u - u Hi < C 3hs \textbackslash\{\}u\textbackslash\{\} H'+i，对于一些整数s和常数C 3，那么这个先验误差估计变为\textbackslash\{\}I-IH\textbackslash\{\} <Cr 2C*C23h2*\textbackslash\{\}u\textbackslash\{\}H.+i\textbackslash\{\}v\textbackslash\{\}H.+i。 这表明函数有限元近似的精度顺序是H 1规范中解u h的精度的两倍。 这种超收敛属性是由于解误差u-u h中的领先阶项与线性泛函的评估中的平滑函数d、e正交。 9.3.6 后方误差分析后方误差分析从I-I h = a(u-u h,v-wh)开始，其中w h将是对偶解v的插值。 将积分分割成每个三角形或四面体上的积分之和，然后按部分积分，得出I - I h = ]T T K，K€Th，其中TK= I (f + V2uh)(v-wh)dV + ]- f I JK I JdK L dn (v-wh)dA，其中\textasciicircum{}\textasciitilde{}是正态导数跨越内面的跳跃，在形成边界dfl的一部分的面上被定义为零。 当使用分段线性有限元时，V 2wh在每个单元格内为零，4\textasciicircum{}-很容易被评估。 插值误差v-w h可以从v的二次导数中估计；这些反过来又可以用第二导数来估计

伴随方程 FOR 误差分析 167 双有限元解的差异 v h。 通过这种方式，可以准确估计细胞误差TK的大小。 对于最佳的网格自适应，策略是细化\textbackslash\{\}TKI较大的单元格，直到误差的边界小到可以接受。 在这个过程中，\textbackslash\{\}TK\textbackslash\{\}将在整个网格中变得相对均匀。 9.3.7 扩展 这里提出的分析是一个简单的p.d.e. （泊松方程）在一个简单的域上（可以精确三角化）和简单的边界条件和线性泛函（对应于有限元空间中的函数）。 分析可以扩展到更难的问题。 Becker和Rannacher推导出了不可压缩Navier-Stokes 方程的后验误差分析[8]，Siili、Giles等人为不可压缩Navier-Stokes 方程[9]开发了先验和后验分析。 这些分析假设了简单的域边界和边界数据。 Siili和Giles即将发表的一篇论文将表明，对于对流/扩散方程，平滑弯曲边界和平滑边界数据可以以不破坏泛函的超收敛属性的方式处理；关键是将边界几何和数据适当投影到有限元空间上。 9.4 一些结论性评论 本文概述了使用适当对偶问题的解来估计CFD计算中近似非线性泛函的误差的方法。 误差估计可用于获得对函数本身的更好的近似值，或用于驱动网格自适应，目的是为给定的计算工作水平实现最准确的答案。 有限体积分析表明，在非结构化网格上，离散守恒对于获得相对于局部截断误差的顺序的精确度至关重要。 然而，概述的分析假设二元解的梯度是有界的。 沿着停滞流线[16]的Euler 方程可能不是这样，因此可能需要额外的分析。 先验有限元误差分析揭示了一个有趣的超收敛属性，表明近似线性泛函的精度顺序是解本身的两倍。 有限体积分析缺乏类似的结果可能表明有限元方法具有显著优势，但这种优势只有在

168 GILES使用比二阶精度更好的方法。 致谢 我非常感谢我的同事Endre Siili对本文的评论以及我们关于误差分析的多次讨论。 这项研究得到了英国工程和物理科学研究委员会GR/K91149研究补助金的支持。 参考资料1。 G。 斯特朗和G。 修理。 有限元法的分析。 Prentice-Hall，1973年。 2. E。 Siili和P。 休斯顿。 双曲问题的有限元方法：后验误差分析和适应性。 数值分析的现状。 我。 Duff和A。 Watson，编辑们。 牛津大学出版社，1997年。 3. T。 声纳和E。 西伊利。 可压缩Euler 方程的双图规范细化指标。 出现在Numerische Mathematik中。 4. E。 西伊利。 后验误差分析和双曲问题的自适应有限元近似的全局误差控制。 第16届数值分析两年期会议记录。 D。 格里菲斯，编辑。 皮特曼，1996年。 5. 我。 巴布斯卡和W.C. 莱茵博尔德。 自适应有限元计算的误差估计。 暹西姆J。 数字。 Anal.，15（4）：736-754，1978年。 6. 我。 巴布斯卡和A。 米勒。 有限元法中的后处理方法-第1部分：位移、应力和其他位移的更高导数的计算。 实习生。 J。 数字。 方法Engrg.，20：1085-1109，1984年。 7. R。 维菲尔特。 事后误差估计和自适应网格细化技术。 J。 计算机。 应用程序。 数学，50：67-83，1994年。 8. R。 Becker和R。 兰纳彻。 在有限元方法中加权后验误差控制。 技术报告，海德堡大学，1994年。 预印编号 96-1。 9. M。 B. 吉尔斯，M。 G。 拉森，J。 M。 Levenstam和E。 西伊利。 黏性流动中升力和阻力有限元近似的自适应误差控制。 技术报告NA97/06，牛津大学计算实验室，1997年。 10. J。 佩雷尔，M。 Paraschivoiu和A。 帕特拉。 椭圆偏微分方程的线性泛函输出的后有限元边界。 关于计算力学进步的Sypmosium，提交给Comp。 冰毒。 应用程序。 Engrg.，1996年。 11. A. 詹姆森 通过控制理论进行空气动力学设计。 J。 科学。 计算机，3:233-260，1988年。 12. A. Jameson，N.A. 皮尔斯和L。 马丁内利。 使用Navier-Stokes 方程的最佳空气动力学设计。 AIAA论文97-0101，1997年。 13. 哦。 Baysal和M.E. 摇摇欲坠。 可压缩Euler 方程的空气动力学灵敏度分析方法。 J。 流体。 Engrg.，113：681-688，1991年。 14. J。 艾略特和J。 佩雷尔。 使用非结构化网格进行实用的3D空气动力学设计和最佳化。 AIAA论文96-4122-CP，1996年。 第六届AIAA/NASA/ISSMO多学科分析和优化研讨会会议记录。

伴随方程 FOR 误差分析 169 15。 瓦克 安德森和V。 Venkatakrishnan。 空气动力学设计优化在具有连续伴随公式的非结构网格上。 AIAA论文97-0643，1997年。 16. M.B. Giles和N.A. 皮尔斯。 CFD中的伴随方程：对偶性、边界条件和解决方案行为。 AIAA论文97-1850，1997年。 17. D.R. 林德奎斯特和M.B. 吉尔斯。 三角形和四边形网格的数值方案的比较。 在D.L. 德沃耶，M.Y. Hussaini和R.G. Voigt，编辑，第11届流体力学数值方法国际会议记录，物理学讲义第323卷，第369-373页，柏林，1989年。 斯普林格出版社。 18. M.B. 吉尔斯。 不规则网格上基于节点的解决方案的准确性。 在D.L. 德沃耶，M.Y. Hussaini和R.G. Voigt，编辑，第11届流体力学数值方法国际会议论文集，物理学讲义第323卷，第273-277页。 Springer-Verlag，1989年。 柏林。

\clearpage
\chapter{在流程计算中增加了耗散}
在流计算中添加了耗散 Robert W. MacCormack 1 10.1 介绍 在过去四分之一世纪中，用于求解可压缩黏性流动方程的高效隐式数值方法已被广泛开发和使用。 这些方法首先用矩阵方程近似控制偏微分流方程，然后用近似因数分解用一系列高效可逆矩阵因子替换原始矩阵。 此外，这些方法通常添加人工耗散来控制数值不稳定性。 要求解的原始方程和实际数值求解的方程之间的差值代表了计算误差。 这个错误是通过1引入的。 有限差分商的导数离散化，2.原始矩阵方程的近似因数分解，以及3.人工耗散的加法。 这些方法通常以比时间精度要求的要小得多的时间步数操作，通常收敛需要数千个时间步程才能达到工程精度。 近似分解或分解引入的误差主要是导致收敛缓慢的原因。 需要相对较短的时间步骤来遏制在流动的初始瞬态期间引入的这些错误，然后是斯坦福大学航空航天系，斯坦福，CA 94305-计算流体动力学的前沿-1998编辑：David A. Caughey和Mohamed M. 哈菲兹 ©1998 世界科学

172 MACCORMACK必须采取时间步骤来清除引入的垃圾溶液。 离散化错误和人工耗散主要负责降低数值解决方案的准确性。 这位作者曾经将作为差分方程的附加项引入的人工耗散与由选择引入的真正耗散，例如向上绕，在原始微分方程组中实际出现的术语的离散化中引入。 前者被认为是较低的形式，因为人类干预可以根据需要增加它或提高其强度，直到获得主观上令人愉悦的结果。 计算流体动力学仍然是一门艺术，而不是一门科学。 另一方面，这位作者将在包含陡峭梯度的区域（例如激波）将耗散一阶精确的迎风差分与高阶导数近似相结合，以控制数值振荡。 然而，今天，这位作者认为即使是这种形式的真正的耗散也是数值解体的癌症。 然而，说完后，数值耗散是数值稳定性和获得解决方案的基本要求。 它必须存在。 本文提出了一个制定和实施它的合理方法。 获得控制可压缩无黏性方程或黏性流动的快速准确解决方案的策略是首先尽可能消除数值误差的来源，然后只重新引入控制不稳定性所必需的，而不破坏控制方程的物理学。 策略如下。 1.通过高阶精度的有限差分近似对导数进行离散化，2.将近似因数分解的误差最小化，以及3.添加与框架无关的人工黏性应力和最小强度的耗散。 在要描述的计算中，三阶准确的上风偏置差分商用于近似所有无黏性空间导数，中央二阶精确近似用于黏性项。 MacCormack [6]描述了最小化近似因数分解误差的数值程序。 它由两部分组成：一个近似因数分解，最初由Bardina和Lombard [1]使用，其分解误差比Briley和McDonald [4]或Beam和Warming [2]方法中使用的分解误差小得多，以及一个额外的程序[7,8]，以消除引入的近似因数分解误差。 由于数值误差最小化，程序既快速又准确。 上述项目1和2将不会在这里进一步讨论。 下一节通过检查笛卡尔和圆柱坐标中的Navier-Stokes应力张量，开始讨论数值稳定性所需的项目3。

174 MACCORMACK 通过比较上述两种描述，我们可以找到相似之处并解释不同之处。 例如，考虑圆柱形方程左侧的两个附加项，-pqg 2/r和+pq rqe/r。 第一个是离心力项，如果流体具有角动量，则总是在径向方向增加动量。 第二个表示“滑冰运动员”项，如果径向流体运动朝向轴线，角动量会增加，否则角动量会减少。 流动的物理学需要这两个项来平衡系统，它们在将动量方程的笛卡尔形式转化为圆柱形系统时自然产生。 方程右侧给出的黏性项也显示了圆柱形系统中的其他项，并且需要平衡流动物理学。 流体动力学方程通常在广义坐标中以数值求解，有些甚至不是正交的，并注意将笛卡尔方程直接或通过有限体积公式转换为计算坐标系，以保留与框架无关的流动物理学。 然而，当需要增加耗散时，无论是直接还是通过选择低阶精确差分近似，很少考虑保持物理平衡。 人工黏度的预期效果是抑制数值振荡，这通常会导致不稳定。 然而，一个意想不到的副作用是破坏整个系统的物理平衡。 这类似于在控制圆柱形方程中添加一个重要的新项，该项直接改变径向动量，并在角动量中引起不必要的显著“滑冰运动员”反应。 在多维曲线坐标系中增加人工黏度，同时保持物理平衡是不可望的--只有一个例外。 例外是使用已经独立的框架的Navier-Stokes黏性应力张量本身作为车辆来添加数值耗散。 与用于在雷诺平均Navier-Stokes 方程中包含湍流混合效应的涡黏性的概念类似，通过增加自然黏度来添加数值黏性可以以独立于框架的方式来控制数值不稳定。 M *\textasciitilde{} \textasciicircum{}物理 + \textasciicircum{}数值（10-1）10.3 添加数值黏性的选择 假设数值过程已被剥离为裸露的数值误差来源，并且该过程需要增加耗散来控制稳定性，我们将说明在激波和亚音速对流的两种流动情况下选择数值黏性。 每个案例都会给出数值结果。

流量计算中的附加消散175 10.3.1 激波 激波可能是可压缩流动最突出的特征，它们在计算流体动力学中得到了最多的关注。 然而，它们仍然是多维流动困难的主要来源。 如果它们很强，高于两个的正常马赫数，并且允许在计算网格上不分青红皂白地切割，就会出现它们的不良表示。 在钝体计算的对称轴附近经常发现令人不安的不规则现象[7]。 网格自适应本身不足以准确表示激波。 需要数值耗散。 添加数值黏性的以下选项用于获得稍后通过球体的5马赫流动的结果。 首先，通过确定正常速度分量是否在流动方向上从超音速变为亚声速流动，以及随之增加的压力，找到每个计算坐标方向上所有冲击点的位置。 考虑£、n和（计算坐标系，点1和2在坐标方向上相邻£和2>£■存在冲击点，如果<? M > ci, <j„ 2 < c2 and p 2 > pi or -qnl < ci, -Q„ 2 > c 2 and pi > p 2 其中q n是垂直于表面的速度£ = 常数，位于点£1和£21之间的中间c是音速，p是压力。 如果满足这些条件中的任何一个\textasciicircum{}numerical i = 1/2 • min \{A\% Ah v,Ah\textasciicircum{}\}i ■ p;c; (10.2) 其中Ah(, Ah v和Ah\textasciicircum{}表示£、n和£方向的网格间距距离，p是密度。 在i = 1和2的两个点上，黏度都增加，如上等式（lO.l）所示。 同样，其他坐标方向n和（）被检查是否有冲击点。 上述程序沿着一条窄带增加自然黏度，两点宽，与激波对齐。 这反过来又会影响沿着四点宽的带子进行计算。 然而，与点1或2相邻的超音速点应受到保护，免受数值黏性的影响。 他们不需要它，超声速流动到数值黏性内信息的上游传播不是物理传播。 10.3.2 亚音速对流 数值不稳定通常始于亚音速区域内对流速度的振荡，它们可以通过自我刺激得到加强。

176 MACCORMACK 添加数值黏性的以下选择用于获得稍后通过椭圆的0.2马赫流量的结果。 再次，考虑相邻的两个点1和2，并检查亚声速流动。 如果km| < ci，和\textbackslash\{\}q n2\textbackslash\{\} < c 2，那么fJ-numericaii = 0.05 • min \{Ahf, A/j„ Ah (\} ■ \{pi km I + /9ak» 2l) ■ kn3l + 3(kml+k„2|) + k„4|+e，其中点3和4与1和2相邻，排序为£ 3 < £1 < £ 2 < £4和e \textasciitilde{} 10 -9被添加，以防止在静止区域中除以零。 商因数以1为界，与8qn ' \textasciicircum{} ' d?成正比 和上述数值黏性的值被添加，如Eq.(lO.l)所示，在points和2。 10.4 计算结果 求解[6,7,8]中描述的可压缩黏性流动方程的方法，加上刚才描述的数值黏性，用于解决两个已知数值困难的基准计算。 第一个是，Blottner [3]建议5马赫的黏性层流经过一个球体，第二个是Pulliam[9]建议的0.2马赫无黏性流动，经过椭圆，轴比为6比为5度角，由Pulliam [9]。 在第一个问题中，只有激波的附加黏度，Eq.（ 10.2），被使用。 自然黏性条件足以保持流动剩余部分的稳定性。 对于第二个问题，只包含亚声速流动，只有人工对流黏度，Eq.（ 10.3），被使用。 10.4.1 Mach 5 黏性流动 过去一个球体 Mach 5 层流过去一个球体基准问题最初由Blottner在1990年选择，用于研究在流动对称轴与弓激波的交点附近遇到的数值困难。 这个区域非常敏感，因为座标奇点在包含因子1/r的术语中引起r -> 0。 例如，参见前面以圆柱坐标给出的Navier-Stokes动量方程。 许多论文都观察到

在该区域的流量计算中增加了177个数值困难。 Blottner在讨论这个数值问题的性质方面做得很好，并使用薄层Navier-Stokes代码提出了出色的计算结果。 弓激波被用作水流的外部边界。 在本文中，提出了完整的Navier-Stokes计算，并捕获了激波作为流动的内部特征。 流在三个不同的网格上得到求解，一个由22x26点组成的路线网格，一个42x50点的中等网格和一个82x98点的精细网格，每个网格都跨越相同的流量。 对于这种流，随着点靠近对称轴，计算难度会增加。 图。 1显示了中等和粗的网格。 初始弓激波位置被选择与网格的网格线对齐。 兰金-胡戈尼奥特冲击-跳动关系用于初始化激波后面的流动，当接近身体表面时，法向体表面的速度分量减少到零。 表面温度保持在98.89开尔文。 萨瑟兰公式用于确定黏度。 基于球体直径的雷诺数为1.89xl0 6。计算从CFL数开始，即实际使用的时间步长大小与显式方法允许的最大值的比率约为6xl0 3。然后，在前52个时间步骤中，CFL数增加到2.5xl0 7，此后保持固定。 这个中等网眼的案例运行了256个时间步骤。 如[7]所述，网眼还动态地与弓激波保持对齐，并在冲击附近进行精炼。 中等网格的压力和马赫轮廓如图所示。 2和所有三个网格的表面压力、皮肤摩擦和传热结果如图所示。 3. 请注意，正如预期的那样，皮肤摩擦随着与对称轴的距离而线性增加，表面压力和传热具有钟形曲线，其最大值在轴上，三个解决方案收敛。 停滞压力比p stag/Poo在三个网格上，从粗到细，为32.5800、32.6228和32.6667，这与Blottner给出的Richardson推断值32.6558相比较。 以千瓦/平方米、103.920、104.628和105.613表示的相应传热值也与Blottner给出的107.360值相比是有利的。 图。 4显示了剩余步骤与时间步骤。 在前64步中，网格通过向1马赫等高线放松所选网格线的位置来跟踪冲击[7]。 在此期间，平均剩余减少率约为0.86，高于同一问题的固定网格上0.80[8]的比率。 从第64步到第96步，网格在冲击附近进行了改进，并跟踪了跟踪，此后网格只跟踪了冲击。 这些结果是在每个单词32位的工作站上获得的，剩余减少限制在大约五个数量级

178 MACCORMACK 10.4.2 Mach 0.2 无黏流动 在1990 T.H.经过一个椭圆 Pulliam提出了突出的结果，即欧拉流量代码无法计算出经过简单椭圆的低马赫数流量。 这里获得的基本结果是任何非对称的网格和/或攻击角度组合的提升解决方案。 然后，他挑战计算流体动力学社区“解释这种不寻常的行为”，他定义了一个基准流动问题，即在0.2马赫无黏流动的5度攻角下轴比为6:1的椭圆，这显示了“令人不安”的结果。 流动应该产生零升力和阻力，但所有已知的数值尤拉解都以变量p、u、v和能量来写成。 Pulliam的结果收敛到1.545的显著升力系数。 近七年来，这个问题在数字上等同于d'Alembert的悖论。 Hafez和Brucker [5]为Euler 方程的稳定不可压缩和可压缩形式获得了这个问题的高度精确的解决方案。 此外，Winterstein和Hafez [10]通过在三角形网格上解决的Euler 方程为这个问题提供了出色的解决方案。 在后一项研究中，他们通过改变位置来固定后部驻点，直到达到零升力的状态。 在本研究中，在停滞点位置没有任何条件的情况下解决了不稳定的Euler 方程。 关于这个无黏流动问题的多种解决方案有很多讨论，该问题没有尖锐的后缘，也没有Kutta条件的规范。 然而，如果流程初始化时没有涡量，并且在计算过程中没有生成任何涡量，则应该会出现对称非提升解决方案。 图。 5显示了这个基准问题的130x66点网格的一部分。 拉伸网的外部边界距离身体大约5000个主轴半径。 自然黏度被设置为零，在身体上实施了滑边界条件。 远场边界保持在0.2马赫恒定流量。 通过突然将身体置于均匀的0.2马赫流量中，冲动地开始计算。 最初的CFL数字是3.0xl0 3，在80步后增加到2.5xl0 9，此后保持固定。 计算运行到2000个时间步骤来测试收敛，但解决方案在更早之前是相当稳定的。 无花果。 6和7显示了压力轮廓和流线的相当对称的模式。 如[9]所示，如果产生显著的升力，后部驻点预计将从后缘顶部向下表面移动。 图。 8显示了沿椭圆表面的压力系数。 同样，它看起来相当对称，尽管不是完全的，正如小升力和阻力所要求的那样。 图。 9 显示提升和拖动与时间步进。 最终值分别为4.00xl0 -3和1.66xl0\textasciitilde{} 3。 100步后，升力下降到7.40xl0-2。 剩余步骤与时间步骤显示在

流动计算中增加的消散179图。 10. 对于“普利亚姆悖论”问题，已经得到了一个相当准确的欧拉解。 这一次，就像d'Alembert的悖论一样，答案似乎是黏度，在这种情况下，黏度太多或框架依赖数值黏性，这可能会破坏Euler 方程的敏感物理平衡。 10.5 结论 提出了一种控制数值不稳定的方法，在概念上类似于用于在雷诺平均Navier-Stokes 方程中包含湍流混合效应的涡黏性的方法，通过增加自然黏度将数值黏性直接添加到流动方程中。 Navier-Stokes应力张量本身是独立于框架的，因此这种方法保留了框架独立性，当在任意曲线坐标系中数值求解时，被认为对控制方程的物理平衡很重要。 在最小化分解误差的流动求解器和数值耗散中，这种方法已被应用于两个困难的基准流问题。 一个由Blottner提出，作为Navier-Stokes求解器的测试，另一个由Pulliam提出，作为Euler求解器的挑战。 计算结果与公认的计算和理论结果很一致。 参考资料1。 巴尔迪纳，J.和C.K. 伦巴第，“使用CSCM隐性上风Navier-Stokes方法的三维高超声速流动模拟”，AIAA论文编号 87-1114，1987年。 2. 梁，R.和R.F. 变暖，“可压缩Navier-Stokes 方程的隐性因子方案”，AIAA杂志，卷。 1978年16月，pp. 293-402。 3. Blottner，F.G.，“高超声速流动在球形鼻尖上的准确Navier-Stokes结果”，航天器杂志，第27卷，第2期，1990年，pp. 113-122。 4. Briley，W.R.和H. 麦克唐纳，“广义隐式方法的多维可压缩Navier-Stokes 方程的解决方案”，《计算物理学杂志》，卷。 24，1977年，pp. 372-397。 5. 哈菲兹，M.和D. 布鲁克，“人工涡量对Euler 方程离散溶液的影响”，AIAA论文编号 91-1558-CP，1991年。

180 MACCORMACK 6。 MacCormack，R.W.，“隐式算法的高效矩阵分解”，在1996年6月24日至28日在加利福尼亚州蒙特雷举行的第15届流体力学数值方法国际会议上发表，会议记录将发表在《物理讲义》上，Spr.-Verlag 7。 MacCormack，R.W.，“快速Navier-Stokes求解器的考虑”，在1997年5月2日至4日，加利福尼亚州戴维斯的流动模拟技术研讨会上发表，会议记录将发布。 8. MacCormack，R.W.，“流体流动的新隐式算法”，AIAA论文编号 97-2100，1997年。 9. Pulliam，T.H.，“计算挑战：椭圆的欧拉解”，AIAA杂志，第28卷，第10期，1990年，pp. 1703-1704。 10. Winterstein，R.和M. 哈菲兹，“使用三角形网格的钝体的Euler解决方案：人工黏度形式和数值边界条件”，AIAA论文编号 9S-3SSS-CP，1993年。

181 流动计算中增加的消散图1。 球体周围的中等和粗的网格。 图2。 超声速流动经过球体的压力和马赫轮廓。

MACCORMACK 0 25 50 图3。 表面压力、皮肤Mction和热传递。lo-':;; 1;! Lva：。 >........... 如果,\textasciitilde{}-y:::;:;::! \textbackslash\{\};I:;;:;;::;;.;:;;:;:\textbackslash\{\};:;;;::;;I;;:;;:;;:;;;;;;;;;;;:;;;:: <:;.............................................................. i.......................................... *.................................................................. -........................................... -................. 0 50 100 150 200 2 50 时间步骤图4。 剩余与时间步骤。

流动计算中增加的消散183图5。 关于椭圆的网格。 图6。 亚声速流动的压力轮廓约为椭圆。

MACCORMACK图7。 关于椭圆的亚声速流动的精简。-................................................................................................................................................................................................................................... 1.0 -................................................................................................................................................................................................................................................................................................................................................................................................. .....................................................,....................................................................... 0.00 0.25 0.50 0.75 1。 OO图8。 压力系数与x/c。

\clearpage
\chapter{用于Euler 方程的四运算器守恒格式}
Euler 方程 Jean-Jacques Chattot1 11.1的四运算符守恒格式 导言 Murman和Cole关于通过跨声速小扰动方程（TSD）的松弛求解[1]的开创性工作，由Murman的续集论文[2]完成，在许多方面都是突破，最重要的是计算效率。 但可能还不够说，该方法还有另外两个功能，这些功能非常理想，在Euler 方程的方案中并不常见：没有调谐参数需要调整，而且该方法非常稳健。 为了为Euler 方程准备基础，现在建立了1-D TSD 势流方程的混合方案和Burgers（非黏性）方程的相应方案之间的联系。 马赫1附近流动小扰动电位的控制方程可以写成：dx\textbackslash\{\} 21 dx J = 0 + 边界条件。 扰动电位（p在节点上定义，并定义两个开关来表征每个点的流动制度：r = <Pi+\textbackslash\{\} \textasciitilde{} <Pi-\textbackslash\{\} c 2Ax 现在考虑Burgers 方程的扰动速度：„ \_ <Pi-<Pi-2 du d dt dx = 0 加州大学戴维斯分校机械和航空工程系，加利福尼亚州戴维斯，95616。 计算流体动力学的前沿-1998年。 编辑：David A. Caughey和Mohamed M. 哈菲兹。 ©1998 世界科学

188 CHATTOT + 首字母和 边界条件。 如果引入以下离散关系：\_q>i-<Pi-l Ac，开关可以用u重写为：M,-1+«,- \_ <Pj-(Pj-2 „. \_ ui+ui+\textbackslash\{\} = <Pi+l-<Pi-l 2 2Ac "。 1：i- 2Ac“。 1 =■ i+- 两种公式之间的对应关系最好在以下草图中描述，a）定义电位和开关，b）速度公式的相应数量：“* u c <Pi-i 9»w <Pi 9M a) 9 1 \textasciicircum{} 1 © 1 \textasciicircum{} b) -9- ■\# -©■ Murman [2]的四个算子混合方案与Burgers 方程完全等价。 然而，将引入一个稍微不同的方案，使用不同的音速点运算元，消除了膨胀冲击。 11.2 Burgers 方程的混合方案考虑以下混合方案：i) u i > 0 u i > 0 1 + -..1+1..«..nj.,.» si. 在Ax ii ii) u i > 0 u 1 < 0 1 i+- 2 2.1+1..n.1 s2,,.« N2，斧。 1 i-,.ns2.,." \textbackslash\{\}2, ",--"", («r+ir/2-( HT)V2 1 ( Mf)z/2-(«r\_ 1)'t/2\_0 Af Ac Ac Af<- Ax III) U l < 0 M i > 0 i- i+- 2 2 M I -U J i- ('+-2 2 2Ac l + £(«? +1 -«？ \_！） = 0，Af<°°

保守的四运营商计划189 iv）u x <0 u。 1 14-,,B+1 \textasciitilde{}"T | (Kf+i)2/2-(Mf)2/2 At Ac 0,Ar< Ac -«. 1！ +- 2 可以注意到，声波点，方案iii），不是保守的离散化，因此它在瞬态中不是保守的，可以证明守恒误差与Ac成正比，然而，它在稳态下是保守的。 这是消除膨胀冲击的代价。 这个方案是局部隐式的，因此没有时间步骤要求，因为它可以写：n+\textbackslash\{\} ui+\textbackslash\{\} \textasciitilde{} ui-\textbackslash\{\} u? +l-ur - + U; = 0 At 2Ac 隐式字符是通过系数实现的，而不是空间导数。 这为声波点带来了无条件的稳定性。 需要就稳定性提出一些意见。 由于控制方程是非线性的，所以使用冯·诺伊曼方法在线性化后对稳定性进行了研究。 超音速点算子i）和亚音速点算子iv）的线性化是直截了当的，特征速度u和u，取为i+-i- 2 2常数。 在声波点，隐含性被识别，空间导数“”仅被视为一个正常数。 对于冲击点，方案2 Ax ii），线性化涉及两个符号相反的恒定波速，如：\{\textasciicircum{})2l2-\{uUfl2 Ax „n ui+\textbackslash\{\} I i+- Ax ■ + u u?-ul\_x Ax 该分析提供了上述稳定性条件，在数值测试中被发现是可靠的[3]。 可以证明，该方案在Burgers 方程的稳态下是二阶精确的，源项对纤细喷嘴中的准一维流动进行建模[4]。 11.3 1-D Euler 方程 一维Euler 方程可以写成：dw | df(w) dt dx，其中：w ■■ fp \textbackslash\{\} pu, /(w) = 'pu pu puH pu2 + p E = -(7-l)p 2 H IP (y-l)p 2

190 CHATTOT 让A = -=-是Jacobian 矩阵，a2 = -是音速的平方。dw p p 与一维Euler 方程系统相关的特征矩阵C有三个不同的特征值：|C| = |A-A/| = 0 X (j) = \{u -a;u;u +a\}, j = 1,2,3 系统是完全双曲的。 左特征向量由以下方式获得：l\textasciicircum{}A = fi¥J\textbackslash\{\}j = 1,2,3 兼容性关系是：rf»:,»\{*:+£) = „,;\_-U,3 这些代表内导数，定义的特征线：f*L) =XU),j = 1,2,3 三个兼容性关系等同于原始系统。 它们构成了下面描述的数值方案的基础。 11.4 特征框方案 考虑节点i周围的两个框[i - l,i], [»',/ +1]。 未知数w”位于节点上。 在每个框中，使用Roe平均值[4]对离散Jacobian 矩阵、特征值和特征向量进行评估。 假设流量从左到右（«" > 0），会出现以下四种情况：0 u" i - a", > 0 u" i - a", > 0 "2 ' 2 i+- i+-2 2 /(3) \textasciicircum{} +1 -\textasciicircum{}, \textasciicircum{}-\textasciicircum{}l A/ Ax = 0 H) M" 1 -a" 1 >0 u" 1 n a。 1 l+- i+- 2 2

保守的四运营商方案 191 r 0 w,n+I -w; n en \_ en en en i | Ji+l Ji, Ji Ji-l At Ax Ax ) z(3) '. 1 我-。 2\_在Ax = 0 in) u n i - a"! < 0 w"! - a" j > 0 i- i- I'H- in- 2 2 2 2."...” B+1 „ AW w w-\textasciicircum{}-i ■ -\textasciicircum{}, ' 2At Af Ax i+± i-± 2 2;(2) «t+l-y»? \textasciicircum{}f？-f！ Li At Ax = 0 iv) «" i -a"! <0 u n i -a”！ <0 我！ ('+- J+- 2 2 2 2 *? +1-»>？ I\&i-rr Af Ax /(2) Z(3), Ar Ax = 0 注意，对于改变符号的特征A w，方案与上述混合方案相似，特别是关于冲击和音速点

192 CHATTOT操作员，因此，该方案是保守的，计算出正确的冲击速度，并通过在声波条件下强制平稳过渡来防止膨胀冲击。 时间导数是从3x3系统的解中获得的，其矩阵由左特征向量组成。 在超音速点，左特征向量的矩阵可以被单位矩阵取代。 11.5 离散特征值和特征向量的使用 上述方案可以通过利用Roe平均值的属性，以等效但略有不同的形式编写，用于兼容性关系CRm和CR°\textbackslash\{\}，如下所示：jU) fi -fi-l = ft) ft) wi \textasciitilde{}wi-l iS Ax r '' i- i- 2 2 Ax 对于声波点运算符，特征值和左特征向量在节点状态下评估。 这确保了该计划不是守恒形式。 在冲击点，有四个特征流向点；'，有两个兼容性关系CRW。 情况与Burgers模型相同，通过将两个周围盒子的空间衍生物相加，从而产生一个不一致但守恒格式，情况被处理得相同。 然而，需要选择特征向量，在上述公式中表示为l，m。 建议将Burgers 方程的类比更进一步，以便Euler 方程能够以与更简单模型相似的方式进行线性化，以两个波速为基本特征。 冲击点方程可以重写为：,.n+l...n n,„n - + A( l) w i+l-wt (1>.. 1 A„。 \_\_\}\_ 2 /, (1) W; W: At t+- 2 Ax Ax i-l i+- i+- Ax 2 «+- Ax One可以选择向量I (i) "2 '"I，使上述方程中的最后两个项抵消。 通过查找线性组合的值6，更容易搜索此类向量：j«=ez<Va-0)' (1\textbackslash\{\} (+- i- W, i \_ = 0 使用特征向量的属性，这简化为条件：efi\textbackslash\{\}u i-i"V). w'"w'-'f(i-e)/",.( A J-AW/).\textasciicircum{}\textasciicircum{} i+-,--,-- Ax,\_\_, + -,- +\_ Ax 有了这个选择，人们可以重写兼容性关系CRm的四个算子方案，其中开关因紧凑性而被省略，以更新形式，如下：

保守的四运营商方案 193 0 / (1>,.wf+1=/« w, f-\textasciicircum{}M- wf-i) a) jp.wTM® m)/P. Wf+1 =Zp。 Af Af w？ \_£L A a> (w » \_ w »)\_£L A a) (wf \_ wf \} Ax,+-L Ax j-i 2 2 Af 2Ax 415« ■W: 1 + \textasciicircum{} (A (D A 0) \} Ax,+i,-\_. L。 V) /« w n+\textbackslash\{\} = f X\textbackslash\{\}。 I+- Vf，f-\textasciicircum{}Vv\textasciicircum{}-wf）斧！ +-对于兼容关系C7？（ 2）和CR 0）方案与case/）相同。 这与Burgers 方程的混合方案相似，因此稳定性条件是相同的。 因为这是一个系统，左特征向量扮演未知向量w的投影算子的角色。 该方案是L“-稳定的，在数值实验中发现非常稳健。 没有自由的参数可以调整。 在稳态下，方案是二阶准确的。 11.6 一些结果 Sod的激波管测试案例在CFL条件编号1,001点的网格上模拟，共500步。 结果如图所示。 1，并且与确切的解决方案完全一致。 方程g(x) = 1 + 4(x - 0.5)2的收敛-发散喷嘴中的窒息流，出口压力p exil = yp"，在发散部分产生压缩冲击。 稳态解与Roe方案进行了比较。 可以在图中看到。 2，罗方案，在声点上没有人工黏度，不会提供平稳的过渡。 最后一个例子对应于超音速风洞启动的物理问题。 高压水库通过双喉管道连接到大气排气管。 正向收敛-发散喷嘴有一个单位喉部面积，在测试部分A、K、=1.25达到最大面积。 收敛-发散扩散器喷嘴有一个可变的喉部，取决于参数a。 随着减少，第二个喉咙增加。 几何方程由以下方式给出：

194 CHATTOT g(x) = -(2x-1)4 -hlx-1) 2 + -,0<x<0.5 4 2 4,!; u)-a\textasciicircum{}(2*-l) 4-|(2*-l) 2\} + |,0.5<x<l,0<a<l 为了使测试部分没有电击，第一次电击必须被第二个喉咙“吞下”。 这需要第二个喉咙来满足不等式[6]：T*\textasciicircum{} A* \textbackslash\{\}.test\textbackslash\{\} A，在这里，这对应于-\textasciicircum{} > 1.1169 A为了模拟启动，初始条件对应于静止空气，左边界点（i = 1）表示储层条件，管道的其余部分处于大气压力。 以无维形式：ui=0 y + 1 M，=0 2y。 „。 Pi=--pexit，i = 2，...，ix y + 1 Pi = Pexit pexil =0.9。 网格有101个点。 代码使用A，I A\{ =1.117运行30,000次迭代，CFL数统一。 风洞几何形状如图所示。 3. 图中描绘了压力随着时间的演变。 每1000次迭代4次。 可以看出，第一次电击需要很长时间才能被第二次喉咙“吞下”。 在测试部分，冲击速度达到最低值，略高于零度。 一旦进入扩散器喷嘴的收敛部分，第一个冲击迅速移动到第二个喉咙的右侧，并与第二个冲击融合，在流量达到稳态时形成一个更强的冲击。 11.7 结论 作者承认E的遗产是公平的。 M。 Murman到计算流体动力学的领域和他的个人课程。 作者于1971年在法国ONERA开始了他的职业生涯，并负责根据Murman的新原创贡献开发一个跨声速代码。 如果有的话，本文试图表明，通过一些进一步的分析，现有的欧拉方案可以转化为更简单、更纯粹的形式，并达到更高的美水平：原作之美。

保守的四运营商计划195参考资料1。 默尔曼，E.M. \& Cole，J.D.，平面稳定超声速流动的计算，AIAA杂志，卷。 9，第1期，1971年。 2. Murman，E.M.，通过松弛方法计算的嵌入式激波分析，AIAA杂志，卷。 12，第5期，1974年，pp. 626-633。 3. Chattot，J.J.，一阶偏微分方程的盒子方案，计算流体动力学的进步，Gordon Breach出版社，1995年，pp. 307-331。 4. 查托特，J.J. \& Malet，S.，Euler 方程的“盒子方案”，数学讲义，1270，1987年，pp. 82-102。 5. Roe，P.L.，近似黎曼求解器，参数向量和差分方案，计算物理学杂志，卷。 43，第2期，1981年，pp. 357-372。 6. 祖克罗，M.J. \& Hoffman，J.D.，气体动力学，卷。 1，John Wiley \& Sons，1976年。

196 CHATTOT图1。 激波管-Sod测试案例。 速度、密度、压力和比能量分布图2。 收敛-发散喷嘴中的窒息流。 使用Roe方案和箱式方案比较压力分布。

四个操作员守恒格式 197 图3。 超音速风洞区域分布。 图4。 超音速和启动 压力的迭代演变。

\clearpage
\chapter{带有分体和铉法兰襟翼的自动阻挡}
带有分体和铅翼的翅膀自动阻挡D。 Scott Eberhardt 1 \& Pratomo Wibowo 1 摘要 提出了一种为具有有限跨度襟翼的机翼生成结构网格的技术。 该方法由四个步骤组成。 它们是：选择前线，确定引导程度，确定前线推进优先级，然后推进前线。 双曲预测器和椭圆校正器用于推进前线。 最后一个块网格都是CO和Cl连续的。 前端在需要时自动创建，在不需要的地方合并。 前连接定义了一个称为引导程度或D的参数。 0. G.，它决定了将要生成的网格的拓扑结构。 单个锋面按照其表面拓扑结构指定的顺序前进。 该方法用于在具有部分跨度分割和铰链襟翼的机翼周围生成结构网格。 所示示例是北美P-51野马机翼上的模型襟翼。 显示了网格生成步骤以及最终网格。 生成的网格仅用于说明目的，并不打算在流动求解器中使用。 1 华盛顿大学，Box 352400，西雅图，华盛顿州 98195-2400 计算流体动力学 - 1998年前沿 编辑：David A. Caughey和Mohamed M. 哈菲兹 ©1998 世界科学

200 EBERHARDT \& WIBOWO 12.1 介绍网格生成是CFD在工程流程中成功应用的重要任务。 网格生成 一直是并且仍然是一项劳动密集型且需要专业知识的任务，特别是对于复杂的配置。 网格生成需要大量人力和劳动力的配置很容易找到：汽车、水下车辆和飞机。 进一步复杂的一个例子是，现代飞行器在起飞和降落时使用高升力装置，如板条和多个襟翼。 围绕这种复杂配置的流场通常被分解为几个相对简单的子域。 一旦确定了子域，用结构化的网格填充每个子域是相对简单和直接的。 子域可以与相邻的子域共享一个公共接口，也可以与另一个子域叠加。 使用具有共同介面的子域的第一个策略是本研究的主题，它被称为多块网格方法。 将物理流场分解成块，即所谓的块生成，是多块网格生成过程中最具挑战性的任务。 自动块生成算法的开发是网格生成领域几个备受追捧的话题之一（[1]-[11]）。 许多现有的块生成方案以拓扑形式生成块边界。 换句话说，一旦用户获得这些块边界，他/她应该沿着这些块边界分配网格点，并为每个块生成内部网格点，以获得最终的多块网格。 人们不能忽视向非结构化网格的推动，它解决了许多网格生成挑战。 作者承认非结构化网格的优势，但一直顽固地不愿意留下结构化网格的熟悉感。 本文概述的方法是生成块，作为网格生成过程的一部分。 表面网格用作输入，并分为“正面”。 关于如何生成前线的描述在[12]和[13]中找到。 本文介绍了一个叫做引导程度的概念，本文将详细介绍。 虽然这个概念在[12]和[13]中使用，但它并不明确详细。 讨论了推进优先级，然后简要介绍了使用的预测校正算法。 最后，介绍了关于北美P-51野马的一些结果。 请注意，这里介绍的方法是生成关于逼真的多元素翅膀的网格的分步过程的一部分。 最简单的襟翼系统首先受到攻击，以便一步一步地解决技术问题。 在这一点上，该方案成功地生成了关于无限小厚度的分割襟翼和铰链襟翼的网格。

翅膀的自动阻塞201图1如何生成前部以保持简单的边缘连接。 12.2 方法 网格生成流程中使用了四个重要元素。 它们包括选择前线，确定引导程度（D. 0. G.）值，确定前面推进优先级，然后推进前面。 正面是向外推进的网格表面，类似于双曲网格生成方案。 对于带襟翼的机翼，正面通常分配给机翼的上下表面、襟翼的上下表面、机翼尖端和机翼/襟翼间隙的表面，如果存在的话。 前推进的一个重要细节是，所有前必须在角落连接。 这意味着前线或块可能必须分成多个前线才能满足这一条件，如图所示。 1. 生成新战线的过程是自动的，不需要用户交互。 对此有更完整的讨论可以在[12]中找到。 12.2.1 指导性前线可以使用“自由推进”或“引导推进”进行推进。 方法的选择取决于D。 0. G.价值。 自由推进意味着表面向外推进，而不考虑相邻的正面（除了确保相邻的正面采取相容的推进步骤）。 自由推进的一个简单例子是未拍打的无限跨度机翼的机翼表面。 引导式前进意味著一条战线的前进明确取决于另一战线。 一个例子是内部的直角角，如图所示。 2. 标准的双曲网格生成器将尝试使两个表面为相同的网格线，这是不自然的。

202 EBERHARDT \& WIBOWO 图 2 内部角落和“自然”网格。 [13]的一个重大创新是发展了指导程度的概念。 在这项工作中，我们完善了概念，并填补了几个缺失的漏洞。 在这里，我们将尝试描述什么是D。 哦。 G. 价值是。 每个前表面都有四个边缘，这些边缘与邻居（或边界）相邻。 例如，考虑一个简单的机翼。 上表面和下表面锋在前缘和后缘相遇。 前缘是圆形的，因此两个正面以零角度相遇。 两条战线可以通过自由推进来发展，但必须保持联系。 后缘有一个非常大的角度（接近180度）。 这些战线应该独立推进（选项将在下文讨论）。 当前方案使用默认的D。 哦。 G. 值取决于有多少表面边缘相接。 表1显示了D. 0. G. 价值观，何时实施，以及哪个数字说明了所选的前沿推进方法。 自由进步D。 0. G.-1和-2的值需要更多的解释。 为了一个D。 哦。 G.值为-1，如果一个正面提前几步D。 哦。 G。 默认角度推进类型图0 -45°至45°自由推进3 -1 45°至135°自由推进4-2 135°至180°自由推进5 -3无邻居自由推进6 1 -135°至-45°引导推进7 2 -180°至-135°引导推进8表1 D的插图。 哦。 G.价值

翼的自动阻挡203图3引导度=0：（a）两个正面之间角度为20度的初始网格块；（b）三次网格推进后的结果；（c）网格推进的最终结果。 图4 引导程度= -1：（a）两个正面之间角度为120度的初始网格块；（b）第一个正面的网格推进；（c）网格推进的最终结果。其侧面将创建第二个正面的延伸，如图所示。 4b。 当这个正面前进时，“角落”被填满了。 A D。 哦。 G.-2的值是人们在机翼的后缘上可能期待的。 如图所示，这两个块独立生长。 5b。 然后，新方块的侧面创建了一个新的前线，然后向前推进，如图所示。 5c。 这个D。 0. G.用于处理机翼后面的唤醒区域。 A D。 哦。 G. 2的值是独一无二的，即产生一条奇异的线。 这个D。 哦。 引入G.值是为了填补分裂的翻盖和机翼之间的区域。 由于两条前线基本上以小角度坍塌，最自然的网格是从奇点辐射的网格。 用户可以覆盖默认值。 例如，在图中。 9 用户可能希望网格缠绕在尖锐的角落。 在图中。 9b，默认的D。 哦。 G. 显示-1的值。 用户可以指定D。 0. G.零值以获得

204 EBERHARDT \& WIBOWO（a）图。 5. 引导程度=-2：（a）两个正面之间角度为150度的初始网格块，（b）每个网格正面独立前进，（c）网格前进的最终结果（a）（b）图。 6. 引导程度=•3：（a）初始网格前没有任何网格块连接到其侧面，（b）五次推进后的网格结果。 （A）（b）（c）图。 7. 引导程度=1：（a）两个正面之间角度为-80度的初始网格块，（b）两次推进后的网格，（c）网格推进的最终结果

翼的自动阻挡205（a）（b）（c）图8引导度=2：（a）两个正面之间角度为150度的初始网格块；（b）三次推进后的网格；（c）网格推进的最终结果。<a>（b）（c）图9强加默认值以外的引导度值：（a）两个正面之间角度为120度的初始网格；（b）使用默认值（DOR = -1）的网格推进结果；（c）使用强加值的网格推进的结果。图所示的网格。 9c。 12.2.2 推进优先级 如果战线按任意顺序推进，可能会出现问题。 在这项研究中，确定前线推进可以优先排序，这可以使用默认优先顺序来完成。 例如，一个空腔应该在其他战线开始前进之前被填满。 因此，一般来说，凹面先前进。 这将包括带有D的正面。 0. G. 2的值和四边带D的值。 0. G. 1的值。 遵循上述凹面，引导前线是最高优先级。 进步的优先事项是拥有最多引导战线的前沿。 最后，一旦引导前推进完成，就可以开始自由推进。

206 EBERHARDT和WIBOWO Fronts可以一步或多步推进。 如果一个人在一个简单的机翼上推进前线，那么向前推进上表面一步，然后推进下表面可能是有意义的，而不是将上表面推进到外部边界，然后推进下表面。 然而，在腔体中，在推进周围前线之前，填充整个腔体是有意义的。 在某些情况下，选择是显而易见的，方案会做出选择。 但是，有时这个步骤并不明确，用户必须进行干预。 这种情况可能发生在机翼尖端。 有时最好先把翅膀向前推进几步，然后再让机翼前向前。 这是一个仍然需要工作才能实现自动化的领域。 12.2.3 前推进 前推进算法使用修改前推进方法作为预测器和椭圆方案作为校正器（MAP方案）的组合。 Kim为2D引入了MAP方案，并将其用于围绕2-D配置的多块网格生成[12]。 MAP方案一次按一个单元格高度推进一组表面网格块。 顾名思义，MAP方案是一种预测校正方法。 正面由预测器推进，由校正器平滑。 每个前线都经过这个预测器-校正器步骤，直到感兴趣的流域被多块网格填满。 预测步骤：这使用修改后的前进前线方法推进前线。 前进前线方法（AFM）最初是为非结构化的三角形网格（[14]-[16]）开发的。 然而，在目前的研究中，这种方法已被修改用于六面体网格生成。 修改后的AFM不仅可以同时生成六面体单元的集合，还可以根据周围情况调整每个网格点的距离和前进方向。 它首先从沿着前边界的这些值中插值前线每个内部点的前进距离和方向。 然后，前面的所有点都向前推进，从而产生一个新的正面，并获得与沿旧正面的四边形单元相同数量的六面体单元。 在提供表面网格作为初始前线时，\{i,j)每个表面块的索引应以一致的方式排序，以便很容易识别其向外方向；例如，使用右侧规则将第三个fc轴作为向外方向（e；x ij = ejt）。 然而，新前线通常携带着旧前线的不光滑性（如果有的话），并放大了它，使进一步的推进变得不可能或毫无意义，特别是对于凹陷地区。 通过采取纠正步骤来避免这种情况。 校正步骤：在校正步骤中，拉普拉斯方程按照科尔多瓦的方法用作椭圆方程[17]。 要应用椭圆校正器，先是旧正面相对于新正面的图像

引入了WINGS 207的自动阻塞。 使用这3个战线，3-D拉普拉斯方程沿着新战线求解，另外两个是固定的边界条件。 作为椭圆方程的解，新前线的网格点分布更平滑。 12.3 从典型的CFD活动的角度来看，Flapped Wings的结果所选的测试问题有点不寻常。 它反映了作者对经典和历史飞机的兴趣，以及对设计和建造自制飞机的学生项目的支持。 机翼是北美P-51野马的机翼。 虽然机翼是原来的几何形状，但襟翼不是。 动机是设计一个新的机翼，为休闲飞行员提供更温顺的行为，同时保留P-51野马的外观。 设计应该是两个座位，比例是原来的75\%。 12.3.1 拆分襟翼 第一个例子是图中所示的拆分襟翼配置。 10a。 在图中。 10b 机翼部分在右侧被切割出来，以显示分体襟翼的位置。 襟翼有两个正面，一个顶部和一个底部，尽管它非常薄。 下部机翼表面前部和上部襟翼前部在有D的边缘连接。 哦。 G. 2的值。 因此，如图所示，首先填充翼和襟翼之间的空隙。 10c。 新填充的间隙块的两侧形成了新的正面。 如图lOd所示，接下来将推进两个襟翼之间的区域。 接下来，正面将推进到翼尖，如图所示，从翼尖到翼尖形成一个光滑的网格。 下一步是推进机翼的上表面和下表面，这些表面现在包括之前创建的尖端到尖端前部的下表面，如图1所示。 请注意，在尖端和尾部区域前进之前，这些前线向前推进到外部边界。 这是用户输入的结果。 接下来生成尖端区域，图lOg，最后是图lOh所示的最终网格的尾迹。 虽然网格太粗糙，无法实际使用，但它说明了网格推进的主要特征。 它还表明，最终的网格是CO和Cl连续的。 12.3.2 铰链襟翼 在逼真的机翼上创建网眼的下一步是引入铰链襟翼。 此时，假定铰链襟翼的厚度为无穷小，就像分体式襟翼一样。 表面网格如图所示。 11a。 自动发电机

208 EBERHARDT \& WIBOWO 图 10 北美 P-51 野马机翼，带分体襟翼，显示 网格生成 的阶段。

机翼自动阻挡 209 图11 P-51机翼带铰链、极薄、铰链襟翼。如图lib所示，选择先填补机翼空隙，因为它的前部边缘最引导。 然后，问题被简化为上一节的分体襟翼案例。 12.4 结论：围绕多元素机翼生成结构网格存在许多技术问题。 试图从简单的机翼/襟翼系统开始解决其中一些技术问题；劈叉和铰链襟翼。 引入了一种使用D指令的前进前线的方法。 哦。 G.价值。 接下来，给出了前线推进的优先权。 最后，对MAP计划进行了简要描述。 展示了围绕带有分体和铰链襟翼的P51机翼的网格生成。 虽然网格并不打算用于流动求解器，但它们显示了一般过程。 生成更逼真的翼/襟翼配置的网格的下一步是引入有限厚度的铰链襟翼。 这项工作此时几乎已经完成。 完成该步骤后，将引入开槽襟翼。 建议处理开槽襟翼的技术是使用[12]和[13]中描述的“桥接”理想，这将使问题减少到前一种情况。 参考资料[1.] Schonfeld，T.和Weinerfelt，P.，通过前进的前线技术自动生成四边形多块网格，Proc。 第三届 数字网格生成会议，CFD，巴塞罗那，1991年6月3日至7日。 [2] Stewart，M。 E。 M.，应用于多元素空投网格的域分解算法，AIAA杂志，卷。 30，不。 6，第1457-1461页，6月，

210 EBERHARDT \& WIBOWO 1992。 [3] 安德鲁斯，A。 E.，基于知识的流场分区，计算流体动力学 '88中的数值网格生成，Pineridge Press，pp. 13-22，1988年。 [4.] Weatherill，N。 P.和Shaw，J. A.，用于飞机配置的组件自适应网格生成，AGARD AG-309，第29-39页，1988年。 [5.] 肖，J。 A.，乔治拉，J. M.和Weatherill，N. P.，Aerodynamcc Geometrees的元件自适应网格的构建，计算流体动力学 '88中的数值网格生成，Pineridge Press，pp. 383-394，1988年。 [6.] 乔治拉，J。 M.和Shaw，J. A.，关于多生物网格生成相关问题的讨论，AGARD CP-464，1990年。 [7.] Allwright，S。 E.，多双色域去电沉积和表面技术网格生成，计算流体动力学 '88中的数值网格生成，Pineridge Press，pp. 559-568，1988年。 [8.] Allwright，S。 E.，多生物拓扑规范和完整飞机配置的网格生成，AGARD CP-464，1990年。 [9.] 丹嫩霍夫，J。 F。 Ill，通用几何的块状结构技术，AIAA论文91-0145，里诺，1991年1月。 [10.] 丹嫩霍夫，J。 F。 Ill，为复杂配置创建网格抽象的新方法，AIAA论文93-0428，里诺，1993年1月。 [11.] 诺布尔，S。 S.和Cordova，J. Q.，结构网格生成的阻塞算法，AIAA论文92-0659，1992年。 [12.] Kim，B.和Eberhardt，S. 用于高升配置机翼的自动多块网格生成，美国宇航局CP-3291，美国宇航局表面建模研讨会，网格生成和CFD解决方案的相关问题，俄亥俄州克利夫兰，1995年5月。 [13.] Kim，B.，关于复杂几何体的自动多块网格生成，博士。 论文，华盛顿大学，1994年。 [14.] 洛，S。 H.，任意平面域的新网格生成方案，Int'l。 J。 数字。 冰毒。 Engr.，卷 21，pp. 1403-1426，1985年。 [15.] Peraire，J.等人，可压缩流动计算的自适应重合，J. 补偿。 物理，卷。 72，pp. 449-466，1987年。 [16.] Lohner，R.和Parikh，P.，通过推进前线方法生成三维非结构化网格。 AIAA Paper 88-0515，里诺，1988年1月。 [17.] 科尔多瓦，J。 Q.，双曲网格生成的进步，第四届。 对Comp的症状。 流体动力学，加利福尼亚州戴维斯，1991年9月。

\clearpage
\chapter{局部预处理：操纵大自然母亲愚弄父亲的时间}
局部预处理：操纵大自然母亲愚弄父亲时间David L. Darmofal 1 \& Bram van Leer 2 13.1 导言 求解稳态方程的常见策略是及时进行相关的非稳态方程，直到解收敛为静止结果。 在可压缩Euler 方程的情况下，这种方法具有重要优势，即将问题从稳态的混合双曲斜问题转变为瞬态阶段的严格双曲问题。 不幸的是，这种方法还将由于不稳定方程而导致的任何刚度引入到稳定方程的收敛过程中。 例如，在接近不可压缩流动的中，声速明显快于本地流速。 因此，声学扰动的传播速度比仅随局部速度传播的对流扰动（如熵和涡量）快得多。 对于显式时间行进代码，不同传播速度的这种刚度可以显著减慢收敛到稳态：快速模式设置最大允许的时间步长，而慢模式设置扰动对流出计算域所需的迭代次数。 为了加速收敛的时间行进方法，可以1航空航天工程，德克萨斯A\&M大学，学院站，TX 77843-3141 2航空航天工程，密歇根大学，安阿伯，MI 48109-2140 计算流体动力学-1998年前沿 编辑：David A. Caughey和Mohamed M. 哈菲兹 ©1998 世界科学

212 DARMOFAL和VAN LEER在非稳态方程中引入了一个预调节器，du -\textasciicircum{}+P(u)r(u) = 0，其中u表示状态向量，r是计算试图驱动到零的空间残差，P是预调节器。 在本综述中，我们将重点介绍二维Euler 方程的局部预处理器。 局部预调节器使用来自空间残差将被预调节的细胞或节点的纯局部信息进行评估。 虽然预调节器修改了非稳态方程，但假设P是一个正定矩阵，分析方程的稳态解是未修改的。 正如我们将讨论的那样，预调节器的离散实现通常确实需要对空间残差进行微妙的修改；幸运的是，这种修改通常具有提高离散化精度的好处。 在本文中，我们将详细介绍为2-D Euler 方程设计局部预调节器的过程。 虽然这些例子将特定于这些方程，但基本的设计考虑和分析技术适用于许多其他方程组。 然后，我们将举例说明局部预处理在实践中取得的成就。 最后，如果局部预处理要充分发挥其潜力，我们将指出仍然需要发展的剩余领域。 13.2 设计考虑 局部预处理的第一次尝试可能与Chorin的[6]人工可压缩性的不可压缩流动方法有关。 虽然Chorin没有将他的方法视为局部预处理，但他的方法将压力的时间导数添加到不可压缩连续性方程中，这相当于局部预处理。 自从最初努力修改非定常流动方程以计算稳定流量以来，局部预调节器的设计不断扩展和完善。 目前，局部预处理器的设计是截然不同的，有时是相互竞争的标准的平衡。 在下文中，我们描述了在为2-D Euler 方程推导局部预处理器时面临的各种设计考虑因素。 通过利用对称变量[35, 1]，大大简化了2-D Euler 方程局部预调节器的分析，这些变量定义为dw = -，dq，dr，dS，。 Pc J

局部预处理 213 其中dq和dr是速度的流向和法向分量的变化，dp是静压的变化，dS = dp-c 2dp与熵的变化成正比，p是密度，c是声速。 使用此基础，预制条件Euler 方程可以写成dw dw \_ \_ dw (13.1)，其中M 1 0 0 1 M 0 0 0 0 M 0 0 M B 0 0 0 0 0 0 0 0 0 0 0 0，马赫数，M = q/c，流向和法向坐标，(\textasciicircum{},77)，q是局部速度。 13.2.1 基本要求 局部预调节器的基本要求是，它不应反转任何波的传播方向。 违反此要求的预处理器将导致边界条件的不正确应用，因为入射波和出波将反转作用。 可以通过强制P为正定（尽管不一定对称）来满足这一标准。 在这种情况下，先入条件的剩余是贡献的正加权总和，这不会逆转时间发展。 虽然积极的确定性是一个基本标准，但它并没有对预处理器施加严重的限制。 确保正确定预处理的典型方法是在满足其他理想的约束后检查它。 当预处理到Euler 方程时，结果也应该是一个良好的双曲方程集。 良好定位的一个必要条件是，随着时间的推进，系统的能量保持约束。 保证有限能量的充分但不是必要的条件是迫使先入条件系统对称化。 如果存在对称的正定矩阵Q，并且QPA和QPB是对称的，则前条件系统是可对称的。 如果存在满足这些要求的Q，那么该系统很容易被证明在规范w Qw中是稳定的，忽略了边界条件的影响。 与P的正定性类似，对称性通常在预调节器设计过程结束时得到验证。 一些重要的双曲守恒系统

214 DARMOFAL和VAN LEER法律，包括Euler 方程，是可进行的。 Godunov[ll]表明，这意味着存在满足额外守恒律的熵函数。 然而，先入为主的Euler 方程已经失去了守恒形式，因此Godunov的结果不适用于这些。 最后，预调节器还应在所有流动制度中满足certsn平滑度条件，特别是在稳定Euler 方程的断点，声波点（M = 1）。 在预调节器的开发中，由于声波传播方向的质变，亚音速和超音速制度必须分开处理。 M = 1的预调节器分支不匹配可能会导致声波点捕获能力差并停滞收敛。 13.2.2 长波时间尺度的均衡 如导言中所述，由于对流模式和声学模式的传播速度不同，Euler 方程在低马赫时具有固有的刚度。 类似的论点也可以对可压缩流动近声条件提出。 局部预处理可用于平衡这些传播速度或时间尺度。 此外，对于黏性流动，耗散时间尺度现在进入，也可以在预处理中进行解释。 请注意，平衡时间尺度的分析和设计过程是在连续偏微分方程水平上进行的。 由于离散数值方案应该正确近似长波行为，我们期望在数值解中再现预调节器的长波属性。 然而，短波或高频通常不能由离散化正确表示。 因此，波速的均衡只涉及长波时间尺度。 分析预条件波属性的第一步是计算形式为w(x, y,t) = w(xcosq!)的平面波解的传播速度，A \textasciitilde{} + ysin4 - At），其中4是波传播方向相对于溪流方向的角度。 将平面波代入方程13.1导致A的特征值问题，其中r是相应的右特征向量。 如果没有预处理，2-D Euler 方程具有特征值。虽然预处理设计可以直接从特征值进行，但Xi，a.必须用这种方法给出警告。 具体来说，波速刚度的一个典型度量是使用最大值与最小值的比率

局部预处理 215大小，即max; max0 |A;(</>)| minimin\textasciicircum{} \textbackslash\{\}\textbackslash\{\}i(<j>)\textbackslash\{\} ' 然而，从2-D Euler 方程可以看出，这种方法是徒劳的，因为每当<j> = 90°时，A2,3 = 0表示无限刚度，无论区域性马赫数如何。 在考虑预制系统时，也存在类似的问题。 为了解决这个问题，Van Leer等人[40]使用了在某个时刻穿过一个点的平面波的包络；这实际上是点扰动发出的波前。 点扰动包络由cos <p - sin ( sin d> cos d> X(<f>) \textbackslash\{\}'(4>)给出，其中\{g\textasciicircum{},g\textasciicircum{})是包络的位置。 事实上，这种方法等同于使用组速度而不是相位速度（即 A）其中（<7£,g,,)是方向<f>的平面波的群速度向量。 群速度的一个众所周知的属性是，扰动的能量在群速度下传播。 因此，优化点扰动包络等同于优化能量传播的速度。 跟随Van Leer等人，我们还可以使用点扰动包络来目视检查系统的波传播特性。 Euler 方程的波前如图所示。 1为M = 0.1、0.5、0.9和1.3。 注意：所有图都已按最大波速缩放。 对应于A2,3的对流波在所有图中都清晰可见为孤立点扰动。 相比之下，声波前以局部对流速度q传播，同时以声速c径向传播。 在低马赫数下，我们可以看到前面提到的快速声学模式和慢对流模式产生的刚度。 在较高的马赫数下，硬度也存在，但在这种情况下，快速/慢模式是声学模式的下游/上游传播部分。 采用群速度大小，g = \textbackslash\{\}lg\textbackslash\{\} + g\%，波速刚度的适当度量是K 9，定义为，\_最大，最大\textasciicircum{} |<? T(</>)| 9 \textasciitilde{} 'mimmin\textasciicircum{} \textbackslash\{\}gi\{4>)\textbackslash\{\} ' 如果没有预处理，2-D Euler 方程的波传播刚度由n g = (M + l)/min(M, \textbackslash\{\}M - 1|)给出。 图中显示了K g的图。 2. 不出所料，K g表明M -¥ 0和M -¥ 1存在刚度问题。 虽然不是第一个可压缩流动的预调节器，但Van Leer-Lee-Roe预调节器[40]是第一个适用于所有马赫数的。 此外，

（a）M = 0.1 DARMOFAL \& VAN LEER（c）M = 0.9（b）M = 0.5（d）M = 1.3 图1 2-D欧拉的波前。 波前已按最大波速缩放。对于对称预调节器类别，P = PT，Van Leer-Lee-Roe可以证明为所有马赫数产生最佳波前[l7]，刚度最小，K\textasciitilde{}。 具体来说，这个预处理由p = dm和7 = min（P，P/M）给出。 对于亚音速流，该预调节器产生的特征值是= - 亚音速PVl的M2 c0s2 4，x2,3 = qcos4

局部预处理 217 图2 波前刚度，K S，与Euler（实线）、最佳预调节器（虚线）、块-Jacobi预调节器（虚线）和Choi-Merkle/Weiss-Smith预调节器（虚线）的马赫数相比。 注意：Van Leer-Lee-Roe预调节器是最佳预调节器。 对于超声速流动，特征值是Ai 4 - c (\textbackslash\{\}/M 2 - 1 cos 6 I sin d>) \textbackslash\{\} r D J'* i ) tor超音速>V; r。 A 2, 3 = qcos<t> J 图中所示的结果波前。 3和图中的刚度。 2. 从这些情节中可以看出，低马赫数的刚度已经完全消除了。 在亚声速流动中，声波前现在是以原点为中心的椭圆。 椭圆的小轴朝向流方向，半宽为q\textbackslash\{\}/l - M 2。 因此，随着声学条件的达到，这种椭圆会坍塌，无限刚度仍然存在，尽管它从无先前条件的情况中显著降低。 在超声速流动中，声学干扰在相对于流动的马赫角度方向对流的两个点上坍塌，速度等于流速q，即完美的预处理。 从分析上看，我们发现M < 1的K g = l/\textbackslash\{\}/I - M 2，M > 1的K g = 1。 显然，局部预处理显著改善了时间尺度的刚度。 Turkel[36]是第一批推广Chorin方法的研究人员之一。 在他的分析中，他提出了利用原始变量p、u、v加上熵5或密度p的预处理。 在对称化中

DARMOFAL \& VAN LEER (a) M = 0.1 (b) M = 0.5 (c) M = 0.9 (d) M = 1.3 图3 最佳预处理, P,,t（包括P,l,）的波前。 波前已按最大波速缩放。变量，Turkel预调节器的一般形式是，其中at和pt是自由参数。 Turkel的预调节器最广泛地用于消除低马赫数的刚度[38]。 在M + 0的极限中，如果在+1和pt + M，那么ng -+ 1。 Turkel预调节器也可以扩展，以在所有马赫数上产生与Van Leer-Lee-Roe预调节器类似的最佳波前。 对于亚声速流动，最佳波前是通过pt = M和at = 1+p实现的：[17]。 不幸的是，当与上风或矩阵耗散方案结合使用时，Pt的最佳形式会产生离散特征值，fd到不稳定的右半部分

局部预处理 219平面用于更高的亚音速马赫数[7，23]。 因此，这种特定形式的预调节器在实践中不能用于跨声速流动。 为了避免这个问题，Turkel等人[38]定义了t和fS t，以返回截止马赫数上方的无预处理情况。 虽然众所周知，P vW产生的波前是对称预调制剂的最佳波前，但没有证据表明该波锋最优地扩展到一般非对称预调制剂。 对于超声速流动，尽管可能存在其他选择，但由于K g = 1，这些波前可以公正地称为最佳。 在亚声速流动中，我们留下了改进的可能性。 然而，所有证据都表明，这些波前在亚声速流动中也是最佳的。 目前，当预调节器产生与P vlr相同的波前时，我们将将其标记为最佳。 由于P vlr和Pt都可以优化，因此最佳预处理显然不是唯一的。 预调节器非唯一性的另一个例证是，给定对称化变量的预调节器P，当预处理通过其移位P T找到相同的特征值。 因此，给定一个最佳的预调节器，它的转置也是最佳的。 最近，我们在推导所有最佳预调节剂的一般形式方面取得了重大进展。 具体来说，存在两个最佳预调节器家族。 一个家庭与Van Leer-Lee-Roe预调节器关系最密切。 这个最优家族的逆数是，A MA - 4-\textbackslash\{\}/l - M 2C2 C 0 G M(C + G) - 0 0 0 0 1，其中可以选择A、C和G来满足其他设计约束。 显然，这种先置调节器的转置也是最佳的。 我们注意到，此先定子的对称形式由C = G = 0和A = (P + 1)/M 2给出，并得出Van Leer-Lee-Roe先定子。 一个有趣的一点是，最佳Turkel矩阵不是这个家族的成员，因此必须是另一个最佳预调节器家族的成员。 不幸的是，目前没有针对这个第二族最优矩阵的干净的代数描述。 其他已开发的预处理器包括D.Choi-Merkle无黏性预制剂]、Y.Choi-Merkle黏性预制剂[5]和Weiss-Smith预制剂]。 这些预处理器是为低马赫数流量开发的。 改造预处理器

220 DARMOFAL \& VAN LEER到对称化变量给出e 0 0 -5 cmve ' 0 10 0 fcm/ws \textasciitilde{} 0 0 10' \_ 8cmi(e -l)00 1，其中e是用户定义的参数。 对于D.Choi-Merkle inviscid预调节器，5 cmi = 1和5 cmv = 0。 对于Y.Choi-Merkle黏性预处理剂，5 cmi = 0和5 cmv = 1。 而且，对于Weiss-Smith预调节器，S cmi = S cmv = 0。 事实上，Weiss-Smith预处理器是Turkel预处理器家族的成员（具体来说，a t = 0和f\}t = v/£）。 对于亚音速无黏性流动，推荐选择是e = M 2。 这三个预处理的特征值是相同的，由A M = i«(l + e) cos 4> T \textasciicircum{}i g 2 (1 \_e ) 2co s 2 0 +e c 2 | \textasciicircum{} rsubgo. c c \textasciicircum{} \textasciicircum{} A2,3 = qcosrf) J给出的波前分析表明，在e = M2的低马赫数中，K 9 -> (VE + l)/(y/E - 1)。 因此，虽然Choi-Merkle/Weiss-Smith预调节器在低马赫数下不能实现最佳调节，但它确实消除了Euler 方程在低速时的无限刚度。 在较高的马赫数下，预调节器会恢复到未预制约Euler 方程的特征值，尽管在约0.4的马赫以上，这些预调节器的波速调节比未预制约Euler 方程更差。 请注意，Choi-Merkle inviscid和Weiss-Smith预调节器成为恒等矩阵，如M ->• 1，但ChoiiMerkle黏性预调节器不适用于rupersonii iiows，我们ctssum.6，e = 1，因此特征值返回到没有预处理的Euler 方程的特征值。 这些结论在图中所示的波前和«值中显而易见。 4和2。 局部预处理的一个更经典的建议是使用由上风或基质耗散离散诱导的Jacobi块[26，27，2，3，29]。 如果使用上风方案来离散化空间算子，则细胞面的通量可以写成，/=2 \{fright + futi) \textasciitilde{} \textbackslash\{\} \textbackslash\{\}Ac\textbackslash\{\} (U right \textasciitilde{} U,«/,), (13.3) 其中A c是保守通量的Jacobian 矩阵的面值，li cI具有与A c相同的特征向量，但其特征值是A c的特征值的绝对值。 块-雅各布预调节器是线性化空间系数矩阵中对角线上发现的块

I,OCAL 预处理 221 (a) \&I = 0.1 (c) M = 0.9 (b) M =0.6 (d) M = 3..3 图4 Choi-Merkle和Weiss-Smith 预处理的波前, PCml,,,, with E -. min(M" 11)。 Wavefronits已按最大波速进行缩放。操作员。 Wr矩形单元格，这个预制因素由其中o是一个因子给出，该因子可用于根据特定的离散化缩放空间足迹的大小。 时间步骤，At，，，，，，为了dirnension一致性而包含，但在实施后，用时间步骤中使用的At，，，取消。 块-Jacobi预调节器的主要优势是其出色的离散特征值聚类（见第13.2.3节），mc1其鲁棒性（见第13.2.4节）。 不幸的是，block-Jacobi预调节器对提高长波传播速度的作用不大。 Pj的波前如图所示。 5. 在低马赫数下，熵模式哈。？ 已被重新定位到（1,0），然而，涡量模式仍然位于原点附近，并且，作为M -+ 0，刚度仍然是无限的。 阿特

DARMOFAL \& VAN LEER (a) M = 0.1............................................. OJ +..... 我.................. > OJ.................................................... “'..”..................... 9）一个.................. 4................... +.................. j.................. +.................. -2 -0.6 0 01 (b) M = 0.5 (c) M = 0.9........................ a i....................................;................................................................................. OJ (. i 4 o*.................. +.................. ji............ 01.................. 6.................. 1........................................ > 0:.:............................................... 4 * 2 -O*..................j................... i............ a..................... ）................................. >................................................................ * - -I 48 0 01 (d) M = 1.3 图5块-Jacobi 预处理的波前，Pj。 波前已被最大的波速缩放。更高的马赫数，块-雅各布系统通过增加上游声波的速度来提高刚度。 虽然在低马赫数下不明显\textasciitilde{}IB，但涡量模式的波前现在可以清楚地看到是非物理的，呈三角形。 相比之下，图中的最佳波前。 3看起来相当合理--这是基于连续体设计的直接结果。 无论如何，虽然最佳预调节器比Euler 方程的波传播有所改善，但在更高的马赫数下，它们并没有比块-雅各布的实质性改进。 13.2.3 离散特征值聚类 虽然预制条件方程的长波属性可以从连续偏微分方程中分析，但短波需要分析离散方程，因为这些小波长通常不是

局部预处理 223由数值方案很好地代表。 分析方案离散属性的常见方法是通过离散空间算子的傅里叶足迹，离散空间算子是其特征值在复平面中的位置。 具体来说，在关于恒定状态的线性化和离散傅里叶模式的替换后，Uj\textasciicircum{} = ue'\textasciicircum{} Sl+k8"\textbackslash\{\}对于任何离散波数组合，（8 X，6 y）都可以找到数值近似的光谱放大矩阵。 离散波数的范围从-n到n。 使用傅里叶足迹来分析预制欧拉和Navier-Stokes 方程的平滑能力，可以在其他几个参考文献中找到[42、19、2、9、20、3、29]。 收敛加速的一个目标是在数值上尽可能快地抑制傅里叶足迹的特征模式。 例如，一个有效的多网格算法需要对随后的粗网格上不存在的所有模式进行有效平滑。 对于全粗多网格，这需要高效平滑高低、低高和高高频组合，其中高频波数范围为7r/2 < \textbackslash\{\}6\textbackslash\{\} < n。 对于半粗化，只有高频元件需要有效的平滑。 正如我们将展示的那样，局部预处理通常可以聚集傅里叶足迹，从而显著增强阻尼。 特别是，我们专注于高频位点，并表明如果没有适当的预处理，许多特征模式的阻尼效果不佳。 当将Van Leer-Lee-Roe、Turkel或Choi-Merkle/Weiss-Smith等连续体预调节器应用于上风方案时，必须修改耗散，使/ = g (fright + fie ft) \textasciitilde{} -M' 1 P' 1 | P (A COS <j> + B sin 4>) | M \{Urigkt \textasciitilde{} Uleff), (13.4) 其中M是从保守状态向量到对称变量的变换矩阵，<f>是网格法线相对于流方向的角度。 这种修改对于稳定性以及提高离散化在低马赫值下的准确性是必要的[40，17]。 对于Van Leer-Lee-Roe先决条件，这个修改后的耗散矩阵太复杂了，无法得出一个简单的分析表示；因此，通常的做法是进一步修改耗散，以/ = 2 (/"« fc< + heft) - \textasciicircum{}M\textasciicircum{}P- 1 (\textbackslash\{\}PA \textbackslash\{\}coscj>\textbackslash\{\}+ Fij(8iii\textasciicircum{}|)M(urijAt-«i e /«).( 13.6）在低马赫数下，上风离散Euler 方程的高高频足迹非常不集群。 图6显示了M = 0.1的傅里叶足迹的高频含量，细胞长宽比为JR = Ax/Ay = 1，流量与网格对齐（</> = 0°）。

224 DARMOFAL \& VAN LEER 作为说明，扩增因子轮廓也叠加了Lynn[18]推导的最佳两级多级阻尼方案。 如果没有预处理，可以看到一组高频傅里叶模式聚集在原点附近，放大因子必须接近一个。 这种聚类是Euler 方程在低马赫数下不同传播速度的直接结果。 相比之下，Van Leer-Lee-Roe预调节器的特征值与原点相去甚。 虽然聚类已经改进，但只需对原始预处理器进行小幅更改，就可以实现更好的聚类。 我们注意到，在原始的P vw中观察到了两个不同的特征值集群。 事实上，这些特征值家族是由于声波和对流波的存在。 Lynn[18]表明，通过缩小声波的传播速度，傅里叶足迹可以完美地聚集于最高频率的波，6 X = 0 y = ±7r。 这种缩放由Lee[15, 14]开发，可以通过将r重新定义为r = fi/(JR q + 0)来实现，其中M \_ T,i=iW\textasciicircum{}k + vAy k\textbackslash\{\} Y,t=i \textbackslash\{\}vAx k - uAy k\textbackslash\{\} 物理上，Al q可以被解释为在流动方向上定义的细胞长宽比，即流向与正常细胞长度的比率。 使用此定义，所有马赫和细胞长宽比组合的高高频模式都聚集在复平面的同一位置。 对于半粗化多网格来说，这种改进的聚类是必不可少的，因为细胞谱比可能非常大。 改进的聚类可以在图中验证。 6. 图中显示了相同条件下但具有细胞长宽比的高频位置，JR = 100。 7. 在这种高长宽比下，与未预制件的欧拉和原始的范利尔预制件相比，改进是巨大的，并证明了r修改的合理性。 Allmaras[2]对block-Jacobi预置剂进行了广泛的分析，发现它在细胞-长宽比、马赫数和流角的所有范围内都具有出色的高频聚类。 块-Jacobi足迹已包含在图中。 6和7。 需要注意的一个重要区别是，块-Jacobi预调器有一个围绕实轴聚集的模式家族。 这种效应也与低马赫数下传播速度的巨大差异有关，而块-雅各布预调节器无法缓解这种差异。 相比之下，Van Leer-Lee-Roe预处理器将这些模式从真实轴上传播开来，使它们与其他错误模式更好地聚类。 13.2.4 健康的特征向量结构 预调节器设计的一个非常重要的方面是保持健康的特征向量结构。 首次研究特征向量的重要性

X，OCAL，PItECONDITIONTNG（a）No 预处理（c）Van Leer与原始T f b）block-Jacobi图6 M = 0.1，= 1，b，= 0的高傅立叶足迹。 放大系数在0到1之间的轮廓，增量为0.1。预凝结器设计由Darmofal和SchmicL[7]执行。 在这项工作中，他们展示了特征向量正交性如何与瞬态放大11的可能性直接相关，即使在dl eigenvdues影响稳定性的系统中，也会出现小扰动。 具体来说，consiler以下线性问题的演化，du -- -t iLu = 0，clt，其中L是N x N矩阵。 在预处理中，这将代表傅里叶变换，先入；ionecI空间算子。 上述eq\textasciitilde{}tstion的解可以使用矩阵指数来紧凑地写成，u(t) = exp (-itL) uo，

DARMOFAL和VAN LEER（a）无预处理（c）Van Leer与原始T（b）块-J-bi（d）Van Leer与修改7图7高高傅里叶足迹，M = 0.1，AR = 100，4 = 0'。 放大系数在0和1之间，增量为0.1的轮廓被叠加。其中uo是初始条件。 给定时间的最大能量放大G(t)可以表示为矩阵指数的规范，其中R是L的右特征向量的矩阵，Cl是特征值的对角矩阵。 放大矩阵可以以所有时间为界，其中特征向量矩阵的条件数K(R)定义为K(R) = IIRIIIIR'lII，wmax是所有特征值的最大正虚数部分。 对于Euler 方程等双曲系统，后者通常是零，因为特征值是严格的红色。 因此，

局部预处理 227最大放大的约束由K(R)给出。 对于正交特征向量结构，K(R) = 1，能量增长是不可能的。 与正交性的更大偏差会导致更大的K(R)和更大的可能的瞬态增长。 如果没有预处理，对称变量是一个正交系统，因为矩阵A和B是对称的。 然而，在预处理之后，PA和PB不再保证对称；因此，一般来说，预处理引入了以前不存在的短暂增长的可能性。 不幸的是，Darmofal和Schmid发现，基于连续体的预调节剂，如Turkel、Van Leer-Lee-Roe和Choi-Merkle，都会受到M -y 0的严重瞬态生长。 事实上，他们发现，t = 0 +的非维化能量增长率与这些预调节器的1/M成正比，因为M -► 0。 此前，许多研究人员发现停滞点极度缺乏稳健性；然而，这种不稳定的确切原因尚不清楚。 例如，Turkel[38]在他的（3 t定义中使用了一个截止值，类似于/3 2 = max（M 2，rjM2eA，其中M Tef是自由流马赫数的一些引用，rj是用户定义的参数，通常在1到3之间。 这有效地避免了j3 t归零，Darmofal和Schmid表明，这限制了非正交性。 然而，这种修复取决于问题，此外，通过引入非局部参数，使预调节器的分析复杂化。 改进先入条件方程的特征向量结构的另一个途径是从设计开始强制执行一些正交度量。 最严格的要求是在所有条件下强制执行正交性。 这相当于要求PA和PB对称。 要求完全正交的预调节器太有限，无法显著改变长波时间尺度。 一个较弱的限制是只要求沿溪流方向进行正交。 Lee和Van Leer[16, 14]遵循了这种方法，并发现了一个具有正交特征向量的最佳亚音速前调制剂。 具体来说，这个预调节器的形式是-Pt>！ 96 = 哪里 2M 2 -/? * \textasciitilde{} 1+ M 2 ' 预调节器被创造为P vi9e，即 Van Leer 96，由Van Leer在1996年夏天开发。 此预调节器以M -y 1的形式与Van Leer-Lee-Roe预调节器顺利连接，因此，f -f / 0 0 -f |+i o 0 0 0/3 0 0 0 0 0 1

228 DARMOFAL和VAN LEER超音速Van Leer预调谐器可用于M > 1。 注意，P„(g 6是P opt家族的先定元；此外，P„06是P opt的唯一正定成员，在溪流方向上是正交的。 13.2.5 精度保持作为M ->■ 0，许多可压缩有限体积技术的一个众所周知的问题是，除了缓慢的收敛[43]外，精度的极度损失。 这个问题可以追溯到低马赫数时不平衡的耗散矩阵[8]。 在这项分析中，对称化的Euler 方程首先通过参考密度和流速进行非维度化。 Gustafsson和Stoor[12]表明，使用这种非维度化，-| = O(M)，因为马赫数会减少。 因此，整个对称状态向量的行为为w = [0(M), 0(1), 0(1), 0(1)] T。 同样，A和B矩阵被发现为以下顺序，A = 0(1) 0(1/M) 0 0 0(1/M) 0(1) 0 0 0 0 0 0 0(1) 0 0 0(1) 0 0 0 0 0(1), B = 0 0 0 0 0(1/M) 0 0 0 0 0 0 0 0 0 结合这些结果很容易表明，对于一阶上风或矩阵耗散方案，类似的分析表明\textasciicircum{}lAI-\textasciicircum{} + A\textasciicircum{}I - [0(1), 0(1/M), 0(1/M), 0(1)]' 因此，除了熵方程外，欧拉和耗散项不匹配为M ->• 0。 当应用局部预调节器时，可以在预调节通量和耗散条件之间实现平衡。 Fiterman等人[8]表明，如果P \_1,P- 1 |PA|,P- 1 |P.B| 0(1/M 2 ) 0(1/M) 0(1/M) 0(1/M) 0(1) 0(1) 0(1) 0(1/M) 0(1) 0(1) 0(1/M) 0(1) 0(1) 0(1) 0(1) 0(1) 0(1) 0(1) 会发生适当的匹配。快速检查表明，为低速收敛加速设计的预调节器通常具有这种精度属性。 具体来说，P„j r，P«，

局部预处理 229 Pcm和Pt,;96都有适当的限制，如M ->■ 0，而block-Jacobi 预处理没有（block-Jacobi 预处理甚至不需要修改耗散）。 最后，我们还注意到，Reed[30，31]进行了更具限制性（但也更直接）的截断误差分析，并表明Choi-Merkle先条件方程的截断误差比某些项没有预处理的截断误差低M 2的系数。 无论如何，总体结果保持不变--局部预处理与相关的修改耗散可以解决常见计算方法的低马赫数精度问题。 13.2.6 分离性预处理在试图将方程分离成独立的块时也很有用，如果时间导数被丢弃，这些方块将保持双曲或变成椭圆。 例如，给定稳定的Euler 方程，，dw du>这些方程可以转换为独立的双曲和/或椭圆部分。 对于超声速流动，稳定的Euler 方程可以分为一组四个对流（即双曲）方程，其中5和H是熵和焓。 前两个方程是沿马赫线的声学扰动对流，而后两个方程是沿流线的焓和熵对流。 在亚声速流动中，焓和熵的对流仍然存在；然而，声子系统现在是椭圆的，等同于Cauchy-Riemann 方程，0 0 0 0" ±1 ' fidp+ pudv \textbackslash\{\} ' 1 0 0 0" ±i ' (3 dp + pu dv 0 p 0 0 ±1 j3 dp - pu dv y 0 0-1 0 0 ±i (3 dp - pu dv 0 0 0 10 1. *l v s y 0 0 0 0 0 M s 1 0 0 0 0 0 " 0 0-1 0 d 0 0 1 ( -P-dp \textbackslash\{\} I pu * H S 0 -1 0 0 0 0 0 0 0 0 0 0 0 0 0 0 0 0 0 0 0 dr] (dr] (dp \textbackslash\{\} r H S = 0. 这种分割的优点是，不同的离散化方法现在可以用于椭圆和双曲部分。 对于不稳定的Euler 方程来说，这种分离是不可能的。 然而，引入局部预调节器提供了分离双曲和椭圆原点的自由，同时保持时间导数并允许使用时间行进算法。

230 DARMOFAL \& VAN LEER Roe [33]开发了一种执行椭圆和双曲部分分离的一般策略。 对于2-D Euler 方程，允许这种分离的预调节器的一般形式是，sep a3p CL2P -a4M 0 0 0 1，其中a*是自由参数。 将此结果与最佳预制件家族P opt相结合，我们可以找到允许椭圆/双波分离的预制件家族，具有最佳波前，sep/opt l+a.i9 2 J\_ M，V1-0 2-；1 M 0 1 0 0 0 a，0 0 0 1 + M，其中a*是一个自由参数。 特别是，我们注意到Van Leer-Lee-Roe矩阵是P sep/opt家族的成员（a* = 1//？）。 基于Van Leer-Lee-Roe预调节器的双曲/椭圆分裂已被用于制定稳定流动的真正多维上风方案[25，24，28]。 与标准网格对齐的上风方法相比，这些方案提高了准确性；然而，它们似乎也容易受到扰动的瞬态放大。 13.3 当前状态 在上一节中，我们重点介绍了本地预处理器的设计，强调了可能选择的不同设计标准。 在本节中，我们将重点介绍局部预处理的现状--具体来说，局部预处理在实践中取得了什么成就？ 13.3.1 长波的刚度去除原则上，长波的不同时间尺度的去除应该会导致加速收敛，因为溶液中的所有误差现在都可以以相同的速度传播，而不会留下缓慢传播的误差模式。 这种刚度的去除首先在低马赫数流动中得到了证明。 Merkle等人[4, 5]使用交替方向隐式算法对各种低速内部流动进行了显著的收敛加速

局部预处理 231 中心差分。 局部预处理在所有马赫数字中的好处的第一个例子是Van Leer-Lee-Roe预调节器[40]。 在这项工作中，在NACA 0012 翼型上获得了明确的上风解，表明在所有情况下，与使用无预处理的欧拉方案相比，应用P„; r会导致显著加快速度。 随后，许多作者调查人员发现了类似的结果[15、10、9、30、14、31]。 在本节中，我们将演示在M\textasciicircum{} = 0.1、0.3、0.5和0.8时对凸起流动的单网格计算的收敛加速度效应。 我们只会测试两个预调节器，P vir和Pj，并将结果与无预调节的欧拉计算进行比较。 基本流动求解器由高解析度的上风方案[39]和Roe的[32]近似Riemann 求解器组成。 时间整合是用Lynn[18]的2级最佳阻尼方案进行的。 如图所示。 8，Moo =0.1和0.8情况的收敛率低于无先发制定欧拉计算的两个中等马赫数。 使用block-Jacobi 预处理，除了MQO = 0.1的情况外，所有情况都收敛得非常相似，在大约1000个周期内下降了大约六个数量级（每个周期包含4个多阶段积分，因此这实际上是4000个多阶段迭代）。 然而，相比之下，Van Leer-Lee-Roe预处理器的性能几乎独立于马赫数，在所有情况下，在大约1000个周期内下降六个数量级。 与未预处理的欧拉结果相比，Van Leer-Lee-Roe预调节器大约快1.5到3.0倍（就收敛六个顺序的周期而言）。 这些结果与第13.2.2节中讨论的波前分析非常一致，该分析表明，块-Jacobi 预处理不有助于低马赫数下的长波传播。 因此，一个设计得当的局部预调节器可以成功地消除长波时间尺度，在实践中，这种效果加速了单网格计算的收敛。 13.3.2 改进的多重网格收敛除了加速单网格方法外，局部预处理还可以对多网格方法产生有利的影响。 事实上，局部预处理在多网格方法中使用时具有双重好处：长波被加速，从而改善了精细网格求解器，空间算子的离散特征值被聚类，从而改善了高频误差模式的平滑。 许多作者展示了局部预处理对多网格性能的好处[34、19、20、18、29、38、21、13、23、22]。 这一领域的一个重要贡献是Tai[34]在1-D和Lynn[18]在2-D的工作，他们表明局部预处理和多网格的组合可以比独立应用的任何一种技术更好地加速收敛。 为了证明这一结果，我们对之前的颠簸流测试案例应用了全和半粗化多网格[26，27]方法。

DARMOFAL \& VAN LEER i 4 O \textasciitilde{}WOOzPmOmI\#O (a) No 预处理 (b) block-Jacobi (c) Van Leer-LeRoe Finre 8 马赫数 独立于收敛的非限制颠簸流，具有64 x 32 ccells的单个网格。 实体：M，= 0.1，破折号：M，= 0.3，破折号：MEQ = 0.5，圆：M，= 0.8。 不，每个周期包含4个twustage intererations。 收敛历史如图所示。 9和10分别用于全粗化md半粗化。 比较非预制和Van Leer-Lee-Roe预制结果，我们看到预制结果至少比非预制结果快五倍，回顾在没有多网格的情况下，观察到的加速在1.5到3.0之间，我们得出结论，多网格和局部预处理的组合增强了两种加速技术的性能。 一个有趣的观察是，对于M, = 0.1案例，block-Jacobi预调节器明显比Van her-Lee-Roe预调节器慢。 这种差异可以归因于Van Leer-Lee-Roe预调制器的长波传播的改进。 图11是解决方案中误差的等高线图。 误差是通过先求解机器精度的方程来计算的。 然后，重新运行模拟和

局部预处理 J m loo IW wm (a) No 预处理 J m IW 60 lm 160 w- wm (b) block-Jacobi (c) Van Leer-Lee-Roe..,. 图9 马赫数 收敛的独立性，具有全粗化muitigrid的非限制凸起流。 细网格64 x 32个单元格；粗网格8 x 4个单元格。 实体：Mw = 0.1，破折号：Mw = 0.3，破折号：Mw = 0.5，圆：M，= 0.8。误差是当前解和最终解之间的差值的大小。 最初，误差分布在凸起上。 在前三个周期之后，block-Jacobi和Van Leer-Lee-Roe的误差相似，两者都沿着凸起的实心壁和下游表现出网格对齐的误差。 在随后的周期中，Van Leer-Lee-Roe预调节器将这些网格对齐的错误模式传播到计算域之外，而block-Jacobi预调节器则不传播。 因此，块-Jacobi预调节器必须依靠平滑来消除这些结果；然而，由于这些模式是网格对齐的，它们无法平滑，并受到融合的影响。

\& VAN LEER -,! -\textasciitilde{}----d--.----.------ J..,L I- J M\textasciitilde{}OI IW 50美元 100 160 cvaat (b) bfock-Jacobi (c) V? I.II Leer-Lee-Roe图10 Mach 11u1\textasciitilde{}1ber独立性，用于converg.cnc:e的rilconfinetl bump f ow与半粗化mnltigrid。 细网格64 x 32个单元格；粗网格8 x 4个单元格。 实体：M，= 0.1，破折号：= 0.3，破折号：M- = 0.5，Cjircles：M，= 0.8。 13.3.3 低马赫数精度 如第13.2.6节所述，由于tllodified耗散矩阵，局部预处理可以提高低马赫数流量计算的准确性。 这种效应首先被观察到\textasciitilde{}ill\textasciitilde{}lzerically 1.q' Van Leer等人[40]，以及，sulseque\textasciitilde{}ltly，l>y 1-na.ziy otller调查人员[37，38，311。 改进的ticcuracy i1.t低马赫数的一个典型例子是A/r，-0.01流a.rour\textasciitilde{}d a。 NACA 0012 翼型在att.a.ck的l.25' a.ngle：如猪所示。 12. 在这些图中，我们将unpreconc:litioned zmtl TiVciss- Snlitli preconditios\textasciitilde{}ed Euler simulztztions与不可压缩的pa.rxef方法进行了比较。 从tlle Cr co\textasciitilde{}ltours和sl\textasciitilde{}rf:i(:e图中都可以看出，61.1（5个未经预先的结果存在显著性\textasciitilde{}lt；不准确。 EIo\textbackslash\{\}vc！ Vc！ R，grecouclitioned结果有ciean Cr轮廓a.和表面压力\textasciitilde{}natclz

局部预处理 235（a）初始条件。 16个误差等高线从0到0.015。 （B）3个周期。 16个误差等高线从0到0.002。 （C）6个周期。 16个误差等高线从0到0.002。 （D）9个周期。 16个误差等高线从0到0.002。 （E）12个周期。 16个误差等高线从0到0.002。 图11使用block-Jacobi和Van Leer预调节器的M„ = 0.1颠簸流的全粗化多重网格收敛的误差轮廓。

236 DARMOFAL \& VAN LEER（a）无预处理 Cp轮廓（c）预处理的Euler C p轮廓（b）表面Cp（d）表面C p图12未预处理的Euler和Weiss-Smith预处理解决方案的低马赫数精度的比较以及与面板解决方案的比较。 NACA 0012，M = 0.01，a = 1.25度。 31 C p等高线从-0.7到0.7绘制。用面板计算。 13.4 未来发展 虽然局部预处理具有重大前景，但唯一似乎在一系列问题上稳健的预处理方法是block-Jacobi。 然而，块-雅各布方法在低马赫数下不能提供有效的收敛加速度或良好的精度。 因此，必须解决的一个主要绊脚石是，马赫数性能良好的低预调节器在停滞点缺乏坚固性。 在工作期间

Darmofal \& Schmid[7]和Van Leer等人[41]的局部预处理 237对稳健性问题的原因有了更好的理解，仍然缺乏充分的解决方案。 虽然我们专注于非黏性流动，但在高雷诺数黏性流量计算中，收敛加速度的需求最强烈。 在将局部预处理扩展和应用到离散化的Navier-Stokes 方程[5、10、2、30、18、44、29、13、14、21]方面已经完成了大量工作。 还需要进一步开发，将结果扩展到其他方程组，如理想磁流体动力学方程组。 致谢 上述工作只有在许多人的合作下才能完成。 具体来说，作者要感谢Dohyung Lee、John Lynn、Barrett McCann和Phil Roe的贡献。 第一作者希望通过NSF职业奖（ACS-9702435）和波音公司感谢NSF的支持。 参考资料1。 S。 Abarbanel和D。 戈特利布。 具有混合导数的二维和三维Navier-Stokes 方程的最佳分割。 计算物理学杂志，41:1-33，1981年。 （另外，ICASE报告编号 80-6，1980）。 2. S.R. 奥尔玛拉斯。 2-D Navier-Stokes 方程的局部矩阵预调节器分析。 AIAA论文93-3330，1993年。 3. S.R. 奥尔玛拉斯。 2-d Navier-Stokes 方程多网格解决方案的半隐式预调节器分析。 AIAA论文95-1651，1995年。 4. D。 Choi和C.L. 默克尔。 将时间迭代方案应用于不可压缩流动。 AIAA杂志，23（10）：1518-1524，1985年。 5. Y.H. Choi和C.L. 默克尔。 预处理在黏性流动中的应用。 计算物理学杂志，105：203-223，1993年。 6. A. J。 合唱。 解决不可压缩黏性流动问题的数值方法。 计算物理学杂志，2:12-26，1967年。 7. D.L. Darmofal和P.J. 施密德。 特征向量对Euler 方程局部预调节器的重要性。 计算物理学杂志，127：346-362，1996年。 8. A. 菲特曼，E。 Turkel和V.N. 瓦萨。 稳态流体方程的压力更新方法。 AIAA论文95-1652，1995年。 9. 交流 戈弗雷。 迈向强大的预处理的步骤。 AIAA论文94-0520，1995年。 10. 交流 戈弗雷，R.W. Walters和B. van Leer。 带有有限速率化学的Navier-Stokes方程的预处理。 AIAA论文93-0535，1993年。 11. S.K. 戈多诺夫。 一类有趣的准线性系统。 多克尔。 阿卡德。 Nauk SSSR，139:521-523，1961年。 12. B. 古斯塔夫森和H。 斯托尔。 Navier-Stokes 方程几乎不可压缩

238 DARMOFAL和VAN LEER流量。 SIAM数值分析杂志，28（6）：1523-1547，1991年。 13. D。 杰斯珀森，T。 Pulliam和P。 布宁。 最近对OVERFLOW的增强功能。 AIAA论文97-0644，1997年。 14. D。 李。 欧拉的局部预处理和Navier-Stokes 方程。 密歇根大学博士论文，1996年。 15. D。 Lee和B. van Leer。 Euler和Navier-Stokes 方程的局部预处理的进展。 AIAA论文93-3328，1993年。 16. D。 Lee、B. van Leer和J. 林恩。 所有马赫和单元雷诺数的本地Navier-Stokes预调制器。 AIAA论文97-2024，1997年。 17. W.T. 李。 Euler 方程的局部预处理。 博士论文，密歇根大学，1991年。 18. J。 林恩。 Euler 方程与局部预处理的多网格解决方案。 密歇根大学博士论文，1995年。 19. J.F. Lynn和B. van Leer。 具有最佳平滑的欧拉和Navier-Stokes 方程的多阶段方案。 AIAA论文93-3355，1993年。 20. J.F. Lynn和B. van Leer。 欧拉和Navier-Stokes 方程与局部预处理的半粗化多网格求解器。 AIAA论文95-1667，1995年。 21. D。 马夫里普利斯。 各向异性非结构网格上的黏性流动求解器的多网格策略。 AIAA论文97-1952，1997年。 22. B. 麦肯。 可压缩Euler 方程多网格解决方案的局部预调节器的评估。 德克萨斯农工大学硕士论文，1996年。 23. B. McCann和D.L Darmofal。 可压缩Euler 方程多网格解决方案的局部预调节器的评估。 AIAA论文97-2028，1997年。 24. L.M. 梅萨罗斯。 非结构化网格上Euler 方程的多维波动分裂方案。 密歇根大学博士论文，1995年。 25. L.M. 梅萨罗斯和P.L. Roe 基于分解方法的多维波动分裂方案。 AIAA论文95-1699，1995年。 26. W.A. 穆德。 解决对流问题的新方法。 《计算物理》杂志，83:303-323，1989年。 27. W.A. 穆德。 基于多网格、半粗化和缺陷校正的高分辨率欧拉求解器。 计算物理学杂志，100：91-104，1992年。 28. H。 Paillere，H。 Deconinck和P.L. Roe 基于Euler 方程的稳定特征的保守上风残余分配方案。 AIAA Paper 95-1700，1995年。 29. N.A. Pierce和M.B. 吉尔斯。 预处理 可压缩流动拉伸网格的计算。 AIAA论文96-0889，1996年。 30. C.L. Lee 低速预处理应用于可压缩Navier-Stokes 方程。 博士论文，德克萨斯大学阿灵顿分校，1995年。 31. C.L. 里德和D.A. 安德森。 将低速预处理应用于可压缩Navier-Stokes 方程。 AIAA论文97-0873，1997年。 32. P.L. Roe 近似黎曼求解器、参数向量和差分方案。 计算物理学杂志，43:357-372，1981年。 33. P.L. Roe 由许多简单问题组成的：对CFD中模型问题作用的反思。 在V中。 Venkatakrishnan，医学博士 萨拉斯和S。 Chakravarthy，编辑，计算流体动力学中的障碍和挑战，第241-258页。 Kluwer学术出版社，1998年。 34. C. H。 泰。 显式代码的加速技术。 密歇根大学博士论文，1990年。 35. E。 特克尔。 流体动力学矩阵的对称化与应用。 计算数学，27：729-736，1973年。

局部预处理 239 36. E。 特克尔。 求解不可压缩和低速可压缩方程的预处理方法。 计算物理学杂志，72：277-298，1987年。 37. E。 图尔克尔，A。 Fiterman和B. van Leer。 预处理和有限差分方案的不可压缩流动方程的极限。 在M。 哈菲兹和D.A. Caughey，编辑，计算未来：计算空气动力学的进步和前景，第215-234页。 John Wiley and Sons，1994年。 38. E。 土耳其，V.N. Vatsa和R。 雷德斯皮尔。 预处理低速流动的方法。 AIAA论文96-2460，1996年。 39. B.范·莱尔。 由Euler 方程控制的空气动力学问题的上风差分方法。 应用数学讲座，1985年22日。 40. B. van Leer，W.T. 李和P.L. Roe 特征时间步进或Euler 方程的局部预处理。 AIAA论文91-1552，1991年。 41. B. van Leer，L. 梅萨罗斯，C.H. Tai和E。 特克尔。 局部预处理在驻点中。 AIAA论文95-1654，1995年。 42. B. van Leer，C.H. Tai和K.G. 鲍威尔。 为Euler 方程设计最佳平滑多阶段方案。 AIAA论文89-1933，1989年。 43. G。 Volpe。 关于低马赫数字下可压缩流动代码的使用和准确性。 AIAA文件91-1662，1991年。 44. J.M. Weiss和W.A. 史密斯。 预处理适用于可变和恒定密度流量。 AIAA杂志，33（ll）：2050-2057，1995年。

\clearpage
\chapter{放松重温--重新审视稳定流动的多网格}
放松重温-重新审视多网格以稳定流动Thomas W. Roberts 1，David Sidilkover 2和R。 C. 斯旺森 1 14.1 导言 1971年出版了一篇具有里程碑意义的计算空气动力学论文，即默曼和科尔[9]。 与许多开创性作品一样，它的意义不在于它所解决的具体问题--小干扰，平面跨声速流动--而在于确定解决技术上重要且理论上困难问题的一般方法。 Murman和Cole工作的主要特点是使用类型依赖微分来正确解释混合椭圆/交错方程的适当依赖域，并引入线松弛来求解稳流方程。 跨声速电位流的所有后续工作都基于这些概念。 Jameson [6]透过两个重要贡献将Murman和Cole的想法扩充套件到全势方程。 首先，他引入了旋转差分模板，该模板将Murman和Cole类型依赖差分算子推广到一般坐标。 其次，他使用了Garabedian引入的解释，将放松作为人工时间的迭代，给1研究科学家，空气动力学和声学方法分部，美国宇航局兰利研究中心，汉普顿，弗吉尼亚州23681-0001 2高级职员科学家，ICASE，汉普顿，弗吉尼亚州23681-0001 计算流体动力学的前沿- 1998 编辑：David A. Caughey和Mohamed M. 哈菲兹 ©1998 世界科学

242 ROBERTS、SIDILKOVER和SWANSON构建了稳定的松弛方案，推广了参考[9]的原始线松弛方法。 20世纪70年代的十年见证了跨声速电位流解法的爆炸性活动，Caughey[4]的评论文章总结了这一问题。 大约在Caughey调查的时候，计算空气动力学研究的主要重点是离开全势方程，转向稳定的Euler 方程的解决方案。 通过与势流方程的松弛方法类比，Euler 方程的解方法可以被认为是伪时间的迭代。 与势流方程的方法不同，到目前为止，最常见的方法是将稳定流作为非稳态方程的长期渐近解求解，而不是直接求解稳态方程。 从概念上讲，这要容易得多，因为不稳定的Euler 方程是一个双曲系统，存在大量的理论。 放弃时间的准确性允许在迭代方案的构建中具有相当大的灵活性。 为了将收敛加速到稳态，使用了各种类型的预处理。 此外，从非定位方程开始，将迭代方法直接扩展到雷诺平均Navier-Stokes 方程，近年来在这方面取得了相当大的进展。 尽管如此，仍然非常需要提高现有方法的收敛速率。 势流方程的线松弛方法以及欧拉和Navier-Stokes 方程的时间迭代方法都受到缓渐近收敛到稳态的影响。 将迭代解释为松弛，自然需要考虑已成功应用于经典松弛方案的收敛加速方法。 这些方法中最重要的是多网格算法。 多网格理论对于椭圆方程是高度发达的，其中0（n）收敛速率是可以实现的，其中n是系统中未知数的数量。 换句话说，获得方程组解所需的工作与未知数的数量成正比。 椭圆方程的经典松弛方案在消除误差的短波长分量方面非常有效，而多网格过程中的粗网格在消除长波长误差方面非常有效。 将多网格加速度应用于欧拉或雷诺平均 Navier-Stokes 方程会导致人们考虑时间积分方法，这些方法也为短波长误差分量提供了良好的阻尼。 基于非稳定方程的多网格方法主要有两类。 一类方法使用上风差分和隐式时间积分作为平滑器[1,8, 17]。 另一种方法是詹姆森最初提出的[7]。 具有显式人工黏度的有限体积空间离散化与Runge-Kutta 时间积分相结合，作为平滑器。 这种方法已成功扩展到雷诺-

新鲜看待多电网243平均Navier-Stokes 方程 [16]。 不幸的是，这些方法导致多电网效率低下。 当应用于复杂几何形状的高雷诺数流量时，收敛速率通常比每个多网格周期的0.99差。 最近，皮尔斯等人展示了显著的改善。 [10]。 然而，当人们考虑到平滑域上的泊松方程收敛在实践中可以达到每周期近0.1的速率时，很明显，现有的求解器有很大的改进空间。 在本文的其余部分，提出了Euler 方程的多网格算法，该算法产生的收敛速率与泊松方程的速率相当。 该算法放弃了稳态的时间行进方法，而是依赖于稳态方程的松弛。 第14.2节概述了该算法的一般原则。 第14.3节给出了当前方法的数学公式，而第14.4节描述了解决方案程序。 第14.5节显示了二维通道和机翼周围的不可压缩无黏流动的结果。 第14.6节简要讨论了将当前方法扩展到可压缩流动方程，其中还显示了与势流求解器的联系。 摘要可在第14.7节中找到。 14.2 多网格方法 根据Brandt [2]的说法，实现对流主导流的理想多网格性能的主要障碍之一是，粗网格仅为平滑误差分量提供所需校正的一小部分。 这个特殊的障碍可以通过设计一个求解器来消除，该求解器可以有效地区分系统的椭圆、抛物线和双曲（对流）因素，并适当地处理每个因素。 这种求解器的效率将受到求解器对系统每个因素的效率的限制。 例如，对流可以通过空间行进来处理，而椭圆因子可以通过多网格来处理。 在本例中，与对流项相关的所有误差组件都在一次扫描中被消除，收敛速率受椭圆求解器的速度限制。 Brandt提出了一种叫做“分配松弛”的方法，通过这种方法，人们可以构建平滑器，有效地区分运算符的不同因素。 使用这种方法，Brandt和Yavneh展示了不可压缩Navier-Stokes 方程的教科书多重网格收敛速率[3]。 他们的结果是一个简单的几何和笛卡尔网格，使用方程的交错网格离散化。 在一个密切相关的方法中，Ta'asan [15]为可压缩Euler 方程提出了一个快速的多网格求解器。 该方法基于一组“规范变量”，这些变量以表示稳定的Euler 方程

244 ROBERTS、SIDILKOVER和SWANSON椭圆和双曲分割[14]。 Ta'asan使用这个分区来引导方程的离散化。 使用交错网格。不同的变量位于单元格、顶点和边缘中心。 在参考[15]中，表明使用身体贴合网格的二维亚声速流动的可压缩Euler 方程可以实现理想的多网格效率。 使用规范变量的一个可能限制是，非黏性方程的分区不能直接适用于黏性方程。 最近，作者[ll]提出了一个替代Brandt的分布松弛和Tarasan的规范变量分解的替代方案。 该方案不需要交错网格，而是使用原始变量的传统基于顶点的有限体积或有限电子发散离散化。 这简化了限制和延长操作，因为所有变量都可以使用相同的运算符。 投影算子被应用于方程组，得出压力的泊松方程。 通过将投影算子应用于离散方程而不是微分方程，压力的适当边界条件直接得到满足。 压力的泊松方程可以用高斯-塞德尔松弛来处理，而动量方程的对流项可以用空间行进来处理。 由于系统的椭圆和对流部分是解耦的，可以实现理想的多电网效率。 与分布式松弛和规范变量方法相比，这种方法非常简单。 14.3 数学]公式 原始变量中的不可压缩Euler 方程，其中u和v分别是s和y方向速度的分量，p是压力。 密度为1。 对流运算符的定义是其中a、、ay是偏微分运算符。 Euler 方程可以写成

新鲜观察多网格245介绍Q的伴数，由Q*(f) = -d x(uf)-dy(vf)定义，(14.3)一个投影运算符P定义为：/ / 0 0 \textbackslash\{\} P = 0 / 0 (14.4) \textbackslash\{\} d t 8 y Q* J 将投影运算符应用于Euler 方程产生/ Q 0 d x\textbackslash\{\} / u\textbackslash\{\} / 0 \textbackslash\{\} Lq = PLq = 0Q3„u+2 0 V 0 A ) \textbackslash\{\} P J \textbackslash\{\} (d XV)(dyU) \textasciitilde{} (d XU)(dyV) J (14.5) A是拉普拉斯。 右侧的矩阵运算符由L的主体部分（即运算符的最高阶项）组成，其余项是次主体项。 这些项的产生是因为运算符Q和Q*中的系数u和v不是恒定的。 需要注意的是，为了构建放松方案，可以忽略子原则术语。 方程组14.5是比原来的Euler 方程 14.2更高阶的系统。 连续性方程是一阶偏微分方程，已被压力的二阶微分方程所取代。 人们可能会想到Eq. 除了墙体流动切线的物理边界条件外，14.5还需要对压力的边界条件，这是Eq.要求的。 14.2. 然而，在域的边界上，14.5的第三个方程采取的形式（-udyv + vdyU + d xp）hx +（udxv - vd xu + dvp）hy = 0，（14.6），其中单位的分量法线在墙上。 这仅仅是垂直于墙壁的动量方程。 因为墙壁上的压力方程采用等式的形式。 14.6，在这种情况下可以被认为是控制方程的兼容性条件，不需要压力的辅助边界条件。 等式左侧的操作员。 14.5是上三角形。 由于压力满足泊松方程，可以使用传统的松弛方法，如高斯-塞德尔，来求解它。 动量方程中的对流算子迎风差分和网格顶点的下游排序允许动量方程的行进。 这里使用了集体高斯-塞德尔方法，其中顶点按流动方向排序。 下一节将更全面地描述这一点。

246 ROBERTS、SIDILKOVER和SWANSON 14.4 解决方案程序 近似L的第一步是离散化Euler 方程 14.2。 与使用交错网格的引用方法[3,15]不同，当前的方法是基于顶点的，所有未知数都存储在网格的顶点上。 四边形结构网格和三角形非结构网格的离散化已经编码。 使用当前方法，动量方程的离散近似形式具有很大的灵活性。 举例说明，考虑非结构化三角形网格上的细胞顶点离散化。 典型的网格单元格0如图所示。 1. 通过使用ft边界周围的梯形规则积分来发现未知数的细胞平均梯度。 例如，对fl上u梯度的离散近似是id\textasciicircum{}u+jdyU = \textasciicircum{} (w 0(yi -2/2)+ "1(2/2 -yo)+u 2(yo -J/i)) \textasciitilde{}2AE ( u°(Xl \textasciitilde{} xz) + u i( x 2 - so) + u 2(x0 - Xi)), (14.7) 其中An是三角形的面积。 上标h用于表示相应微分运算符的差分近似值。 同样获得v和p的梯度。 这些梯度用于近似等式。 三角形上的14.2。 通过将单元平均动量方程残差适当地分布到每个三角形的顶点，获得网格顶点Q的上风近似值。 目前的方案不与任何特定形式的上风离散化绑定。 用于获得第14.5节中提出的非结构化网格结果的一个选择是Giles等人的对流方案。 [5]。 在此方案中，残差是图1非结构化网格的Traingular cell

FRESH LOOK AT MULTIGRID 247分布到Q的顶点。使用权重Arii Wi 1- vjl i = 0,l,2，其中l n是Q投影到横流方向的长度，An；是第i个顶点对面边缘横流方向长度的分量。 作者目前正在开发的另一种上风离散化是基于Sidilkover的多维上风公式[12]。 一旦计算了连续性和动量方程的单元平均残差，将投影算子P应用于这些离散方程，以获得方程的压力泊松方程的残差。 14.5. 让R Pi是顶点i的剩余压力方程，P的应用可以用积分形式写成，RPi= L I / ( \textbackslash\{\}Qhu + d\textasciicircum{}\textbackslash\{\}-u\textbackslash\{\}(d\textasciicircum{}u + d\textasciicircum{}j\textbackslash\{\} ) dy三角形\textbackslash\{\}g\textasciicircum{}。 (A*u + a*t>)| )dx| (i4.8) -( Qhv + <9> 其中A;是以第i个顶点为中心的控制体的面积，上标h用于表示三角形上相应微分算子的离散近似。 求和在与第i位顶点相邻的所有三角形上。 图中的虚线。 1是位于细胞Q中的Ao、A\textbackslash\{\}和A 2边界的部分。 14.8中的框项是x和y动量方程和连续性方程的单元平均残差。 通过评估Eq，找到细胞平均残差在Q上对R p的贡献。 14.8在A\{位于ft的边界的适当段上，取等式中的框项。 14.8在细胞上是恒定的。 以这种方式在离散水平上应用投影算子P，而不是从微分方程14.5开始并将它们离散化，有两个重要的优势。 首先，方程的离散近似。 14.8在边界顶点减少到Eq。 14.6，自动为压力提供正确的边界条件。 其次，如果动量和连续性方程以守恒形式在三角形上离散化，就有可能得到一个完整的守恒格式。 这对于带有冲击的可压缩流动尤为重要。 多网格算法使用一系列网格GK、GK-\textbackslash\{\}、■ ■ ■、Go，其中GK是最好的网格，Go是最粗的。 将离散近似调用到fc-th网格L\&上的运算符L，并让q*是解

248 ROBERTS、SIDILKOVER和SWANSON那个网格。 该系统的形式为Ljtqjt = f*，其中LA的条目是3x3块矩阵，在每个网格顶点的未知数（«，v，p）T上操作。 通过将运算符LA写成lik = Mfc-Nfc来构建一个通用迭代方案，其中选择拆分，使MA很容易反转。 词汇高斯-塞德尔是通过将MA作为忽略lik的对角线块上方的项而产生的块下三角矩阵获得的。如果MA的对角线块仅包含与运算符主体部分相对应的条目，则可以获得进一步的简化。 因为Eq中的运算符。 14.5是上三角形，MA的对角线块将是3 x 3上三角形矩阵。 让q£是fc-th网格上解的ra-th迭代，放松迭代是MAq£ +1 =f* + NAq;j。 运算符LA是非线性的，因此MA和NA是qjj和q£ +1的函数。 让£q£ = qjj +1 - <\$■>迭代可以重写为Mjfcfctf = fA - Uqt（14.9）因为MA是块下三角形的，<5q£通过正向替换找到。 在每个顶点，必须反转一个3 X 3的上三角矩阵。 如果对对流算子Q的离散近似是全上风，并且网格点是按照流动方向排序的，那么NA的3x3块将在前两行中为零。 在这种情况下，词典格高斯-塞德尔松弛等同于对流项的空间行进。 对流误差在一次松弛扫描中被有效地消除，系统的收敛速率成为压力泊松方程的收敛速率。 有可能获得理想的多重网格收敛费率，因为错误的每个组成部分都得到了适当的处理。 直截了当的全近似方案（FAS）多网格迭代应用于方程组。 让LA-I是粗网格运算符，I\textasciicircum{} -1是细到粗网格限制运算符，/£ \_1是粗到细网格延长运算符。 如果qA是网格fc上的当前解，则该网格上的残差为TA = fA - LACJA- 这导致粗网格方程LA-iqA-i = ffc-i = l£\_ xrk + LA-I (l k-Ak) ■ (14.10) 求解qt-i的粗网格方程后，通过qr w <- qA + I\textasciicircum{} (q*- - ti-Ak) ■ (14.11) 方程14.10通过应用用于求解精细网格方程的相同放松程序求解。 多网格递归地应用于粗网格方程。 在最粗的网格上，执行许多松弛扫描，以确保方程完全解。 使用传统的Vcycle或W-cycle。

FRESH LOOK AT MULTIGRID 249 14.5 结果 基于第14.3和14.4节理论的非结构化网格和结构化网格流求解器都已编写。 这些代码在参考[11]中进行了描述，其中提出了广泛的解决方案。 这里展示了说明该计划效率的结果。 两个求解器都获得了一个通道中的不可压缩无黏流动的解决方案。 通道几何形状和边界条件如图所示。 2. 0 < x < 1之间的下墙形状是y(x) = r sin 2 -KX。 对于此处显示的计算，厚度比r为0.05。 流量角和总压力在入口指定，压力在出口指定。 流动切线条件u-h = 0在通道的上壁和下壁强制执行。 在准均匀的四边形网格上得到了解。 在通道的中心部分使用了简单的剪切变换，以获得符合边界的网格。 对于非结构化网格求解器，通过沿对角线划分每个四边形单元格来对网格进行三角测量。 通过在每个坐标方向上将细网格粗化两倍，得到了一系列嵌套粗网格。 在下图所示的所有情况下，最粗的网格是7x3的顶点。 使用了词典学高斯-塞德尔松弛，网格顶点从通道的左下角到右上角排序。 这导致了动量方程的下游松弛。 使用了V（2,1）多网格周期；即，在限制到粗网格之前，在每个网格上执行两次放松扫描，在将粗网格校正添加到微网格溶液中后执行一次放松扫描。 图中显示了97 x 33顶点网格上的计算压力。 3用于非结构化网格流动求解器和图中。 4用于结构化网格求解器。 图中显示了不同网格密度的收敛速率的比较。 5和6分别用于非结构化和结构化网格流求解器。 显示了压力方程残差的L\textbackslash\{\}规范；图2通道几何。

250 ROBERTS、SIDILKOVER和SWANSON图3压力，轮廓增量Ap = 0.01，用于97 X 33顶点的非结构化网格。 图4压力，等高线增量Ap = 0.01，对于97 x 33个顶点的结构化网格。动量方程残差显示相同的行为。 每个流动求解器使用的最佳网格包含385 x 129个顶点，共7个网格级别。 最精细网格上非结构化网格求解器的收敛速率每个多网格周期的残差减少约为0.190。 结构化网格结果略好，每周期0.167。 这些速率与泊松方程每周期0.125的理想速率相当。 结构化网格求解器的性能更好很可能是由于更好的限制和延长运算符；非结构化流动求解器仅使用细网格顶点的位置和包含该顶点的粗网格单元的三个顶点执行双线性插值。 最重要的是，这些数字显示了几乎理想的多重网格收敛速率，与电网间距无关。 这表明收敛是在有序操作中实现的。 对于复杂的几何形状，生成一系列嵌套的非结构化网格可能不切实际，多网格求解器的性能可能

新鲜看待多电网251图5非结构化电网上收敛速率的比较。 图6结构化网格上收敛费率的比较。

252 ROBERTS、SIDILKOVER和SWANSON图7通过扰动49 X 17网格的顶点生成的网格。 图8 压力，等高线增量Ap = 0.01，97 X 33顶点的随机扰动非结构化网格。预计会恶化。 为了显示当前方法的鲁棒性，对一系列非嵌套粗网格运行了三角网格求解器。 这些是通过独立随机扰动每个嵌套网格上的顶点位置来生成的。 扰动的49 x 17网格如图所示。 7. 图中显示了具有5个网格级别的扰动97 x 33细网格的计算压力。 8和收敛速率如图所示。 9. 压力轮廓非常平滑，没有显示缺乏网格平滑的迹象。 渐近的收敛速率已经退化到每周期仍然值得尊敬的0.24。 使用结构化网格求解器获得了对称Karman-Trefftz 翼型上的非升降流的解决方案。 从共形映射中生成了一个385 X 193顶点的精细O网格，粗网格通过递归地消除每个坐标方向上的所有其他顶点来嵌套。 选择网格间距是为了获得单位长宽比网格单元格。 外部边界距离翼型大约13个弦长。 Farfield 边界条件由分析解给出。 在流入点

新鲜看多网格253图9 收敛速率，随机扰动无结构网格，97 x 33个顶点，5个网格级别，V（2，1）循环。沿外边界指定了总压力和流动倾角。 对于流出点，规定了压力。 在翼型表面上，强制执行切线条件。 为了获得理想的多重网格收敛速率，有必要按下游顺序对顶点进行排序，以便动量方程中的对流项进行行进。 通过沿着径向网格线从域的外边界到域的前半部分的翼型表面，以及域的后半部分从翼型表面到外边界的径向网格线放松，这很容易完成。 对于每个案例运行，最粗的网格由13 x 7个顶点组成。 卡曼-特雷夫茨翼型周围非升力流动的计算和分析表面压力系数的比较如图所示。 10. W(2,1)多网格周期用于这些计算。 计算解决方案与分析解决方案非常一致，除了后缘的再压缩。 请注意，这个区域没有网格聚类，这加剧了问题。 图中显示了三个网格密度的压力方程残差的收敛速率的比较。 11. 观察到收敛速率随着网格加密的增加而略有恶化：在385 x 193网格上，速率为每周期0.153。 然而，与通道流量结果一样，收敛速率几乎独立于电网，非常接近每周期0.125的理想速率。

254 ROBERTS、SIDILKOVER和SWANSON图11非起重Karman-Trefftz 翼型的收敛率比较。 图10 表面压力系数，非升降Karman-Trefftz 翼型 193 x 97网格。

新鲜看多网格255 表1列出了最佳网格上的收敛速率的摘要。 显示两组结果：每个多网格周期的收敛速率和每个工作单元的收敛速率。 为了讨论的目的，一个工作单元（WU）被理解为在最好的网格上的一个高斯-塞德尔放松扫荡。 这本质上是在最佳网格上进行一次剩余评估的成本。 实际收敛费率与理想的收敛费率进行比较，其计算方式如下。 让y是放松方法的平滑率。 对于V（m，n）或W（m，n）周期，理想的收敛速率是nm+n。 泊松方程的词汇表高斯-塞德尔具有平滑率\textbackslash\{\}i = 0.5。 对于V(2,1)或W(2,1)周期，这给出了0.5 3 = 0.125的理想收敛速率。 为了计算每个工作单元的收敛速率，我们使用以下公式。 对于每个周期，总共有m + n个精细网格放松扫描。 等式检查。 14.10表明，细到粗网格限制要求对细网格进行一次残差评估，并对粗网格进行额外的残差评估。 粗网格残差评估是细网格残差评估成本的1/4。 由于放松扫描的大部分成本都在于剩余的评估，因此每个周期在最精细的网格上总共需要（m + n + 1 + 1/4）工作单位。 忽略了在电网水平之间插值和余量和解决方案的成本。 对于V(m, n)-循环，我们有WU \_\_ « (m + n + i +! ）（我+！ + x +...) = §(m + n+f)。 由于W-cyc\textbackslash\{\}e每个周期涉及两个粗网格解，我们有（m + n + 1 + i）（l + I + 1 + • • •）W-cycle = 2（m + n + |）。 对于F周期，这些数字产生了每个工作单元/i3(m+n)/(4(m+ n+5/4))的理想收敛率，对于jy周期，每个工作单元Aj('"+n)/(2(m+n+5/4))的理想收敛率。 V（2<1）周期被认为每个周期需要5 2/3 WU。 W（2,1）贵50\%，每个周期需要8V2 WU。 相比之下，一个V（2,1）周期只比最佳网格上5级Runge-Kutta方案的单个时间步骤略多。 词汇学高斯-塞德尔的理想收敛速率为V(2,1)周期每WU 0.693，W(2,1)周期为每WU 0.783。 表1所示的实际利率似乎非常接近理想。 表1中的收敛费率可用于估算获得精细网格上离散化误差水平的解决方案所需的工作。 让p是离散算子的近似顺序，让hk是fc-th网格上的网格间距参数。 对解决方案的初步猜测

256 ROBERTS, SIDILKOVER \& SWANSON 案例周期 收敛速率 案例周期 每个工作单元的案例周期 理想实际理想实际通道，385 x 129，非结构化通道，385 x 129，结构化通道，结构化V(2,l) V(2,l) 0.125 0.125 0.190 0.167 0.693 0.693 0.746 0.729 翼型, 385 x 193，结构化W(2,l) 0.125 0.153 0.783 0.802 表1 通道和翼型流的精细网格上多网格求解器的收敛速率摘要，并与理想速率进行比较。在精细网格上，GK是通过插值GK-I上计算的解决方案获得-假设GR-I上的解决方案已获得 离散化 误差 T K-I = 0(h p ) 在该网格上。 多网格循环用于减少从T„\_到r„的误差。 让fi w是每个工作单元的收敛速率，从GK-I上的初始解决方案获得GK解决方案所需的工作量WK = logp -log w •K-l P log fi log w -K-\textbackslash\{\} 全多网格（FMG）周期从最粗网格Go上的解决方案开始，并使用上述策略在更精细的网格上递归生成改进的解决方案。 对于此处考虑的嵌套网格，网格间距参数由hk-i = 2/ijt相关，每个网格上的工作量由Wk-\textbackslash\{\} = Wk/4相关。 离散化是二阶精度，即p = 2。 这为我们提供了获得精确到T K为W的解的总工作的估计，总L\_ l°gM W fcs(i)(i + i + \& + -0 \_8 1OR2 3 log M W (14.12) 使用表1最后一列中n w的值，我们看到通道流解可以使用FMG周期在大约6.4 WU内获得离散化误差水平。 翼型溶液可以在大约8.3 WU中获得。 这些估计值通常很低，实际上比单个FMG周期的工作要小（通道和翼型病例分别为7.6和11.3 WU）。 使用Eq计算的工作。 14.12也没有解释在粗网格解插值到细网格时引入短波长误差。 尽管如此，Eq. 14.12是多电网方案预期性能的有用指南。

FRESH LOOK AT MULTIGRID 257 14.6 扩展到可压缩流动 这里介绍的方案直接扩展到可压缩欧拉方程。 在原始变量中，方程为Lq = ( Q 0 0 0 pQ 0 0 0 pQ \textasciicircum{} o pdx Pdy o \textbackslash\{\} ( u \textbackslash\{\}p) = 0, (14.13) 其中p是密度，c是音速，s是熵。 该系统的投影算子P是P = / / 0 0 o \textbackslash\{\} 0 / 0 0 0 0 / 0 \textbackslash\{\}o dx dy Q* ) (14.14) 运算符Q和Q*在方程中定义。 14.1和14.3。 将此应用于Eq。 14.13 和以前一样忽略次元项，产生PLq = 0 0 PQ 0 0 0 PQ 0 0 \textbackslash\{\} I s \textbackslash\{\} d, u dy V M2a32 ) \textbackslash\{\}P ) + s.p.t., (14.15) 其中M是马赫数，d a是沿流方向的偏导数，“s.p.t.”是次元项。 可压缩方程和不可压缩方程之间最重要的区别在于，类似Prandtl-Glauert的算子作用于压力。 请注意，这个系统接近于消失马赫数极限的不可压缩方程的系统。 对于亚声速流动，可压缩方程可以通过与不可压缩方程相同的松弛方案求解。 与时间行进方法不同，收敛速率不会随着马赫数接近零而退化。 Prandtl-Glauert算子在压力方程中的出现很重要。 实际上，求解压力方程的问题与求解全势方程的问题没有区别。 这是一把双刃剑。 一方面，在跨声速情况下放松压力方程的困难正是放松跨声速全势方程的困难。 这是任何直接在稳态流动方程上工作的方法都面临的根本困难。 另一方面，人们可以期待在解决潜在流动方面的丰富经验可以直接应用于可压缩Euler 方程的当前方案。 对流术语的处理与不可压缩的案例没有本质区别。

258 ROBERTS、SIDILKOVER和SWANSON 14.7 结论 Murman和Cole将类型依赖的微分和松弛方法引入到计算空气动力学中；其结果是求解非线性流动方程的实用和高效的方法。 在随后的几年里，重点转向了基于非定方程的迭代方法。 本文表明，通过将观点从不稳定方程转变为稳定方程，可以大大提高流量求解器的效率。 由于Murman和Cole引入了类型依赖性差异化，因此目前的方法依赖于离散化，该离散化区分了系统的椭圆和双曲部分。 正如Murman和Cole使用松弛来求解稳态方程一样，目前的方法将多网格的松弛应用于稳态方程。 这种方法为稳定的Euler 方程提供了教科书式的多网格效率。 这是一个特别简单的方法；可以使用控制方程的传统有限差分或有限体积离散化，允许在选择基本数值方法时具有灵活性。 与时间行进方法不同，但与势流方法一样，该方法的收敛速率不会对低速流动进行降解，并恢复正确的不可压缩极限。 最后，根据Sidilkover和Ascher的想法，这种方法可以应用于不可压缩的黏性流动[13]。 在像这里显示的简单问题一样有效地解决完整的雷诺平均Navier-Stokes 方程之前，仍有很多工作要做。 本工作只涉及问题的一个特殊但仍然重要的方面，即控制方程的适当离散化。 如果要实现可压缩黏性流动的教科书多网格效率，它可能需要按照这里介绍的一般大纲的方法。 致谢 如果没有Jerry South的倡导和鼓励，本文中介绍的工作就不会开始。 作者感谢他在这项研究过程中的兴趣、热情和鼓励。 参考资料1。 安德森，W。 K.，托马斯，J. L.和Whitfield，D. L.，“通量分割Euler 方程的三维多网格算法”，美国宇航局技术论文2829，1988年。

重新看MULTIGRID 259 2。 Brandt，A.，“多网格技术：1984年流体动力学应用指南”，GMD-Studie 85，GMD-FIT，1985年。 3. Brandt，A.和Yavneh，I.，“加速多重网格收敛和高雷诺再循环流”，SIAM J. 科学。 统计学家。 计算，卷。 14，不。 3，pp. 607-626，1993年。 4. 考伊，D。 A.，“跨声波势流的计算”，Ann。 牧师。 流体机械，卷。 14，pp. 261-283，1982年。 5. 吉尔斯，M.，安德森，W. K.和Roberts，T. W.，“上风控制音量：一种新的上风方法”，AIAA论文90-0104，1990年。 6. 詹姆森，A.，“翼型和机翼上的跨声速流动的迭代解，包括1马赫的流动”，Comm。 纯Appl。 数学。，卷。 27，pp. 283-309，1974年。 7. 詹姆森，A.，“通过多重网格法为二维跨声速流动的Euler 方程解决方案”，应用。 数学。 计算，卷。 13，不。 3和4，pp. 327-355，1983年。 8. Mulder，W.，“Euler 方程的多网格放松”，J. 计算机。 物理，卷。 60，不。 2，pp. 235-252，1985年。 9. 默尔曼，E。 M.，Cole，J. D.，“平面计算，稳定跨声速流动”，AIAA J.，卷。 9，不。 1，pp. 114-121，1971年。 10. 皮尔斯，N。 A.，Giles，M.，Jameson，A.，Martinelli，L.，“加速三维Navier-Stokes计算”，AIAA论文97-1953，1997年。 11. 罗伯茨，T。 W.，Sidilkover，D.，Swanson，R. C，“Steady Euler 方程的教科书多网格效率”，AIAA论文97-1949，1997年。 12. Sidilkover，D.，“可压缩Euler 方程的真正多维上风方案和高效的多网格求解器”，ICASE报告94-84，1994年。 13. Sidilkover，D.和Ascher，U. M.，“使用压力泊松配方的稳态Navier-Stokes 方程的多网格求解器”，Comp. 应用程序。 数学卷。 14，不。 1，pp. 21-35，1995年。 14. Ta'asan，S.，“多维稳定无黏流动的规范形式”，ICASE报告93-34，1993年。 15. Ta'asan，S.，“稳态Euler 方程的规范变量多重网格法”，ICASE报告94-14，1994年。 16. Vatsa，V.，Wedan，B. W.，“3D Navier-Stokes方程的多网格代码的开发及其在网格细化研究中的应用”，计算机与流体，卷。 18，不。 4，pp. 391-403，1990年。 17. 沃伦，G。 P.和Roberts，T. W.，“上风偏置数据重建的多网格属性”，第六届铜山多网格方法会议，美国宇航局会议出版物3224，第2部分，1993年。

\clearpage
\chapter{并行计算机上的航空航天工程模拟}
并行计算机上的航空航天工程模拟K。 摩根，1 N。 P。 Weatherill，1 O。 哈桑，1 P. J。 布鲁克斯，1 M。 T。 曼扎里1和R。 Said 1 15.1 导言 航空航天公司正在采用基于非结构化四面体网的解决方案技术来模拟复杂几何形状上的稳定可压缩非黏性流动[4]。 该方法在各种计算机平台上的实现表明，对于这一类问题，通常可以通过使用由数百万个元素组成的网格来实现可接受的准确性[18]。 然而，该方法的计算机内存要求现在阻止了其扩展到需要使用明显更大的网格的模拟。 当考虑涉及湍流流的问题，甚至是存在移动边界的瞬态无黏性流动的问题时，这一点尤为明显。 当模拟航空航天车辆在现实波频率下的电磁散射时，情况会更糟，因为可用的数值算法将要求使用由数亿个元素组成的网格。 在此背景下，鉴于当前计算机技术的发展，我们已将研究基于使用并行处理来解决计算大问题的技术。 然而，为了使这种方法在英国斯旺西SA2 8PP的威尔士大学土木工程系中取得成功。 Frontiers of 计算流体动力学 - 1998 编辑：David A. Caughey和Mohamed M. 哈菲兹 ©1998 世界科学

262 MORGAN, WEATHERILL, HASSAN, BROOKES, MANZARI \& SAID 在上下文中，它不仅需要对正在使用的基本方程求解器的并行实现，还需要对网格生成和区域分解的预处理阶段采取新方法。 在本章中，我们将考虑航空航天工程应用的基本非结构网格求解器的并行化，并概述我们正在遵循的方法，以便我们能够满足预处理阶段的计算机内存需求。 与可视化计算解决方案相关的额外困难正在相关工作[6]中得到解决，这里没有考虑。 15.2 解算法 考虑以dU \&F 3\_\_ 8\& dt dxj dxj的形式表示的偏微分方程组的解，其中j取隐含求和中的值1,2,3。 在封闭的空间域中寻找解决方案，受制于适当的初始和边界条件。 有了F 3、G 3和S的正确定义，这种一般形式允许使用适当的湍流模型[8]和电磁散射[10]对不黏性流动[12]、湍流流动进行建模。 开发基于使用非结构化网格的解决方案方法的起点是用弱变分公式取代问题的经典陈述。 空间解域被离散化为四面体元素的一般集合，使用具有自动点创建的Delaunay程序[17]。 当需要适合湍流 黏性流动 [3]或电磁散射[11]的网格时，此基本离散化程序会被修改。 假设空间中近似解的分段线性变化，然后使用Galerkin近似变分公式导致每个内部节点/的半离散方程MIJ-\textasciicircum{}- = M uSj + RI（15.2），其中隐含求和延伸到网格中的所有节点J。 当网格以基于边缘的数据结构[13]表示时，根据R ' = E Ck [(FJi + Fl) - (G»; + Gj.)]，通过将单个边缘贡献相加来评估右侧向量R轴的分量。 （15.3）其中节点I假定由边缘e直接连接到节点Ie。 在这个方程中，与边缘e和方向Xj相关的权重Cj T，

在平行计算机上模拟263表1性能统计（RSB）。 在黏性流动中稳定。 NoP MinEd MaxEd MinPo MaxPo RT(min) RSU 4 469 771 484465 66 523 70130 650 1 8 230459 244 827 33 262 37 357 345 1.88 16 114139 125 040 16 631 19 674 159 4.09 32 55144 64 926 8 316 10 609 92 7.07 64 26 760 33 475 4158 5 930 44 17.77 128 12 826 16 980 2 079 3 239 27 24.07 256 6189 8 762 1040 1822 14 46.42 表2 2 性能统计（RSB）。 湍流。 NoP MinEd MaxEd TNIE RT(min) RSU 4 579 277 579 291 48 069 2 097 1 8 289 623 289 651 63 502 1053 1.92 16 144 774 144 827 98 312 548 3.84 32 72 311 72417 157 324 271 7.73取决于局部网格几何[13]。 这些重量是预先计算的，然后存储。 通过用适当的数值通量函数替换每个边缘的实际对流通量函数来实现稳定和不连续捕获[5]。 采用有限差分程序来离散化时间维度。 就目前而言，使用欧拉或多级时间推进来推进解决方案。 对于稳态模拟，通过将[20]方程（15.2）中的质量矩阵M块来求解生成的方程系统，而显式迭代用于真正的瞬态问题[1]。 对于涉及移动边界的瞬态模拟，采用空间/时间变异公式[9]。 15.2.1 并行实现 此解决方案算法的高效并行实现需要一个有效的程序来获得网格的初始分解，以确保平衡处理器之间的计算负载，并优化必要的处理器间通信。

264 MORGAN, WEATHERILL, HASSAN, BROOKES, MANZARI \& SAID 表3 绩效统计数据（RSB）。 电磁散射。 NoP MinEd MaxEd TNIE RT（sec） RSU 4 549429 569 181 47577 4516 1 8 271 198 294234 79 792 2232 2 16 132911 148268 112764 1174 3.8 32 64814 76339 156772 613 7.4 64 32 166 38441 209080 342 13.2 128 14971 20659 268356 176 25.7 256 7 142 10701 344607 136 33.2 15.2.1。 I 区域分解 有许多不同的分解方法，在序列中，给定的非结构网格 12）。 在我们的初步研究中，递归光谱平分（RSB）是选择的方法[15]，因为它通常产生平衡的子域和低通信要求。 这些属性应该会导致平行方程求解器的良好性能。 15.2.1.2 数据结构 RSB \textasciitilde{}rocedure将全局网格细分为一个数字。 NoP，子域的子域透过对网格中的节点著色。 连接两个相同颜色的节点I的边缘是子域I的内部边缘，两个节点都被视为该子域的内部节点。 连接颜色I节点和颜色J节点的边缘。其中I < J，将是子域I中的接口边缘。 颜色I的节点和颜色J的节点将分别是子域I和J的内部节点。 在这种情况下，颜色J的节点被复制为子域I中的接口节点。 应该注意的是，边缘没有重复。 每个子域中都使用顶点、元素、边缘和边界面的局部编号。 通信阵列是实现子域之间信息传输所必需的，在域分区阶段进行评估。 15.2.1.3 求解器并行化 求解算法的并行实现使用标准的PVM或MPI例程进行消息传递，并采用单个程序多个数据模型。 在时间步骤开始时，接口节点从接口边缘获得贡献。 然后，这些部分更新的接口节点贡献被广播到相邻的相应内部节点

在PARALIJEL计算机上模拟265图1具有网格生成器（实线）、RSB（虚线）和方程求解器（虚线）网格大小的内存回收的变化\textasciitilde{}。s\textasciitilde{}ibclorr\textasciitilde{}ains。 在接收接口节点贡献并随后更新所有内部节点值后，在内部eclges上循环。 将更新的mlt\textasciitilde{}es发回接口节点完成程序的tt时间步骤。 该程序的实施方式如此之大，以至于它试图允许计算和通信同时进行。 15.2.2 示例Thc：通过考虑来自悬崖应用领域的示例来演示方程求解器的pardlel性能。 So1ve1.s之前是f\textasciitilde{}\textasciitilde{}lly vdicbatecl \{8,10,12]，因此本章不再考虑vdiction。 15.2.2.1 稳定无黏流动 稳定无黏流动求解器的pardlel版本的pe\textasciitilde{}*fosmance特性被证明为由1612 174 dements, 266092组成的网格

266个MORGAN、WEATHERILL、HASSAN、BROOKES、MANZARI和SAID节点和1 912 170个边缘。 表1说明了RSB方法的负载均衡功能，还包含CRAY T3D上结果实现的性能信息。 在这里，MinEd、MaxEd、MinPo、MaxPo分别表示任何子域中的最小边缘数、最大边缘数、最小点数、最大点数，RT是计算1000个时间步骤所需的运行时间，RSU是与使用四个子域时所需的时间相比获得的相对速度。 对于这种尺寸的网格，观察到使用256个处理器实现了72.5\%的并行效率。 为了说明使用k-ui 湍流模型[19]模拟湍流时可以获得的性能，生成了一个由1 966 731个元素、331 685个节点和2 317 145个边缘组成的网格。 表2显示了网格分区统计数据，以及CRAY T3D上并行流动求解器的结果性能信息。 在这里，TNIE表示接口边缘的总数。 在这种情况下，由于流动求解器执行的计算量增加，并行性能通常更好。 15.2.2.2 电磁散射 展示了由1897 844个元素、310813个节点和2 242682个边缘组成的网格的并联电磁散射模拟能力的性能特性。 表3显示了使用RSB实现的网格分区统计数据，以及CRAY T3D上并行求解器结果性能的信息。 请注意，RT现在表示计算500个时间步骤所需的运行时间。 在这种情况下，并行性能通常不是那么好，因为数值通量函数[10]采用简化形式，以及由此减少求解器执行的计算量。 15.3 内存要求 分析师使用非结构网格方法对航空航天工程中的现实问题进行模拟，很快就遇到了可用计算机资源对可以使用的网格大小施加的限制。 在CRAY T3D等机器上，上述方程求解器的实现可以在单个处理器上容纳大约400000个元素。 这意味着，如果512个处理器可以同时访问，就可以处理约2亿个元素的网格。 然而，网状生成和区域分解流程需要大量的

在并联计算机上模拟 图2 内存要求与RSB（虚线）和RBM（实线）的区域分解网格大小的变化。内存量通常以串行方式执行，这意味着在实践中可以使用的网格要小得多。 图1说明了这一点，图1显示了网格生成和区域分解的内存要求如何随典型串行计算机上的网格大小而变化。 为了感兴趣，尽管在这种情况下没有直接相关，但流动求解器的串行内存要求也显示出来。 通过将讨论限制在特定的计算机平台上，可以最好地证明这个数字的意义，在这里我们选择一个CRAY YMP/EL，它具有大约200 MWords的可用内存。 从图中可以看出，在这台机器上，可以生成和分解的最大网格将由大约800万个元素组成。 因此，对于涉及更大网格的模拟，必须采用替代区域分解和网格生成策略。 15.3.1 替代区域分解策略 从图1可以看出，区域分解程序对可用内存的要求最大。 因此，一个显而易见的第一方法是调查内存需求，以及随后的

268 MORGAN, WEATHERILL, HASSAN, BROOKES, MANZARI \& SAID 表4 绩效统计（RBM）。 电磁散射。 NoP MinEd MaxEd TNIE RT(sec) RSU 4 550181 571164 57 215 4610 1 8 274 369 291416 92 073 2 325 2 16 134 371 147 006 133 994 1070 4.3 32 64 296 74412 173 396 578 8 64 31607 37 956 224 545 332 13.9 128 15176 19 607 286 736 191 24.1 256 7 220 10154 361971 145 31.8性能，其他区域分解策略的表现。 作为代表性的替代方案，还考虑了使用基于递归带宽最小化（RBM）的简单直接分区算法[2]。 这种方法减少的内存需求在图2中显而易见，图2将RSB与RBM对不同尺寸网格的内存要求进行了比较。 图中显示，RBM方法可用于分解由CRAY YMP/EL上多达1600万个元素组成的网格。 为了说明可以实现的并行性能，表3中详述的电磁散射模拟中使用的网格现在由RBM分解，获得的相应结果显示在表4中。 对于这个问题，在使用这两种区域分解方法后产生的并行性能是相似的。 区域分解的其他方法尚未被研究，因为很明显，每种方法都有自己的网格大小限制，因此任何方法都不可能允许直接分解这里感兴趣的非常大的网格。 15.3.2 替代网格生成策略 已经研究了提供生产更大网格能力的三种方法。 前两种方法涉及将标准/i-精炼技术应用于生成的网格[7]。 通过h-细化，将一个新节点添加到网格中的每个边缘，然后将现有的四面体细分，并通过节点点的适当重新连接形成新元素。 这种方法的一个缺点是，用户失去了确定最精细网格上可用节点位置的能力，而另一个复杂的是，当考虑边界边缘时，添加的节点需要位于边界表面上。 第三种方法是解决问题的新方法，其中采用了网格生成的并行方法。

在平行计算机上模拟269图3通过长度18A的完美导电飞机对波长A的平面波的散射-散射电场15的x\textbackslash\{\}成分的计算轮廓。 S.Z.I第一方法 第一种方法允许在较大的网格上进行模拟，是通过生成的网格的\textasciicircum{}-细化来生成网格，并使用RBM作为区域分解策略，然后受到RBM程序的内存需求的限制。 作为通过遵循这种方法可以合理建模的问题类型的说明，考虑了完美导体对电磁波的散射。 散射器被理解为长度为18A的完整飞机，其中A是入射波的波长。 遵循此配置直接生成的网格的/i-refinement，使用的网格由15182 752个元素、2个553495个节点和17 872347个边缘组成。 该解决方案在36个入射波周期后输出，在使用256个处理器的CRAY T3D上计算大约需要3个小时。 飞机表面散射电场X\textbackslash\{\}分量轮廓的计算分布如图3所示。 应该注意的是，经常表达的对/i-refinement产生的网格质量的担忧在当前情况下似乎无效，因为总是采用网格质量增强技术[10]。 图4显示了这一点，其中180万个直接生成的网格和210万个元素的h精炼网格的元素二面体角的分布相当相似。

270 MORGAN, WEATREIZILL, HASSAN, BROOKES, MANZAlEI \& SAID 0 20 40 60 80 100 120 140 160 角度图4 在大约相同大小的h精炼网格上直接生成的atld上元素二面体角度变化分布的比较。 15.3.2.2 第二种方法 另一种可以在拉丝网格上进行模拟的方法是最初使用可以在可用计算资源规定的限制内生成和分解的网格。 这个网格被分解成相同数量的区域，因为有处理器可用于执行方程解。 然后，通过对每个分解区域进行11次细化，产生所需尺寸的xnesh。 在这个waby中，不需要进一步的域解-nposition，因此，相关的内存约束被删除。 从理论上讲，根据ea.rlier提到的h-精炼技术的抽包，这种方法可用于生产任何所需尺寸的nleshes。 由于最终的lnesh永远不会完全组合在一起，因此在评估方程（15.3）的边缘权重时需要特别注意。 采用不同的策略来克服这个问题，但是，通过在子域边界复制接口边缘，将pa.rtia1权重与每个接口边缘相关联，从每个subdorna.in单独累积边缘贡献，以及tlien sumn\textasciitilde{}ing适当的子域贡献\textasciitilde{}\textasciitilde{}s来解决。 通过考虑在飞机配置上模拟稳定无黏流动，证明了使用这种方法的可行性。 使用R.SB将初始网格细分为32个子域。 配置的更精细的网格，由1 809 651个元素t\textasciitilde{}nd 327 090个节点组成，是

在并行计算机上模拟271图5在飞机配置计算的压力轮廓上稳定无黏流动，通过在每个子域内使用单个级别的/i-细化产生。 流动条件以0.85的自由流马赫数和2度的攻角来定义。 飞机表面压力轮廓的计算分布如图5所示。 15.3.2.3 第三方法 一种通用方法，它允许在任意大网格上模拟问题，是将平行策略应用于网格生成阶段[16]。 采用的策略采用单一程序多数据模型和经理/工人结构，其中工人执行网格生成器的子程序。 从域边界的离散化开始，产生域的初始Delaunay离散化。 然后，使用体积标准的元素聚集来将域的几何划分为一组NoP子域。 然后，每个子域都可以使用网格算法的串行版本独立网格化。 管理员使用PVM或MPI消息传递库将子域网格数据分发给工人。 在实践中已经证明了

272 MORGAN、WEATHERILL、HASSAN、BROOKES、MANZARI和SAID表5平行3D 网格生成，有16个子域。 子域元素点面 1 1243 837 200733 34 216 2 1065 724 174316 37724 3 778 201448 40 230 4 1225 162 198 944 38 240 5 1311018 212815 40 342 6 1098 964 180 691 42 812 7 983 377 163178 43 730 8 1180 311 192 392 39 774 9 9 977565 161659 41476 10 487 476 176100 40 280 11 921642 154756 47 868 12 1198 717 195286 39 768 13 13 04 066 212249 42 522 14 1200 54 54266 15 61266 106012 43 986 16 动态负载均衡的使用动态负载均衡，其中需要网格的子域多于可用的处理器多，导致高效的计算实现。 应该注意的是，该方法的一个副产品是，可以在任何计算机平台上生成任意的大网格，即使是仅由一个处理器组成的平台。 这是因为初始问题可以简化为一组网格问题，每个问题都可以由任何给定的处理器处理。 通过展示位于两个半球之间的区域离散化时产生的结果来说明了这种方法的潜力[14]。 该区域边界的三角测量由279 256个三角形和139 630个点组成。 这些点的Delaunay三角测量产生了一个初始体积网格，其中包含425 092个四面体，并将其细分为16个子域。 使用并行计算机的8个处理器，每个子域都使用标准的Delaunay方法离散化。 生成的最终网格有171235333四面体，表5提供了有关子域网格大小的信息。 从边界的更精细的初始离散化开始，第二个网格，包含49个168881四面体也通过相同的过程产生。 从表5中可以看出，在实际使用这些分解之前，需要均衡预测的计算负载。 目前正在研究实现这一目标的方法，但希望通过纳入更好的方法来改善情况

在并联计算机上模拟 0 20 40 80 80 100 120 140 180 角度图 6 序列生成、并行生成和大小大致相同的h精炼网格上元素二面体角度变化分布的比较。构建域的初始分解，并使用负载均衡技术对生成的子域进行后处理。 网格是并行生成的事实预计不会对网格质量产生任何重大的不利影响。 图6证实了这一点，该图显示了序列生成、并行生成和大致相同大小的h精炼网格上元素二面体角度变化分布的比较。 16.4 结论在非结构化四面体网格上求解守恒律形式的方程组的显式程序已被并行化，并应用于航空航天工程中感兴趣的许多不同问题领域的模拟。 已经提出并实施了允许将该方法应用于涉及非常大网格的模拟的方法，并确定了悬而未决的困难。 对于某些类别的问题，在设计环境中广泛使用该方法之前，需要进一步开发解决方案算法技术。

274 MORGAN, WEATHERILL, HASSAN, BROOKES, MANZARI \& SAID 致谢 作者感谢英国工程和物理科学研究委员会在研究补助金GR/K42264下提供对爱丁堡并行计算机中心的CRAY T3D的访问，并支持研究补助金GR/J12321、GR/J91234和GR/L18860下的并行化活动。 作者还感谢英国航空航天、空中客车和军用飞机以及法恩伯勒的DERA为黏性流动建模提供的支持。 R。 Said要感谢K提供的部分支持。 R。 S基金会。 参考资料1。 多尼亚，J.，朱利安尼，S.，拉瓦尔，H. \& Quartapelle，L.，有限元对流扩散问题的时间准确解，应用力学和工程中的计算机方法45，123-146，1984年。 2. 格林霍，C。 \& Fowler，R. F.，非结构有限元网格的分区方法，报告RAL-94-092，Rutherford Appleton Labora tory，Didcot，1994年。 3. Hassan，0.，Morgan，K.，Probert，E. J。 \& Peraire，J.，三维黏性流动的非结构化四面体网格生成，国际工程数值方法杂志39，549-567，1996年。 4. 希尔斯，D。 P.，数值空气动力学：从工业角度看过去的成功和未来的挑战，在J.-A. Desideri等人，编辑，应用科学计算方法'96-受邀讲座和特殊技术会议，John Wiley \& Sons，Chichester，166-173，1996年。 5. Hirsch，C，内部和外部流动的数值计算-第2卷，Wiley-Interscience，Chichester，1990年。 6. 琼斯，J。 W。 \& Weatherill，N。 P.，ViPar：大型数据集的并行可视化，提交给国际流体数值方法杂志，1997年。 7. 洛纳，R.，摩根，K. \& Zienkiewicz，0。 C，可压缩Euler 方程的自适应网格加密，在I中。 Babuska等人，编辑，有限元计算中的精度估计和自适应细化，John Wiley，Chichester，281-297，1986年。 8. 曼扎里，M。 T.，Hassan，0.，Morgan，K. \& Weatherill，N。 P.，3D非结构化网格上的湍流计算，分析和设计的有限元素，1997年（印刷中）。

在平行计算机上模拟 275 9. 摩根，K.，贝恩，L. B.，Hassan，O.，Probert，E. J。 \& Weatherill，N。 P.，在M-0中模拟具有移动边界的3D非稳无黏性可压缩流。 Bristeau等人，编辑，21世纪的计算科学，John Wiley，Chichester，347-356，1997年。 10. 摩根，K.，布鲁克斯，P. J.，哈桑，0。 \& Weatherill，N。 P.，电磁波散射问题模拟的并行处理，应用力学和工程中的计算机方法，1997年（印刷中）。 11. 摩根，K.，哈桑，0。 \& Peraire，J.，模拟分段均匀介质中电磁散射的时间域非结构化网格方法，应用力学和工程中的计算机方法134，17-36，1996年。 12. 摩根，K.，佩雷尔，J.，佩罗，J. \& Hassan，0.，使用非结构化网格计算三维流动，应用力学和工程中的计算机方法87，335-352，1991年。 13. Peraire，J.，Peiro，J. \& Morgan，K.，欧拉流经过安装的航空发动机的有限元多网格解，计算力学11，433-451，1993年。 14. 说，R.，Weatherill，N。 P.，Morgan，K. \& Verhoeven，N。 A.，用于非常大的网格的分布式Delaunay网格生成，提交给计算力学，1997年。 15. Simon，H.，用于并行处理的非结构化问题的分割，计算系统工程2，135-148，1991年。 16. Verhoeven，N.，Weatherill，N. P。 \& Morgan，K.，2D并行Delaunay网格发生器中的动态负载均衡，在A中。 Ecer等人，编辑，Parallel 计算流体动力学：使用并行计算机的实现和结果，爱思唯尔科学，阿姆斯特丹，641-648，1996年。 17. Weatherill，N。 P。 \& Hassan，O.，具有自动点创建和施加边界约束的高效3D Delaunay三角测量，国际工程数值方法杂志37，2005-2039，1994年。 18. Weatherill，N。 P.，Hassan，0.，Morgan，K. \& Marchant，M. J.，非结构化网格上的大规模计算，在F中。 Benkhaldoun和R。 Vilsmeier，编辑，复杂应用有限卷会议记录，Hermes，巴黎，77-98，1996年。 19. 威尔考克斯，D。 C，湍流 CFD建模，DCW工业公司，加拿大，1993年。 20. Zienkiewicz，0。 C. \& Morgan，K.，有限元和近似，John Wiley，纽约，1983年。

\clearpage
\chapter{优化CFD代码和算法，以便在Cray计算机上使用}
优化CFD代码和算法，以便在Cray Computers Laurence B上使用。 Wigton 1 16.1 导言 在本文中，我们将讨论波音CFD研究人员生活中相当忙碌的一年。 在这一年里，我们学到了很多关于使用新的Cray T90计算机以及实施新的算法和编码技术来提高我们两个主力代码的性能，即TRANAIR [10, 4]和TLNS3DMB [8, 9]。 在T90计算机上进行了广泛的测试，包括用CAL（Cray汇编语言）运行代码，以便我们可以将性能特征与前几代Cray计算机进行比较。 我们关于T90硬件和软件的许多发现和观察都传回了Cray Research/SGI，这样他们将来就可以为我们和整个CFD社区提供更好的支持。 去年，我们还大幅更新了TRANAIR设计的线性代数例程，因为“设计灵敏度计算”主导了TRANAIR设计/优化运行的成本。 特别是，我们添加了带有特征向量放缩的GMRES块，使用锯齿状对角格式改进了矩阵向量乘法，并使用充分利用T90上新硬件功能的特殊CAL代码加快了三角形矩阵求解。 1 波音商用飞机集团。 Frontiers of 计算流体动力学 - 1998 编辑：David A. Caughey和Mohamed M. 哈菲兹 ©1998 世界科学

278 WIGTON 对于TLNS3DMB，我们引入了一个新的算法，以更快地计算湍流模型所需的距离函数。 我们还创建了TLNS3DMB的非核心版本，我们称之为NSOCR，以利用Cray计算机上的SSD。 NSOCR也可以在核心中运行，速度与原始TLNS3DMB一样快。 NSOCR和快速距离计算程序被送回给美国宇航局，这样它们将来也能更好地支持我们。 16.2 包括T90在内的Cray计算机 由于所有重点都放在分布式内存并行计算机上，有趣的是，传统的Cray计算机仍然是许多商业/工程环境中首选的计算平台。 其原因是它们可靠，可以同时运行许多大型作业，具有良好的软件支持，特别是随着T90的推出，它们实际上是物有所值的计算机。 与其他计算机相比，Crays具有快速的CPU、出色的记忆体频宽以及足够大的记忆体和SSD，以便它们能够满足手头的问题。 尽管其他计算平台取得了所有进步，但Crays仍然是评判所有其他平台的机器。 16.2.1 传统Cray计算机概述型号T90 C90 Y/MP时钟速度440Mhz 240Mhz 160Mhz管道/CPU 2 2 1 Memspeed 34或56 23 17管道：每根管道有3个内存端口和一套完整的分段功能单元。 在矢量模式下设置了一条管道。 在每个时钟周期的初始化成本后，管道可以做：2个获取+1个加+1个乘法+1个存储。在Cray语言中，管道在时钟周期内完成的操作集被称为“钟”。 Memspeed：向量初始化成本的主要部分是“memspeed”，即从内存中获取第一个操作数所需的时钟周期数。 小型T90为34，大型T90（具有更多的内存板电路）为56。 请注意，内存速度跟不上CPU速度。

优化CRAY计算机的代码和算法 279 16.2.2 SSD除了高速内存外，Cray计算机还可以支持SSD，它就像一个非常高速的磁盘（在我们的T90上，SSD传输速率为800Mw/sec vs. 条纹磁盘为30Mw/秒）。 多年来，固态硬盘变得更大、更快、更便宜。 SSD允许我们解决Cray上的大问题。 然而，充分利用SSD需要一个核心外代码，这增加了编码的复杂性。 16.2.3 关于T90的初步猜测 人们可能会认为T90和C90一样，只是它的时钟速度快1.8倍。 这并不完全正确，因为正如我们所看到的，T90的向量启动速度较慢（由memspeed控制）。 此外，T90具有内存带宽限制，对指令排序的依赖性更大，对所谓的CMR指令的灵敏度更高。 我们将讨论这些要点中的每一个。 16. Z.S.I 内存带宽模型 T94 C90 YMP/8 numcpus 4 16 8 Pipes 8 32 8 Memreq 24 96 24 banks 64 1024 256 bankbusy 7 6 5 Memavail 9 170 51 这里的“Memreq”是所有管道可能需要的内存（Memreq = 3 * Pipes），“banks”是数字内存库。 访问特定内存库后，在“银行繁忙”的时钟周期过后，无法再次访问它。 对于内存密集型操作，时钟期间可用的内存字数为：Memavail = banks/bankbusy 对于C90和Y/MP Memavail > Memreq。 对于T90，我们有Memavail = 9和Memreq = 24，因此对于内存密集型操作，T90可以减慢9/24的系数。 严格基于CPU性能，人们可能会认为T90每个CPU的速度应该比Y/MP快6倍左右。 毕竟，T90的时钟速度是3倍，每cpu有2个管道。 然而，事实上，内存带宽

280 WIGTON限制将此优势减少到：（9/24）6 = 2.25。 这接近于我们用T90取代Y/MP时观察到的生产工程代码的平均性能提升。 这项协议是偶然的，因为平均生产作业不会以峰值内存密集型模式运行。 然而，T90上的内存库冲突比Y/MP上的内存库冲突更多，即使是平均性能也会受到影响。 增加矢量启动时间，提高对指令排序和CMR的敏感性也起到了一定的作用。 16.2.3.2 指令排序 在练习为BLAS常规SAXPY编写CAL代码时，注意到CAL指令的简单反转对T90的影响比对C90的影响要大得多：2个版本的SAXPY机器T94 NASA C90 SAXPY1 639 674 SAXPY2 1208 697 以下电子邮件中向Cray Research发送了有关此的说明。 Frank Chism：（cc：Ecale、Hilmes、Whitaker、Heroux）似乎T90上的尾随并不总是坏的。 我刚刚发现，对“坏”的asymc代码进行小幅更改，使其在T90上运行速度翻倍。 在“坏”代码不对称中，我使用了像乘法+添加+获取+获取+存储等链，如：V2 S7*RV0;SA*X V3 V2+FV1;SA*X+Y AO A2+A5 VO,A0,1;获取更多X' = look 256 forward AO A3+A5 VI,A0,1;获取更多Y's look 256 forward AO A3,A0,1 V3;存储回Y，如果我将商店移到获取之前，乘法+添加+存储+获取+获取，如：

优化CRAY计算机的代码和算法 281 V2 S7*RV0;SA*X V3 V2+FV1;SA*X+Y AO A3,A0,1 V3; store 1 AO A2+A5 VO,A0,1; Fetch i AO A3+A5 VI,A0,1;Fetch i;获取更多X's look 256 forward 获取更多Y's look 256 然后我们不是在T90上以639兆浮点运行，而是以1208兆浮点运行。 这几乎和f90代码一样好。 似乎商店会释放你，这样你就可以开始获取链中正在使用的收银机。 在移动的C90上，商店说明也帮了一点忙（697兆浮点对674），但效果远没有T90那么显著。 一些硬件大师将不得不解释这一点。 拉里·威格顿16.2。 S.S完整内存参考（CMR）在从内存中获取时，Cray编译器担心被获取的数字可能已通过之前发出的存储命令发送到内存位置，但可能仍在内存网络上传输中。 每当有这种可能性时，编译器就会生成一个完整内存引用（CMR）指令，以确保所有内容都已按照之前发出的存储命令到达内存中。 不幸的是，编译器倾向于生成不必要的CMR命令。 与前几代Cray计算机相比，可疑的CMR命令对T90的负面影响更大。 考虑SAXPYI的以下程序：子例程saxpyi.test(m,a,x,indx,y)实数x(m),y(*)整数indx(m) CDIR\$ IVDEP do 10 i = 1, m k = indx(i) y(k) = a * x(i) + y(k) 10 继续返回结束

282 WIGTON 根据C90和T90上的Cray文档（CF90命令和指令参考手册SR-3901 1.0，第37页），指令：CDIR\$ IVDEP等同于：CDIR\$ IVDEP [SAFEVL=128]，128是可以分配给SAFEVL的最大值。 因此，cf77编译器担心一组128 y可能会占用上一组y使用的一些内存位置。 这迫使编译器在每个组被向量指令处理后发出CMR指令。 当然，用户真的意味着SAFEVL=无限，所以不需要CMR说明。 人们可以获取cf77编译器生成的CAL代码，并手动删除不必要的CMR指令。 生成的代码在T90上快50\%，在C90上快约10\%-15\%。 CMR问题也存在于f90的早期版本中，但似乎在f90版本2.0.3 16.2.4中已经修复了。为什么我们仍然喜欢T90，尽管T90有各种明显的弱点，但它仍然比C90更快地运行平均工作。 此外，它比C90便宜得多。 此外，Cray正在迅速改进T90硬件和软件。 特别是，当我们提请他们注意各种缺陷时，Cray正在调整编译器。 16.2.5“杀手零”如以下程序所示，在Cray计算机（不仅仅是T90）上，有可能有：0 + 2 = 0程序奇怪c引入杀手零的八进制表示：等价（ia，a）ia = 0'0411650000000000000000' b = 2.0

优化CRAY计算机的代码和算法283 c = a + b打印*，"a,b,c= ",a,b,c c 现在以八进位制（*,1）a,b,c 1格式打印所有内容(lx,"a,b )c=",/,(025))停止结束 cf77下程序"奇怪"的执行：triton<35> cf77 weird.f -o weird triton<36> weird a,b,c= 0.,2.,0. a,b,c= f90下的程序「weird」的执行：triton<37> f90 weird.f - weirdo triton<38>奇怪的a,b,c= 0。 E+0，2.，0。 E+0 a,b,c=“杀手零”一词是由Cray/SGI的Frank Chism创造的。 请注意，Killer Zero在Cray计算机上打印为0，但在功能单元中不像0。 我们第一次在TLNS3DMB的核心外版本中遇到了Killer Zero。 Killer Zero是由Cray内存分配例程hpalloc生成的。 如果用户正确初始化了hpalloc分配的内存，这不是问题。 我们的问题可以追溯到TLNS3DMB原始版本中的错误，该错误由Veer Vatsa在后来的代码版本中修复。 杀手零是未归一化数字的特殊情况。 在Cray计算机上，如果浮点数不是完全为零，那么曼蒂萨的第一个位必须是1，否则该数字不会归一化，并会造成问题。

284 WIGTON 16.2.5.1 跟踪非规范化数字 理想情况下，当用户尝试使用非规范化数字时，硬件应该停止执行。 如果没有这种硬件更改，用户很难追踪非规范化数字的使用（Cray在处理不定数方面要好得多）。 不幸的是，对于Cray来说，实施所需的硬件更改将非常昂贵。 我们已经要求Cray Research/SGI实施软件改进：1. 编译器中的调试模式，该模式生成检查未规范化数字的指令。 2. 更改格式化的写入语句、打印语句和符号转储，这些转储将标记未规范化数字。 16.2.6 Cray仍然是代码调试之王，尽管存在“Killer Zero”的问题，我们仍然认为Cray是调试代码的领导者。 Crays具有不常见的功能，包括将核心预设为不定的能力，以及在出现错误时生成符号转储的能力。 此外，与之前的Cybers不同，Cray计算机一旦发生浮点异常就会停止执行，这使得追踪代码中的错误变得更加容易。 16.3 TRANAIR TRANAIR [10, 4]是波音的通用几何全电位加边界层代码。 此代码是核心外的；它在带有SSD的Cray上运行得最好。 TRANAIR使用Drela的2D+扫/变整体边界层代码。 全势流方程在自适应精炼的octree笛卡尔网格上求解。 全势方程和边界层方程被耦合在一起，并使用近似牛顿迭代求解。 这是使用GMRES和ILUT 预处理完成的。 TRANAIR设计/优化使用直接而不是伴随的配方。 成本以求解与线性化设计灵敏度相关的多个右侧（MRHS）的线性方程占主导地位（我们计算每个设计变量调整时整个流场的变化）。 我们现在将描述TRANAIR解决MRHS问题的方法的改进。 16.3.1 块GMRES 当用MRHS求解线性方程时，块GMRES结合了与每个右侧解相关的所有Krylov空间。 这

优化CRAY COMPUTERS 285的代码和算法会自动增加子空间的维度，计算出每个方程的最佳解。 Saad的书[6]中出现了一个很好的讨论。 与标准GMRES相比，块GMRES需要更多的Krylov向量存储空间。 当然，还需要更多的正交（SDOT和SAXPY）操作。 然而，块GMRES收敛得更快，因此需要更少的预条件矩阵-向量乘法。 此外，正如后面所讨论的，额外的好处来自处理内环内的多个右侧，包括减少内存引用和SSD I/O。 16.3.2 通货紧缩的GMRES通常小特征值慢GMRES 收敛。 当执行GMRES迭代时，具有通货紧缩的GMRES（DGMRES）计算小特征值和相应的特征向量。 这些特征向量被添加到Krylov空间，以加快收敛的速度。 16.3.3 dbgmr 根据合同，我们让Yousef Saad和他的研究助理Andrew Chapman制作dbgmr。 这是具有特征向量通货紧缩的GMRES块。对于“艰难”的测试问题，dbgmr通常比标准GMRES有一个数量级的改进。 对于“简单”的问题来说，改进并不那么显著。 用于TRANAIR的ILUT 预处理非常有效，因此它确实倾向于使TRANAIR矩阵易于解决。 然而，对于跨声速TRANAIR案例，dbgmr将解决问题所需的迭代次数减少了2.5倍（GMRES块本身提高了1.9倍）。 除了引入dbgmr外，我们还改进了矩阵向量乘法例程和与ILUT 预处理相关的三角矩阵求解。 16.3.4 稀疏矩阵向量乘法 TRANAIR中用于将稀疏矩阵乘以向量的旧方法本质上是Blelloch等人开发的分段和（SEGMV）程序。 [2]。 SEGMV是由TRANAIR团队独立发现的，与通常用于在TRANAIR中存储矩阵的压缩行格式配合得很好。 虽然SEGMV比矩阵向量乘法的天真逐行实现更快，但它存在某些问题。 首先，需要CAL代码来捕捉SEGMV和TRANAIR团队的全部优势

286 WIGTON并没有为此目的使用CAL。 更重要的是，SEGMV覆盖了矩阵的内容，这使得它不适合多个右侧问题。 因此，我们决定使用锯齿状对角线（JAD）技术，萨德的书[6]中也有描述。 JAD要求重新格式化矩阵，这对TRANAIR来说尤其不小，因为矩阵存储在核心之外。 然而，这种重新格式化成本很容易在正在进行的大量矩阵向量乘数中摊销。 JAD可以在FORTRAN中高效地实现，正如下一节所述，它允许我们在内圈内处理MRHS。 总体而言，基于SEGMV的实施受损，JAD被证明比旧的TRANAIR代码快2.3倍。 16.3.4-1 内环内多个右侧处理MRHS（多个右侧）在发送给Yousef Saad的电子邮件中讨论了内环内：优素福·萨德：让我再评论一下将多个右侧放在Cray的内环内。 我们总是说Cray是一台矢量机器，但它实际上是一台管道机器。 它能够管道：从内存中获取2次，1次加法，1次乘以1次存储到内存，在启动此管道的初始化成本后，它能够在一个时钟周期内执行所有这些操作。 Cray指的是所有可以作为钟声一起管道的操作，他们喜欢用钟声来衡量执行a.向量循环的执行成本。 因此，如果我们看一下循环：做10 i=l，n k=ip(i) y(i)=y(i)+a(i)*x(k) 10继续在一个钟声中，Cray可以获取ip(i)和a(i)。 此时，它不能进行任何乘法或加法操作，所以这就是第一个钟声期间发生的一切。 在第二个钟声中，Cray将获取y(i)和x(k)（2个获取），将a(i)乘以x(k)，加到y(i)并存储到y(i)。 在第二个钟声中，它实际上做了2个获取乘法、加法和存储。 底线是右手

CRAY COMPUTERS 287侧环最佳化程序和演算法，在2个钟声中执行2个浮点操作。 多个右侧循环：做20 i=l,n k=ip(i) y(i,l)=y(i,l)+a(i)*x(k,l) y(i,2)=y(i,2)+a(i)*x(k,2) y(i,3)=y(i,3)+a(i)*x(k,3) y(i,4)=y(i,4)+a(i)*x(k,4) 20继续，我们在5个钟声中执行8个浮点运算。 因此，同时处理4个右侧的多个右侧循环是单个右侧模式的1.6倍。 如果我们考虑到“a”和“ip”的矩阵信息的SSD传输，那么我们必须为每个循环添加2个钟声（一个钟声读取a，一个钟声读取ip，我们将在核心中存储x和y）。 在这种情况下，do 10循环在4个钟声中执行2个浮点运算，而do 20循环将在7个钟声中执行8个浮点操作。 在这种情况下，多个右侧的效率更高，系数为（8/7）/（2/4）=32/14，这就是为什么我说你的想法值2的系数。 拉里·威格顿评论：除了减少固态硬盘的使用外，在实际实践中，处理内环内的MRHS会导致CPU时间减少1.3倍。 简单的钟声论证过于乐观；如果Cray在每根管道上有更多的向量寄存器，它会更准确（而且计算机会更快）。 JAD的FORTRAN版本本质上比SEGMV的FORTRAN实现快1.8倍左右。 因此，对于MRHS来说，新的TRANAIR矩阵向量乘法速度加快了1.3 * 1.8 = 2.3。 16.3-4.2「PreferVector」指令我们想编码多个右侧回圈，例如：做30 irhs = l.nrhs do 20 i = l,n k = ip(i) y(i,irhs) = y(i,irhs) + a(i) * x(k,irhs) 20 继续 30 继续

288 WIGTON，但这效果不好。 根据Cray Research，我们应该能够像这样编码多个右侧循环：CDIR\$ PREFERVECTOR do 20 i = l,n k = ip(i) fac = a(i) do 30 irhs = l.nrhs y(i,irhs) = y(i,irhs) + fac * x(k,irhs) 30 continue 20 continue 但到目前为止，这也没有奏效。 如电子邮件中所述，手工展开循环取得了最佳性能。 16.3.5 三角矩阵求解 三角矩阵求解，也称为正向/向后置换，是TRANAIR最昂贵的部分。 通过引入列格式而不是行格式（在Cray计算机上，短的SAXPYI比短的SDOTI更快），在内部循环中处理MRHS，删除FORTRAN编译器生成的虚假CMR指令，最后编写利用T90上新的双收集硬件指令的CAL代码，取得了巨大的改进。 在T90上，三角形矩阵解码比旧代码快4.2倍。 有趣的是，T90上的三角形矩阵求解比C90快3倍。 这是一个重要的例子，说明T90比C90快3倍，尽管时钟的速度只有1.8倍。 1v6.3.6 TRANAIR Design Cray配置文件的配置文件结果在T94上运行，用于具有61个设计变量的新旧TRANAIR代码。 新代码一次解决大约10个右侧。 旧代码新代码三角形解决8,748,100 844,060矩阵-向量乘以2,803,201 481,792 SAXPY+SDOT 187,156 643,813 总计11,738,457 1,969,655

优化CRAY计算机的代码和算法289评论：•配置文件采样率为512微秒。•dbgmr将迭代（因此将三角形求解和矩阵向量乘法的数量）减少了2.5倍•三角形求解现在快4.2倍。 •矩阵向量乘数快2.3倍。t更大的Krylov空间诱导更多的SDOT和SAXPY活动。 •整体改善是6的因素，这对我们来说是一个很大的胜利。 16.4 TLNS3DMB TLNS3DMB [8, 9]是NASA Langley开发的3D块结构匹配网格Navier-Stokes代码。 它采用詹姆森型技术，包括多网格、Runge-Kutta时间行进和隐式残余平滑。 TLNS3DMB尽可能推进解决方案，就好像网格在一个块中一样。 所有区块都分为相同数量的多网格级别。 跨接口边界无缝完成多网格。 Runge-Kutta时间行进的每个阶段都访问了所有街区的每个网格点。 只有隐式残余平滑是逐块完成的。 16.4.1 核心外TLNS3DMB代码（NSOCR）为了充分利用SSD，编写了TLNS3DMB的核心外版本，我们称之为NSOCR。 大部分转换工作都是使用用FORTRAN编写的自动化工具完成的。 然而，在块之间传递幽灵细胞信息的算法被更改（手动），以允许更高效的核心外实现。 运行测试用例来验证NSOCR生成的结果是否与TLNS3DMB生成的结果是二进制相同的。 NSOCR将数组放在数据库上。 小阵列保留在核心中，而大型阵列则根据需要进出固态硬盘。 用户可以控制一个名为“incore\_size”的参数，该参数决定了阵列在移动到SSD之前必须有多大。 通过将incorejsize设置为非常大的数字，大于任何数组尺寸，那么代码以纯核心模式运行。 在这种情况下，NSOCR和TLNS3DMB一样快。 通过将incore\_size设置为较小的值，例如10000，那么NSOCR主要以非核心模式运行。 在这种情况下，访问SSD所需的时间约为CPU进行计算所需的时间的23\%，这被认为是非常可接受的开销。 由于NSOCR在核心内和核心外模式下都运行良好，我们计划使其成为TLNS3DMB的唯一版本。 NSOCR被还给了美国宇航局。

290 WIGTON 16.4.2 距离计算 最近开发的湍流模型，如Baldwin-Barth（[1]）和Spalart-Allmaras（[7]），要求用户计算从字段网格中每个点到所考虑的配置的距离。 对于涉及数百万网格点的计算，即使在Cray C90级计算机上，计算距离函数的天真方法也很容易消耗数小时的CPU时间。 这些距离计算非常昂贵，以至于一些代码开发人员选择避免执行适当的距离计算，从而影响其代码的准确性。 在本节中，我们希望讨论一种计算距离函数的有效方法。 天真算法：对于每个NF场网格点，天真算法只需计算到每个NS表面点的距离，并选择这些距离的最小值。 成本：NF * NS快速r算法：构建大致\textbackslash\{\}/NS框，每个框包含大致y/NS表面点。 对于每个字段网格点，计算到每个框的距离。 选择最近的框，并计算到框中每个表面点的距离。 如果与最接近的方框中包含的点的最小距离小于与其余方框的距离，我们就完成了。 否则，检查第二个最近的盒子等。 经验表明，平均而言，人们必须看1.5个盒子。 成本：NF * y/NS 盒子的建造：从一个包含所有表面点的大盒子开始。 以最长的方向划分方框。 选择分割平面，使一半的表面点位于每边。 递归地进行。 当盒子有y/NS或更少的表面点时停止。 作为最后的装饰，一旦建造了一个盒子，我们就会把它缩小，这样它几乎只能容纳里面的所有表面点。 盒子的二维示例如图所示。 1. 当然，这个概念在3D中同样有效。 图1 距离计算中使用的方框示例

优化CRAY计算机的代码和算法291快速距离功能已返回美国宇航局，现在在TLNS3DMB [8, 9]和CFL3D [5]中可用。 该方法大大节省了时间。 事实上，一个在C90上需要3个小时的案例现在只需要5分钟。 在实践中，该方法已被证明比其他方法更快（基于octrees等），因为它完全矢量化。 当然，它也是平行的。 最初，并行代码在cf77和f90下都失败了。 Cray放弃了对cf77的支持，但他们立即修复了f90中的错误。 致谢 在这项相当紧张的研究工作中，与我不知疲倦并随时准备帮助“支持小组”交换了数百封电子邮件。 Yousef Saad教授和他的研究助理Andrew Chapman提供了升级TRANAIR线性代数包所需的关键要素。 我在Cray Research/SGI的联系人，包括Frank Chism、Dave Whitaker、Mike Heroux和David Ecale，在处理我关于Cray计算机和高效编码技术的许多问题时反应非常灵敏。 参考资料1。 鲍德温，B。 S.和Barth，T. J.，“高雷诺数壁界流量的单方程湍流运输模型”，AIAA论文91-0610。 2. Blelloch G.E.，Heroux M.A.和Zagha M.，“向量多处理器上稀疏矩阵计算的分段运算”，CMU-CS-93-173，卡内基梅隆大学，1993年8月。 3. Chapman，A.和Saad，Y.，“放气和增强的Krylov子空间技术”，数值线性代数与应用，卷。 4（1），1997年，pp. 43-66。 4. Johnson，F.T. “TRANAIR矩形网格方法解决复杂配置的非线性全势流方程”，Sadhana，第16卷，第2部分，1991年10月，pp. 165-177。 5. Rumsey，C.L.，Biedron，R.T.和Thomas，J.L.，“CFL3D：其历史和一些最近的应用”，MIS'； TM-112861，1997年5月。 6. Saad，Y.，“稀疏线性系统的迭代方法”，PWS出版公司，1996年。 7. 斯帕拉特，P。 R.和Allmaras S. R.，“空气动力学流动的单方程湍流运输模型”，AIAA论文92-0439。 也发表在La Recherche Aerospatiale，第1期，1994年，第5-21页。 8. Vatsa，V.N.，Sanetrik，医学博士和Parlette，E.B. “开发灵活高效的多电网多块流动求解器”。 AIAA文件编号 93-0677，1月。 1993年。 9. Vatsa，V.N.和Wedan，B.W. “3D Navier-Stokes方程的多网格代码的开发及其在网格细化研究中的应用”计算机与流体，卷。 1990年18月，pp. 391-403。

292 威格顿 10。 Young，D.P.，Huffman，W.P.，Melvin，R.G.，Bieterman，M.B.，Hilmes，C.L.和Johnson，F.T. “设计优化中的不精确性和全球收敛”，AIAA论文94-4386。

\clearpage
\chapter{使用Nsmb结构化多块求解器在空气动力学中的最新应用}
最近在空气动力学中的应用与NSMB结构化多块求解器Carlos Weber 1，Cyril Gacherieu 1，Arthur Rizzi 2，Anders Ytterstrom 2，Jan Vos 3，Nathalie Duquesne 2和Loic Tourrette 4 17.1介绍使用结构化多块网格的Navier-Stokes求解器已成为研究航空领域空气动力学实际问题的标准CFD技术。 该程序的一个非常常见的变体是基于以细胞为中心的概念的有限体积方法。 这种方法在工业中得到了青睐，因为多块方法促进了围绕复杂几何形状的网格生成，并自然地将问题划分为较小的单元进行并发处理。 这种求解器的最终用途是空气动力学设计。 但1 Centre Europeen de Recherche et de Formation Avancee en Calcul Scientifique, F-31057 Toulouse, France的工程师。 2 航空系，KTH皇家理工学院，S-10044，斯德哥尔摩，瑞典。 3 液压机械和流体力学研究所（IMHEF），瑞士联邦理工学院（EPFL），CH-1015，瑞士洛桑。 4 Aerospatiale Avions，F-31060 图卢兹，法国。 Frontiers of 计算流体动力学 - 1998 编辑：David A. Caughey和Mohamed M. 哈菲兹 ©1998 世界科学

294 WEBER、GACHERIEU、RIZZI等行业尚未例行模拟三维超音速湍流流，因为CPU时间大、内存大成本高，周转时间长。 Navier-Stokes网格在黏性层中放置大量网格点，以便正确解析这些区域，即使使用强大的矢量计算机，这导致周转时间也比一夜之间长得多。 大规模并行计算机的出现为减少此类流模拟的经过时间提供了希望，但只有当软件特别适应并行架构时，才能以高效的方式利用这些超级计算机的巨大潜力。 因此，欧盟成立了EUROPORT项目，以调查这一潜力是否以及如何实现[19]。 我们参与了这个项目，总体结论是，即使现在和未来的高性能超级计算机，对现实工业设计案例的Navier-Stokes模拟仍然对CPU时间要求很高。 对于这种类型的问题，显式方案的收敛速度不够快，需要新的方法来加速收敛。 在某些情况下，多网格加速确实有帮助，但大多数时候，复杂的几何形状及其在黏性区域的高度拉伸细胞以及湍流模型效果使这项技术无效。 因此，需要一个更强大的算法来加速收敛。 从EUROPORT的角度来看，我们现在认为，只有当三个相互关联的领域取得进展时，工业应用的Navier-Stokes模拟才会取得进展：1）更准确的物理模型，即 湍流，2）更快的收敛和更强大的算法，以及3）在矢量/并行超级计算机上实现更好，性能更高。 不幸的是，在一个领域的进步通常在其他领域会适得其反。 例如，一个更复杂的湍流模型增加了并可能改变要求解的方程组的特征，增加了要执行的计算数量，并使给定的并行处理范式的编码复杂化。 同样，具有更高精度和更好收敛属性的算法可能难以实现，并且并行性能可能较差。 因此，任务是平衡所有这些因素，并在所有三个领域实现整体改进。 本文的目的是描述我们在一个致力于控制这三个项目的学术工业合作伙伴财团中正在做的事情。 NSMB代码是在三个研究机构（洛桑的EPFL、斯德哥尔摩的KTH和图卢兹的CERFACS）和两个工业合作伙伴（图卢兹的Aerospatiale和林凯平的SAAB）的联合项目中开发的。 从一开始，NSMB代码就使用有限体积法和空间离散化的中心差分，以及显式多阶段Runge-Kutta解决了可压缩的Navier-Stokes方程

空气动力学295时间积分方案中的应用。 为了加速收敛，使用了局部时间推进，并通过隐性残差平滑和多网格程序引入了一些隐性。 数值方案用人工耗散项进行了增强，以防止奇数/偶数振荡并提高冲击分辨率。 几个湍流模型实现了，从简单的代数Baldwin-Lomax模型到双方程模型和代数雷诺应力模型。 最近，我们为时间积分实施了隐式LU-SGS（下-上对称高斯-塞德尔）方案。 Jameson和Turkel[12]最初将此方案与中心空间方案一起开发，但我们发现收敛在使用高单元宽高比的网格时退化了。 因此，实施了二阶罗氏方案，确保了对角线主导地位。 已经实现了两种变体，原始的标量版本使用光谱半径，以及全矩阵版本，其中数值通量尽可能准确地线性化。 因此，代码的用户现在可以从不同的湍流模型、时间积分方法和空间离散化方案的库中进行选择。 我们说明了一些不同的选择在许多计算示例上产生的影响，并比较了它们在几个矢量和并行机器上的计算性能。 这些案例包括A-翼型上的低马赫数 湍流，三角翼上的跨声速流动和Aerospatiale的AS28G飞机配置周围的跨声速流动。 17.2 物理模型 Navier-Stokes 方程是这里考虑的流体动力学模型。 这些方程中的湍流波动是时间平均的，系统由使用Boussinesq近似的适当模型关闭。 过渡没有建模。 该流程要么被视为完全湍流，要么先验指定过渡位置。 17.2.1 控制方程 在三维笛卡尔坐标（x,y,z）中，以保守形式表示的完美气体的可压缩雷诺平均纳维尔·斯托克斯方程是ot ox oy oz，其中U=（£Ti）（17-2）

296 WEBER、GACHERIEU、RIZZI ET AL T-fK?-??'\} Q-（QTJ-Q7 S\textbackslash\{\} H=（Knf-K？ F) (173) / 0 \textbackslash\{\} (17.4) \textbackslash\{\}Stm/ 术语U mf表示平均流动的五个保守变量，平均流动的非黏性和黏性通量项分别是J\textasciicircum{}、Q™f、U™ f和T™ f、Q™ f、U™ f 这些向量的成分是众所周知的，在这里不重复。 Reynolds平均系使用湍流 Reynolds应力的Boussinesq近似和仅代数方程或一个或多个偏微分方程的湍流模型闭合。 模型的变量是U tm，这些模型中的黏性通量和黏性通量项分别是F\&、Gj™、H\textbackslash\{\}\%和T\textbackslash\{\}m、£j m、U\textbackslash\{\}m最后，模型中的源项成分是S tm。 湍流模型向量在以下章节中指定。 17.2.2 湍流模型 NSMB中提供各种不同复杂度的湍流模型，从代数模型到一方程和二方程模型。 这里使用了四种不同的湍流模型：代数Baldwin-Lomax [2]和Granville [10]模型，Spalart和Allmaras的单方程模型[18]以及Chien [6]的双方程k-£模型。 前三个模型已由Aerospatiale在代码中实施，而后者已由KTH的航空部实施。 Baldwin-Lomax模型已以通用方式实现多块计算[7]。 17.2.2.1 Granville Granville模型[10]基于Baldwin-Lomax模型，并使用相同的管理雷诺平均 Navier-Stokes 方程（17.1）公式。 Granville在Baldwin-Lomax配方中引入了Cki和Ccp的替代表达式，其中包括Coles Wake因子n，以表示有利和不利的压力梯度，直到分离。 S=

空气动力学中的应用 297 Granville假设了湍流边界层 Cfc的外部相似性定律；= 9iT\textasciicircum{}-（17 -5）通过消除Coles Wake因子CCP = \textasciicircum{}FTTT,- Qn.., \textasciicircum{}3\textasciicircum{}（17 -6 ）2CH\{2 - 3C kl + C\textbackslash\{\} t）在平衡压力梯度的情况下，Coles Wake因子在沿溪流方向上是恒定的，并且可以与克劳瑟压力梯度参数经验相关。 Granville提出了一个修改后的克劳瑟压力梯度参数p：- -'■- nr（17.7）适合Dx最后，格兰维尔得出了一个Cu的显式公式，作为修改后的克劳瑟压力梯度参数的函数（3：2 \_ 0.01312\_ 3 0.1724 + /? V'内部涡黏性也必须修改，以包括压力梯度。 线性变化（fty）需要修改，因为墙附近的总剪应力不再恒定。 根据Galbraith等人[8]的论文，墙附近的混合长度变为：linneT \textasciicircum{} K y\{-)\textasciicircum{}[l - exp\{- y-)\}。 （17.9）Tw A拒绝使用（it\textbackslash\{\}w\textbackslash\{\}）来确定诱导非线性行为的r，Thomas和Hasani [20]提出了一个r（y）的代数插值公式，该公式在整个边界层中是合理准确的：- \textasciitilde{} 1 + £r? - (3 + 2£>7 2 + (2 + O\textasciicircum{} 3 (17-10) Tw，n = | and f = -\textasciicircum{}P-。 最近，Granville [10]建议，最平滑的合并（与对数定律截距及其斜率）是通过以下阻尼常数变化实现的：A\textasciitilde{} \textasciicircum{}-r (17.11) (1 + \&P+) 5和

298 WEBER、GACHERIEU、RIZZI ET AL P+ = -\textasciicircum{} (17-12) puT3 ax，如果p+ > 0，则b = 12.6，如果p+ < 0，则b = 14.76。 17.2.2.2 Spalart-Allmaras Spalart和Allmaras的单方程模型[18]使用与上述代数模型相同的质量、动量和能量的平均流方程，以及湍流黏度的额外传输方程，以获得涡黏性的表达式。 这给出了湍流模型状态向量（17.2）中的U tm = v，以及Eq中的其他项。 （17.1）是Tg = UP，g\textbackslash\{\}™ = vv，？ C = ™ Ttm \_ v + g gg Qtm \_ v + g di> ytm \_ v + g gg c«m g-, C i2 (V£) 2 f ( \textbackslash\{\}( \textasciicircum{}给出以下运输方程的湍流运动黏度，j£ = ctiSv +\textasciicircum{}[v-((u + i>)VD)+c b2(Vi>)2(Vi>)2'\textbackslash\{\}-cwlful(r)(\textasciicircum{}\textbackslash\{\} \textasciicircum{}\textasciicircum{}\textasciicircum{} production v s s-*™ -\textasciicircum{} coni-ecijon 消散（17.13) 涡黏性被定义为Ht = pvfvi = pvt (17.14) 为了确保v等于对数层、缓冲层和黏性子层中的Kyu T，阻尼函数/„i被定义为U = j\textasciicircum{} (17.15)作为全局部变量A A = - (17.16)

空气动力学中的应用 299 Spalart-Allmaras建议使用涡量 \textbackslash\{\}u\textbackslash\{\}大小的平方来定义生产函数S。 必须修改该项，以保持其对数层行为（5 = U T/（Kh））一直到墙S=S 1/2 + 7T\textasciicircum{}2/»2（17.17）（An），这是在函数f v2 U = 1 - IT \textasciicircum{} r（17.18）的帮助下完成的，销毁项应在边界层的外部区域消失。 Spalart-Allmaras提出了函数1 + r 6, 1/6 \textasciicircum{}r)=9\textbackslash\{\}-f-\textasciicircum{}-\textbackslash\{\} (17.19)，参数r r=idr§ (17-20) r和fw在对数层中都等于1，在外部区域都等于。 函数g只是一个限制器，可以防止/„,, g = r + cw2(r6 - r)的大值。 （17.21）Spalart-Allmaras模型中的常数为cbl = 0.1355，c b2 = 0.622，c w2 = 0.3，c vl = 7.1 2 \_ CM，（l + c 62）it2 v = T；，C WI = 7j -I，c w3 = 2，A = 0.41。 17.2.2.3 Chienk-e 双方程模型的起点是Boussinesq近似，可以写成：,du, Buj 2dui,,2 - -puru/1 = M\textasciitilde{} + g\textasciicircum{}' 3 Q\textasciicircum{}Sij) \textasciitilde{} \textasciicircum{}pkSij，其中，根据法夫尔平均法，u;是瞬时速度的质量平均部分，u是波动部分。

300 WEBER, GACHERIEU, RIZZI ET AL 双方程湍流模型使用两个额外的传输方程来闭合法夫尔平均纳维耶·斯托克斯方程（Eq. （17.1））。 对于k-e模型，第一个方程确定湍流动能k（=g-ujwj），第二个方程确定湍流动能e的耗散，并与k一起使用来定义湍流长度尺度。 因此，在状态向量中增加了两个额外的项：U<m=(\textasciicircum{}) (17.22) 这个和低阶模型之间的基本区别在于湍流黏度的定义。 k-e模型的湍流黏度不是与代数长度尺度相关联，而是由关系确定的：r Pk2与C„ = 0.09。 Boussinesq近似|/5fc中的最后一个项，对应于湍流运动引起的湍流压力，并被新增到动量和能量方程中非黏性通量中的压力项中，产生p* = p + \textbackslash\{\}\textasciitilde{}pk--由于与k和e方程的耦合，能量方程略有修改，并解释了额外的扩散项<7fc，-它也出现在k方程中。 用前面的符号写成的k和£的最后两个方程是：pue \} ' "* \textasciitilde{} y pve J ' in \textasciitilde{} \textbackslash\{\} pwe J是通量，对于黏性通量，扩散项写成0>i = /i + JH\_\textbackslash\{\} dk\_ Prk) dxi JH\_\textbackslash\{\} de\_ Pre ) dxi，其中Prk和Pr£类似于湍流 Prandtl-Schmidt数。 这个基本的双方程模型通常是“高雷诺数”模型，它不考虑湍流和湍流之间的相互作用

空气动力学中的应用301流体黏度，因此不适用于墙壁附近的区域。 为了允许湍流方程一直集成到墙壁上，许多研究人员提出了不同的黏性校正。 在这项研究中，使用了Chien [6]的低雷诺数版本。 在这个模型中，为耗散率的各向同性分量求解了一个传输方程。 由于耗散的各向同性分量，而不是耗散本身，k和e的源项通过引入校正项略有修改，然后写入：Stm= U) = Uj4P*-C, 2/J±-§exp(-0.5y+)) \textasciicircum{}\textasciicircum{}使用生产项。对于墙附近的正确处理，使用以下阻尼函数：/„ = l-exp(-0.0115y+) /i = l / 2 = l-0.22exp\textasciicircum{}- 3(;其中湍流 雷诺数, Re t = \textasciicircum{}--和y + - pun\textasciicircum{}n是两个无维变量，x n是距离最近的墙的距离。 然后涡黏性是：fJ-t = C\textasciicircum{}fuP\textasciicircum{}/e，Chien模型中使用的典型常数是：CV = 0.09，C £l = 1.35，C £2 = 1.8，Pr k = 1.0，Pr e = 1.3 k和e墙上的边界条件是：k = 0，£ = 0对于本文中执行的计算，k和e方程通过显式Runge Kutta积分求解。

302 WEBER, GACHERIEU, RIZZI ET AL 17.3 数值模型NSMB使用以细胞为中心的有限体积（FV）近似值作为空间离散化，部分原因是在冲击等不连续性中守恒的重要性。 FV方法还有可能处理网格中复杂几何形状和奇点周围的拉伸网格。 细胞面的通量向量由非黏性部分和黏性部分组成，其中黏性部分是使用偏移控制体上的梯度定理计算的。 NSMB的用户可以选择带有二阶和四阶人工耗散的詹姆森中心差分，或者空间离散化的基于Roe的迎风差分。 通常的显式多级Runge-Kutta 时间积分一直是时差的标准方法，加上詹姆森隐性残差平滑和多网格加速到稳态解。 我们在这里没有描述它。 相反，我们确实描述了代码的新功能，其中包括上风空间差异和隐式时间差异。 在本文中，隐式方案仅应用于平均流方程，即仅实现代数湍流模型。 17.3.1 空间离散化 17。 S.1.1 詹姆森的中心方案 二阶中央方案使用二阶和四阶差分的组合，用詹姆森等人的人工耗散模型增强。[13] 事实证明，这种自适应耗散模型可以在没有振荡的情况下再现不连续性。 四阶人工耗散项抑制奇数/偶数振荡，这些振荡不会被方案本身阻尼。 细胞i和i + 1之间界面的数值通量由以下方式给出：Fi+1/2 = Tin ( y ' +1 2 +y ') - <W (17.26) 其中Tin是对流通量，di+1/2 1S tne耗散通量定义为：di+1/2 = e|+i /2 (\textasciicircum{}+i - U.) - e\textasciicircum{} 1/2(Ui+2 - 3U i+1 + 3U t - U t\textasciicircum{}) (17.27) 系数e\textasciicircum{} 2'和e (4\textasciicircum{}用于局部适应耗散通量，并由Jacobian 矩阵 £的光谱半径r(A)进行定向缩放！ +l/2 = fc(2)K\textasciicircum{})i+l/2\textasciicircum{}+l/2

空气动力学中的应用 303 ejj 1/a = max(0.0.fcWr(A) j+1/a -e<\} 1/2 ) (17.28) 传感器变量J\textasciicircum{}+1/2控制激波附近的二阶耗散。 它是使用压力归一化二阶差的绝对值构建的：\& = 然后将传感器视为Pi+l \textasciitilde{} 2Pi +Pi-1 Pi+i +2pi +pi-i (17.29) ui+1/2 = max (\textasciicircum{}, m +1) (17.30) (17.8.1.2 Roe的上风方案 这个方案是Roe方案的二阶总变减小（TVD）版本，应用守恒定律的单调上风方案（MUSCL）推[24]。 单元格侧i + 1/2的数值通量为：-F.+ 1/2 = \textbackslash\{\} yFin\{U i+1i2) +Fin\{U i+1i2) J - \textbackslash\{\} p(tf\&i/2.tf£i/ a)| (\textasciicircum{}1/2-\textasciicircum{}1/2) (17.31)，其中U-,J 2和U\textasciicircum{}\textbackslash\{\}\_ 1,2是保守变量的单元格侧外推值。 上标L和R分别指相应接口的左侧和右侧。 矩阵A（U L，U R）是Roe矩阵，基于Roe的近似Riemann 求解器。 单元格接口的左右状态定义为[4]：/l + \$ \textasciitilde{} 1 -\$ \textasciitilde{} Ut+i/2 = \textasciicircum{},+ -A i+1/2 + -A,., /I + \$ \textasciitilde{} 1 - <3> \textasciitilde{} \textbackslash\{\} U,R-1/2 = Cr,-\textasciicircum{}-J-A i\_1/a + - r -A i+Va J (17.32) 定义AUi +1/2 = fi+i - Ui，限制斜率使用minmod函数计算：A i+1 / 2 =minmod(A£/ I-+i/2,wAE/j\_ 1/2) (17.33) A,\_ 1/2 =minmod(At/i\_ 1/2,o;A[/ J+1/ 2) (17.34)，其中u是以下范围内的压缩参数：3- \$ 1<\textasciicircum{}< =• (17.35) 1 \_ \$

304 WEBER, GACHERIEU, RIZZI ET AL 方程中的精度参数\$。 （17.32）定义了一系列高精度TVD上风方案。 使用\$ = - 1产生二阶全迎风方案，\$ = 1/3产生基于斯卡尔对流方程的三阶方案的方案。 作为minmod功能的替代方案，Van Leer或Superbee限制器可用于全上风方案。 在实践中，限制器被指定在特征变量上，而不是保守变量的斜率。 其优点是更好地考虑了流动的传播特性，并且可以将不同的限制器用于不同的特征场。 17.3.2 收敛加速 17.3.2.1 多网格 多重网格法是将收敛加速到稳态的强大方法。 通过在一些较粗的网格上迭代后结合解决方案，通常可以大幅节省计算时间。 这个想法是由Ni[17]和Jameson[14]为可压缩流动计算开创的。 在NSMB中实现的多网格版本是Brandt的FAS算法[11]，它可以以递归方式编写。 多网格可以与显式Runge-Kutta方案或隐式LU-SGS方案相结合。 对于湍流计算，涡黏性仅在精细网格上计算，其在较粗网格上的值是通过限制获得的。 17.3.2.2 D-ADI 使用常数系数的隐式残差平滑对欧拉计算非常有效，但其对纳维尔·斯托克斯模拟的有效性较低。 对于高雷诺数流量，有必要解决固体边界附近的薄剪切层，这需要具有非常大的细胞长宽比的网格。 众所周知，显式方案很难在这样的网格上收敛。 为了克服这些问题，Caughey [3]开发了一种对角线交替方向隐式算法。 对于曲线坐标系，Navier Stokes方程写成r\textbackslash\{\} r\textbackslash\{\} r\textbackslash\{\} r\textbackslash\{\} dlu + di (jr,n \textasciitilde{} Tv) + d\textasciitilde{}\textasciicircum{} Gin \textasciitilde{} gv) + ac (<Hin \_ w ") = ° ( 17-36) 控制体上的积分，离散化使用隐式方案生成：\{I + yM An \textasciitilde{} L") + Sv(B" \textasciitilde{} Mn) + h\{C n - N n)]\}AUn = Rn 17.37

空气动力学305中的应用，其中Ajjn \_ jjn+l \_ Vn和（显式）残差等于：JT = -\textasciitilde{}\textbackslash\{\}h\{FL \textasciitilde{} 3\textbackslash\{\}) + *,(\$?„ - Gnv) + hWn - «?)] （17-38）An，Bn，Cn（resp. L n,Mn,Nn）是非黏性（异常黏性）雅各布矩阵，5\textasciicircum{},8\textasciicircum{}和S\textasciicircum{}是曲面法线。 为了使隐式方法成为有效的平滑算法，人工耗散的第四差分包含在隐式算子中，这导致每个一维因子的五对角线系统的反转。 为了避免求解块五对角线系统的高成本，使用相似性变换在每个点对角线方程，见Chaussee和Pulliam [5]。 这具有解耦合方程的效果，并要求为三维问题T\textasciicircum{}[I+\textasciicircum{}S\textasciicircum{}]Tf1Tri[I+\textasciicircum{}5vAr,\}T-1T<[I+\textasciicircum{}A<]T<-1AUn = R n (17.39)解方程组的五个标量五对角线系统，其中A\textasciicircum{}[I+\textasciicircum{}S\textasciicircum{}]Tf，A<]Tf是包含非黏性雅各布矩阵A、B、C、A = iAiT\textasciitilde{}\textbackslash\{\} B = T„AVT-1、C = T CA(T\textasciicircum{} (17.40)和T(,T V,T(是包含相关右特征向量的矩阵的矩阵。 由此产生的方法具有良好的高波数阻尼，因此它是与多重网格法结合使用的良好算法。 它也具有计算效率，因为只涉及标量系统。 确定局部相似性变换以及执行残差以及中间和最终校正的矩阵乘法所需的额外计算工作只占求解块系统所需的工作的一小部分。 上述算法已被扩展，通过将黏性雅各布的光谱半径添加到Tysinger和Caughey [3]之后的无黏性隐式因子中来求解Navier-Stokes 方程。 例如，在£方向上，A（矩阵被替换为：A（= As + p（Lt）（17.41），其中黏性Jacobian 矩阵的光谱半径由以下方式给出：\textasciicircum{}）= \textasciicircum{}-（，+ \textasciicircum{}LiJl！ （17.42）

306 WEBER, GACHERIEU, RIZZI ET AL 17.3.3 隐性时间 离散化 17.3.3.1 标量LU-SGS方案 以下小节介绍了最初由Yoon和Jameson引入的LU-SGS方案[25]。 之所以选择它，是因为它的数值努力低、内存要求低和收敛速度合理。 此外，事实证明，该计划非常稳健。 该方案基于下上因式分解和对称高斯-塞德尔松弛。 控制方程在空间和时间上分别离散化。 这确保了稳态解决方案将独立于时间离散化程序，因此独立于时间步骤。 将时间水平n的残差线性化导致以下方程：与/是恒等矩阵。 方程（17.43）表示一个大型稀疏线性系统，必须在每个时间步骤中求解。 术语dR/dU象征性地代表通量线性化产生的雅各布矩阵。 可以直接方法来求解线性系统，该系统需要大稀疏块带状矩阵的反转。 然而，使用这种方法，数值成本和存储要求是禁止的。 相反，对线性系统本身应用迭代方法和/或进行近似。 起点是等式的对角优势形式。 （17.43）为了满足松弛方法的稳定性要求。 Yoon和Jameson使用以下形式：（A7 + rA）AU'+l \textasciicircum{}-\textasciicircum{} J )'+i/2 AC/ '+i-\textasciicircum{} (A+r AI)i\_1/2 At/,-! = -R(U n)i (17.44)，AAU = AF，r\&是Jacobian 矩阵 A的光谱半径。 LU-SGS方法通过将线性系统分解获得如下：\{E + D) D- 1 (F + D) AU = -R n (17.45) 其中E仅包含下三角部分，F包含上三角部分，D包含隐式算子的主对角线。 线性系统现在通过在平面上的网格中向前和向后扫荡而反转，i + j + k = const: (E + D)AU* = -R n \{F + D) AU = D AU* (17.46)

空气动力学中的应用 307 通过在斜平面上扫荡，对角线外术语EAU*和FAU分别变得众所周知，并被添加到右侧。 因此，只有块对角矩阵必须反转。 如果引入使用光谱半径的近似值，隐式算子甚至可以简化为标量对角矩阵。 对于黏性计算，黏性雅各布的光谱半径被添加到无黏性光谱半径中。 17.8.3.2 全矩阵LU-SGS方案 发现，使用高单元宽高比的网格模拟黏性流量时，标量LU-SGS方案的收敛速率会显著降低。 其原因是由于使用光谱半径，系统过度稳定。 收敛速率可以通过使用基于Roe的上风方案的全矩阵LU-SGS方案来改进[23]。 假设Roe矩阵是恒定的，基于全上风方案的数值通量可以线性化，从而导致以下方案：At' ' 2 I + - a (\textbackslash\{\}A 1+1/2\textbackslash\{\} + |A,\_ 1/2 |) AUi + \textbackslash\{\} a (A(U\& l/a) - |I, +1/2 |) AU i+1 - \textbackslash\{\} a (A(U, L\_1/2) + |I,\_ 1/2 |) AU \{-i = -R(U n)i (17.47) 因子a是忽略限制器时向上风外推到接口的结果，并设置为1.5。 为了增强收敛和方案的稳定性，包括了黏性雅各布矩阵。 17.3.4 时间精确方案 可以使用明确的Runge Kutta方案进行时间精确计算，其中时间步骤等于所有网格单元格的最小时间步骤。 对于具有高度拉伸网格的纳维尔·斯托克斯计算，最大允许的时间步进变得非常小，因此，詹姆森[15]提出的双时间推进技术被实施。 双重时间步进背后的想法是使用完全上风方案的时间精确时间步进的外部时间步进循环，以及一个具有虚构时间步进的内部时间推进循环。 局部时间推进、多网格和其他收敛加速技术可用于在虚构时间内将内部时间推进回路收敛到“稳态”。 通过为外部时间步骤使用完全隐式方案，可以进行大型时间步骤。 双时间步进方法的问题在于找到外部时间步进的最佳值，因为外部时间步进的值太大需要大量的内部时间步进才能收敛到稳态，而值太小

308 WEBER、GACHERIEU、RIZZI等人的外部时间步骤，与Runge Kutta方案相比，双时间推进方法的增益很小。 17.4 程序结构和多块实现NSMB是在名为MEM-COM [16]的数据库系统上开发的。 MEM-COM是一个面向对象的数据管理系统，用于内存和内存到磁盘数据处理。 使用数据库系统的主要优势是，对于大规模的多块流模拟（即100多个块和超过100万个网格点），访问独立块的速度非常快，并且几乎独立于块数。 从用户的角度来看，数据库文件显示为单个UNIX文件，因此与模拟相关的所有信息都可以轻松地保存在存档系统上，而不会丢失信息。 MEM-COM库包括一个动态内存管理器（DMM），它提供了在运行时分配NSMB中阵列的必要存储的可能性[22]。 在分布式内存计算机上运行时，DMM 会在计算机的每个节点上分配内存。 NSMB中使用的显式Runge-Kutta方案的单块或多块计算几乎没有区别。 主要在边界条件的合并层面上出现了差异。 LU-SGS 隐式方法使用块之间的显式耦合实现，这意味着隐式时间推进过程是逐块独立执行的。 显式耦合不会引入任何额外的计算成本，但隐式方案的全局性质会降低。 当前版本的NSMB假设网格线跨块接口是连续的，这意味着接口算法没有必要。 计算块两侧的虚构幽灵细胞用于将信息从边界条件（包括块连接边界条件）传输到用于更新内部点的算法，有关更多详细信息，请参阅[22]。 17.5 高性能计算策略 现代矢量和并行平台，例如 如果代码适应架构的特殊功能，NEC SX-4和富士通VX、IBM-SP2、Cray J932和T3D现在为用户提供了高性能。 我们描述了我们为通过NSMB实现这一目标而采取的步骤。

空气动力学应用309 17.5.1 RISC/矢量优化NSMB专为在矢量计算机上运行而设计，编写了最耗时的例程，以在块中的所有网格点上具有矢量化的do循环。 这在NEC SX4上产生了超过1100 MFlops的出色性能，超过单个SX4处理器峰值性能的50\%。 在RISC架构上，NSMB的初始版本表现相当差，达到峰值性能的10\%到15\%。 性能不佳的原因是内存和缓存争用问题，因为缓存中加载的数据无法重复使用。 当问题大小减少时（即当使用较小的块时），性能会提高，因为所有数据都可以放入缓存中。 此外，当在所有网格点上使用do-loops时，在幽灵单元格中执行不必要的工作。 在矢量计算机上，这不是一个问题，在所有网格点上使用矢量循环的增益在很大程度上补偿了这种开销，但在RISC架构上并非如此。 研究了几种技术来提高RISC架构上最耗时的例程的性能[1]。 一些工作阵列的互换数组索引通过对代码进行了小幅修改，就取得了很好的改进。 17.5.2 并行实现 基线并行NSMB版本的设计在可移植性、代码更改最小化和未来维护的便利性方面受到重要限制。 在开发并行NSMB时做出了两种设计选择，首先在并行NSMB执行之前执行域分区，其次NSMB中的并行实现基于主/从范式。 后一种选择允许将NSMB轻松移植到不同的并行平台。 17.5.2.1 域分区工具 MB-Split 区域分解中最重要的方面之一是在并行机器上分配块，以便在不同处理器之间实现良好的负载平衡。 分块簿记的复杂性需要自动的负载均衡和分块工具。 这样的区域分解工具MB-Split[26]是在KTH在Parallel Aero项目中开发的。 MB-Split是一个用C++编写的程序，它可以读取具有任意数量块的MEM-COM数据库，并生成具有新数量块的新MEM-COM数据库，以便在并行计算机或计算机集群上获得平衡计算。 MB-Split自动纠正现有块上的所有边界条件，并生成边界条件信息

310 WEBER、GACHERIEU、RIZZI等新创建的区块。 使用递归边缘分割或贪婪负载均衡算法拆分块。 还有一个Greedy 负载均衡算法的修改版本，称为GreedyXtra 负载均衡，它考虑了在ghostcells中完成的工作。 也可以根据不同块的实际时间进行负载平衡。 不同块的定时由NSMB输出。 后一种算法不拆分任何块，它只为给定数量的处理器重新分配它们。 对分配给节点的块数量没有限制，每个节点也不需要有相同数量的块，这意味着可以对网格进行负载平衡，其原始块数量比可用处理器多，反之也相反。 除了拆分块外，MB-Split还包含将块合并回原始块结构的选项。 通常在工业CFD模拟中，使用较粗的网格解决方案来初始化精细的网格计算。 最有可能的是，粗网格计算将在不同数量的处理器上进行，因此块分解与精细网格计算不同。 决定使用原始的块拓扑来从粗网格到细网格的溶液插值。 这意味着MB-Split不仅可以拆分网格，还可以拆分保存在MEM-COM数据库中的解决方案。 17.5.2.2 并行NSMB的主从实现 主从实现的主要特点是：1. NSMB的并行版本中使用的消息传递原语使用最便携的通信原语，如发送、阻止和非阻止接收。 对消息传递库的调用是通过透明的通信层进行的，IPM[9]是一组M4宏，在编译时生成PVM或PARMACS指令。 2. 在并行NSMB中，主站执行对MEM-COM数据库的所有访问。 在计算阶段（即时间推进循环），除了打印收敛历史记录外，从站和主机之间没有通信。 3. 从属节点以二进制树结构组织。 根节点的任务是分配和控制流程模拟（检查收敛，评估全局时间步骤......）。 在计算过程中收集信息也是通过树结构完成的。 4. 在Runge-Kutta阶段开始时，每个处理器都会将块连接边界条件的数据发送到持有相邻块的处理器。 同一处理器上的块在不使用网络的情况下交换数据。 一旦收到来自相邻块的所有数据，块中的下一个Runge-Kutta阶段的计算就会立即开始。

空气动力学中的应用 311 同样的程序也适用于LU-SGS方案。 在解决线性系统之前执行所有数据交换。 5. 每x个时间步骤（通常x = 10）一次，湍流模型所需的墙和唤醒信息在每个块上计算，被组装并重新分配到每个处理器。 6. 主从可执行文件是使用CPP预处理器指令在编译时生成的。 NSMB的串行版本以类似的方式获得，因此只有一个源代码，可以从中生成三个可执行文件（主、从和串行）。 17.6 计算结果 计算了三个不同的测试案例：翼型，称为A-翼型，65度扫描三角翼，以及全配置运输机，称为AS28G。 17.6.1 A-翼型 A-翼型是一个二维亚音速测试案例，通常用于测试湍流型号。 它有据可查，并且有很多来自风洞实验的数据。 流量是低马赫数流量，M = 0.15和Re = 2.1 x 10 6的雷诺数。 攻角是a = 7.2°。 17.6.1.1 Baldwin-Lomax模型 在第一次计算中，带有Roe差异的LU-SGS方案与Baldwin-Lomax 湍流模型一起使用，它非常适合这种没有或只发生轻微分离的流动类型。 这也是通过将收敛与显式Runge Kutta方案进行比较来证明LU-SGS方案性能的第一个测试案例。 流动条件假设一个固定的过渡点。 使用的网格是一个单块C网格，具有256 x 64 = 16384个单元格，用于最好的网格级别。 靠近壁的非常小的细胞和高的细胞长宽比会导致小的时间步，因此在使用显式Runge-Kutta 时间推进方案时，收敛速度很慢。 使用了五级Runge-Kutta方案，CFL数为1.0，将误差标准减少了30年。 然后将其提高到2.0，以加速收敛。 使用隐式方案时，CFL数最初设置为10。并在时间推进期间增加，因此在大约170次迭代后达到10 6。 没有发生稳定性问题，并且证明该算法非常稳健。 在图中。 1，显示了硅图形功率挑战中迭代和CPU时间方面的收敛历史记录。 显式方案收敛速度极慢，需要

312 WEBER、GACHERIEU、RIZZI等人大约50000次迭代，剩余减少50年。 完整矩阵（三阶上风）和标量隐式（中央）版本分别需要550次和7000次迭代。 CPU时间可以分别减少32倍和8倍。 在前四十年里，标量隐式版本表现良好。 然后它继续缓慢收敛并变平。 这种行为是由于高细胞长宽比造成的。 使用完整的矩阵版本，这种收敛退化不会发生。 边界层在翼型上的两个位置的速度曲线如图所示。 2. 中央和三阶上风方案都与x/c = 0.7位置的实验数据非常一致。 X/c = 0.96时速度曲线的差异是由于Baldwin-Lomax 湍流模型未能以适当的方式预测再循环。 沿弦的压力系数可以在图中看到。 3. 三阶上风方案和中央方案给出的结果几乎相同，并且都与测量的压力很一致。 17.6.1.2 k-e Chien模型 使用明确的五级Runge-Kutta方案将k和e方程与平均流量方程整合在一起。 使用的网格是一个单块C网格，具有256 * 64 = 16384个单元格，已被证明足够精细，可以给出与网格无关的结果。 墙距y +的平均值约为0.4。 使用了五阶段Runge-Kutta方案，CFL等于0.5。 显式方案收敛得非常缓慢，A-翼型的剩余图1 收敛历史需要大约70000次迭代。

空气动力学中的应用313

图2 A-翼型上两个位置的速度剖面。减少了50年。 表1比较了与NSMB中包含的各种湍流模型和实验测量计算的力系数CL和Cp。 零方程模型高估了升力系数，而斯帕尔特·阿尔马拉斯模型低估了实验值。 阻力系数被大多数模型低估了，而k-e对阻力系数和升力系数都显示出很好的一致性。 使用k-e模型获得的边界层的速度曲线与前面描述的Baldwin Lomax模型在同一位置绘制，如图。 2. 数值结果和实验结果的比较表明，在流动连接的第一个位置，墙附近有很好的一致性（见图。 4）。 在第二个位置，实验结果表明了分离的流动，而fc-e以及Baldwin Lomax的计算表明了附着的流动。

314 WEBER, GACHERIEU, RIZZI ET AL 图3 压力系数 沿着A-翼型 17.6.2 三角翼的弦，虽然三角翼的几何形状很简单，但由于多个涡流的相互作用，激波与涡流的可能相互作用，最常见的流动是湍流，在高攻角下模拟机翼上的涡流是一个具有挑战性的问题。 这里展示了具有圆形前缘的65°扫三角的流量结果。 基于根翼弦，Moo = 0.85，a = 12°和Re/m = 9 x 10 6的湍流计算结果呈现。 使用此测试案例有3个原因：•评估空间离散格式对结果的影响cL cD Baldwin-Lomax模型1.067 0.0152 Granville模型1.046 0.0148 Spalart-Allmaras模型1.002 0.0143 Chien k-e模型1.032 0.0155实验结果（F2风洞ONERA）1.033 0.0155表1用不同湍流模型计算的力系数与在A-翼型上测量的力系数的比较

空气动力学中的应用 315 图4 A-翼型上两个位置的速度曲线 • 评估所选湍流模型对结果的影响 • 研究不同的收敛加速方法 除了Re/m = 9 x 10 6的湍流计算外，还对Re/m = 72 x 10 6进行了一次计算，并进行了一次层流计算。 对于所有计算，使用了总共675'840个网格点的8块网格。 墙附近的最大y +值在0.2到0.5之间，墙和自由流边界之间有80个点。 使用LU-SGS隐式方案和/或多网格收敛加速的计算直接在最好的网格上开始。 对于所有其他计算，使用具有84480点的粗网格来初始化精细网格计算。 17.6.2.1 空间离散格式的影响 计算使用Baldwin Lomax 湍流模型进行，具有以下空间离散化方案：尽管使用不同参数进行了多次测试，但似乎不可能使用矩阵耗散的中心方案获得收敛解。 图5显示了这六个空间离散化方案的X/C = 0.4时的C p。 使用矩阵耗散的结果不是收敛的，并且涉及中间解决方案。 如图所示。 5，空间离散格式对压力吸力峰值有很大影响。 在NSMB中实现的标准Jameson 人工耗散是各向异性的，因此耗散性比Martinelli 人工耗散低。 三阶上风方案耗散性最小，并预测涡流的最小宽度。 使用中央方案+矩阵耗散的两种计算都接近三阶上风方案的结果，但如前所述，这些计算没有收敛。 在三角翼的迎风一侧，所有的计算

316 WEBER、GACHERIEU、RIZZI等图中的图例。 5 空间离散格式 cp-ntf-re9e6-std中心方案+标准詹姆森耗散cp-ntf-re9e6-mar中心方案+马蒂内利消散cp-ntf-re9e6-mxu中心方案的缩放+基于原始变量的矩阵耗散cp-ntf-re9e6-mxw中心方案+基于保守变量的矩阵耗散cp-ntf-re9e6-upw上风方案cp-ntf-re9e6-3upw三阶上风方案表2使用不同空间离散化方案在65英寸三角上进行Baldwin Lomax计算的计算的计算摘要。完全在一起，除了使用非收敛计算 中央方案+矩阵耗散。 图6显示了使用Baldwin Lomax模型和Spalart-Allmaras和Edwards-McRae的单方程模型，以及层计算，在X/C = 0.4时计算出的C p。 所有计算都是使用中央方案+詹姆森耗散进行的。 如图所示，层压计算预测了二次分离，这在湍流计算中不存在。 涡流的宽度是Spalart-Allmaras模型中最小的。 与Baldwin Lomax模型相比，一个方程湍流模型预测了较低的吸力峰值。 雷诺数的效果如图所示。 7，它显示了两个雷诺数的计算结果，以及Re/m = 9x 10 6的层次计算。 这三种计算是使用三阶上风方案进行的。 如前所述，层级计算可见二次分离。 Re/m = 72 x 10 6的计算产生了最高的吸力峰值，这里也可以看到一个小的次级分离。 雷诺数的效果是吸力峰值增加，涡核向对称平面小幅移动。 图8显示了Re/m = 9 x 10 6时湍流计算的机翼上方的粒子痕迹，涡流清晰可见。 使用Baldwin Lomax 湍流模型和不同的收敛加速方法进行了几次计算。 L2残量作为时间步骤的函数如图所示。 9，计算总结在表3中。 所有计算都是在曼诺瑞士国家计算中心的NEC SX4上进行的，使用本地时间推进。 可以看出

空气动力学中的应用 317 图5 C p at X/C = 0.4使用不同的空间离散化方案图。 9，使用两种上风方案与全矩阵LU-SGS隐式方案相结合的计算在时间步数方面收敛速度最快。 对于中央方案，使用多网格和标量LU-SGS隐式方案的计算产生了最快的结果（同样在时间步数方面），其次是使用标量LU-SGS隐式方案的计算。 表3总结了所做的不同计算，包括有关所需计算资源的信息，以及在NEC SX4上获得的性能。 在计算每个时间步骤的成本时，使用隐式标量LU-SGS方案的中央方案+Martinelli 人工耗散的计算似乎最便宜，而中央方案+Martinelli和带有多网格的显式Runge Kutta是最昂贵的。 使用每个时间步骤的成本，为不同的方案确定了将残留物减少到10-2的成本。 这表明，带有隐式标量LU-SGS方案的中央方案+Martinelli 人工耗散需要约2400 CPUsec，而带有显式Runge Kutta方案的中央方案+Martinelli 人工耗散需要约5200 CPUsec。 使用隐式矩阵LU-SGS方案的上风和三阶上风方案分别需要5100和5760 CPU秒。 所有其他时间积分计划都贵得多。 值得一提的是，上风方案比人工耗散的中央方案更贵，这部分解释了差异

318 WEBER, GACHERIEU, RIZZI ET AL 图6 计算结果 65° deltawing使用不同的湍流模型。在中央和上风方案之间。 从表3中可以看出，使用矩阵LU-SGS隐式方案的上风方案的计算比使用显式Runge Kutta 时间推进的中央方案或使用标量LU-SGS方案的中央方案多出约45\%的内存。 17.6.3 AS28G全飞机 AS28G测试案例代表了全飞机配置，其中包括黏性效应轴很重要的发动机/机身集成，特别是在发生流动分离的塔架上。 最好的网格包含62个块，有350万个网格点。 图10显示了此配置的表面网格。 这里显示的流动条件代表了塔塔优化的巡航条件：Re/m = 11.16 x 10 6，a = 2.2°和M = 0.8。 17.6. S.1 流分析和与实验的比较在这里使用了中心空间差异，二阶和四阶人工耗散的系数分别设置为0.8和0.015。 通过将壁法人工耗散通量乘以局部马赫函数，人工耗散在边界层中进行了阻尼

319 空气动力学中的应用 图7 计算结果是65° deltawing在两个雷诺数。数字在壁上趋于0，在边界层的边缘趋于1。 如[7]，[21]所示，这种阻尼对边界层和激波位置的正确预测有很大影响。 使用标量LU-SGS方案时，计算是从自由流初始解和CFL数1开始的。 经过大约700次迭代，CFL数字已线性增加到10 9。 经过5000次迭代，使用Cray T3D上的109块网格、32个IBM SP2处理器上的75块网格和富士通VX上的原始62块网格，残差（p）的Li范数减少了四个数量级。 这相当于T3D的总经过时间约为30小时，SP2的总经过时间约为55小时。 即使从自由流初始解决方案开始，并为人工耗散使用非常小的值，也没有遇到稳定性问题。 Lift和阻力系数历史记录如图所示。 11，并在大约4000次迭代后显示收敛值。 对于使用128个处理器在Cray T3D上进行计算，网格已分为276个块。 由此产生的块数是，每个处理器上至少有两个块。 为了获得合理的负载平衡，测量了实际的CPU定时，并使用这些定时重新分配块。 此过程要求每个处理器至少两个块。 然而，当使用276个区块时，具有区块之间显式耦合的标准LU-SGS方案导致了分歧。 然后对方案进行了修改，因此在更新块时执行了几次扫描

320 WEBER、GACHERIEU、RIZZI等人图8颗粒痕迹描绘了65°三角翼的前缘涡流。每次扫描之间的连接性边界条件。 在迭代方法的收敛，将获得与1块网格相对应的更新AU。 这种方法需要更多的沟通，但随着几次扫描的进行，计算工作也会增加，因此通信时间与经过时间的比率只是缓慢增加[23]。 以上述方式进行三次扫描，经过大约2000次迭代后可以获得收敛解决方案。 使用相同算法的276块和109块计算的收敛速率如图所示。 12. 在前三十年的剩余减少中，这两种计算收敛于一致。 那么276块计算的收敛似乎有些退化。 然而，经过2000次迭代，残留可以减少40年。 在查看压力系数历史记录时，使用276块时，错误似乎减慢一些。 不可能使用显式的Runge-Kutta方案获得收敛解决方案，但IBM SP2上使用了Diagonal-ADI方案。 然而，ADI方案对湍流模型中的参数值非常敏感，大于0.1-0.2的CFL数字会出现发散。 ADI和LU-SGS方案的收敛率如图所示。 13. 使用ADI方案的计算在大约6万次迭代后停止。 到那时，残差几乎没有减少二十年，升力和阻力系数，但Cx仍然比正确值低20\%左右。 表4给出了基于风轴的计算升力（CL）和阻力（Co）系数，并与不同计算机和块分割的风洞测量进行比较。 所有计算都是使用Baldwin-Lomax 湍流模型和LU-SGS方案进行的。 二阶人工耗散的系数降低到0.4

空气动力学中的应用 321 空间 离散化 时间积分 步 CPU（秒）内存（Mb） MFlops 中央 + 标准显式 Runge Kutta 9000 34420 334 810 中央 + 矩阵显式 Runge Kutta 6000 30864 339 941 中央 + Martinelli 显式 Runge Kutta 6000 23091 334 810 中央 + Martinelli 显式 Runge Kutta + 多网格 3000 47203 448 824 中央 + Martinelli 隐式标量 LU-SGS + 多网格 2000 25406 476 724 中央 + Martinelli 隐式标量 LU-SGS 5000 15056 339 714 上风隐式矩阵 LU-SGS 2000 23970 481 520 三阶上风隐式矩阵 LU-SGS 2000 22163 481 表 3 使用不同的时间推进方法进行Baldwin Lomax计算的结果摘要。在富士通上，在T3D上为0.25，而最初为0.8。 有人认为，降低这个系数将使机翼顶部的冲击力进一步向前移动，从而更好地预测升力和阻力。 然而，这种变化对激波位置没有影响，对CL和Cry也没有影响--正如预期的那样，块的分裂对最终结果的影响很小。 图16至14绘制了压力系数，-C p，在机翼的不同位置。 使用LU-SGS和Baldwin-Lomax的NSMB结果与风洞测量进行了比较。 从y/(2b) = rj = 0.29的站2开始，到y/\{2b) = rj = 0.57的站4结束。 半跨度为b/2 = 3.704米。 机翼下侧的C p分布对所有站都预测得很好，但激波位置略落后于正确位置，特别是对于站点3（图。 15）靠近塔（在rj = 0.35时）。 靠近翼尖，在第4站（图。 14），对激波位置的预测要好得多。 因此，塔似乎使冲击-边界层相互作用进一步复杂化，以至于简单的代数湍流模型失去了准确性，无法精确预测激波位置。 鲍德温-

322 WEBER、GACHERIEU、RIZZI等图9 收敛历史不同的时间推进方法，结果使用Baldwin Lomax 湍流模型进行65°三角化。 众所周知，Lomax 湍流模型很难正确计算冲击边界层相互作用，在这种情况下，通常预测其正确位置的下游激波位置，这与此处显示的结果相对应。 更好地预测冲击-边界层相互作用将需要更先进的湍流模型，因此，由于每个时间步的浮动操作更多，收敛速度较慢，因此需要更长的流逝时间。 塔内侧计算出的皮肤摩擦线与图中的油流图片进行了比较。 17. 在图片的上部可以看到机翼，在右下角可以看到发动机，在塔架之间。 数值模拟正确再现了定性流量行为。 17.6.8.2 HP C性能使用IBM SP2上的32和64处理器以及Cray T3D上的64处理器的并行计算的定时与NEC SX4和Fujitsu VX两台矢量计算机上使用单个处理器的定时进行了比较。 表5显示了不同计算机的每次迭代的定时。 富士通

空气动力学中的应用 323 图10 AS28G飞机配置的表面网格。正时适用于原始62块网，而SP2的定时适用于使用32个处理器的75块网格和使用64个处理器的106块网格。 Cray T3D和NEC SX4正时适用于109块网。 32个SP2处理器的定时适用于ADI和LU-SGS方案，而64个处理器的定时仅适用于ADI方案。 第二列给出了没有通信的计算时间。 时间步进程序的计时和总经过时间之间的整个差异不能全部视为沟通，而是沟通的很大一部分。 所有计算都是使用Baldwin-Lomax 湍流模型进行的，所有并行计算都是使用PVM消息传递库进行的。 NEC SX4计算的性能率约为500 MFlops，可以与其他计时进行比较。 LU-SGS和ADI方案每次迭代所需的时间大约相同，如表5中32个处理器SP2案例所示。 SP2的性能率非常低，这可以解释为SP2上公共域PVM的低通信带宽（小于10 MB/s）。 由于要发送的数据量更大，这对于精细的AS28G网眼比A-翼型更明显。 在富士通上使用NSMB的性能率，对于矢量长度有利的外壳，可以高达900 MFlops，NEC SX4略高，“薄”SP2处理器约为40 MFlops。 这将需要大约20到25个SP2处理器才能获得与富士通和NEC相同的性能，而不考虑通信，但负载均衡和高效并行化的问题使它变得困难

324 WEBER，GACHERIEU，RIZZI ET AL CL cD IBM SP2，75块 0.5234 0.0307 Cray T3D，109块 0.5231 0.0307 Cray T3D，109块，减少Artdisss2 0.5300 0.0305 Cray T3D，109块，减少Artdiss2，3扫荡0.5282 0.0302 Cray T3D，276块，减少Artdisss2，3扫0.5289 0.0302 富士VX，62块 0.5247 0.0306 NEC SX4，62块 0.5233 0.0303 富士通 VX，减少Artdiss2 0.5297 0.0303 Windtunnel 0.5145 0.0249 表4 AS28G的升力（CL）和阻力（CD）系数，使用NSMB中的LU-SGs和Baldwin-Lomax。 计算机Colfim\textasciicircum{}f° n，总通过时间富士通（LU-SGS）0 58 NEC SX4（LU-SGS）0 40 32 P。 SP2（ADI）44 1100 32 P。 SP2（LU-SGS）51 57 64 P。 SP2（ADI）48 860 64 P。 SP2（LU-SGS）63 42 64 P。 T3D（LU-SGS）33 30.5 表5 不同计算机上使用ADI和标量LU-SGS方案的精细AS28G网格的定时。

空气动力学中的应用 图11 升力（CL）和阻力（Co）系数作为迭代的函数，用于AS28G网格，使用LU-SGS和Baldwin-Lomax。以接近理论性能，这就是为什么与现代并行wmputer相比，像富士通和NEC这样的矢量计算机实现了良好的价格/性能。 表6给出了Cray T3D上LU-SGS方案的时间，执行3次扫描，每次扫描之间更新块边界条件。 乍一看，与128和64处理器相比，大约2.5的加速令人吃惊。 这是因为64个处理器的负载平衡不佳。 由于使用64个处理器时的内存限制，网格无法以最佳方式分割。 使用128个处理器，可以在10.5小时后获得收敛溶液。 对于空气动力学设计工作来说，这种运行时间可以被认为是可以接受的，因为这种流动模拟必须在夜间运行中完成。 17.7 总结和结论 我们介绍了我们最近在Navier-Stokes模拟的三个方面的工作，即单方程Spalart-Allmaras 湍流模型、矩阵LU-SGS隐式方案及其在IBM SP2、Cray T3D、NEC SX4和Fujitsu VX上的实现。 我们已经调查了代码的这些进步在三个不同的测试案例中的表现。 对于飞机配置的测试案例，我们已经表明，Baldwin-

326 WEBER、GACHERIEU、RIZZI等人

在更新块连接边界条件时执行3次扫描。 计算机Coi\%\$£mtion。 总通过时间 64 Procs 35 24.7 128 Procs 26 10.5 表 6 使用LU-SGS方案进行3次扫描和块连接边界条件更新，在Cray T3D上精细AS28G网格的定时。 Lomax 湍流模型是350万网格点的Navier-Stokes模拟，以2000个时间步进收敛，可以在带有128个处理器的Cray T3D上运行过夜。 计算结果与实验一致，在升力约2\%和阻力约4\%以内。 标准公共域PVM被用作在IBM SP2上执行计算的消息传递库，对于32个处理器，其通信时间与计算时间一样大，因为SP2上公共域PVM的带宽小于10 MB/s。 由于网络更快，在Cray T3D上使用64个处理器进行计算的通信时间要低得多，但负载平衡不佳，尽管网络速度快，但并行性能较差。 进一步对域进行分区可能是一种解决方案，但由于处理器的内存限制，这种情况不适合计算机。 为了缓解这一限制，尝试了128个处理器的分区，但导致图12使用109块和276块比较收敛历史记录

空气动力学中的应用 图15 基于密度的L2残差，用于LU-SGS方案和对角线-AD1方案。 两者都使用Baldwin-Lomax和32 SP2 procamm。由于块之间的显式耦合，数值不稳定。 在LU-SGS方案中执行三次扫描，而不是一次，增加了耦合，被证明是稳定的，并在2000个时间步进中导致收敛。 此外，在矢量机器SX4和富士通VX上进行了运行。 单个NEC SX4处理器上精细（3.5 Mpoints）网格的计算时间大约需要50个T3D处理器和SP2上的45个处理器。 富士通VX的相同数字分别为34和32，矢量处理代码相当优化。 这使得使用现代矢量计算机在价格/性能方面成为有竞争力的选择。 隐式方案收敛迅速且稳健，运行的CFL数为10亿，但目前只有代数湍流模型包含在该方法中。 致谢 NSMB的平行\textasciitilde{}ebion的开发是在ESPRIT I11平行航空项目的框架内完成的。 瑞典参与该项目由瑞典国家工业和技术发展委员会NUTEK资助。 瑞士参与该项目由瑞士教育和科学部（OFES）资助。

WEBER、GACHERXEU、RIZZI等人......，...... I;'.,................... -*.,,........... "a6; 0 0.1 0.2 0.3 0.4 0.5 0.0 0.7 0.8 0.9 1 \$Jc 图14 AS28G机翼上的Cp分布，在站1，半跨度的57\%。 瑞典KTB平行计算机中心（PDC）对计算机的广泛使用表示感谢。 感谢CCERFACS的计算机支持小组对Cray T3D的帮助。 Aerospatide，Toulouse非常感谢ac1cnowledged提供AS28G网格和windtuanel数据。 参考资料1。 Amestoy P.R.和Dayde M.J.，在高性能计算机上移植工业代码，在：Ruid Dynamics中的高性能计算，编辑：P. Wcssding，Iiluwer Acadenzic Publisshcr，BRCOPTAC系列，1996年3月。 2. Baldwin B.S.和Lomax I-I.，分离湍流犁的薄层近似和代数模型，AIAA论文。 78-257，1978年。 3. Caughey D.A.，Euler 方程的对角线隐式多网格算法，AXAA杂志，卷。 26）不。 7，1988年，pp. 841-851。 4. Chakravarthy，S.R.：Navier-Stolces方程的飞行分辨率上风Wrmulations，在：VICI LS 1988-05，Con\textasciitilde{}putational Ruid U\textasciitilde{}ynamics，1988年。 8. Chaussee D.S.和Pulliam 'I'。 I.f.，使用对角线隐式算法的二维入口模拟，AIAA Joum\textasciitilde{}ul，Vo1。 19，l\%b。 1981年，\textasciitilde{}JP。 153-159。 6. Chien I<.-Y.，用低雷诺数* 湍流模型预测通道和边界层流动。 AIAA杂志，卷。 1982年1月20日，pp. 33-38。 7. Qacherieu C., Etude cl'un mod8 de 湍流 algdbrique 31): Applicatio\textasciitilde{}\textasciitilde{} au cefcul Navier-Stokes de I'dcoulement autour d'une instaflation motrice d'avion de transport, PIID thesis No. 1248）Imtitut National Polylechniy.c\textasciitilde{}e de Xoulozrsc，1996年。

空气动力学中的应用 -0pn\textasciitilde{}y\textasciitilde{}1som.wbn-a41m ASBG机翼上的图分布，在站点3,42\%半跨度。 8. 加尔布雷斯R.A. 麦当劳。 S.A. Sjolander和M.R. 头部，湍流墙壁区域的混合长度，飞行员。 Quart.，卷，27，pt。 2，1977年，pp. 97-110。 9. Giraud L.、Noyret P.、Sevault E.和Van Kemenade V.，IPM 2.3，用户指南和参考手册，CERFACS，1994年11月。 10. Granville P.S.，压力梯度中湍流边界层的Baldwin Lomax因子，AIAA杂志，卷。 25，1987年，pp. 16241627。 11. Hackbusch，W.，多网格方法和应用，Springer-Verlag，柏林，海德堡，1985年。 12. Jameson，A.和Turkel，E.，隐含方案和LU分解，数学汇编，卷。 37，不。 156，1981年，pp. 385-397。 13. Jameson，A，Schmidt，W.，Turkel，E.：使用Runge-Kutta 时间推进方案的有限体积法对Euler 方程的数值解，AIAA论文81-1259，1981年7月。 14. 詹姆森，A，Euler 方程的二维跨声速流动的SoIution，由多重网格法，Appl。 数学。 玉米。，卷。 13，1983年，pp. 327-355。 15. 詹姆森，A，使用多网格的时间依赖性计算，应用于通过机翼和假发的不稳定流动，AIAA论文91-1596，1991年6月。 16. Merazzi S.，MEM-COM集成内存和数据管理系统-MEM-COM用户手册6.0版，SMR TR-5060，SMR公司，P.O. Boa：4 1，CH-2500 Bienne，1991年3月。 17. Ni，R.-H.，解决Euler 方程的多网格方案，AIAA杂志，卷。 20，1982年，pp. 1565-1571。 18. Spalart，P.R.和Allmaras，S.R.，空气动力学流动的单方程湍流模型，AIAA论文9.2-0439，1月。 1992年。 19. 斯图本，。 K.，Mierendorff，H.，Thole，C.A.和Thomas，0.：带有真实代码的工业并行计算，并行计算杂志，将出版，1996年。 20. Thomas L.C.和Hasani S.M.F.，“流动”的补充边界层近似，流体力学杂志，卷。 111,1989，pp. 420-427。 21. Tourrette L.，DLR-F4机翼/机身配置周围的跨声速流动的湍流模型评估，AIAA论文96-2034，1996年。

空气动力学中的应用 331 图17 塔内侧的皮肤摩擦线

\clearpage
\chapter{航空航天应用及其他领域的不可压缩的Navier-Stokes计算}
航空航天应用和超越D的不可压缩Navier-Stokes计算。 夸克1，C。 亲属2，J。 Dacles-Mariani3，S。 罗杰斯尔和S。 Yoon 1 摘要介绍了黏性不可压缩流动计算程序开发的最新进展。 检查了两种不同的流动求解器，一种使用人工可压缩性方法，另一种使用压力投影方法，然后应用于工程问题。 计算的例子包括收缩通道中的脉动性流动、机翼翼尖涡的形成和传播、液体火箭发动机的高级叶轮分析以及心室辅助装置的开发。 最后一个例子是将这种计算技术扩展到非航空航天应用。 从算法开发和应用角度讨论了涉及这些应用程序的数值问题。 18.1 导言 从空气动力学家的角度来看，不可压缩流动可以被视为可压缩流动的极限情况，因为与声速相比，流速接近明显较低的值。 在航空航天和其他领域，有大量具有实际重要性的流动问题属于这一类别，例如1美国宇航局艾姆斯研究中心高级计算方法分部，加利福尼亚州莫菲特菲尔德2 MCAT公司，加利福尼亚州山景城■*加州大学戴维斯分校计算流体动力学-1998年。 编辑：David A. Caughey和Mohamed M. 哈菲兹。 © 1998 世界科学

334 KWAK、KIRIS、DACLES-MARIANI、ROGERS和YOON可以被视为本质上不可压缩。 代表这些流动的不可压缩的Navier-Stokes 方程提出了一个满足质量守恒方程的特殊问题，因为它没有与动量方程耦合。 在物理上，这些方程的特点是压力波的椭圆行为，其速度是无限的。 过去已经开发了各种方法，可以根据公式、变量或算法的选择，以多种方式进行分类。 由于涉及复杂几何形状的三维应用是我们的首要兴趣，因此在本研究中选择了原始变量公式。 与流函数或涡量等推导量相比，原始变量，即压力和速度，可以在实几何中轻松定义。 因此，为了方便和灵活性，美国宇航局艾姆斯研究中心使用原始变量公式开发不可压缩的Navier-Stokes代码（INS3D代码系列）。 在本文中，讨论了INS3D两个版本的一些算法功能，然后是计算结果的演示。 这里介绍的解决方案程序是在结构化网格框架内。 在Kwak等人中可以找到对我们工作的总体审查。 [1、2]。 在过去的几年里，大量关于CFD的评论文章和书籍都包括了关于不可压缩流动计算的讨论。 为了更全面地回顾不可压缩流动的一般计算方法，读者可以参考这些材料，即Hirsch [3]、Hafez和Oshima [4]。 18.2 解法 本节回顾了INS3D开发中使用的两种解法。 将首先给出控制方程，然后讨论两种方法。 18.2.1 定性密度的三维不可压缩流动由以下Navier-Stokes 方程控制：£ = 0 (18.1) du，du-Uj \_ dp d\%是dt dxj dx ( dxj，其中t是时间，x（笛卡尔坐标，M，相应的速度分量，p是压力，T-黏性应力张量。 所有变量都已通过参考速度和长度尺度进行非维度化。 在广义的曲线坐标中，（f，rj，£），这些方程可以写成

不可压缩的NAVIER-STOKES 335 dt H> (ei-evi) + s = \textasciitilde{}r \_d\_ (u-i\$)，其中|, = |,h或z for i=l, 2, or 3 „ 1 X\$t)xP + »U t:-VI,. V|, d] urfe) r+fe) x»+fe)/ +teL (18.3) (18.4) (18.5) 7 = 变换的雅各布 u = 运动黏度 s = 源项 源项 s 用于在稳定旋转参考系中表示离心力和科里奥利力。 对于大多数流应用程序，此项设置为零。 18.2.2 基于可压缩流动算法的方法：可压缩性方法人工：0（18.6）CFD与可压缩流动计算一起取得了重大进展。 因此，能够使用其中一些算法是非常有趣的。 为此，可以使用Chorin [5]的人工可压缩性方法。 在这个公式中，通过添加压力项的时间导数来修改连续性方程，结果为1 dp du t，其中/\}是人工可压缩性参数。 与非稳动量方程一起，这形成了一个双曲-波形类型的与时间相关的方程组。 因此，可以实现为可压缩流开发的隐式方案。 需要注意的是，在这个公式中，t不再代表真正的物理时间。 从物理上讲，这意味着有限速度的波被引入不可压缩流动场，作为分配压力的介质。 对于真正的不可压缩流动，波速是无限的，而这些伪波的传播速度取决于人工可压缩性的大小

336 KWAK、KIRIS、DACLES-MARIANI、ROGERS和YOON参数。 在真正的不可压缩流动中，压力场受到流量扰动的瞬息影响，但使用人工可压缩性，流量扰动及其对压力场的影响之间存在时间滞后。 理想情况下，人工可压缩性参数的值应在特定算法选择允许的尽可能高，以便快速恢复不可压缩性。 这必须在不降低实现的数值方法的准确性和稳定性属性的情况下完成。 另一方面，如果选择人工可压缩性使这些波传播太慢，那么伴随这些波的压力场变化非常缓慢。 这将干扰黏性边界层的正常发育。 在黏性流动中，边界层的行为对流向压力梯度非常敏感，特别是当边界层分离时。 如果存在分离，以有限速度传播的压力波将导致局部压力梯度的变化，这将影响流动分离的位置。 分离流量的这种变化将反馈到压力场，可能阻止收敛保持稳定状态。 当黏性效应对整个流场很重要时，就像在大多数内部流动问题中一样，伪压力波和黏性流动场之间的相互作用尤为重要。 人工可压缩性放宽了在每一步中满足质量守恒的严格要求。 然而，为了利用这个便利的功能，从物理和数学上了解人工可压缩性的性质至关重要。 Chang和Kwak [6]报告了人工可压缩性的详细信息，并提出了一些选择人工可压缩性参数的指南。 据报道，从这个概念演变而来的各种应用用于获得稳态解决方案（例如，[7-9]）。 为了使用这种方法获得与时间相关的解，可以在每个物理时间步骤中应用一个迭代程序，以便满足连续性方程（见[10-12]）。 关于人工可压缩性方法的进一步讨论可以在参考中找到。 [13,14]。 18.2.2.1 稳态公式组合方程（18.6）和动量方程给出以下方程组：i°=-\textasciicircum{}-\textasciicircum{})+s=-k（i8 -7），其中R是动量方程的右侧，可以定义为稳态计算的残差，其中\textasciitilde{}P\textasciitilde{} J J E.= i V-te），E。 =（18.8）

INCOMPRESSIBLE NAVIER-STOKES 337 通过将前一个时间步骤的通量向量线性化，可以获得非因数隐式方案。 对于欧拉隐式情况，这可以写成/AT KdDj n J（这个方程在伪时间内迭代，直到解收敛到稳态，此时原始不可压缩Navier-Stokes 方程得到满足。 方程的直接反转（18.9）将成为稳态解的牛顿迭代。 然而，在三维中，非因数方案的大块带状矩阵的直接反转是不切实际的。 因子方案 为了克服非因子隐式方案的困难，许多研究人员设计了间接方法。 Beam and Warming [15]和Briley和McDonald [16]的交替方向隐式方案（ADI）通过三个一维算子的乘积近似了无因数方案中的隐式算子。 很难将ADI方案应用于方程（18.9）的完整矩阵形式。 注意到在稳态下，方程（18.9）的左侧接近零，黏性项的简化表达式可以在左侧使用。 为了保持解决方案的准确性，整个黏性术语在右侧使用。 反转被简化为一系列一维问题。 这个方案在二维上是无条件稳定的。 在三维中，通过使用数值耗散，该方案变得有条件稳定。 ADI方案引入了一个因数分解误差，降低了收敛的速率。 然而，这种方法也许引入了CFD的第一个隐式方案，因此黏性流动计算变得更加可行。 在Jacobian的对角化降低成本的帮助下[17,18]，该方案已被用于许多求解器的开发。 在美国宇航局Ames，使用这种方法开发了第一个广义三维坐标中的不可压缩Navier-Stokes求解器，INS3D代码（Kwak等人，[19]；Chang和Kwak，[6]）（有关用户指南，请参阅Rogers等人。 [20]）。 INS3D代码被广泛用作航天飞机主发动机（SSME）热气歧管升级设计的主要工具（见，[9]）。 这是CFD对火箭推进系统开发的早期贡献之一。 为了提高ADI方案的稳定性和收敛，已经开发了几种双因素方案。 LU-SGS方案结合了LU因数分解和对称高斯-塞德尔松弛的优势，由Yoon和Kwak在人工可压缩性公式中实现[21]。 在这种方法中，黏性通量的导数使用二阶中心差来近似。 对流通量项使用中心差分离散，这需要数值耗散项来稳定。 通过选择不同的数值耗散模型和雅各布矩阵，可以构建各种方案。

338 KWAK,KIRIS,DACLES-MARIANI,ROGERS \& YOON 非因数方案 为了消除因式分解错误，可以形成类似于MacCormack [22]使用的线松弛隐式方案，而不是通过左侧矩阵的因式分解，而是通过迭代解过程。 方程（18.9）左侧矩阵的离散形式是一个带状矩阵，使用迭代方法近似求解。 三个计算方向中的一个被选择为隐式方向，而扫荡域则在另外两个方向上进行。 该算法的实现是，可以为扫描方向选择任何或所有三个计算方向。 最佳方向和扫描次数在很大程度上取决于问题。 该算法的经验表明，对于大多数问题，最好使用墙法线方向作为隐式方向，并且应该使用10次扫描的东西。 Rogers等人用这种方法开发了基于人工可压缩性的新流动求解器，即INS3D-UP代码。 [10]。 由于在人工可压缩性公式中，控制方程被更改为双曲-贬角线类型，因此为可压缩流动方程（Roe，[23]）开发的迎风差分方案被纳入流动求解器中。 这导致了一个更对角的主导系统。 18.2.2.2 时间精确配方 由于控制方程的椭圆性质，不可压缩流动的时间依赖性计算特别耗时。 从数值上讲，这意味着在每个时间步骤中，压力场必须经历一次完整的稳态迭代。 在瞬态流动中，物理时间步骤通常必须很小，每次时间增量期间流动场的变化可能很小。 在这种情况下，获取无发散流场的每个时间步骤的迭代次数可能不如获得常规稳态解所需的次数高。 然而，时间准确计算通常比稳态计算耗时一个数量级。 Rogers等人[10]使用人工可压缩性开发了一种时间准确的方法。 在这个公式中，动量方程中的时间导数使用二阶、三点、向后差公式3u"+,-4u"+un-] An+1 m =\textasciitilde{}r (18 -10)进行微分，其中上标n表示时间t = nAt的数量，r是方程（18.3）中给出的右侧。 为了求解在（n +1）时间水平下发散自由速度的方程（18.10），引入了一个伪时间水平，并用上标m表示。 方程是迭代解的，当U"+Um+*的发散接近零时，w" +l,"'+1接近新速度K" +1。 为了将这个速度的发散推到零，引入了以下迭代关系：

INCOMPRESSIBLE NAVIER-STOKES 339 „n + l,m+l \_n+l,m P \textasciitilde{}P AT -j8V.fi” (18.11) 将方程（18.11）与动量方程相结合，并在m + l伪时间水平上线性化残差项，结果为以下方程为delta形式A \textbackslash\{\}\#i+l.m I An+l,m+l An+l,m\textbackslash\{\} 其中7，r是一个对角矩阵，由l„ = diag = -\&»*■"• -Ia-h5Dn+]m -2D" +0.5D"-') (18.12) At \textbackslash\{\} ' J\_ L5 L5 L5 At' At' At' At' At (18.13）可以看出，该方程与（9给出的稳态公式非常相似。 两个方程组都需要相同残差向量R的离散化。 尽管这种配方已被成功使用，但针对时间依赖性问题，加快流动求解器是可取的。 在下一节中，将讨论一种基于分数步骤方法的方法；这已被发现是获得时间准确解决方案的有效方法。 18.2.3 基于压力投影的方法 1965年，Harlow和Welch [24]发表了第一个使用泊松压力方程的原始变量方法。 在这种方法中，称为标记和细胞（MAC）方法，压力被用作映射参数，以满足连续性方程。 通过取动量方程的发散，得到压力的泊松方程：dx，dt dx，（18.14）其中h，=-dU-U，dt::dx，dx-通常的计算程序涉及选择当前时间步的压力场，以便在下一个时间步骤中满足连续性。 原始的MAC方法基于二维笛卡尔网格上的交错排列。 交错网格以自然的方式保存质量、动量和动能，并避免了常规网格中遇到的压力的奇数-偶点解耦合（见[25]）。 尽管原始方法使用了明确的欧拉求解器，但各种时间

340 KWAK、KIRIS、DACLES-MARIANI、ROGERS \& YOON推进计划可以通过此公式实施。 自推出以来，已经设计了许多MAC方法的变体，并进行了成功的计算。 MAC方法可以被视为投影方法的特殊情况（即 Chorin [26]）或分步法（见Yanenko [27]，Marchuck [28]）。 在这种方法中，在每个步骤中获得发散自由速度场的正确压力的严格要求可能会大大降低整体计算效率。 为了满足网格空间的质量守恒，泊松方程中第二导数的差分形式必须与离散动量方程一致（见夸克[1]）。 要求解稳态解，只有当解收敛时才需要正确的压力场。 在这种情况下，压力的迭代程序可以简化，这样它只需要在每个时间步骤中迭代几次。 使用这种方法最著名的方法是用于压力链接方程的半隐式方法（SIMPLE）[29]（有关最近的进展，见Chen等人[30]）。 这种方法的独特之处在于估算速度和压力校正的简单方法。 这个功能简化了计算，但将经验主义引入了方法。 尽管有经验主义，但该方法已成功用于许多稳态计算。 本文无意评估这种方法，对这种方法感兴趣的读者请参阅上述参考文献。 计算程序 在这里，时间演变可以通过几个步骤来近似，其中可以通过将动量方程视为对流、压力和黏性项的组合来完成算子分割。 这种方法的常见应用需要两个步骤。 第一步是使用动量方程求解辅助速度场，其中压力梯度项可以从上一个时间步骤的压力中计算出来。 在第二步中，计算压力，可以将辅助速度映射到无发散速度场上。 笛卡尔坐标中的以下示例说明了此过程：第1步：计算辅助或中间速度，«r，通过\textasciicircum{}，I(3*;- Sr,)-|: + i\_Lv\textasciicircum{} +,,) (1,15）其中H,■ =---u\{u,, Re =雷诺数 Sxi 在这里，对流项由二阶Adams-Bashforth方法推进。 压力梯度项被添加到程序中，以尽量减少每个时间步骤中的压力校正。

INCOMPRESSIBLE NAVIER-STOKES 341 第2步：解决压力校正问题。 在第二步中，动量方程可以写成，\textasciicircum{}£A-\textasciitilde{}T-(r1-r) (i8.i6) 在2 ox t，该方程与连续性方程相结合，得出以下用于压力校正的泊松方程。v2(0» +1 -\textasciicircum{}) = AA 5 (18.17) v ' At Sx：一旦计算出压力校正，新的压力和速度计算如下：p"+'=p"+(0« +'-0-)--iLv 2(f,+, -0") (18.18) „"+1=M"-\textasciicircum{}-i\_(0"'\_0») (18.19) 在广义坐标中使用分数步进方法开发了成功的方法（Rosenfeld等人，[31-33]，Wessling等人。 [34]）。 这种方法的一个特殊方面需要特别注意，是中间的边界条件。 Rosenfeld等人[31,32]设计了一个通用方案，物理边界条件可以在中间步骤中使用。 与其他基于压力的方法一样，分数步进方法的效率取决于泊松求解器。 多网格加速在物理上与椭圆场一致，是提高计算效率的可能途径之一[33]。 Rosenfeld等人[31]定义了因变数，即压力和体积通量分别在原发细胞的中心和表面。 这种选择等同于交错网格上的有限差分公式，选择缩放的逆变速度分量作为未知数。 黏性项通过近似因式分解处理。 由此产生的求解器成功验证了需要小物理时间步骤的依赖性问题。 为了放松三维CFL数限制，并提高求解器的网格依赖性鲁棒性，Kiris和Kwak [35]实施了一种放松方案，其中对流和黏性项都隐式处理。 第一步是求解辅助速度场的动量方程。 三点后向微分公式用于时间离散动量方程，如下所示：-U3«,. -4K," + «r') = -\$f- + hXii,) (18.20) 2Atv ' ax: 运算符拆分的公式是为了实现整体时间精度。 生成的求解器在具有度量不连续性的网格条件下进行验证。 发现可以使用大的CFL数字而不会造成不稳定

342 KWAK、KIRIS、DACLES-MARIANI、ROGERS和YOON问题。 当时间步骤不受问题的物理限制时，这是非定常流动计算的理想特征。 使用这种方法开发了一个时间准确版本INS3D-FS代码[35]。 18.2.4 人工可压缩性与 时间精确计算的压力投影方法 为了讨论计算时间精确流量的两种方法的特点，选择了收缩通道中的脉动流量。 计算使用两个版本的INS3D进行，即INS3D-UP和INS3D-FS分别代表人工可压缩性和压力投影方法。 几何形状与Park [36]的实验设置一致，如图所示。 1. 收缩的高度由a = 0.57给出。 这是从通道顶壁到收缩中最低点的距离。 收缩上游通道的长度由L„ = 7给出。 收缩长度isLc = 4.66，通道的下游部分由实验的流入边界给出，位于收缩上游的100个通道高度。 然而，计算流入边界被放置在收缩的上游的七个通道高度，具有抛物线速度曲线，这样质量流与实验设置的流相匹配。 脉动流入速度由图中给出的形状函数给出。 2. 这个问题之前是计算过的（见Wiltberger等人，[37]），为了比较两种算法，在这里重复。 为了获得时间准确的解决方案，INS3D-UP在每个时间步骤中都进行小迭代，直到满足不可压缩性，而INS3D-FS在每个时间步骤中都是时间准确的。 如图所示，所得解决方案是可比的。 3和4。 在实验中，测量了沿底壁的涡流中心的位置，定义为B-涡流。 B涡流紧挨着收缩后生长，并向下游脱落。 在图中。 3，从这些计算中生成的流线轮廓以周期0.1的时间增量显示。 这两个代码的结果显示了相同的现象，INS3D-FS在整个周期内产生更大的涡流结构。 然后根据图中的实验数据绘制B涡流的位置。 4. 这两个代码都与实验结果相比较。 为了比较时间准确的程序，在一个时间步骤的推进中详细研究了迭代过程。 在INS3D-UP代码中，时间精度是通过在每个步骤中迭代来实现的。 这种分级是用修改为双曲-半斜线类型的控制方程完成的。 因此，从出口边界传播的上游压力波必须有足够的经过时间来平衡通道内的黏性效应。 完成此操作所需的迭代次数与人工可压缩性参数的大小直接相关，人工可压缩性参数之前定义为f\$。 这种现象如图所示。 5a。 另一方面，INS3D-FS使用泊松方程来测量压力，将流场映射到一个

不可压缩的NAVIER-STOKES 343在新时间水平上的无发散速度场。 泊松迭代在这里使用GMRES 收敛加速度的点高斯-塞德尔迭代完成，没有像人工可压缩性案例那样表现出波样现象。 这如图所示。 5b。 与泊松迭代相比，在INS3D-UP中进行隐式反演通常更昂贵。 然而，对于稳态解决方案，INS3D-UP可能需要比INS3D-FS更长的时间步骤。 在我们的计算中，INS3D-FS需要1.5到2倍于稳态解决方案的CPU时间。 在需要小时间步骤来解决物理问题的时间依赖性情况下，INS3D-FS的运行速度至少比INS3D-UP快三倍或更大。 这两种方法的算法特征将在未来的报告中进一步讨论。 18.3. 计算结果 在前几节中，简要回顾了基于原始变量方法的解决方案方法和相关代码。 这些求解器在过去被应用于许多工程问题。 在本节中，从航空学的外部和内部流动问题中选择计算的例子，然后是生物流体问题的应用。 18.3.1 翼尖涡流的形成和传播 翼尖涡旋的研究在流体工程的许多领域都非常重要。 它的重要性可以从飞机着陆和起飞期间的间距、叶片/涡旋涡相互作用对旋翼飞机性能以及翼尖涡空穴对船舶螺旋桨性能等问题中看出。 尽管在翼尖涡上以理论、实验和计算研究的形式做了大量工作，但真正涉及近场细节的工作很少。 本研究的重点是近场物理学，它将为模拟中场和远场区域的尾波传播提供流入条件。 翼表面的形成 Dacles-Mariani等人[38,39]使用ENS3D-UP代码与Chow等人的实验研究一起对形成和卷起过程的计算研究。 [40]。 实验和计算都对初始汇总过程和近场细节进行了非常详细的研究。 在计算研究中，计算域由翼风洞壁几何组成，如图所示。 6a，这与美国宇航局艾姆斯32英寸x48英寸低速风洞的实验设置很接近。 计算域包括一个带有NACA 0012 翼型截面的矩形半翼、一个圆形翼尖和周围的边界。 机翼的长宽比为0.75，以10度的攻角安装在风洞内。 足够的

344 KWAK,KIRIS,DACLES-MARIANI,ROGERS \& YOON流是湍流，根据和弦长度，雷诺数为460万。 使用实验值规定了速度曲线的流入边界条件，并根据使用一维黎曼不变量的特性方法计算流入压力。 在固体表面（风洞壁和模型表面），速度被指定为零。 单一网格方法（见图。 6a）用于这项研究。 解决这个问题的物理需要一个非常密集的网格。 因此，CPU要求很大。 Cray C-90的250万个网格点需要28 CPU小时的总运行时间，残差下降了五个数量级。 这项研究主要侧重于解决翼尖涡计算的基本问题，而不解决减少网格点总数的需要。 然而，人们认识到，为了更实际的实施，必须解决电网密度要求。 如图所示。 6b，涡旋在由涡量从尖端边界层供电的机翼尖端形成。 机翼上表面和下表面之间存在的压差驱动边界层中的流体围绕尖端和机翼的吸侧。 形成的离散涡旋变得高度三维和复杂。 随着涡流向下游移动，它卷起越来越多的机翼尾迹，直到其循环名义上等于机翼的循环。 与进场路径上的飞机分离相比，这种卷起距离很小，但与相互作用的起重面之间的距离相比，可能不一定小。 因此，近场的流动本身就很重要，并提供了控制远场涡流的可能手段。 随着水流向下游前进，涡核的最大横流速度增加。 轴向压力梯度反过来发展，加速了涡核中的流体。 Phillips和Graham [41]表明，对于湍流涡流，假设v' e2 = v'2，我们得到dr r dr这个方程表明，导致径向压力梯度的主导项是一个包含v e的项，这反过来又影响了峰值轴向速度。 然而，请注意，方程右侧的第二个项不能被忽视，因为Zilliac等人在实验中测量的近场核心湍流强度很高。 [42]。 这表明，核心属性的分辨率不仅取决于周向速度的梯度，还取决于湍流应力项之一。 如果这个术语没有正确考虑，测量值和计算值之间可能无法进行很好的比较。 近场传播 一旦形成，由于其在核心区域的欧拉性质，尾涡几乎可以长期保持其强度。 这也需要特殊的湍流建模

不可压缩的NAVIER-STOKES 345作为核心区域的高网格分辨率。 首先，分数步进代码INS3D-FS用于调查空间差分和湍流模型引起的误差水平。 然后比较了INS3D-UP和INS3D-FS的唤醒速度曲线。 在这项研究中，使用了Baldwin-Barth [43]单方程湍流模型。 对于INS3D-FS计算，计算域包括从机翼的后缘（x/c=1.0）到机翼下游0.673弦长度的区域，具有H-H网格拓扑结构。 广泛的实验数据[44]在x/c=1.0、1.12、1.24、1.447和1.673处提供。 x/c=1.0站的实验速度曲线用作流入边界条件。 边界处的压力分布是根据兼容性条件计算的。 计算使用尺寸为36x82x82的相对精细的网格进行。 该解决方案在1000次迭代中收敛到机器精度。 此计算所需的CPU时间约为3.5 CRAY-C90小时。 最初，计算是使用粗网格进行的。 发现涡核速度峰值预测不足。 通过将k和/方向的网格点数量从36x42x42的网格尺寸增加一倍到36x82x82，提高了网格分辨率。 对涡核峰值的预测有所改善，但没有实质性。 数值结果表明，涡核的数值耗散在下游时过多。 这与使用人工可压缩性的翼尖涡研究结果一致（见[38,39]）。 当使用应变率张量的范数修改湍流模型中的生产期限时，涡核的过度数值扩散减少了（见Kiris和Kwak，[35]）。 在无花果中。 7a和7b，通过显示沿涡流核心线流动量的轴向进展来比较湍流模型和三阶和五阶对流微分方案的影响。 当使用三阶通量差分时，数值耗散的数量很大。 为了减少这个数值耗散，人们可以使用更精细的网格，或者在上风偏置差异中增加模版。 由于提高差分精度的成本比增加网格大小的成本低得多，因此使用了第五阶迎风差分。 应该指出的是，尽管对流项使用了五阶迎风差分，但该方法的整体空间精度是二阶的。 这是因为体积和表面积向量被评估为二阶精度。 简单平均用于半点位置的度量项，二阶中心差分用于黏性项。 然而，与低阶微分相比，增加迎风差分中的模版对减少数值耗散量有显著效果。 图中比较了使用三阶和五阶方案的结果。 7a和7b。 随着尾迹向下游传播，核心区域仍然表现出过度耗散，主要是由于湍流黏度高。 为了研究湍流模型中生产期限的影响，对涡量大小和应变率的各种组合进行了实验[35]。 图中可以清楚地看到湍流模型在核心区域的影响。 7a和7b。 通过修改

346 KWAK、KIRIS、DACLES-MARIANI、ROGERS \& YOON 核心生产术语，以纳入核心区域的应变率，将误差减少到2\%以下。 最后，图中绘制了三个不同尾迹位置的速度幅度和交叉流速度。 8. 如图所示，两个代码的结果相当接近，并且都显示了很好地捕获涡核的能力。 18.3.2 液体火箭发动机的应用 直到最近，高性能泵的设计过程与三十年前没有显著区别。 在此期间，大量的实验和操作经验表明，泵流的许多重要特征在半经验设计过程中没有充分考虑。 在同一时间段内，计算机、数值算法和物理建模取得了巨大的进步。 在将原始的INS3D应用于SSME热气歧管的重新设计后，提高先进火箭发动机中涡轮机械部件的性能变得非常重要。 这促使开发了用于火箭泵流量模拟的CFD程序。 感兴趣的液体燃料和氧化剂泵正在以恒速运行。 因此，首先寻求旋转稳态解决方案，从而在稳旋转参考系中开发使用INS3D-UP代码的泵模拟程序。 火箭泵涉及全叶片和部分叶片、尖端泄漏和扩散器的出口边界。 除了几何复杂性外，涡轮泵流还遇到了各种流动现象。 这些包括湍流边界层分离、唤醒、过渡、翼尖涡、三维效果和雷诺数效果。 为了在设计过程中使用CFD，必须验证计算流程分析工具，以便设计师能够定义这些工具的准确性和变化。 泵应用的CFD程序的验证侧重于火箭感应器。 由美国宇航局元帅太空飞行中心组织的泵技术团队进行了广泛的计算验证（见，Garcia等人，[44]）。 然后，由此产生的计算程序被应用于通过SSME高压燃料涡轮泵叶轮的流动和先进的泵叶轮的开发（Kiris和Kwak，[45]）。 这里介绍了高级叶轮流量分析的结果。 在图中。 9，先进叶轮的横截面图以示意图显示。 先进泵的计算模型包括叶轮和出口腔区域。 图10显示了叶轮枢纽区域附近的计算网格。 叶轮设计流量为1,205加仑/分钟，设计速度为6,322 rpm。 此计算的雷诺数是每英寸181,273。 在图中。 11，子午线速度，即圆周平均轴向速度，显示在叶轮放电处。 相对x距离是从裹尸布到轮毂测量的，其中x=1.0是轮毂。 合并了经度速度，Cm

不可压缩的NAVIER-STOKES 347沿着每个恒定x位置的径向条，它们由249.5英尺/秒的轮尖速度进行非维度化。 还绘制了出口裹尸布腔5\%和10\%再循环的子午线速度分布。 当出口护罩腔泄漏到叶轮眼时，叶轮出口处的速度峰值向b2宽度的中心移动，其中b2被定义为叶轮出口处的叶片高度（见图。 9）。 然而，在r/rtip= 1.0275时，裹尸布泄漏对溶液只有轻微影响（图。 11）。 在图中。 11，符号代表Brozowski的实验数据[46]，线条代表叶轮出口处无叶片空间流动的Cm分布。 测试数据显示，峰值更接近b2宽度的中心。 计算结果和实验数据之间的差异部分是由于轮毂腔中的再循环流。 轮毂腔的泄漏导致更强的再循环区域，该区域将速度峰值转移到b2宽度的中心。 由于CFD分析不包括轮毂腔的泄漏，因此无叶片空间中预测的再循环区域没有实验研究中那么强。 图12显示了叶轮出口处叶片到叶片的速度分布。 叶片到叶片的速度分布说明了叶轮出口流失真。 符号代表实验数据，线条代表计算结果。 在两个子午位置都捕捉到了喷气醒模式，该模式在扩散器叶片上产生不稳定的负载。 总体而言，数值结果与实验数据相当好。 18.3.3 心室辅助装置的设计 当自然心脏无法提供足够的血流时，心室辅助装置（VAD）是一种救生工具。 1989年，美国宇航局和贝勒医学院（BCM）的德巴基心脏中心开始开发一种新的植入式VAD系统（图。 13）。 该VAD基于磁驱轴向泵，在100毫米汞柱压力下需要每分钟5升的血流率（图。 14）。 为了使其可植入，该设备必须变得小，这需要非常高的转速。 必须解决两个主要问题，才能使这一设计投入使用。 首先，红细胞损伤必须保持在低于可接受的水平，其次，必须防止血液凝固在任何局部区域形成，这将阻止泵旋转。 此外，血液必须正确冲洗，因为血凝块的形成可能出现在停滞区域。 由于设备尺寸小，操作条件恶劣，流量测量仪器极其困难。 因此，有必要使用计算技术来观察流程。 INS3D-UP代码已被用于解决这些问题，并使泵符合临床测试的要求。 将流动求解器应用于火箭发动机开发的经验教训被重新应用于VAD分析和设计中。 首先，在稳定旋转的参考框架中模拟了通过VAD叶轮基线设计的流动。 本分析使用了区域多块网格，共有35万个网格点。 这个叶轮的设计流量是每5升

348 KWAK,KIRIS,DACLES-MARIANI,ROGERS \& YOON分钟，设计速度为每分钟12,600转/分钟（rpm）。 这个问题由管径（0.472英寸）和叶轮尖端速度非维度化。 当最大残差下降至少五个数量级时，该解决方案被认为是收敛的。 每次迭代每个网格点所需的计算机时间约为1.5 x 10英寸*秒。 这些计算所需的计算机总时间约为6到8个单处理器Cray-C90小时。 进行了参数化研究，以优化叶轮叶片形状和尖端间隙。 最初，分析了三种不同的叶轮叶片设计，尖端间隙为0.009英寸。 然后，图中所示的设计。 用两个尖端间隙分析了15a；与0.009英寸的尖端间隙相比，0.0045英寸的尖端间隙在效率和头部系数方面显示出更好的流体动力学性能。 使用这种尖端间隙为0.0045英寸的设计作为基线叶轮设计，引入了火箭推进的想法，以开发一种新的可植入VAD（图。 15b）。 研究了由基线叶轮和感应器组成的新设计。 基线叶轮和新叶轮的轮毂和叶片表面，由非维度压力着色，如图所示。 15. 压力由pV 2非维度化，其中p是密度，V是叶轮尖端速度。 显示了由于离心力作用导致的叶片上的压力梯度，以及从流入到流出的压力上升。 这里显示了沿叶轮叶片高度的子午线速度分布，用于各种设计（见图。 16）。 最终设计基本上消除了早期设计中存在的回流。 潜在血液凝固的关键区域之一位于旋转和非旋转组件之间的轴承区域附近。 凝血可能由于高剪切或停滞而导致枢纽区域的凝血，这取决于该区域的间隙和配置。 分析了几种腔体形状和间隙宽度。 图17显示了腔体宽度为b的原始基线设计和间隙宽度为8b的最终设计的速度向量。 原始设计显示，旋转轮毂面附近有非常高的剪切应力，腔体下部有一个非常停滞的流体区域。 将腔体宽度增加到8 b表明腔体中的再循环增加。 这个和其他修改被纳入了最终配置，它通过了VAD两周临床使用的要求。 表1给出了原始设计和新设计之间的性能比较。 表中的临床结果是由BCM获得的。 本表中报告的溶血指数显示了泵产生的血红蛋白量，以每100升的克为单位。 红细胞的破坏导致血红蛋白的释放。 与基线设计相比，新设计显示出显著的性能提升。 新旧设计之间的效率提高了22\%。 感应器为流量提供足够的压力上升，以防止叶轮叶片的空化。 除了提高泵送效率外，VAD的设计还要求通过减少停滞区域，在实心墙附近进行良好的墙面清洗。

不可压缩的NAVIER-STOKES 349 原始设计新设计泵效率0.25 0.33溶血指数0.02 0.003所需功率12.6瓦9.8瓦旋转（转速）12,600 10,800血栓形成是否测试运行时间2天30天以上表1。 NASA DeBakey VAD的原始和新设计的性能比较。 通过使用计算方法，VAD的性能得到了极大的提高。 此外，详细的流量分析导致开发了非凝血泵。 1991年7月，医学研究所估计，北美每年约有25,000至60,000名患者可以从高效的左心室辅助装置中受益。 因此，通过使用CFD工具实现的改进设计将对管理有关键需求的患者的健康产生深远的影响。 18.4 结论 在本文中，讨论了主要为三维流动模拟设计的不可压缩Navier-Stokes求解器。 由于原始变量公式在设置边界条件时造成的复杂性较少，因此讨论仅限于原始变量公式。 计算结果用于说明数字程序。 尽管计算机速度和内存最近大幅提高，但流动求解器的速度和内存要求仍然是影响周转时间的主要因素。 INS3D-UP是基于人工可压缩性方法的上风有限差分代码，正被应用于稳态、时间准确和旋转稳态解决方案的广泛应用。 INS3D-FS基于在交错网格上使用有限体积离散的分数步进方法，旨在解决与时间相关的问题。 对于稳态解决方案，INS3D-FS需要比INS3D-UP多1.5到2倍的CPU时间。 对于不稳定的计算，INS3D-FS是时间准确的，而INS3D-UP需要在每个时间级别进行分级，因此INS3D-FS预计会很经济。 由于这些代码适用于各种问题，我们希望更详细地量化这些代码的优点和缺点。 尽管数值解程序在算法速度、准确性和网格依赖性方面存在局限性，但当正确理解其适用范围时，这些求解器对现代流设备的开发人员来说具有重大价值。

350 KWAK,KIRIS,DACLES-MARIANI,ROGERS \& YOON REFERENCES 1. Kwak，D.，“黏性不可压缩流动的计算”，冯·卡曼流体力学研究所，1989-04年系列讲座。 还有美国宇航局TM 101090，1989年3月。 2. Kwak, D., Kiris, C, Dacles-Mariani, J., Rogers, S., and Yoon, S., “三维稳定和非定常流动模拟的不可压缩Navier-Stokes求解器”，计算流体动力学 Reviewl997, Hafez, M. and Oshima, K., ed. John Wiley and Sons，1997年。 3. Hirsch，C，内部和外部流动的数值计算。 John Wiley \& Sons，1988年。 4. Hafez，M.和Oshima，K.，ed。 计算流体动力学。 John Wiley and Sons，1995年和1997年。 5. 乔林，A。 J.，“解决不可压缩黏性流动问题的数值方法”，J. 补偿。 物理，卷。 2，第12-26页，1967年。 6. 昌，J。 L。 C和Kwak，D.，“关于数值求解不可压缩流动的伪可压缩性方法”，AIAA论文84-0252，AIAA第22届航空航天科学会议，内华达州里诺，1984年1月9日至12日。 7. 斯特格，J。 L.和Kutler，P.，“计算涡旋的隐性有限差分程序”，AIAA J.，卷。 15，不。 4，pp. 581-590，四月 1977年。 8. 郭，D.，张，J. L。 C，Shanks，S。 P.和Chakravarthy，S.，“使用原始变量的三维不可压缩Navier-Stokes 流动求解器”，AIAA J，卷。 24，不。 3，第390-396页，3月。 1986年。 9. Chang，J.L.C.，Kwak，D.，Rogers，S. E.和Yang，R-J，“不可压缩流动的数值模拟方法和对缅因州航天飞机发动机的应用”，Int. J。 流体中的数值方法，卷。 8，pp. 1241-1268，1988年。 10. 罗杰斯，S。 E.，Kwak，D.和Kiris，C，“不可压缩Navier-Stokes 方程的稳定和非稳态解”，AIAA J. 卷。 29，不。 4. 603-610，四月ll.Merkle，C。 L.和Athavale，M.，“基于人工可压缩性的时间不稳定不可压缩流动算法”，AIAA论文87-1137，1987年。 12. Belov，A。 Martinelli，L.和Jameson，A.，“用于非稳定不可压缩流动计算的多网格的新隐式算法”，AIAA论文95-0049，1995年13。 Temam，R.，Navier Stokes方程。 修订版，北荷兰，1979年。 14. Rizzi，A.和Eriksson，L.-E.，“旋转的Inviscid 不可压缩流动的计算”，J. 补偿。 物理，卷。 153，第275-312页，1985年。 15. 梁，R。 M.和Warming，R. F.，“可压缩Navier-Stokes 方程的隐式因子方案”，AIAA J.，卷。 16，pp. 393-402，1978年。 16. Briley，W.R.和McDonald，H，“广义隐式方法的多维可压缩Navier-Stokes 方程的解决方案”，J. Comp Phys Vol 24 No. 4，第372-397页，1977年。v y '

不可压缩的纳维-斯托克斯351 17。 普利姆，T。 H.和Chaussee，D. S.，“隐式近似因数分解算法的对角形式”，J. 补偿。 物理，卷。 39，pp. 347-363，1981年。 18. 罗杰斯，S。 E.，Chang，J. L。 C和Kwak，D.，“伪可压缩方法的对角算法”，J. 补偿。 物理。 卷。 73，不。 2，pp. 364-379，1987年。 19. 郭，D.，张，J. L。 C，Shanks，S。 P.和Chakravarthy，S.，“使用原始变量的三维曲线坐标系中的不可压缩的Navier-Stokes 流动求解器”，AIAA论文84-0253，1984年1月。 20. 罗杰斯，S。 E.，Kwak，D.和Chang，J. L。 C，“INS3D-广义三维坐标中的不可压缩导航-斯托克斯代码”，美国宇航局TM 100012，1987年11月。 21. Yoon，S.和Kwak，D.，“具有源项的三维不可压缩Navier-Stokes 方程的LU-SGS隐式算法”，AIAA论文89-1964，1989年。 22. 麦科马克，R. W.，“Navier-Stokes 方程数值解的现状”，AIAA论文85-0032，1985年。 23. 罗，P。 L.，“近似黎曼求解器、参数向量和差分方案”，J. 补偿。 物理，卷。 43，第357页，1981年。 24. 哈洛，F。 H.和Welch，J. E.，“自由表面时间依赖黏性不可压缩流动的数值计算”，物理。 流体，卷。 8，不。 12，pp. 2182- 2189，1965年。 25. 格雷肖，M。 P.和Sani，R. L.，“关于不可压缩Navier-Stokes 方程的压力边界条件”，Int. J。 流体中的数值方法，卷。 7，pp. 1111-1145，1987年。 26. Chorin，A.J.，“Navier-Stokes 方程的数值解”，计算数学，卷。 22，不。 104，745-762，1968年。 27. Yanenko，N.N.，分数步骤的方法，Springer-Verlag，柏林，1971年。 28. Marchuk，G.M.，数值数学的方法。 Springer-Verlag，1975年。 29. 帕坦卡，S。 V.和斯帕尔丁，D. B.，“三维抛物面流中热、质量和动量传递的计算程序”，Int. J。 热量和质量转移，卷。 15，pp. 1787-1806，1972年。 30. Chen, Y.S., Shang, H.M and Chen, C.P., “统一CFD算法与基于压力的方法”，第6届Int'l. Comp研讨会。 Fluid Dyn.，1995年9月4日至8日，内华达州太浩湖。 31. Rosenfeld，M.，Kwak，D.和Vinokur，M.，“广义坐标系中非稳不可压缩Navier-Stokes 方程的分数步解法”，J. 补偿。 物理，卷。 94，No.l，第102-137页，1991年5月。 32. Rosenfeld，M.和Kwak，D.，「移动座标中黏性不可压缩流的时间依赖解」，Intl。 J。 数字。 流体方法，卷。 13，第1311-1328页，1991年。

352 KWAK,KIRIS,DACLES-MARIANI,ROGERS \& YOON 33. Rosenfeld，M.和Kwak，D.，“广义曲线坐标系中分数步求解器的多网格加速”，AIAA J. 卷。 31，不。 10，第1792-1800页，1993年10月。 34. Wessling，P.，Segal，A.，van Kan，J.J.I.M.，Oosterlee，C.W.和Kassels，C.G.M.，“交错网格上一般坐标中不可压缩Navier-Stokes 方程的有限体积差异化”，计算流体动力学杂志，卷。 1，第27-33页，1992年4月35。 Kiris，C.和Kwak，D.，“使用分数步法对不可压缩Navier-Stokes 方程的数值解”，AIAA论文96-2089，1996年6月。 36. Park，D.K.，“动脉骨流的生物流体力学”，硕士论文，LeHigh大学，宾夕法尼亚州伯利恒，1989年。 37. 威尔特伯格，N.L.，罗杰斯，S. E.和Kwak，D. “两种不可压缩的Navier-Stokes算法的非稳内部流动的比较”，美国宇航局TM 108794，1993年11月。 38. Dacles-Mariani，J. S.，S. 罗杰斯，D。 夸克，G。 Zilliac和J。 周，“翼尖涡流场的计算研究”，AIAA论文93-3010，1993年7月。 39. Dacles-Mariani，J.，Kwak，D.和Zilliac，G，“近场区域翼尖涡流场的准确性评估”，AIAA论文96-0208，1月。 1996年。 40. Chow，J.S.，Zilliac，G.G.和Bradshaw，P.，“湍流翼尖涡流的近场形成”，AIAA 93-0551，第31届航空航天科学会议，内华达州里诺，1月。 11-14。 1993年。 41. Phillips，W.R.C.和Graham，J.A.H.，“湍流拖曳涡流中的Reynolds应力测量”，J. 流体机械，卷。 147，第353-371页，1984年。 42. Zilliac，G.G，Chow，J.S.，Dacles-Mariani，J.和Bradshaw，P.，“湍流近场翼尖涡流的结构”，AIAA论文93-3011，1993年7月。 43. Baldwin，B.S.和Barth，T.J.，“高雷诺数壁边界流动的单方程湍流运输模型”，AIAA论文91-0610，1991年。 44. Garcia，R.，McConnaughey，P.和Eastland，A.，“MSFC泵阶段技术团队的活动”，AIAA论文编号 92-3232，1992年。 45. Kiris，C.和Kwak，D.，“推进流分析的不可压缩Navier-Stokes计算的进展”，美国宇航局CP 3282，第二卷，先进的地球到轨道推进技术，1994年。 46. 布罗佐夫斯基，L。 A.，Ferguson，T.V.和Rojas，L.，“叶轮流场激光测速仪测量”，第五届国际运输现象和旋转机械动力学研讨会论文集，1994年5月9日至11日，毛伊岛宽普卢。 47. Kiris，C，Kwak，D.和Benkowski，R.，“左心室辅助设备的计算流量分析”，第六国际。 对Comp的症状。 Fluid Dyn.，内华达州太浩湖，1995年9月4日至8日。

不可压缩的NAVIER-STOKES计算353 igure 2 流入速度幅度相对于收缩通道的一段时间。 Re = 131.9，基于流入速度。 图1 收缩通道计算模型的物理尺寸。 计算网格大小：241 x 63。

354 KWAK, KIRIS, DACLES-MARIANI, ROGERS \& YOON 图4 B-漩涡与周期的位置。 图3 0.1 T到1.0 T的流线型：（a）INS3D-Up结果：（b）INS3D -f3结果。

不可压缩的NAVIER-STOKES计算355图5一个时间步骤内的压力校正：（a）INS3D-UP再迭代；（b）INS3D-FS泊松迭代。 （b）图6机翼翼尖涡研究；（a）风洞测试部分和计算域的示意图；（b）显示翼尖涡卷起和核心运动的粒子痕迹；Re = 4.6 X 106。

356 KWAK, KIRIS, DACLES-MARIANI, ROGERS \& YOON 图 7 翼尖涡附近流动细节；（a）沿涡流核心线速度幅度的轴向进展；（b）沿涡流核心线速度的轴向进展。

不可压缩的NAVIER-STOKES计算357涡流在三个不同位置。 图8 跨尾流的速度幅度和交叉流速度的比较

358 KWAK, KIRIS, DACLES-MARIANI, ROGERS \& YOON igure 9 先进泵叶轮截面的示意图。 图10轮毂表面的高级泵叶轮计算网格。

不可压缩的海军-斯托克斯计算359平均Cm vs. 相对X：R/fitip - 1.0275 图11 图12中圆周平均子午线速度的比较 叶轮出口处叶片到叶片的子午线速度的比较。

360 KWAK、KIRIS、DACLES-MARIANI、ROGERS \& YOON图13轴流心室辅助装置（VAD）植入示意图。

不可压缩的NAVIER-STOKES计算361 F igure 14 S ch ematic of the NASA - De Ba key VAD。 图15 原始和新VAD叶轮的轮毂abd叶片表面因压力而阴影。

362 KWAK、KIRIS、DACLES-MARIANI、ROGERS \& YOON igure 16 沿各种设计的叶轮叶片高度的子午线速度分布。igure 17 原始和新轴承几何形状内的速度向量。

\clearpage
\chapter{翼型优化的利弊}
翼型优化的利弊 Mark Drela 1 19.1 介绍 长期以来，优化一直被认为是以正式和一般的方式解决空气动力学设计问题的手段。 Hicks、Murman和Vanderplaats[1]的早期工作调查了跨声速翼型流动的这种可能性，结果令人鼓舞，但显示了一些相当意想不到的结果和困难。 这些早期努力的特点是使用的设计参数相对较少，主要是由于黑匣子梯度计算通过有限差分的计算成本和可用的计算机资源有限。 参数梯度计算方法的最新进展，如伴随方法[2]和基于牛顿的直接方法[3]，以及计算机速度和内存容量的不懈增加，在很大程度上消除了对可以使用的设计参数数量的限制。 乍一看，这应该允许计算出“真正优化”的设计，从而解决优化的有效性问题。 作者最近的经验表明，情况并非如此，因为随着参数数量的增加，出现了不可预见的困难和令人惊讶的结果。 本文的目的是调查存在相对大量自由设计参数的约束优化解决方案的行为。 这些示例将仅限于二维黏性翼型优化。 即使在如此看似简单的1航空航天系中也出现了困难，麻省理工学院，剑桥，马萨诸塞州02139 计算流体动力学-1998年前沿编辑：David A. Caughey和Mohamed M. 哈菲兹 ©1998 世界科学

364 DRELA空气动力学设计问题很好地说明了形式优化作为一种空气动力学设计方法的优势和缺点。 19.2 方法摘要 本文侧重于优化本身的有效性，而不是分析和优化算法。 因此，本节将仅限于用于应用示例的分析和优化方法的简要摘要。 19.2.1 分析方法 目前的设计/优化方法使用黏性/非黏性MSES代码作为基础分析求解器[4]。 整体方程组R(U; a,M,G k) = 0 (19.1)由内部稳定Euler 方程、边界层方程以及必要的耦合和边界条件组成，通过直接应用牛顿方法，将流动解U求解为一个全耦合系统。 SU = - [dR/dU]' 1 R, U «- U + SU (19.2) 牛顿循环收敛后，分解的雅各布被重新用于计算对设计参数Gk和流量参数a, M的流场灵敏度。dUldGk = -[dR/dU]- 1 \{dR/dG k\} dU/da = -[dR/dU]' 1 \{dR/da\} (19.3) dU/dM = -\{dR/dU\}- 1 \{dR/dM\} 这种重复使用只涉及反向替换，并允许以微不足道的成本计算大量设计参数的流场灵敏度。 数百个设计参数可以在简陋的工作站上进行交互式计算。 19.2.2 几何参数化 翼型形状的合适参数化，以翼型两侧的分数弧长s/s s;de为单位定义，是正弦基函数c/jt的总和，该函数扰动翼型的距离为垂直于其当前表面。 设计参数Gk是模式振幅。 K 1 An(«) = y\textasciicircum{}GkgkU), g k\{>) = r siri(kns/s s\textbackslash\{\}Ae) (19.4)

翼型优化365 1/fc缩放系数使所有基函数具有相同的最大斜率，这在理论上不影响最佳解，但它似乎确实大大改善了优化下降序列的行为。 当然可以定义其他几何基础，事实上，许多几何基础更适合特定问题。 正弦基的有用特征是有保证的相互正交性，以及通过添加模式均匀和可预测的几何分辨率增加，使其非常适合本文中的参数计数调查。 19.2.3 优化方法 MSES牛顿求解器生成的体积解灵敏度输出应用于交互式LINDOP程序中的优化问题[3]。 这允许设计者使用灵敏度信息以交互方式尝试各种目标函数、约束和设计参数集，并生成由显式参数变化、强加压力分布或目标函数下降产生的线性化预测。 线性化预测会显示所有正在考虑的操作点，因此在优化遇到麻烦时发出视觉警告。 一个MSES/LINDOP周期构成了设计空间中的一个下降步骤。 众所周知的BFGS方法[5]用于生成下降方向的顺序。 下降继续，直到目标函数拒绝在Co中进一步下降到0.00005以内-所需的下降次数通常与自由设计参数的数量相当。 19.3 低雷诺数 翼型应用DAE-11 翼型由作者于1987年专门为Daedalus人力飞机设计，采用传统的逆方法和直接几何操作技术。 这是第二代翼型，由其前身机翼重新设计，其性能在飞行测试中得到了部分验证[6]。 因此，它成为通过数值优化进一步改进的有趣候选者。 人力飞机翼型的关键要求是在设计飞行升力系数时实现最小阻力。 它受到钢架位置的结构厚度要求的限制，以及影响机翼二级结构重量的其他一些次要几何要求的限制。 在这里，重点将放在主要要求上。

366 DRELA图1基线DAE-11 翼型，以及部分优化的无约束翼型。 19.3.1 一点优化 以下优化问题体现了低阻力要求。最小化！ F（Gki ex）= Co（19.5）使用40个没有几何约束的Gk DOF（只有CL是固定的），在采取一些优化下降步骤后不久，物理上无法实现的翼型结果，如图1所示。 翼型变得非常薄，第一个麻烦迹象是重新进入的后缘的外观。 当然，这种来自无约束优化的结果是完全可以预料的。 必须根据结构考虑施加适当的厚度限制，在这种情况下，后缘和前缘角度也必须明确限制。 经过多次优化尝试，发现以下约束适合此问题。 （19.6）CM约束也被发现是必要的，因为优化器倾向于强烈地使其更负，然后对人力飞机的机翼结构重量产生巨大不利影响。 指定的角度、厚度和CM与起始DAE-11 翼型的相同。 前缘角度6 = 180°约束只是与翼型鼻头顶部和底部表面之间的表面坡度相匹配，这是防止出现锋利的凿子型前缘所必需的。 Ci约束有效地消除了a作为自由度，而CM约束和四个几何约束中的每一个都删除了一个几何CL = 1.25 0TE \textasciitilde{} 6.25° (Vc)o.33c = 0.128 CM = -0.133 eLE = 180° (*/ c)o.90c = 0.014

翼型优化.."...... 我11.-...... Opt1 10 OOFs +,,, \_........ -..... “...，Opt1 20 ow5 -..... - Opt1 YO OOFs 图2 单点优化的翼型几何形状和目标函数与设计自由度数。自由度。 因此，自由设计参数自由值的有效总数为K-5。 通常需要使用100个几何自由度来生成一个几乎任意的翼型，尽管如果起始翼型与最终设计相当接近，则较少的数字可能就足够了，就像这种情况一样。 通过使用拉格朗日乘数增强目标函数 3，在LINDOP中明确施加了所有约束。 净效果是将设计空间的所有更改投射到可接受的约束子空间上。 图2显示了使用5、10、20和40净自由度的优化翼型 CD和几何形状。 翼型形状和CD看起来都是渐近的，表明40-DOF结果实际上是所提出的优化问题的真正解决方案，至少在局部最小值意义上是这样。 CD从原始DAEll 翼型中的0.0994减少到优化的翼型中的0.0836--大幅减少13\%。 然而，每个翼型的整个阻力极性的计算讲述了一个非常不同的故事。 随着自由度数量的增加，在越来越窄的CL范围内实现了阻力降低。 图3显示，极性曲线呈现出尖尖的形式，因此收益范围缩小到几乎为零。 因此，“随着自由度数量的增加，优化的机翼实际上在实际意义上变得相当糟糕。 回想起来，这种戏剧性的行为并不令人惊讶。 翼型设计--事实上大多数空气动力学设计--从根本上是由权衡驱动的，这几乎总是包括设计外的性能。 如果不考虑这种权衡，就像这个1点优化示例一样，非常差的结果几乎肯定会偶然发生--差翼的数量远远超过好机翼的数量！ 在这种情况下，优化器在表面上增加一个凸起来“填充”天位分离气泡，有效地减少了气泡发生过渡和重新连接时发生的混合和相关阻力惩罚[6]。 然而，气泡位置随着CJ的变化而变化，因此这种Uoi>timized“翼型形状仅对优化器看到的采样CL点有效。 对于较低的CL值，剪切

368 DREXA........ HSES V 3.0 2.0............................ 图3 klars for origina1 DRIZ-11 翼型 arld 1点优化机翼。lthyer不会在tlle bulnp上经历trt\textasciitilde{}11sitio11，而是在lasl\textasciitilde{}inar状态下分离，导致一个比碰撞发生的更大、更有损失的分离气泡。 对于更高的CI，值，过渡运行forwa.rd allead tlne bump更快t;ha\textasciitilde{}n，它将非常迅速地加速皮肤摩擦阻力的迅速上升，然后导致瓷砖下层边界层迅速增厚，并迅速加速失速。 真正的缺陷Inere不是tine优化技术，它给出了detnonst;ral)ly正确的答案，而是优化问题本身的公式ti011。 简单的1点阻力最小化，即使有许多真正的几何约束确定的llry相当大的试验和错误，仍然不符合翼型的真正设计要求。 不幸的是，这个缺点一开始并不明显，“为了弥补1点设计中的明显缺陷，2点optirstion通过替换目标fui\textasciitilde{}ctior\textasciitilde{}（19.5）来提供2点，以便预期操作Cb范围的两端现在sa'lnplecl。 两个操作点之间的Tlne 1:2 weigl\textasciitilde{}tilng有bce\textasciitilde{}i deterrr\textasciitilde{}ned是necwsa.ry，这样tlne上部：阻力极点不会被不太重要的下部过度连接。 Sanns geo\textasciitilde{}xletric

翼型优化.... 2'0 RlrFoil.... 2.0...... *..................................... - OPU 60 WF................. “...................... CL --.... MU 20 WF...... ---. DAE I1 ----- -- i,.::.: 4;: CL,-.-............................................................................................. （......，...... >...... 我............ 0......，...... >...... (..............,.................................. “。” '！ >我'...... “...(""........... 0.0.... 100 2004 -2 0 2 11 6 8 10 1Oq=Co a 原始DAEll 翼型和\%点优化机翼的图4极性。使用了约束（19.6），但当然两个操作点使用了两个单独的CL约束。 与1点情况相比，目标函数 IIOW随着自由度数的减少速度较短，需要60次自由度才能合理地渐近到真正的最佳解决方案。 图4显示了阻力极地随着自由系数的增加而演变。 显而易见的显著特征是，峰值局部优化极性形状持续存在，但现在这发生在两个采样的CL值上。 翼型几何体现在有两个凸起，每个凸起都位于采样操作点的气泡位置。 颠簸不像1点示例中那样明显，但它们仍然显示出远离采样CL值的显著阻力惩罚，如图4中的扇形阻力极点所示。 同样，“优化的翼型在实际意义上不如启动的DAEll 翼型，尽管其缺陷不如1点优化的翼型严重。 19.3.3 六点优化 进一步推进多点优化概念，定义了6点目标函数，以便采样的CL点0.8...1.6现在略高于翼型的预期1.0...1.5工作范围，以提供一些余量

图5六点优化的翼型几何形状和目标ft\textasciitilde{}nction与ntr\textasciitilde{}nber ttf tlesign POPS.结束。 添加了rnodes的tlie dcsig\textasciitilde{}i的进化现在甚至是Illore graclnal，alcl 90模式是recluireci到neasly asyiii\textasciitilde{}ptote to tlie最佳设计，如图5中的sbowli。 这个6点优化翼型的Tlle极点是cornpa.rcd，起始DAB-11极点如图6所示。 每个s;\textasciitilde{}.nlplcxi点周围的峰值beliavior仍然存在，但程度比以前少得多。 这款翼型的geol-netry更引人注目，然而，图7显示了1-point、Zpoint和6-point的geonletry；最佳化的翼型，在sanlpled操作点的分离气泡位置下降1、2和6 buxnps。 Tlie i\textasciitilde{}lviscid C，图8中6点翼型的分布显示了其六个凸起的严重程度。 另请注意，由于表面曲率轻微不规则，无黏性Cl在大多数翼型上也“振荡”。 这在许多几何自由度问题中都很常见，因为这种smdl-scde不规则性几乎没有空气动力学惩罚，因此tlie优化器是看不见的。 在1点、Z点、a和6点优化exalnples J\textasciitilde{}ove中，optimiaer在最同步的物理scde处操纵几何，这对目标函数有显著影响\textasciitilde{}\textasciitilde{}--在这种情况下是过渡分离气泡。 如果预置了witli suffit:ie\textasciitilde{}it几何DOPs，那么优化器执行这种ma.\textasciitilde{}iigulation的程度是令人震惊的。 带有六个明显颠簸的6点翼型肯定不是预期的，但t11e结果\textasciitilde{}nalclres完美的感觉在ret;rosl>ect。 图9所示的黏性GI分布在CI处的sam\textasciitilde{}>led poiilts之一，= 1.5表明l;he larninar区域的bu\textasciitilde{}nps被浸没在自由剪切1:\textasciitilde{}yer下。 Ea.cii buml>从表面突出\&刚好够ba.rely到达剪切层，并定位在表面，以便它“ca.tches”s11e.d\textasciitilde{}层，因为它不进行过渡，并使其重新连接，而没有明显的阻力--在剪切层下产生停滞流体的混合。 计算出的速度曲线

翼型优化loq=C，原始DAEll 翼型和6点优化机翼的图6极。在凸起处，图10清楚地显示了这一点。 在采样点之间的Cz值中间，没有凸起来抓住剪切层，剪切层必须通过混合山谷中的层状气泡来重新连接，这会带来阻力惩罚。 这是图6中阻力极地的iiscallopingn的原因。 当然，即使是最有经验的空气动力学家也无法预见这种相当愚蠢的多点最佳解决方案，并说明了一个拥有大量设计自由度的优化器所表现出的几乎狡猾的聪明才智。 如图6所示，鉴于CL范围内更密集的操作点采样可以消除局部颠簸，优化似乎比极地下部的基线DAEll 翼型略有改善。 远离颠簸，图7所示的机翼明显收敛成独特的形状。 然而，这种形状的几何细节与人力飞机中使用的结构技术并不特别兼容。 高度弧形的后缘、接近x/c = 0.7的上表面的明显凹面以及底面拐点都很难实现。 由于一开始就没有办法预测这些特定特征的出现，因此可能必须构建并解决具有适当约束的新优化问题。 结论是，在工程环境中，对翼型设计进行优化仍然是一种迭代的切割和尝试。 但与传统的逆向技术相比，切割和尝试不是在几何上，而是在优化问题的精确公式上。

372 DRELA 图7 优化机翼的几何形状，显示表面凸起。 图8 6点优化翼型的Inviscid C p分布。 19.4 Transonic 翼型应用 在本例中，从参考[7]的基线案例13a开始，优化了众所周知的RAE-2822 翼型。 这种情况部分进入阻力上升马赫范围，但不足以冲击引起的分离。 因此，它处于低效但并非不切实际的飞行状态，在优化重新设计后可能会有显著改善。 优化计算都是在相同的CL = 0.733和Re = 2.7 x 10 6下进行的，就像开始的RAE-2822情况一样，相同的M = 0.74用于1点优化的情况。

翼型优化373图9抽样点C L = 1.5的6点优化翼型的黏性C p图10 CL = 1.5的6点优化翼型的速度曲线，显示第二个凸起“捕捉”重新连接剪切层。 19.4.1 单点优化 降低阻力的目标由以下合理的优化问题体现，再次使用正弦模式系数Gk和攻角a作为设计参数。 前缘和后缘角度以及局部厚度作为合理的几何约束被强加。最小化F\{Gk,a) = C D (19.9).... CL = 0.733 6TE - 8.62°受\{f/c)o\textasciicircum{}\textasciicircum{} = 0 1206 \textasciicircum{} = m ° (19.10) 图11显示了与实验数据相比的基线分析解，以及1点优化的翼型。 虽然减少了40\%的阻力，并且保持了局部厚度，但翼型不适合喷气式运输应用，因为平均厚度超过了典型的

374 DRELA -2 0 arr "3 0 OPtl t/c HaM - D.7W Re - 2.100.10" -1.5 RIT1 - 1.810 CL - 0.1330 CP CO - 0.01014 I 0 CH - --O.liZB LID - 68.23 -0 5 0.0 0.5 图11 基线RAE-2822 翼型的Cp分布，以及1点优化的翼型，局部厚度限制在.u/c = 0.35。宽的spar box大幅减少。 因此，局部厚度约束并不代表真正的结构要求，尽管一开始并不明显。 代替（19.10）中的局部tlc约束，r.ni.s.应变E,,,（每单位弯曲力矩和单位材料模量）被施加到Inore，以现实地解释宽梁箱的要求。 积分围绕翼型周长进行计算，并由局部slcin thiclcness t加权。 局部弯曲相关单元伸展应变e由局部翼型表面点y和局部弯曲惯性和弯曲重心位置yo来定义。 皮肤厚度t仅在spar箱0.15 < z/c 5 0.65的弦范围上选择为非零。 实际上，r.m.s.应变约束在spar箱上强制执行大致平均翼型厚度，而翼型的remdinder在结构上无关紧要。 图12显示了使用净40自由度优化的翼型。 与局部厚度约束相反，r.rn.s.应变约束现在在spar箱的范围内强制执行合理的厚度，但在spar前面会产生严重的底面凹面，在后面会产生较小的凹面凹。 虽然它增加了空气动力学上有利的底部载荷，但从实际意义上讲，前凹面可能不可行，

翼型优化-2.0 a。 选择Ib.,I I(d. O.NO Re。 2.7W.ld -1.5 Rlh - 2. Y28 CL - 0.73\% c，CO - 0.01117 -1.0 C)I. -0.152 -0.5 0.0 0.5 1.0 RRE 2822 -OPT Ib图12 C，1点优化翼型的分布，r.m.s.应变约束超过0.15 5 x/e < 0.60，以及优化翼型与基线RAE2822 翼型相比。因此提出的优化问题可能仍然无法体现设计问题的实际要求。 同样，这个缺点一开始并不明显。 无论如何，优化的翼型都显示出完全等熹流动，没有激波和波阻。 净CD已从0.0178降至0.0112，下降了37\%。 然而，这部分是虚幻的。 图13中的整个马赫数扫描表明，阻力降低主要在M = 0.74的取样操作点附近实现，尽管这里的区域性最佳化程度不像低雷诺数 翼型案例那样极端。 阻力在较低的马赫数下会大大增加，这种行为在无冲击超音速机翼中是众所周知的。 导致这种行为的几何特征是以x/c = 0.50为中心的轻微凸起，它在超音速区的后部产生逐渐的再压缩，从而消除了冲击。 与颠簸相关的是一个扁平的前上表面，在较低的马赫数下产生强烈的膨胀和强烈的冲击，从而增加了阻力。 在高于采样M = 0.74的马赫数时，流量通过颠簸强烈膨胀，导致强烈的冲击和快速的阻力上升。 本例中优化器的行为在某些方面与低雷诺数优化情况下的行为相似。 优化器利用最小的几何尺度的流动特征，该特征可以通过可用的几何扰动模式来解决。 在这种情况下，被操纵的物理尺度是冲击/边界层相互作用区域。

376 DRELA图13马赫扫荡，用于1点优化和基线RAE-2822机翼。 19.4.2 两点优化 为了减少1点设计的点优化性质，通过替换目标函数（19.9）来定义2点优化，以便现在对一些合理的马赫数范围进行采样。 使用相同的CL和r.m.s.应变，以及几何角度约束。 结果翼型如图14所示，整个马赫扫描如图15所示。 一些“本地化”优化是显而易见的。 翼型上现在有两个预压缩凸起--一个在两个采样的马赫数字中每个位置的shoclc脚位置。 高于M = 0.74采样点没有改善。 19.4.3 四点优化 为了改善2点设计的本地化性质，4点优化问题定义如下。 1 1 F（Gk，a）7 CDIM=\textasciitilde{}。 \textasciitilde{}\textasciitilde{} + 5 CDIM=\textasciitilde{}。 \textasciitilde{}\textasciitilde{}现在使用了Mac11数字范围内的更密集的采样，在范围的上部应用了更大的权重，以反映其更大的重要性。 另请注意，现在抽样了更高的M = 0.76 Mach，以降低之前取样的最高M = 0.74点的阻力。

翼型 优化 -2.0 m5 OPT 2a Y 3.0 CL - 0.7330 Re - 2.70W.IOo -1.5 \textasciitilde{}\textasciitilde{}JLL\textasciitilde{} 0.6W 2.600 0.01069 -0.0- 0.1UO 2.063 0.01118 -0.106 -1.0 -0.5 0.0 0.5 RRE 2822 - OPT 2a 图14 C，两个采样马赫数下2点优化翼型的分布，以及优化几何形状与基线RAE2822 翼型的比较。 除了前底面凹陷、x/c > 0.9的相当薄的后缘和更负的CM外，图17中的结果表明，这种4点优化的翼型似乎比基线RAE2822有吸引力的改进。 阻力发散马赫数的增加是最重要的。 虽然图16中的翼型轮廓在上表面似乎相当光滑，但有限数量的采样点仍然导致四个采样点激波位置中的每一个都出现明显的凸起。 这生动地显示在图18所示的马赫波图上，一个膨胀风扇从三个前凸处发出，这些凸起完全包含在M = 0.76采样的第四个操作点的超音速区域中。 优化器使用的这些冲击管理颠簸显然比低雷诺数 DAE-11示例中的重新连接管理颠簸要轻得多。 尽管如此，它们显然是有限数量的采样操作点的人工制品，在马赫数以几乎连续的方式采样的理想情况下不会存在。 19.5 结论 所展示的例子清楚地说明了许多陷阱，这些陷阱很容易出现在看似简单的翼型优化问题中。 可以通过增加采样操作点的数量来抑制低雷诺数 翼型和跨声速翼型示例中小规模几何不规则的出现，以便在中间位置出现额外的凸起，最终全部融合成光滑的表面。 这个论点支持了这样的假设，即对于一个

图98 2点optil\textasciitilde{}iizecl和基线RAE-2822机翼的马赫扫描。平滑几何形状，有必要有\#操作点采样=8（\#dcsign参数），因为设计参数数量的增加会减少优化器可以利用流量的l\textasciitilde{}lgth尺度，然后必须与采样点数量的比例增加相匹配，以控制这种利用。 当然，可以施加几何约束，如表面曲率，来控制这种利用，但这实际上是自由几何设计参数的有效数量的减少。 由于成本随着采样操作点的数量而线性增加，因此使用大量设计自由度进行优化的能力是一个非常昂贵的命题。 它强调了将设计自由度的有效数量减少到绝对最低限度的必要性。 这些结论不仅对2-D 翼型优化有直接影响，而且对具有多种设计自由度的大规模3D优化问题也有直接影响。 没有理由希望3-D问题能免受这种困难的影响。 根据目前的优化示例和作者之前通过优化进行翼型设计的经验，提供了以下观察。 19.5.1 翼型 Optkization - 缺点 r 如果没有密切相关问题的广泛经验基础，目标函数和有效体现实际设计问题的约束从一开始就无法实现。e 如果呈现足够的设计模式分辨率，优化器将很容易（并且令人讨厌地）操纵和利用如何在最小显著

翼型 OPTIMIZATION OPT Ua RRE 2822 - OPT la 图 16 Cpoint优化翼型在四个采样马赫数下的Cp分布，以及优化几何形状与基线RAE2822 翼型的比较。物理刻度存在。 这些例子在过渡性分离泡沫的尺度和冲击/边界层相互作用区的尺度上表现出这种操纵。 最小规模的操纵往往只有在采样操作条件附近才能产生性能的提高。 点优化的翼型经常显示出设计外性能可能严重下降。 多点优化可以阻止在最小尺度上操纵。 因此，分离气泡和减震器等小型移动流功能似乎被最佳化器「抹上」，而且被利用得不那么严重。 增加几何自由度的净数似乎需要相应地增加目标函数采样的操作点数量。 在具有大量自由度的一般翼型设计问题的理论极限中，可能需要对操作空间进行近乎连续的采样--这是一个非常昂贵的命题。 在多点优化中实际采样的最合适的操作点并不明显。 根据有限的经验，稍微超出预期操作范围的取样似乎是最好的。 多点优化中使用的点权重是任意的，如果没有事先的经验，很难估计其适当的值。 优化的空气动力学形状通常是“有异的”，通常需要事后平滑。 19.5.2 翼型 优化-当传统逆向技术的直觉耗尽时，优化可以快速指示可能的改进方向，或者当

翼型优化381图18马赫波用于M = 0.76的4点优化翼型，在较低的马赫数条件下，在激波位置有指示的颠簸。 6. M。 德雷拉。 麻省理工学院Daedalus原型的Low-雷诺数 翼型设计：案例研究。 飞机杂志，25（8）：724-732，1988年8月。 7. P.H. 库克，医学博士 麦当劳和M.C.P. 坚定。 气翼RAE 2822压力分布以及边界层和唤醒测量。 在计算机程序评估的实验数据库中，AR-138。 阿加德，1979年。

\clearpage
\chapter{迈向工业实力的Navier-Stokes守则--重温}
迈向工业实力 Navier-Stokes代码-重访 Wen-Huei Jou 1 20.1 导言 我们在1992年撰写的论文[1]评估了Navier-Stokes技术。 该论文的重点是Navier-Stokes代码的准确性。 从那时起，CFD技术得到了改进，应用范围也扩大了，因此重新审视这些主题是合适的。 根据我们的经验和预计的应用程序，我们可以帮助识别算法问题，以便进一步研究。 我们经常听到CFD的研究人员吹捧他们的代码被行业使用。 至少在商业飞机公司，似乎需要澄清CFD代码的应用。 在工业中使用CFD代码有两种方式。 两种模式中的一种的首选特征截然不同。 当然，代码最重要的用途是在飞机的开发过程中。 在这里，在应用之前，必须根据用户的舒适程度广泛验证工程任务所需的代码的准确性。 在很短的时间内扭转结果的能力也很重要。 在激烈的战斗中，满足时间表是非常重要的优先事项。 需要数百小时来准备输入和/或在数百个CPU小时内运行以获得一次分析的代码，在开发飞机的过程中没有地位。 CFD代码的第二次使用是在研究中。 对于1波音商用飞机集团，P。 哦。 Box 3707，MS 67-LL，西雅图，华盛顿州98124-2207。 Frontiers of 计算流体动力学 - 1998 编辑：David A. Caughey和Mohamed M. 哈菲兹 ©1998 世界科学

384 JOU 整理一些复杂的空气动力学问题的因果关系，CFD可以在“和平时期”期间以研究模式使用，以获得可以在飞机开发过程中间接应用的知识。 代码解决感兴趣的空气动力学现象的能力可能是该应用程序的首要因素。 就时间表而言，要求较低。 然而，为了解决技术问题，可能需要调查许多案例，因此计算成本和周转时间仍然是此模式下使用的程序的重要问题。 最后，也有业内人士会从研究人员那里拿一个代码来“玩”。 这些努力甚至不应该被视为工业应用。 本文探讨了CFD代码的工业使用的两个方面，并更加关注第一种模式。 我们审查了空气动力学工程几个关键领域的应用状态，并尝试确定代码和算法改进的机会。 由于真正有效的CFD工具几乎总是三维代码，我们没有花太多精力审查二维代码，尽管我们认识到在二维中工作通常有助于解决算法问题。 20.2 巡航设计 对于商用飞机的开发，谨慎的巡航配置设计是迄今为止最重要的空气动力学任务。 它的成功是满足有效载荷和航程要求、提供运营经济性和支持航空公司性能保障的关键。 对于亚音速飞机来说，机翼的形状主要由跨声速巡航条件决定，而高升力机翼的形状或多或少必须符合巡航形状给出的几何约束。 巡航机翼配置分两个步骤确定：平面形式定义和随后在平面形式定义施加的限制下进行详细的空气动力学塑造。 第一步涉及非常复杂的多学科权衡和系统考虑，第二步本质上是空气动力学功能。 目前，CFD代码主要用于第二步。 由于其重要性，过去二十年来，跨声速巡航翼的设计和分析一直是CFD开发的重点。 流程主要是附著的。 因此，势流近似与适当的边界层和Navier-Stokes方法同样适用，并在波音和许多其他商用飞机公司中实施。 定义机翼形状以实现设计目标的能力是CFD的独特能力，任何风洞测试方法都无法取代。 因此，它是空气动力学最抢手的能力

工业强度NAVIER-STOKES代码385商业飞机公司的工程社区。 主要基于压力匹配方法开发了几代设计工具[2]。 最近，工业CFD随着使用TRANAIR代码[3]的多点、受限优化的成功达到了顶峰。 这种新功能允许在支柱、机舱和整流罩等其他飞机部件的情况下进行机翼设计。 通过多点优化，新方法同时考虑了对飞机运行最重要的几个设计条件。 它结合了结构和制造约束作为优化问题的一部分，而不是压力匹配方法中的后处理步骤。 当使用压力匹配方法时，在多个设计点之间的手动迭代和设计几何体的后处理，通常会导致性能降低和实现设计闭合的流动时间长。 当配置很复杂，而先进技术翼型的简单扫描近似值不能近似三维流量时，情况尤其如此。 该方法还用于优化超音速运输飞机的巡航翼和机身的形状，从而大大提高了性能。 毫不夸张地说，优化方法是超音速运输机发展的关键赋能技术之一。 该技术是新的，对如何利用这种方法的探索才刚刚开始，但它已经显示出缩短设计周期时间和提高产品性能的潜力。 多点、受限优化正在成为设计巡航配置的首选方法。 使用Navier-Stokes方法的机翼设计实践仅限于使用DISC方法的压力匹配方法[4]。 基于詹姆森科技（J-T）的TLNS求解器[5,6]集成了内部网格生成代码[7]，在产品开发阶段被波音公司广泛使用，因为它使用我们内部的Triton计算机快速周转。 代码的多块版本也经常用于测试前分析[8]。 在过去的两年里，使用过集网格策略的OVERFLOW代码[9]已经发展成为一种有效的工具。 这两个代码用于研究环境，以解决CFD的准确性问题。 无论是在分析模式还是设计模式下，CFD巡航空气动力学代码最理想的特征是其预测阻力的能力。 出于工程目的，预测因几何变化或雷诺数等流参数变化而导致的阻力增量是一项非常有用的能力。 绝对阻力数的预测将受到高度重视，但更难实现。 为了初步验证此能力，最好使用一个简单的几何形状，可以生成密集网格，以了解它们对阻力等敏感量的影响。 因此，选择了一个简单的跨声速翼体组合

386 JOU这个目的。 图1显示了我们用来调查TLNS3D阻力预测准确性的网格配置。 网格在机翼体配置周围具有C-H拓扑结构。 在这个网格配置的约束下，生成了一个297 x 97 x 97的密集网格。 如果我们通过减去25个阻力计数来转移阻力极性，计算出的极性形状与测试数据非常一致。 这种转变非常大，无法用任何可能的实验误差来解释。 计算中一个可能的误差来源是由于网格拓扑结构，机身前部附近的分辨率不足。 在前连接点附近，流量变化迅速。 C型网格为机翼提供了良好的分辨率，但可能无法在机身机头附近提供足够的分辨率。 重叠网格技术为C-grid拓扑的网格生成图1机翼体提供了灵活性

工业强度NAVIER-STOKES代码387，使我们能够进一步调查这个问题。 图中所示的重叠网格。 生成了2个，并得到了一个OVERFLOW溶液。 这种网格配置覆盖了比C-H网格密度更高的附着区域，而机翼上仅提供略高的网格密度。 测试数据的阻力极移减少到大约10个计数，仍然超出了可能的实验误差范围。 将网格进一步细化50，将阻力差异减少到8.3计数。 由于网格变化，阻力变化的最大来源是机身上的压力阻力。 需要进一步研究，以了解是否可以通过提高数值准确性来进一步缩小这一差距。 这个例子告诉我们，我们没有先验知识，不知道在哪里集中电网来获取某些工程信息。 很快将测试数据和计算结果之间的差异归咎于风洞测试、湍流建模以及数值精度以外的因素往往太诱人了。 具有适当自适应标准的自动网格自适应功能肯定会提供更一致的结果。 原则上，我们需要准确的阻力预测才能使用CFD代码进行性能优化。 然而，如果误差是系统性的，并且是正在进行机翼形状详细设计的配置的固定值，那么只有预测形状变化导致的增量才重要。 设计决策通常通过检查basehne的增量来做出。 到目前为止，我们的经验表明，由于配置更改，CFD对阻力增加的预测似乎相当准确。 这如图所示。 3，它显示了两个机翼体配置之间的计算和测量增量。 该协议足以为工程师做出设计决策提供基础。 为了预测飞机的绝对性能水平，包括支柱、发动机和所有尾部，需要一个大的网格尺寸。 据估计，分析襟翼向上配置可能需要1500万到2000万个网格点。 有了如此大的网格，需要快速收敛方案才能构建在产品开发环境中使用的有效代码。 Navier-Stokes代码可以准确地预测机翼表面的压力分布。 图4显示了多块Navier-Stokes计算的翼体支柱机组配置的结果，在16个块中有510万个网格点。 CRAY C-90的执行时间为15 CPU小时，以获得具有足够工程精度的压力分布。 该代码还用于在形状设计期间监控边界层的运行状况，以防止边界层过早分离。 在这里，根据给定的雷诺数，估计了边界层内在飞机表面各个位置的网格间距分布。 如果能够开发出可靠的解决方案自适应方法，计算结果的置信度将大大提高。

388 JOU（a）（b）图2 重叠网格用于机翼体

工业强度NAVIER-STOKES代码389图3两个机翼体配置之间的阻力增量.65 20.3高升机翼分析和设计高升机翼配置的定义可能是仅次于巡航设计的第二大空气动力学任务。 今天的三维Navier-Stokes代码尚未应用于飞机开发过程中的这项任务。 甚至它们在研究环境中的使用也非常少。 CFD技术根本还没有准备好。 与巡航翼设计一样，高升翼设计过程涉及两个步骤：平面配置的定义以及前缘和后缘设备的细节塑造。 在许多情况下，第一步涉及离散配置更改，例如不同的前缘和后缘设备。 基于配置连续变化的优化方法不适用于这种情况。 因此，我们设想CFD在高升机翼的初始应用是评估不同配置的相对优点。 这将需要一个可以非常快速地周转的代码。 对于流程的第二步，需要有限的设计能力。 由于巡航配置的重要性，高升翼几何形状受到高度限制。 在这个阶段，只允许更改机翼表面的几个地方和几个几何参数。 此外，襟翼支撑机制的运动学和其他因素进一步限制了机翼的形状。 设计代码必须处理这些限制。 对于高升翼来说，估计在3500万到5000万之间

390 JOU 图4 可能需要机翼-车身-支柱网格点上的压力分布来获得具有工程价值的解决方案。 根据这个数字以及JT-TLNS和OVERFLOW代码的当前性能，在Cray C-90计算机上需要大约400个CPU小时才能完成一项分析。 例如，对于新飞机是使用板条还是克鲁格前缘设备的配置决定，至少需要32次运行。 这相当于Cray C-90的大约13,000小时。 此外，任务需要在一个月内完成，以便及时做出配置决定。 人们还必须认识到，这只是要做出的众多计划决定之一。 因此，从吞吐量的角度来看，今天的CFD代码完全不足以执行预期功能。 当然，计算速度和收敛加速度并不是CFD应用到高升翼的唯一问题。 流量分离始终存在于高升翼中，需要建模。 不幸的是，在连续的不利压力梯度下，来自光滑表面的流动分离对流动条件的微小变化非常敏感。 这种灵敏度不仅给CFD带来了问题，也给飞行中的车辆的风洞建模带来了问题。 在雷诺平均Navier-Stokes（RANS）公式的背景下对流动物理学进行建模的主题本身可能值得一篇完整的论文。 我们的论点是，除非我们有一个能以数值方式很好地解决流程的代码，否则任何关于任何RANS模型缺陷的结论充其量都是暂定的。 例如，一个方程湍流模型，如Spalart-Allmaxas模型[10]，是否可以充分预测高升力的三维分离流动

在我们确定数值不准确不会污染结果之前，无法解决工业强度NAVIER-STOKES代码391配置。 像高升翼一样复杂的几何形状总是对网格生成工艺构成挑战。 使用当前的结构化网格生成系统，大约需要200天才能生成超集或多块网格。 对于类似配置的变体，网格生成时间可以缩短到几周。 这个网格生成问题的可能解决方案可能是使用非结构化网格算法，这对开发自动网格生成过程的问题更少。 20.4 预测处理质量和控制有效性 空气动力学稳定性学科涉及飞机开发过程中工程的两个方面。 一个是预测飞机的操控特性，这样就可以定义一个好的配置。 另一个是提供一个数据库，规模约为20万个负载、飞行控制和飞行模拟器。 如果打算使用CFD代码生成所需的数据库，代码的吞吐量必须能够在几个月内处理大量案例。 当前的代码和CFD技术不符合这些要求。 很少有CFD的成熟到足以在飞机开发过程的这一阶段使用。 在研究环境中，只有代码预测一些关键特征的能力被探索。 在这里，飞机俯仰特性的预测作为一个研究例子。 图5显示了机翼-机身组合的俯仰特征。 带有Spalart-Allmaras 湍流模型的Navier-Stokes代码似乎能够合理地预测电击诱导的流动分离。 下一个问题是，是否可以准确预测整个飞机（包括尾翼）的跨声速音高。 图6显示了与风洞测量相比，尾部位置的预测向下冲刷是机翼-机身组合的攻角的函数。 对于这种计算，没有解决方案自适应能力来捕捉机翼的尾迹，机翼后面的粗网格细胞在很大程度上扩散了尾迹。 似乎唤醒系统的纵向成分的循环通过数值方案得到了守恒，并且以惊人的准确性预测了下冲。 因此，我们似乎有很好的机会使用CFD来预测整个飞机的俯仰。 随后的问题是，我们是否可以预测涡流发生器对音高特性的影响。 这在短期内是雄心勃勃的，但在飞机开发过程中需要减少涡流发生器配置的测试。 开发的最佳技术方法

392 JOU 这样的能力仍然不清楚。 如果要避免解决每个涡流发生器，则必须开发一个包含涡流发生器的跨度平均效应的“湍流模型”。 这不是一件容易的事，甚至可能是不可能的。 如果我们决定解决每个涡流发生器，它将大大增加所需的网格点数量，并且肯定会需要网格自适应能力来捕获自由涡流。 即使这是可能的，我们仍然需要一个适用于浸入边界层的涡流的湍流模型。 一旦一个人清除了CFD的下栏，标准就会立即提高！ 20.5 结论 CFD用于巡航应用的高度发达。 相比之下，应用于分离流的Navier-Stokes代码正处于起步阶段。 对于工程应用来说，算法研究的两个方面将非常有价值。 这些是收敛加速和复杂几何形状的网格自适应能力。 对当前结构化网格技术的评估表明，在块状结构或覆盖方法的背景下，网格自适应非常难以实现。 看来，非结构化网格方法是实现这些新算法的更好工具。 当前CFD代码的光谱半径大于0.99。 有了这个图5翼体的俯仰特征

工业强度NAVIER-STOKES代码393 5 c 图6在收敛的水平尾部位置速率下向下清洗的机翼体的预测，需要大约400次迭代周期才能使残留物减少四阶。 多网格和其他迭代方法的最新进展表明，有大幅改进的潜力。 可以合理地预期，算法的改进将在未来5年内将CPU小时减少5倍。 我们的预测是，计算机硬件的进步将使我们的吞吐量每5年增加5倍，同时保持成本不变。 这是一个合理的期望。 有了这个预测，未来十年的吞吐量将增加125倍。 然后，高升力机翼分析将需要不到3小时才能计算，并可以成为产品开发环境中的有效工具。 为了确保分离流动的Navier-Stokes解决方案的工程准确性，需要网格自适应到压力梯度，以及壁边界和自由剪切层。 对于高雷诺数，这意味着定向自适应算法。 复杂三维流的适应标准构成了算法研究的一个重要主题。 在实现网格自适应、曲面几何描述、网格生成和求解器可以非常紧密地连接，并且需要在代码开发中尽早采用集成方法。

394 JOU 致谢 本文借鉴了波音商用飞机集团研究团队的工作，特别是博士的工作。 Steve Allmaras、John Bussoletti、Hoa Cao、Dinesh Naik、Philippe Spalart、Bernard Su、Larry Wigton、Jong Yu和Mr. T。 J。 考。 作者特别要为博士们。 Jong Yu和Venkat Venkatatrishnan，以及Dinesh Naik提供了数字，审查了草案，并提出了有价值的建议。 参考文献 Wen-Huei, Jou, Laurence B. 威格顿，史蒂文·R。 Allmaras、Philippe Spalart和Jong Yu，《走向工业强度的航海家-斯托克斯代码》，在《空气动力学流动的数值和物理方面》V中，ed。 T。 Cebeci，Springer-Verlag，1992年1月。 金锤，M。 我。 和Steinle，F. W.，使用CFD和非常高的雷诺数风洞测试的先进跨声速机翼的设计和验证，ICAS会议记录，1990年，pp. 1028-1042。 Huffman，R.P.，Melvin，B.G.，Young，D.P.，Johnson，F.T.，Bussoletti，M.B.，Bieterman，M.B. \& Hilmes，C.L.，计算流体动力学中的实用设计和优化，AIAA论文93-3111，1993年7月。 于，N。 J。 和坎贝尔·R。 L.，使用Navier-Stokes代码的Transonic 翼型和机翼设计，AIAA论文92-2651，1992年。 瓦萨，V。 N。 和Wedan，B。 W.，为3-D Navier-Stokes 方程开发高效多网格代码，AIAA论文89-1791，1989年。 瓦萨，V.N.，萨内特里克，M. D。 \& Parlette，E.B.，开发灵活高效的基于多网格的多块流动求解器，AIAA论文93-0677，1993年。 Kao，T.J.，Su，T.Y. \& Yu，新泽西州，Navier-Stokes用机芯和支柱进行运输机身配置的计算，AIAA论文93-2945，1993年。 宇，新泽西州，苏，T.Y. \& Wilkinson，W.M.，使用Navier-Stokes求解器进行复杂配置分析的多块网格生成流程，AIAA论文96-1995，1996年。 Buning，P.G.等人，OVERFLOW用户手册，版本1.6ap，美国宇航局艾姆斯研究中心，莫菲特菲尔德，加利福尼亚州，1994年斯帕拉特，P。 R。 \& Allmaras，S。 R.，空气动力学流动的单方程湍流模型，La Recherche Aerospatiale，No.l，1994年，第5-21页。

\clearpage
\chapter{我们从计算流体动力学火车空气动力学研究中学到了什么？}
我们从计算流体动力学火车空气动力学研究中学到了什么？ Kozo Fujii1和Takanobu Ogawa 2 21.1 介绍 在过去三十年里，航空航天飞行器的应用加速了计算流体动力学技术。 现在，这些技术正在扩展到其他应用领域。 其中一个应用是训练空气动力学。 近年来，火车的速度迅速加快。 现在，一些高速列车以300公里/小时的速度运行，大约是马赫数 0.25。 在这个马赫数范围内，可压缩性效应可能不是太重要，但当火车进入隧道时，它变得至关重要。 图1显示了物理1空间和航天科学研究所的示意图，相模原，神奈川县，229，日本2清水公司，越中岛3-4-17，江户区，135，日本计算流体动力学-1998年。 编辑：David A. Caughey \& Mol Fig 1 轰动噪声生成机制

396 火车进入隧道时发生的富士和小川现象。 当火车进入隧道时，就像活塞一样，它会增加隧道内的压力。 这种增加的压力在火车前面产生压缩波，并传播到隧道出口。 由于压缩波的非线性效应，该波的梯度随着传播而变得越来越陡峭。 因此，在隧道出口附近产生了快速的压力变化。 如果隧道足够长，可以建立激波。 当这个压缩波从出口出去时，它会变成一个脉冲状波。 它被称为“轰轰烈烈的噪音”，被认为是与高速列车发展相关的重要社会问题之一。 这种“轰动噪音”是一种低频噪音，它强烈地震动了隧道出口附近房屋的窗玻璃。 小泽[10]澄清说，轰鸣噪音的强度与火车在隧道出口产生的压缩波的压力梯度成正比。 因此，通过减缓隧道入口处的压力增加，可以有效地缓解轰鸣噪音。 为了了解入口处的压力是如何增加的，并找出如何减少梯度，有必要调查火车入口产生的隧道内的瞬态流场。 过去，这种流场主要通过实际列车的现场测量来调查。 没有讨论瞬态流场的细节，因为大多数场测量只给了我们有限的信息。 关于实验室实验，将轴对称物体发射到管道中，并研究了管道内的压力增加。 这些类型的实验隐含地假设管道内的压力增加只取决于横截面面积的变化，三维效应并不重要。 使用最近的CFD技术，可以非常容易地进行流场的调查。 使用移动重叠网格方法，改变火车几何形状是一项简单的任务，可以获得三维瞬态流场的详细数据。 在过去的几年里，我们一直在使用这项技术调查流场[5-9]。 在本文中，我们想总结一下我们所做的，并尝试展示火车进入隧道的空气动力学的基本重要特征。 基于流量模拟，开发了一个简单的理论来估计隧道内产生的压缩波。 开发过程可能是如何使用CFD技术分析工程问题的一个很好的例子。 21.2 数值方法 21.2.1 基本方程 基本方程是用广义坐标系书写的三维可压缩Navier-Stokes 方程。

CFD用于火车空气动力学dtQ+d4E+dnF+d（G=Rs-1 d？ Gv（21.1）其中E=- J pU puU+£xp pvU+Zyp pwU+£zp (e + p)U-£ tp F=- J pV puV+r]xp pvV+r]yp pwV + T]zp (e + p)V-r] tp G=- J pW puW+£P pvW+pwW+Czp (e + p)W-\{ lP imxul-+\{3l p)nv£xW£xW£(3/ti)m2£z IJm3+(3/ H)m2m2 xu+yv+\textasciicircum{}yv+\textasciicircum{}zw) (21.2 where and 和V + \textasciicircum{}yV + \textasciicircum{}W V=t]t+Tixu+rizw mi=C 网格的运动表示为术语£„ T] t和£ t。 在一些计算中，轴对称Euler 方程代替三维Navier-Stokes 方程。 基本方程是相似的，除了一些项被简化，并且不包含黏性项。 这些方程是及时解决的，但基本方程被修改，以允许

398 富士和小川在每个网格区之间进行信息交流。 本研究中使用的接口方法称为“强化解算法[2]，下一节将对其进行简要解释。 21.2.2 强化解决方案算法-区域接口方法 流场是“\textasciicircum{}\textasciicircum{}\textasciicircum{}使用重叠网格方法将入口分解为几个区域。 虽然每个案例的详细区域结构都略有不同，但图中显示了三维计算的典型网格结构。 2. 火车有自己的网格，可以以火车的速度移动。 它被称为“火车区”。 对于几何处理，为列车下方的流场准备了另一个运动网格。 这被称为“底区”。 隧道有自己的网格，入口区域有自己的网格。 由于隧道入口的角落有一个拓扑奇点，所以准备了一个修补的网格来避免奇点。 它被称为“领区”。 准备了一个覆盖火车网格的中间网格区域，以节省从一个区域传输到另一个区域的数据插值所需的计算机时间。 这个网格区域具有x=常。网格部分，并与火车网格一起移动。 火车网格和这个中间网格之间的插值系数不会随着时间的推移而变化，因为两个网格以相同的速度移动。 由于隧道网格和中间网格都有x=常量网格部分，中间网格和隧道网格之间的插值本质上是二维的，使用这个中间网格区域可以节省计算机搜索和插值的时间。 这个中间网格区域也可以增强局部网格分辨率。 为了在每个区域之间传输信息，Eq. （21.1）被修改为包括源术语，dTQ+d\textasciicircum{}E+dnF+d!。 G=R\textasciicircum{}dl-Gv+\textbackslash\{\}x\textbackslash\{\}(Qf-Q) (21.4) 右侧的第二个项是强制项。 当绝对值X被设置为足够大时，解被强制为Qj。 当x的绝对值为零>我求解普通Navier-Stokes 方程时。 强化解决方案算法（FSA）方法使用x函数作为标志来设置图2流场和区域网格结构

CFD用于火车空气动力学399区域，解决方案应强制（强化）为给定的值。 在当前的区域计算中，具有更高优先级的网格区获得的解为g/，\%值在该区域被设置为足够大。 FSA区域方法的详细信息可以在参考中找到。 7. 计算网格可以进入身体几何。 由于身体内网格点的物理值没有任何意义，因此它们不应该用于计算。 在当前的FSA方法中，\%被设置为身体内的负大，程序会自动丢弃不应使用的无意义数据。 在本研究中应用覆盖区域运动网格系统有几个优点。 显然，几何复杂性得到了缓解，结构化网格的解码可以很容易地应用于复杂的流程配置。 当前类型的问题经常需要参数化研究。 在覆盖区域方法中，当考虑不同的列车配置时，只需要更改列车网格。 通过将精细网格区域放在适当的区域，网格点总数的增加较少，可以轻松增强局部网格分辨率。 21.2.3离散化LU-ADI算法[4]用于时间积分，对流项由Roe通量差分割离散化[11]。 使用原始变量的MUSCL插值[13]实现了高阶精度。 黏性术语被中心差分离散化。 Baldwin \& Lomax [1]的代数湍流模型被使用，因为雷诺数在107位，预计不会有强烈的分离。 21.2.4 边界和初始条件 由于很难模拟整个隧道区域，并且已经表明轰鸣噪音的强度取决于入口附近压缩波的坡度，因此在压缩波完全建立的地方切割隧道。 在这个边界处，施加了一维非反射边界条件，以避免非物理波反射。 为了避免火车加速造成的初始干扰，火车最初应放置在离隧道入口足够远的地方。 然而，火车在隧道外移动需要大量的计算机时间，这不符合我们的兴趣。 我们采取的方法是计算列车上空的稳态流场，并将该解决方案作为列车附近流场的初始条件。 通过在计算区域的右端设置边界条件来大致考虑隧道入口的影响，用于这种稳态计算[9]。 对于黏性计算，不滑边界条件强加在火车车身表面，因为流量分离可能会改变有效的火车堵塞，从而可能会改变压力波的强度。 换句话说，在火车身上发展的边界层已经解决了。 另一方面，滑移边界

400 富士和小川图。 3 Maeda et si的实验装置。[13]即使在黏性计算中，条件也强加在隧道壁和地面表面。 对于轴对称情况的非黏性计算，滑移边界条件强加在所有边界表面上。 通过求解正态动量方程获得固体壁上的压力。 21.3 结果和讨论 21.3.1 轴对称计算-验证和参数研究 关于隧道进入问题有几项分析和数值研究。 山本证明了压缩波的梯度与U3成正比，其中U是列车的速度[15]。 实验室实验证明了这种高达200公里/小时的理论是合理的。 然而，列车的速度在300公里/小时的范围内，日本铁路中心正在开发的MAGLEV（磁悬浮）列车的速度将为550公里/小时。 在这个速度范围内，非线性效应可能会诱发偏离理论。 到目前为止，大多数数值工作都使用了一维方程，其中火车入口被建模为隧道截面的时间变化。 事实证明，一维模拟可以很好地预测压力上升的水平。 然而，压缩波的梯度不一定被很好地预测（稍后将展示）。 在本节中，进行了轴对称计算，以确定现有理论的适用性。 有一个由Maeda等人进行的实验室实验。 [13]。 图3显示了他们的实验设备。 轴对称物体被射入管道中，并测量了管道表面的压力变化。 与实验相对应，考虑了三种鼻子形状：椭圆形、抛物面和圆锥形，如图所示。 4. 实验中使用的条件总结如下。 火车与隧道的阻塞比，Rt =0.116，火车速度，U = 230公里/小时（马赫数 0.188），长宽比，a/b（见图。 4），=5.0。 请注意，日本典型子弹头列车的阻塞率为0.22。

CFD用于火车空气动力学401使用轴对称Euler 方程进行相应的计算。 黏性效应将在后一节中讨论。 初步计算显示，压缩波的非维空间尺度为5级（注意隧道高度为单位），解析压缩波的梯度至少需要20个网格点，x方向的网格大小（火车运动方向）为0.059，远小于5/20。 表一总结了网格点的数量。 由于只完成了尤拉计算，因此使用了两步Runge Kutta 时间积分方案，而不是LU-ADI隐式时间积分方法。 在下面，非维时间计算为0.0是火车机头进入隧道的时间。 在检查空间和时间网格分辨率的影响后，时间步骤确定为0.005。 图5显示了与实验相比测量压力的时间历史。 实线是计算结果，标记是实验数据。 测量位置为6.80非图。 5 隧道墙压力的时间历史-与实验的比较-

CFD用于火车空气动力学到目前为止，结果似乎表明，一维分析可以...............预测压力梯度对两个......................重要参数的依赖性；火车速度和堵塞比。............. 现在，为了看到1 I的一-D-D-nal.................. i...................... l............... i......... 模拟的能力，两个问题是...................... 6................;.................................... （................. -.........，考虑。 一个是阻滞比0.1 0.2 0.3 0.4与R的影响，图中横截面的梯度。 7 压缩波的最大时间导数。与隧道堵塞比，火车要保持不变。 一维分析将火车的运动视为隧道的面积变化，因此3 I：Rk0250压力增加是由于隧道截面的时间相关变化而发生的。 3类型'：Rk0J16假设火车速度相同，-n类型3的梯度3：Rt=0.060压缩波将只有U取决于3 T\textasciitilde{}\textasciitilde{}\textasciitilde{}4：\textasciitilde{}k0.030列车横截面的时间相关变化。 图8显示了考虑的四种类型的火车。 图。 8 图中绘制了8 火车形状的阻塞比不同，但相同的梯度截面。 9. 可以看出，所有四列列车的；车头区域的；行截面面积的导数都是相同的。 计算结果如图所示。 10. 绘制了轴对称和一维模拟的结果。 在一维分析中，压力是恒定的，直到火车进入隧道，因为当有效横截面开始变化时，压力开始增加。 如前所述，尽管最大压力水平不同，但所有四种情况的压力增加的时间历史都是相同的。 换句话说，压缩波的梯度在所有情况下都是相同的，并且不依赖于阻塞比。 而在轴对称分析中，在火车进入隧道之前就已经发生了压力上升，这在物理上是合理的。 如图10所示，压力增加（波的轮廓）的时间依赖行为永远不会被一个-准确模拟

404 富士和小川图。 9 火车头横截面面积的分布图。 10 隧道壁压力的时间历史：x = 6.8 图。 Shinbo等人[12]尺寸分析测量隧道墙的11个位置。 此外，阻塞比不影响压缩波的梯度。 实际上，阻塞比是压缩波梯度的重要参数之一，多维分析对于准确评估压缩波梯度是必要的，尽管一维分析对于粗略估计仍然有用，如图所示。 10表示。

列车空气动力学CFD 405 21.3.2 三维分析-黏性流动实用列车配置模拟-在本节中，根据与称为“Nozomi”新干线列车300系列的现场测量的比较，讨论了实际列车配置的黏性流动模拟计算结果。 Shinbo等人[12]于1991年12月在日本的“片仓”隧道进行了现场测量。 隧道长1266米，双轨。 这列火车叫300 Fig。 12 计算网格系列的一个视图（Nozomi，16辆汽车）在左侧轨道上以270公里/小时（0.221马赫）的速度进入这条隧道。 在入口附近隧道墙的几个位置测量了不稳定的压力。 图11显示了位置。 测量详情可在参考中找到。 12. 由火车速度和隧道高度（7.6米）定义的雷诺数作为代表速度和长度是4.0 x 107。 计算网格如图所示。 12. 由于我们对入口附近压缩波的形成感兴趣，因此将隧道长度定为40.0（312.0m），并在正确的边界应用了非反射边界条件，以避免非物理波反射。 火车长度也缩短了，但保持在20.0（156米），这样端鼻不会影响压缩波的形成。 表二总结了网格点的数量。 最初，火车位于入口前3.0。 一次模拟大约需要30个小时在富士通VPP500向量平行（b）后部（a）前鼻图的1个加工元件上。 13 近表面流线

406 富士和小川图。 14 压力轮廓图超级计算机的时间演变。 当火车机头达到90米（52.0非维时间）时，计算停止。 图13显示了火车进入隧道前的近地表流线。 在前鼻中没有观察到分离区域。 在后部，可以清楚地看到流动分离。 然而，火车很长，分离对压缩波网格点数表II形成的影响是三维模拟小的。 火车区101x67x25底部区101x27x15中间区101x41x41隧道区141x35x21入口区51x35x21领区35x35x9共532,000

CFD用于火车空气动力学407图14（a）-（d）显示了火车进入期间压力轮廓图的顺序。 如图所示，在火车进入隧道之前，隧道内的压力略有增加。 14（a）。 前鼻的压力增加，前鼻产生更强的压力波。 一旦波浪从火车区域分离，它就会变成平面波浪，并向出口传播。 由于停滞区域的影响，在这个过程中左墙的压力更高，尽管在这个图中很难识别。 压力增加在一定时间后停止，火车附近的流场在隧道内达到稳定状态。 图15显示了隧道壁上三个周角点的压力分布的时间顺序。 左侧的高峰和低尖峰是火车周围的压力变化，压缩波的形成可以看作是右侧的增加。 最初，火车前的压力增加，然后一旦达到最大值，压缩波的梯度就会变得更陡峭，并从火车上传播。 为了缓解压缩波的梯度，应在压缩波完全形成之前进行处理。 在这种情况下，这种距离在非维尺度上是2.3。 该图还显示，除了火车经过时，所有三个点的压力几乎相同。 它表明压缩波是一维的。 然而，应该注意的是，火车附近的流场是完全三维的，其影响对压缩波的强度和梯度很重要。 图。 15 隧道墙压力的时间历史 实线：左墙，破折线：右墙和上墙

用于火车空气动力学409的CFD图。 18清楚地显示了区别。 在中心运行的火车的最大坡度要高得多。 结果表明，轨道的位置改变了隧道内的流场，也是压缩波梯度的一个重要因素。 因此，轴对称流模拟不足以准确估计实际列车配置的压缩波。 原因将在下一节中讨论与新的一维估计方法相关。 21.4 一种新的准一维预测方法 迄今为止的研究表明，无黏流动假设是有效的，因为列车表面没有发生强烈的分离。 因此，在开发新方法时，我们假设无黏性流。 图中示意图描绘了瞬时流线的时间变化。 19. 我们可能会认为隧道墙接近火车，而不是火车靠近隧道入口。 当隧道靠近火车时，它缩小了火车附近的溪管。 让我们考虑一下沿着隧道顶壁的溪管。 随着隧道墙的靠近，溪流管变窄。 可以理解为压力上升是由于这种变窄的溪管而发生的。 在下面，我们根据这个想法开发了一个新的一维理论。 考虑面积变化的非稳一维Euler 方程写成，图，19在与火车d,(QA)+dx(FA)=G (21.5) p\textasciitilde{} pu 0 pu, F= p+pu2, G= -pdxA e \_(e + p)u\_ -pdtA\_ (21.6) 这些方程被重写为三个特征变量的三个独立方程。dt w0 + (u+c)dxw0 = 0

410 FUJII \& OGAWA dtw+ + (u+c)d,w,= (21.7) -.,... c DA d,w\_+(u-c)drw\_ = ' V; x A Dt Here, 5w0=8p+-E, Sw +=Su+-, 8w\_=Su-- (21.8) c pc pc 我们对火车前产生的压缩波感兴趣。 因此，只有正特征值的方程，即方程的第二个方程。 （21.7）被考虑。 我们认为与稳态相比，流动变量的变化相对较小，并引入了小扰动公式。 经过一些操作后，这个方程变成了，r DA' dtw' + +(u+c)dxw' + =--- (21.9) 在这里，小扰动项被定义为附着素数。 我们假设火车横截面面积的空间导数很小，忽略了高阶项。 从迄今为止的研究来看，要形成的压缩波是一个平面波。 因此，速度的扰动与压力的扰动直接相关，作为方程，Su=\textasciicircum{}-（21.10）pc，然后特征变量8w' +的扰动可以被压力的扰动所取代，5w\textbackslash\{\}=-8p'（21.11）pc将此放入Eq中。 （21.9），d.p +（u+c）drp = pc V ' x 2 H A Dt 1 ypDA 2 A Dt -LIE M 2 A " dx (21.12)

用于火车空气动力学的CFD 411 x+dx 图。 质量守恒方程（21.12）的20 控制体是计算时间压力变化的方程，它告诉我们压力变化与溪管的时间面积变化直接相关。 换句话说，如果给定了溪管在时间上的面积变化，我们可以使用这个方程计算压力变化。 由于火车的速度是恒定的，右侧的实质导数已被空间导数所取代。 如图所示。 20，我们考虑流管的控制体，并应用质量守恒。wan - (puA) = \$pu-ndl (21.13) 在这里，多维性的影响包含在右侧的积分中。 方程（21.13）可以大致重写为，d，dA A -\textbackslash\{\}pu )+pu =puL dx dx（21.14）在这里，L是积分隧道的长度，u n是垂直于隧道壁的速度分量。 通过额外的合理假设[8]，得到了以下方程。\textasciicircum{}dA \_ T u = u n L (21.15) 将此放入等式中。 （21.12），

412 FUJII \& OGAWA dj+iu+cpj/J-\textasciicircum{}A*" \textasciicircum{}- X'U') (21.16) [o \{t<-xlU t) 在这里，我们将垂直于隧道壁的向外向速度分量的积分定义为v waU。 该方程表明，压力增加，换句话说，压缩波的梯度与图中的示意图所示的v waU成正比。 21场演出。 我们称这个新开发的理论为“vwaii理论”。 V waU可以根据隧道外列车上空的稳态流场来计算。 使用该理论，可以通过计算列车上空稳态流场中与假想隧道壁位置垂直的速度分量来预测预期压缩波的梯度，而无需复杂的流模拟，包括隧道效应。 前田等人[3]对图中所示的四个轴对称体进行了实验。 22. 每个身体的区别是圆锥形鼻子的截断区域。 在实验中，表明压力的最大梯度几乎是恒定的，直到截止区域在一定限度内。 如果截止区域变得太长，梯度会突然增加，即案例7（c）。 图23显示了该理论计算的dp/dt。 实验结果也显示出来。 达成了显著的协议。 该理论清楚地捕捉到了案例7（c）的dp/dt的突然增加。 还对不同的身体配置进行了额外的计算，所有结果表明，v waU理论可以准确地预测dp/dt（现在位置是固定的完全导数）分布。 应该注意的是，v waU理论只需要对自由流中的身体进行稳态计算。 此类计算需要少量的计算机时间。 Q-案例7（a/b=7.0）案例7A案例7B案例7C "0.809 a"' '"\textasciicircum{}0.597 a*' 0.401 a -9- -9- 图。 22 四个鼻子配置

CFD用于火车空气动力学413 wa| Iheoiy O ° A O: Experimaits (case7,7A,7Band7C) 图。 23 隧道内的压力梯度历史（x=6.8）0.010 接下来，我们试图找到实际三维问题的vwall理论的准确性。 在三维情况下，v waU被定义为沿隧道壁的法线分量的表面积分的平均值。 一个横截面的速度分布代表如图所示。 24. 考虑的案例是300系列“Nozomi”的案例，左轨道和中心轨道都计算了。 Dp/dt分布如图所示。 25. 实线是图中已经显示的图。 18是包括隧道在内的昂贵模拟的结果。 这些标记是vwaU理论的结果。 协议相当不错。 然而，值得注意的是，v waU理论清楚地捕捉到了左轨道情况的更高dp/dt。 为了了解差异的原因，图中绘制了vwall的局部值。 26. 由于靠近火车的隧道墙附近的v waU更大，因此左轨道情况的平均v wau会更高。 这一结果表明，有效缓解轰鸣噪音的处理应该放在隧道的近侧。 V waU只需要稳态计算，这需要少量的计算机时间。 上述两种情况只需要一次计算，因为v waU分布需要没有隧道效应的稳态解，并且假设隧道的位置获得。 上述结果表明，火车形状是蓬勃的Fig的最佳。 24个代表性的Vwan分布

414 富士和小川噪音可以通过最小化v waU的最大值来设计。 基于这一想法，为轴对称流开发了一种基于面板方法的简单设计算法。 被认为是轰鸣噪音的良好形状的抛物伞形状被认为是最初的列车配置，并设计了新的列车配置。 研究表明，通过这种简单的设计方法，压缩波的导数降低了30\%。 算法和结果的详细信息写在参考中。 8. 图。 25 压力梯度的时间历史-Vwaii和计算的比较 21.5 结论 计算流体动力学最近的技术已被应用于火车进入隧道的空气动力学问题。 对轴向和全三维流场进行计算。 与现有理论、实验和现场测量的比较表明，目前的方法有助于理解隧道内压缩波的形成机制。 结果表明，现有一维理论和多维分析的应用极限对于准确预测压缩波梯度是必要的。 全三维配置的计算表明了轨道位置的影响，图26左行列车和中央行列车的V壁分布

CFD FOR TRAIN AERODYNAMICS 415到压缩波的梯度。 基于模拟结果，开发了一种叫做理论的新理论。 该理论只能根据没有隧道效应的列车的稳态流动字段准确预测压缩波的形成，因此非常高效，尽管它仍然需要三维流动模拟。 通过与实验和模拟的比较，该理论的能力是合理的。 这个新开发的理论可以成为缓解高速列车轰鸣噪音的良好工程工具。 CFD技术被认为处于第一个成熟阶段。 随着计算机性能的快速发展，任何有经验的人都可以模拟复杂的流场。 然而，要使CFD处于第二个成熟阶段，既需要了解流动场，也需要得出基本理论。 这将使CFD成为真正的工程工具。 我们提出了与隧道相关的火车空气动力学分析，作为此类努力的一个例子。 参考资料1。 鲍德温，B。 S.和Lomax，H.，分离湍流流的层近似和代数模型，AIAA论文78-257，1978年。 2. Fujii，K.，基于强化解算法的统一区域方法，计算物理学杂志，卷。 118，pp. 92-108，1995年。 3. 前田，T.，松村，T.，井田，M.，中谷，K.和宇智田，K.，火车头形状对火车进入隧道产生的压缩波的影响，铁路和磁悬浮车辆加速技术国际会议，日本横滨，1993年11月。 4. Obayashi, S., Matsushima, K., Fujii, K. and Kuwahara, K., 使用LU-ADI分解算法提高Navier-Stokes计算的效率和可靠性，AIAA论文86-0338，1986年。 5. Ogawa，T.和Fujii，K.，隧道中移动的火车诱导的可压缩流动的数值模拟，AIAA论文93-2951，1993年7月。 6. 小川，T.和藤井，K.，火车进入有导轨隧道的空气动力学计算方法，铁路和磁悬浮车辆加速技术国际会议，日本横滨，1993年11月。 7. 小川，T.和藤井，K.，进入隧道的火车诱导的可压缩流动的数值模拟，计算流体动力学杂志，卷。 3，不。 1，pp. 63-82，1994年4月。 8. 小川，T.和藤井，K.，高速列车进入隧道产生的轰鸣噪音的预测和缓解，ECCOMAS 计算流体动力学会议，1996年9月。

416 富士和小川 9. Ogawa，T.和Fujii，K.，对进入隧道的火车引起的三维可压缩流动的数值调查，计算机和流体，将于1997年出现。 10. 小泽，S.，关于隧道出口轰鸣噪音的研究，JR研究报告编号 1121，1979年，（日语）。 11. 罗，P。 L.，Euler 方程的基于特征的方案，流体力学年度评论，pp. 337-365，1986年。 12. Shimbo，Y.和Hosaka，S.，高速列车的稳压和非稳压测量，铁路和磁悬浮车辆加速技术国际会议，日本横滨，1993年11月。 13. 托马斯，J。 L.，van Leer，B.和Walters，B. W.，Euler 方程的隐式通量分割方案，AIAA论文85-1680，1985年。 14. 渡边，R.，藤井，K.和东野，F.，进入隧道的火车诱导的压缩波的一维数值模拟，日本机械工程学会杂志B，第61卷，第 592，1995年，（日语）。 15. 山本，A.，火车和隧道的空气动力学，JR研究报告，第 1230，1983年，（日语）。

\clearpage
\chapter{与CFD一起追求价值}
与CFD Paul E.一起追求价值。 Rubbert1摘要 本文的关键发现是，CFD的价值与它在飞机设计过程中对学习率的贡献直接相关。 更高的学习速度导致更好的设计。 学习率由两个项的乘积组成，即（i）每个设计周期的学习，乘以（ii）在给定时间内可以执行的设计周期数量。 CFD的早期发展倾向于关注前者，而忽略或忽略后者。 但20世纪90年代的教诲更加注重后者，其结果是，今天用于设计飞机的流程正在以前所未有的速度改进。 22.1 导言 在解决CFD对价值的追求时，首先必须弄清楚如何定义价值的含义。 这最好通过检查航空公司判断飞机价值的方式来完成。 一个常见的价值衡量标准是一个单一数字，称为TAROC（飞机相关运营成本总额）。 通用宽体飞机的典型TAROC明细如下：百分比项目24利息所有权14折旧成本39\%1保险12燃料9机组人员现金5机身维护DOC 29\%3发动机维护'波音公司。 计算流体动力学的前沿-1998年。 编辑：D。 A. 考伊和M。 M。 哈菲兹。 ©世界科学

418 RUBBERT 9 机舱乘务员 8 一般和行政。 4 地面处理飞机 4 着陆费相关 3 控制和通信 IOC 32\% 3 Gr。 Prop Main \& Dep 1 APU Fuel 可以观察到，除了空气动力学效率之外，还有很多东西都具有价值。 在过去的一年里，我有幸听到了联合航空公司高管戈登·麦金齐的演讲，关于他们从波音、空中客车和道格拉斯的宽体飞机产品中选择的过程。 他们将评估分为11个领域，按排名顺序列在下面。 等级加权1。 任务能力2。 内饰3。 维护12 4. 设施10 5。 货物6。 地勤7。 飞行操作8。 训练9。 产品支持10。 技术问题11。 驾驶舱5 总计100 信息似乎再次表明，价值主要由与空气动力学几乎没有明显关系的东西组成。 然而，空气动力学设计过程与这些其他因素之间的联系比显而易见的要强得多。 这源于这样一个事实，即设计和将飞机推向市场首先是大规模系统集成的练习。 设计和生产像现代喷气式运输车一样复杂的车辆是一个令人费解的挑战。 这意味着与各个技术学科（如空气动力学）相关的价值概念必须远远超出传统的、特定学科的措施，如阻力、升力等。 与特定学科相关的更广泛的价值衡量标准也增长到包括这些学科和学科流程为大规模系统整合实践做出有效贡献的能力。

关于CFD 419的价值追求，空气动力学也是如此。 我们生活在一个信息时代，所有严肃的竞争对手都可以获得相同的技术。 我们都使用相同种类的电脑和风洞。 我们使用相同的湍流型号。 以及许多相同的计算算法。 业务中所有严肃的竞争者都能够设计到可比级别的任务性能。 这被视为“被原谅”。 这只是一个人有资格进入这个行业的赌注。 留在业务中和繁荣的能力越来越取决于能否很好地完成其他事情，大规模系统集成接近榜首。 因此，定义CFD价值的一种有意义的方式是衡量它如何促进公司最好地执行大规模系统集成过程的能力。 这就是我在本文中选择定义CFD值的方式。 我认为实现最先进的空气动力学性能的能力是必然的。 波音知道如何做到这一点。 空中客车知道怎么做。 道格拉斯知道怎么做。 和其他人。 那不是关键问题。 价值的真正衡量标准在于我们如何做到这一点，以我们在更大规模的系统集成和业务实践的更广泛框架中整合空气动力学设计过程的有效性来衡量。 22.2 空气动力学设计过程必须努力实现什么？我编写了一份需求、愿望和目标的简短清单，这些需求、愿望和目标定义了空气动力学工程对大规模系统集成实践的支持有多好或有多差。 列表如下：22.2.1 提供提供最先进的性能的空气动力学线路。 这是必然的。 22.2.2 提供多种空气动力学设计解决方案。 仅仅设计一组实现空气动力学目标的线条是不够的。 人们需要能够提供所有可能的空气动力学设计解决方案，以便做出最能支持整体飞机系统集成的设计决策。 例如，在机翼/机舱集成中，空气动力学家希望能够通过机翼变化、机机机改变、支柱变化或其任何组合来消除干扰。 通过这种方式，有可能达成一种在空气动力学、结构、燃料量、制造成本、重量和系统之间取得最佳平衡的设计。 22.2.3 在性能、可制造性、可维护性、非经常性成本、上市时间之间提供交易的快速可见性和适应性，

420 RUBBERT 22.2.4 快速确定约束成本。 22.2.5 由一些非常快速的设计周期组成。 这就是允许积极追求最大性能和“突破设计范围”的原因，它提供了一种从失败中快速恢复的手段。 当可用的设计周期数量太少时，失败是不可能的。 恢复是不可能的。 这迫使保守的设计和为确保不会发生失败而做的事情的额外成本。 22.2.6 提供全面飞行性能的完美预测。 22.2.7 敏捷地适应新的和不同的设计限制和机会。 22.2.8 非常迅速地完成上述所有工作，因此空气动力学设计不是业务收购过程中或飞机设计和整合的实际过程中的起搏专案。 上述构成了一项艰巨的任务，其中仍然存在着巨大的改进机会。 清单上所有项目中隐含的共同主题是，无论正在做什么，都需要迅速获得非常高水平的空气动力学知识。 人们想知道所有合理的设计解决方案、大量交易的空气动力学后果以及各种约束的空气动力学成本，所有这些都以下降规模的飞行性能来准确表达。 那是很多知识，而产生这些知识的时间非常短。 我们今天做不到这一切。 但对我来说，越来越清楚的是，设计解决方案的质量在很大程度上取决于在可用时间内可以学到的数量。 22.3 导致价值的总体指标 如果设计解决方案的质量在很大程度上取决于在可用时间内可以学习的数量，那么总体指标必须是学习率。 在所有其他条件相同的情况下，最好的设计将由最大限度地提高学习速度的过程产生！ 因此，CFD在当今飞机设计世界中的价值可以通过它对工程设计团队了解他们试图做的事情的速度的贡献来表达。

追求价值与CFD 421 22.4 管理设计的方程 前面的想法可以以以下方程的形式表示，首次出现在参考1. fAmount\textasciicircum{} = TLearning\textasciicircum{} x fCycl\textasciicircum{} x \textasciicircum{} \textasciicircum{}Learned J [ 周期 ) \textbackslash\{\} Time) (22 1) 业务目标要求我们：（i）充分学习，以便能够达到非常高质量的设计，以及（ii）最大限度地减少分配的时间，从而减少成本和上市。 实现这些目标的唯一方法是努力最大化方程右侧出现的前两个项的乘积（22.1）。 为了非常清楚地表明这一点，可以通过以下方式重新排列术语。f 金额\textasciicircum{} (质量 \textbackslash\{\} 学习率 (22.2) [学习] /'循环 \textasciicircum{} 学习 J \_ \textbackslash\{\} 产品 [ 周期 J A V 时间) ( 时间 \textbackslash\{\} = 7 时间到S \{分配的 J \textasciicircum{} 市场 J 这个方程的左侧将学习率表示为两个项的乘积，即（i) 设计周期内学习的数量和（ii）在给定时间内可以执行的周期的周期数。 这就是最重要的。 这不是一个或另一个。 它是两者的乘积，定义了与给定设计过程相关的值。 这个方程还为我们提供了一种量化过程增量变化值的方法，例如使用新的CFD代码可能产生的变化。 让我们定义：（。 学习\textasciicircum{}每个周期的添加信息或学习是给定变化f提供的\textbackslash\{\}周期J。 周期\textasciicircum{}在给定时间J的给定时间内可以执行的额外周期数，由给定的变化产生，然后给定变化的值可以表示为：

422 RUBBERT（学习+A学习\textasciicircum{}\%（学习+A周期\textasciicircum{}=\textasciicircum{}\textasciicircum{}+AValue）\textasciicircum{}周期周期J\textasciicircum{}周期时间，）（22.3）忽视高阶项和非维化，这变成了学习\textasciicircum{}（周期化时间周期（周期时间A值学习周期i周期V时间j值（22.4）一个常见的困境是，引入增加学习/周期的新CFD工具，通常伴随着执行周期所需的时间的增加，从而减少周期。 这可以将AVALUE的符号从正值变为负值！ 22.5 CFD的影响或价值 1990年之前 2 当我进入飞机业务时，没有CFD。 定义一组候选空气动力学线的过程并没有花很长时间，因为没有很多分析可以完成。 风洞模型的制造也没有花太多时间，因为几乎不需要表面压力仪器（没有CFD可以与之比较）。 模型表面定义要求也更加宽松，因为线定义过程无法区分表面保真度微小变化的空气动力学后果。 因此，设计周期被迅速执行。 我记得早期的飞机计划在80年代有模型机翼编号，这意味着已经执行了许多设计周期。 现在回到方程（22.1），早期的设计师在给定的周期中没有像我们现在学到的那么多，因为我们现在有CFD和广泛的风洞模型仪器可以学习。 他们的学习/周期，方程（22.1）右侧的第一个项，很低。 但他们（周期/时间），他们在给定时间内可以执行的周期数量非常高，因此那些早期设计师能够在给定时间内学到很多东西。 DC-8s、DC-9s、707s、第一代737s，是的，仍然宏伟的747的早期版本，都是以这种方式设计的，几乎没有CFD，但有一些快速执行的空气动力学设计周期。 分配的时间，方程右侧（22.1）的最后一个项，在那些早期也较少。 747在大约3年内投放市场，而不是我们最近开发777的5年。 我觉得这相当谦卑。 在CFD过去的大部分历史中，发生的情况是，每次我们再次取得重大进展时，我们都会增加每个周期可以学习的数量。 方程（22.1）右侧的第一个项变大了。 2 本节几乎逐字引用自参考1。

关于CFD 423的价值追求，但是，CFD的每一次新进展都会带来周期/时间的减少。 右侧的第二个术语变小了。 我清楚地记得20世纪70年代初的时代，当时我参与了一个设计非常快速、区域规则的亚音速运输的项目。 这是波音公司第一个在机翼设计中使用CFD的飞机项目。 结果是，我们在定义线路所需的时间上增加了大约六周，我们增加了几个月的风洞模型制造时间，因为我们需要广泛的表面压力测量来校准CFD。 当时，我认为我们取得了长足的进步，因为我们可以“开箱即用”生产出具有良好空气动力学的设计，但我忽略了一个事实，即我们花了6个月来设计和测试一个配置。 在同一时期，博士 美国宇航局LaRC的Richard Whitcomb是我们教我们能做什么的真正领导者，他专注于增加方程（22.1）右侧的第二项。 他开发了一个在几个小时内而不是6个月内亲自使用文件和更改（重新设计模型）风洞模型的过程。 有据可查，他制作的最终设计在时间上领先于我们，并且在空气动力学上表现同样好。 回想起来，我现在意识到他的学习速度比我们快得多，他是领导者，而我们是追随者，就证明了这一点。 22.6 20世纪90年代的ERA 近年来，该行业采用了一种新的范式，该范式认识到周期时间、变化和敏捷性至关重要。 为了应对这种新范式，我们和其他人重新调整了对CFD的投资，专注于那套新的指标。 这在那些没有用方程（22.1）-（22.4）表达的术语中思考世界的人中引起了很多的沮丧和手。 但足够长的时间已经过去，我们现在可以开始观察这种重定向的后果。 我观察到的后果是，通过各种措施，进展和改进都以前所未有的速度发生。 最近设计的空气动力学性能非常出色。 我们预测了777和下一代737的全尺寸飞行性能，其准确度几乎令人难以置信。 但更重要的是，我们开展了空气动力学设计过程，这显著有助于波音在创造新飞机产品时出色地完成大规模系统集成的能力。 最终结果是其市场占主导地位的波音777产品线，以及首次交付前打破所有销量记录的下一代737产品线。

424 RUBBERT 22.7 关键推动因素 在我看来，在20世纪90年代贡献最大价值的关键推动因素包括以下项目和成就。 22.7.1 对学习/循环的贡献•通过CFD优化 在我们的空气动力学设计流程中部署优化方法，为我们每个周期提供了更多信息，特别是在敏感性和约束成本方面。 •将可制造性考虑因素与基于CFD的设计相结合，我们不再事后发现设计的形状无法制造。 我们积极研究空气动力学性能和制造成本之间的交易。 •使用CFD进行阻力分解，我们不断提高计算阻力增量的能力，这明显提高了CFD的价值和有效性。 •不确定性管理CFD在使我们能够定义和有效管理与预测全尺寸飞行性能相关的不确定性方面发挥了重要作用。 • 学科 仅开发可以在可用时间内使用的CFD工具 22.7.2 对周期/时间的贡献 • 快速和一般CFD的早期战略投资决策。 20世纪70年代和80年代做出的关键决定是：（i）用于“标准”配置安排（如机翼/机身）的网格生成将完全自动化。 （二）使用采用解决方案适应性的矩形、非表面拟合网格来处理不涉及大面积分离流动的流动中复杂和任意的几何安排。

关于CFD 425的价值追求，这些战略决策对于为我们提供可以在可用时间内使用的CFD工具至关重要。 对于复杂的黏性流动，我们还没有可比的快速设置能力，我们目前无法支持CFD在巡航制度中享有的有效性来支持高升力襟翼系统的设计，这反映了这一事实。 •将可制造性考虑因素与基于CFD的设计相结合。 通过减少在空气动力学线定义和可制造性之间迭代的需要，我们缩短了设计周期。 •不确定性管理 这最大限度地减少了在将设计投入生产之前等待风洞数据的需要，从而缩短了巡航配置设计的周期时间。 在获取最终线路的风洞数据之前，我们越来越多地观察到设计致力于生产。 如今，高速设计的自行车几乎完全由CFD完成。 •六天风洞模型 对周期时间最有效的改进之一是开发在一两个星期内而不是六个月内设计和建造风洞模型的能力。 去年，我们利用了这一能力，在单个风洞测试进入的范围内循环了整个机翼设计。 综上所述，过去开发产生的价值分为增加单个设计周期内可以学习的数量和增加给定时间内可以执行的周期数量。 在我看来，除了六天风洞模型能力外，CFD一直是所有主要的启用技术因素。 22.8 未来方向 再次参考方程（22.1），我观察到，就学习量而言，我们最终将达到回报递减的点。 一旦一个人学到了足够的知识来做出正确明智的设计决策，获得更多知识的回报就会减少。 我们还没有到那个点，但我们正在迅速接近巡航制度的设计。 我们现在面临的真正挑战是需要大幅减少分配的时间。 如果我们能把新飞机的上市时间从5年缩短到2.5年，我们的产品在市场上的价值将大大提高。 我们只能通过f学习\textasciicircum{}（周期\textbackslash\{\}\{周期J[时间J）的相应增加来实现这一点

426 RUBBERT 我们必须抵制减少周期数量的诱惑。 是的，我们也许可以在一个周期内设计出一个相当不错的机翼。 但这样做，我们剥夺了自己从正在发生的学习中受益的能力，因为没有第二次机会将这种学习纳入另一个机翼的设计中。 单周期设计方法的一个更严重的后果是不允许失败。 它必须第一次就运作良好，这迫使设计方法更加保守，并随之而代之的是各种昂贵和耗时的“保险政策”。 这让你很容易受到一个快速学习的竞争对手的伤害，新西兰队在上届美洲杯比赛中向游艇赛车兄弟会的其他成员传授了这一教训，当时他们在44场比赛中向竞争对手展示了他们的船黑魔法船的船尾！ 以下引文总结了他们是如何做到的。 “新西兰团队快速设计周期的速度是使其迭代过程和探索数千个选项成为可能的关键因素。 加快设计周期的制造公司同样可以击败市场竞争。 但他们也可以跟随Kiwi的领导，利用额外的时间提高其产品的质量、性能和客户契合度。” 当我检查我的水晶球，以深入了解我们将如何实现所需的学习速度的提高时，我看到了以下图像：•CFD对于低速飞行制度、与结构负载、稳定性和控制相关的飞行包络的周长，以及与颤动和噪音相关的不稳定现象，它快速而有效。 •结合CFD和实验方法，让我们有了更多的知识和理解 唤醒涡流，以及更好地处理与唤醒涡旋相关的问题的能力。 •优化方法的更多贡献。 •一种更外科的风洞测试方法，我们只测量我们当时需要的东西，并放弃“测量所有东西以防万一”的旧习惯。 •对基础设施和流程的系统性变化。 我们从6天风洞模型项目中学到的一件事是，我们是自己最大的敌人。 解决方案不在于发明新技术以进一步改善周期时间。 相反，解决方案主要在于改变管理我们决策流程、预算流程、奖励和认可系统等的官僚机构。

追求价值与CFD 427 这些是当今世界大部分周期时间的原因。 •学习如何利用全球24小时工作。 我们的首席执行官Phil Condit挑战我们走向全球。 一旦全球视野的元素描绘了我们每天24小时、每周7天、跟随太阳工作的事情。 想象一下，在我们一天结束时，把我们的工作交给那些刚到工作的同事，他们离世界只有三分之一的路程。 反过来，他们会把它交给其他环游世界的2/3的人。 这样，当我们第二天早上到达工作时，我们会发现前一天交出的任何东西都增加了16个小时的进度。 不难想象这对周期时间有什么帮助。 不难想象信息系统将使这成为可能。 他们已经存在了。 主要困难在于学习如何发展人际关系类型，以建立使其发挥作用所需的信任和理解水平。 22.9 结束演讲 通过CFD追求价值，导致学习率被公认为价值的关键衡量标准。 学习率由两个组成部分的产物组成。 一个是学习/循环。 另一个是周期/时间。 在CFD的大部分历史中，我们主要关注CFD对学习/周期的贡献。 每个新发展都因其揭示以前没有揭示过的东西的能力而被吹捧。 研究人员倾向于忽视或忽略对周期/时间的影响，这往往是负面的。 但在20世纪90年代，我们面临着一个新的价值体系，这个体系认识到了周期时间的重要性。 这导致了对CFD（和风洞过程）的投资，旨在改善周期时间。 结果是，飞机设计过程现在正以前所未有的速度改进，以该过程产品在全球市场上的价值来衡量。 参考资料1。 Rubbert，P. E，“CFD在飞机设计中的使用”，CFD 97加拿大CFD协会第五届年会，1997年5月25日至27日，不列颠哥伦比亚省维多利亚。 2. Rubbert，P. E.，“AIAA莱特兄弟讲座：CFD和不断变化的飞机设计世界”，ICAS-94-0.2，1994年9月。

\clearpage
\chapter{十字路口的CFD：行业视角}
十字路口的CFD：行业视角 Pradeep \textasciitilde{}aj* 23.1 介绍 计算流体动力学（CFD）在从六十年代中期到九十年代中期的三十年间取得了重大进展。 两个因素在很大程度上导致了这一进步：数字计算机的速度和内存的惊人增长了四个或更多，以及数值算法和软件的惊人进步。 CFD现在是所有科学和工程学科的组成部分，流体动力学相互作用在其中发挥着重要作用。 对于科学研究来说，CFD的重要性从它在更好地理解湍流 [l-31. 从工程角度来看，CFD现在被广泛认为是一种关键的赋能技术，不仅支持固定翼飞机的设计--本文的主要重点--还支持旋翼飞机、航天器、涡轮机械和汽车的设计。 CFD的工程应用的广度和范围确实很大。 然而，在作者看来，我们正处于一个阶段，CFD对飞机设计的fir11好处很容易证明，但更难实现。 CFD正处于十字路口。 过去三十年来，我们沿着CFD开发道路的渐进步骤不足以满足未来设计挑战的需求。 我们需要仔细检查可用的途径，并积极追求最有前途的途径，以便CFD能够更有效地“洛克希德·马丁航空系统技术研究员，玛丽埃塔，佐治亚州30063-0685，美国。 @ 1997 洛克希德·马丁。 所有权利重新seed Frontier of Computorional Flurd enomics - 1998。 编辑：David A. Caughey和Mohamed M. 哈菲兹 01998 世界科学

430 RAJ满足未来的飞机设计需求。 本文的其余部分讨论了与这些观察有关的问题。 这篇论文按照以下思路排列。 关于CFD在飞机设计中的作用的一些想法在第23.2节中。 第23.3节概述了与CFD飞机设计有效性有关的问题。 第23.4节总结了过去三十年取得的进展。 随后，第23.5节讨论了设计过程中发生的转变及其对CFD的要求。 第23.6节给出了一些总结意见。 23.2 CFD在飞机设计中的作用 本节简要概述了CFD在飞机设计中的作用。 强烈鼓励读者阅读关于这个主题的其他有趣和发人深省的文章。 它们包括但不限于米兰达关于CFD在飞机设计中的应用的调查论文[4]，汉考克的兰彻斯特纪念讲座[5]，米兰达[6]关于CFD在战斗机设计中的挑战和机遇的论文，科斯纳[7]关于CFD航空航天应用问题的讨论，以及拉伯特[8]关于CFD在不断变化的飞机设计世界中的作用的莱特兄弟讲座。 典型的设计过程通常分为三个顺序阶段：（i）概念阶段，（ii）初步阶段和（hi）生产阶段。 在大多数情况下，有一个概念前阶段，客户要求被映射到设计空间中，并为描述整体尺寸、形状、性能、重量、成本等的设计变量设定目标。 然后在概念阶段生成几个候选配置。 在权衡研究之后，通常在初步设计阶段选择单个配置进行进一步开发和验证。 设计通常在概念和初步阶段进行多次修改。 设计团队的目标是创造满足所有客户要求的“最佳”设计。 在生产设计阶段，在发布制造设计之前，会进行最终布局和更广泛的验证。 在所有设计阶段发生的无数活动可以分为综合或分析。 综合涵盖定义、完善和更改概念和配置；分析包括产生有用数据来评估概念和配置的方法、工具和专业知识。 在大多数飞机设计项目中，CFD主要被视为一种工具--就像风洞一样--用于生成空气动力学数据。 准确估计空气动力学数据对每个飞机设计项目都至关重要。 需要力和力矩数据来评估性能和飞行质量；表面压力为结构设计提供输入；流场数据促进了系统集成，例如发动机或武器与

CFD AT A CROSSROADS 431机身。 为了生成所需的空气动力学数据，风洞和CFD的明智组合现已发展成为大多数设计项目的首选方法；Bangert等人[9]关于F-22战术战斗机设计的论文就是一个例子。 时间表和成本限制是飞机设计所有阶段的核心，这一点怎么强调都不为过。 同样重要的是要注意，在设计的早期阶段做出的决定会对最终设计的生命周期成本产生深远的影响。 据各种估计，飞机生命周期成本的70\%至90\%是在设计的早期阶段锁定的。 当然，使用CFD和风洞生成的空气动力学数据有助于设计团队做出决策。 23.3 CFD 飞机设计的有效性 只有当CFD能够满足设计团队的愿望和期望时，它才能成为飞机设计的有效工具。 团队自然希望在时间表和预算内提供最忠实的工程数据。 为了充分有效，CFD必须同时满足三个关键要求：（i）快速周转（或短的分析周期时间），（ii）可靠的准确性和（iii）负担得起的成本。 这些要求背后的理由由作者在参考文献10中提出，这里不会重复。 然而，为了完整起见，包括了一个简短的梗概。 快速周转。-周转定义为从初始出发到数据最终交付的时间跨度。 当Bangert等人[9]最近审查CFD对F-22设计的应用时，“最小化CFD分析的日历时间”出现在列表的首位。 典型的CFD申请流程由三个步骤组成。 第一步涉及获取几何形状和建立计算模型（或网格生成）。 第二步是运行流动求解器。 第三步也是最后一步是从流动求解器输出中提取所需的空气动力学数据。 如下一节所述，CFD 程序的周转时间差异很大。 可靠的准确性。-虽然缩短周转时间至关重要，但产生可靠准确的空中动态数据同样重要，如果不是更重要的话。 可靠准确的数据必须具有已知且可接受的误差水平。 大多数CFD分析在提供误差估计方面都不足。 数值方案和物理模型都会导致CFD解决方案的误差。 必须注意的是，当CFD在设计环境中行使时，分析师通常没有通过不同的数值方案或物理模型来进行广泛的参数研究来评估其预测的准确性水平。 负担得起的成本。-负担得起的CFD分析的要求对飞机设计团队来说非常重要。 费用包括人工和电脑费用。 人工费用与几何获取/建模、网格生成、流动求解器运行的监控和流动求解器输出的处理有关。 高速的

通常需要432台具有大内存的RAJ计算机来支持设计数据需求；通常需要进行数百次分析运行来生成所需的数据量。 此类机器可能相当昂贵，因为计算机成本与处理器速度和内存密切相关。 23.4 CFD 进展：过去三十年图1。 四个级别的CFD方法 在过去的三十年里，开发了各种各样的CFD方法。 它们在支持飞机设计需求方面的有效性差异很大。 这些方法可以大致分为图1所示的四个级别。 这种分类主要基于数学公式、能力和局限性，以及将方法引入飞机设计的时间框架。 众所周知，基于Navier-Stokes（N-S）方程的IV级代码可以模拟几乎所有感兴趣的流动现象，而连续性假设是有效的。 （基于气体动力学理论的玻尔兹曼方程对分子流动进行建模需要；这里不会涵盖相关的数值方法。） 然而，在20世纪60年代，没有足够的计算机能力和高效的数值算法来求解N-S方程。 这迫使研究人员根据N-S方程的非黏性近似来探索替代方案；前三个级别对应于基于非黏性近似层次的代码。 20世纪60年代末引入的最低级别代码现在被广泛使用和接受；20世纪80年代末引入的最高级别代码的使用也在迅速增加。 本节重点介绍了每个级别的 CFD 方法的基本特征及其有效性。 一级：线性势。-线性势法基于Prandtl-Glauert或Laplace方程，这些方程构成了N-S方程的最低水平非黏性近似。 大多数代码采用边界积分方法。 控制偏微分方程（PDE）和边界条件使用格林定理以曲面积分形式铸造。 解决方案是通过将几何形状离散化为小元素，并为每个元素分配一个奇点（源、双元或涡流丝）来构建的。 奇点强度是通过在每个元素的控制点满足非正常流动条件来确定的。 根据曲面离散化中使用的近似值（平均曲面或实际曲面）以及奇点的类型和功能形式（常源、恒定双字、线性双字等），可以开发具有不同特征的代码[4]。 最简单的代码，广为人知

CFD AT A CROSSROADS 433涡流晶格方法，使用几何和涡流丝的平均表面表示作为奇点。 VORLAX代码[11]是一个代表性的例子。 当实际的表面几何形状离散化时，这些方法通常被称为面板方法。 在QUADPAN[12]中使用了每个元素上恒定的低阶奇点分布，在PANAIR[13]中使用了高阶奇点分布，线性或二次。 线性势码的数学公式的简单性本质上限制了其有效性，仅限于纯粹的亚音速和超音速附着流。 尽管有这些限制，但这些程序在设计工作中被广泛使用，因为它们的易用性、计算效率和建立在多年的使用基础上的相对较高的信心。 经验丰富的用户可以在几个小时内设置计算模型，即使是相对复杂的配置，如完整的飞机。 图2显示了P-3猎户座飞机的QUADPAN模型，带有用于空中预警和控制的旋转圆顶[14]。 I级代码的计算时间从超级计算机上的几秒钟到工作站上的几分钟不等。 然而，用户专业知识和经验对于确保正确解释解决方案至关重要。 涡流晶格方法和面板方法通常为稳定飞行条件提供良好的力、力矩和分布式气载估计。 数据构成了设计早期阶段性能和权重估计的基础。 一些代码还提供了设计选项，可用于确定一组规定的空气动力学参数的几何特征（如机翼的扭转和弧）。 为了满足空气弹性和颤动学科的空气动力学数据需求，双晶格方法[15]的版本是首选代码。 线性电位代码于20世纪60年代末首次引入飞机设计环境，整个代码类别在80年代初达到了高度成熟度。 除了振荡空气动力学程序[16，17]之外，目前很少在研究和开发此级别的程序方面投入的努力。 第二级：非线性电位。-非线性电位方法基于跨声速小扰动（TSP）方程或全势流方程（FPE）。 与I级代码相比，他们用冲击对跨声速流进行建模的能力是最重要的优势。 然而，这种好处是以增加复杂性为代价的，这源于需要诉诸现场方法来解决非线性PDE。 现场方法需要一个区域图2。 围绕给定配置的P-3 AEW\&C QUADPAN模型

434 RAJ被分成小的基本体积；仅仅分割表面已经不够了。 在20世纪70年代，在开发实用的跨声速流分析方法方面取得了相当大的进展[18]。 进展起源于Murman和Cole的里程碑式论文[19]。 Boppe的TSP代码[20]和Jameson和Caughey的FPE代码的FLO系列[21]在20世纪70年代末被广泛使用。 在实践中，由于边界条件处理的不同，TSP代码比FPE代码更容易使用，特别是对于复杂的几何形状。 TSP方法允许基于在平均表面应用无法流条件的简化处理。 相比之下，FPE方法需要在实际表面应用。 因此，笛卡尔场网格系统对于TSP代码来说就足够了，而FPE代码需要符合边界的网格。 与符合边界的网格相比，笛卡尔网格更容易设置。 当然，TSP代码在几何和流动条件类别方面受到限制，它们可以准确地建模--这是其简化边界条件处理的直接结果。 在70年代，新发现的计算跨声速流动能力的前景和兴奋是如此强烈，以至于在分析方法仍在发展时，甚至开发了机翼设计程序[22]。 由于跨声速流动特别容易受到与冲击/边界层相互作用相关的黏性效应，因此在将无黏性TSP和FPE代码与边界层代码耦合方面进行了大量研究。 还开发了基于TSP配方[23]的空气弹性分析能力。 尽管二级代码提供了跨声速流建模所需的能力，但它们不如一级代码有效。 各种因素促成了这种情况。 例如，网格生成技术水平不够成熟，不足以支持FPE方法的常规应用到任何比机翼或机翼体更复杂的东西。 正如人们可能怀疑的那样，代码的应用证实，对于含有强激波或涡量大区域（例如，前沿涡流）的流动，溶液的准确性下降了。 因此，代码的效用受到严重限制，特别是在战斗机设计方面，因为典型战斗机的性能受到前缘涡流的极大影响。 在作者看来，非线性势码基本上被八十年代初欧拉码的快速进步所接管。 TRANAIR代码[24]是这一趋势的例外。 TRANAIR采用一种非常规的混合方法，将处理复杂几何形状的面板方法的灵活性与在笛卡尔网格上实现的FPE公式处理跨子流的非线性的能力相结合。 HI级：欧拉。-构成III级代码基础的Euler 方程，代表N-S方程的最高级非黏性近似。

CFD AT A CROSSROADS 435 它们适用于整个亚音速到高超音速飞行制度。 再加上它们自动捕获旋转流区域（如机翼后面的尾迹和从战斗机翼锋利的前缘发出的涡流）的能力，不需要对此类区域进行明确的先验定义，这使得它们比一级或二级代码更有用。 然而，增强的能力是以解决至少四个和通常五个耦合一阶PDE而不是一个二阶PDE的额外计算成本为代价的。 在八十年代之初，两个因素说服了大多数研究人员将重点转移到Euler 方程上：预计计算机能力的增长和开发更高效的数值算法来解决Euler 方程[25，26]。 此外，随着有限体积概念的引入，将流动求解器与网格映射分离，边界符合网格生成技术的加速步伐，对于实现CFDer例行分析完整飞机几何形状的梦想具有相当大的希望。 这里介绍了所取得的令人印象深刻的进展的概要；详细信息可以在包括参考文献27和28在内的许多出版物中找到。 尤拉代码可以确定两条不同的发展路径：一条基于六面体结构网格，另一条基于四面体非结构网格。 在八十年代初期，大多数研究人员将精力集中在结构化网格方法上，但从八十年代中期开始，推动力转向了非结构化网格方法。 意识到非结构化电网在处理复杂几何形状方面提供了更大的灵活性，并承诺将电网生成过程“自动化”，这促使了这一转变。 结构化网格的倡导者追求多块网格，以克服在处理复杂几何形状时遇到的困难；基于修补或覆盖的多块网格的代码演变为高度复杂。 结构化和非结构化电网方法现在都可以使用高性能计算机在几个小时内为飞机配置提供稳态解决方案。 这里展示了两个示例，说明了典型的补丁多块尤拉代码TEAM的能力，以支持F-22设计需求。 第一个涉及完整的飞机载荷预测[9]。 分析了近370个案例，涵盖了六个马赫数字、几个攻角和一些偏航角。 使用表面蒸腾图3模拟了前缘和后缘襟翼偏转、水平尾部偏转和方向舵偏转的影响。 F-22表面压力概念[29]。 计算和相关性的相关性

436 RAJ图4。 图3显示了F-22入口锤击超压模拟测量的0.9 马赫数和8度攻角的表面压力。 第二个例子涉及超音速飞行条件下入口管道中锤击过压的时间史预测[30]。 图4显示了管道压力和前体周围的马赫场的快照。 尽管结构化网格生成技术取得了相当大的进步，但为复杂的几何形状构建多块网格仍然是一项劳动密集型和耗时的任务。 非结构化网格生成尚未充分自动化，尽管它肯定需要更少的时间和精力。 最近笛卡尔网格方法的复苏[31，32]提供了一个有吸引力的替代方案，因为它们基本上消除了网格生成的困难，大大减少了时间和精力。 图5显示了笛卡尔网格SPLITFLOW代码应用于F-16存储-carriage/分离建模[32]的示例。 自由流马赫数为0.95，攻角为12度。 尤拉代码开发的另外两个方面值得一提。 首先，减震捕获而不是减震装置已成为首选方法；自适应耗散方案的上风和中央差速器都取得了巨大成功。 其次，即使用于建模稳定流，大多数代码也能解决Euler 方程的时间依赖形式。 收敛加速技术，如局部时间推进和多网格，用于以计算高效的方式获得时间渐近稳态解。 显式和隐式时间行行方案都得到了有效利用。 由于使用了与时间相关的方程，扩展代码以对非稳定流进行建模相对简单。 这一方面已被用于开发动态空气弹性分析方法[33]。 最近开发逆向设计[34]和空气动力学设计优化[35]方法的尝试也值得注意。 图S。 F-16 SPLITFLOW 虽然欧拉码在建模强冲击方面优于非线性势码，但解决方案不一定更准确地模拟跨声速流，因为它们无法模拟冲击/边界层相互作用效应。 一些研究人员将欧拉代码与边界层代码相结合来纠正这一缺陷。 欧拉代码确实比潜在流动方法有优势

CFD在十字路口437捕捉前缘涡流。 但是，在次级和/或三级涡流产生相当大影响的情况下，初级涡流的位置和强度可能不准确。 此外，代码无法提供光滑表面的总阻力（包括皮肤摩擦）或模型流动分离的估计值。 因此，毫不奇怪，N-S代码的开发一直在积极地进行并行追求。 第四级：Navier-Stokes.-Navier-Stokes代码与Euler代码有很多共同点。 在实践中，单个代码通常满足求解欧拉方程和N-S方程的需要。 这直接源于两组方程之间的相似之处。 消除扩散项很容易将N-S方程转换为Euler 方程；它们都共享一组共同的对流项。 然而，这种看似微不足道的差异的实际影响是巨大的。 为了准确解析扩散项，需要靠近固体表面（以及黏性应力较大的其他区域）的高度聚类网格。 N-S分析对网格生成、数值准确性、计算资源等构成了重大挑战。 有了适当的网格聚类，人们可以使用N-S方程以相对直接的方式模拟层流。 但使用这些方程直接模拟即使是简单的湍流流，也会将当前的超级计算机拉到极限[3]。 目前，雷诺平均Navier-Stokes（RANS）方程几乎完全用于模拟飞机配置上的湍流流量。 对于大多数问题，采用RANS方程的薄层近似来将问题缩小到可管理的大小。 但这些简化造成了沉重的成分；我们现在需要湍流型号！ 近年来出现了各种湍流模型，从相对简单的代数模型到更复杂的雷诺应力模型。 这些模型已被应用于使用多块结构网格方法产生令人印象深刻的结果，包括修补[36]和覆盖[37]。 用于湍流流分析的非结构化网格方法也在加速发展[38]。 方法范围从各向异性四面体网格到混合网格，结合了靠近固体表面的棱镜和其他地方的四面体或笛卡尔网格。 总的来说，模拟湍流流量的经验是好坏参半的。 在使用简单模型的湍流流模拟中取得了很多成功，而使用更复杂的模型也有很多失败。 完善现有模型和开发新的、改进的模型的尝试继续有增无增无阻[39]。 还致力于开发层向湍流过渡的模型，这是精确的黏性流动模拟至关重要的领域。 虽然正在取得进展，但CFD从业者的困境非常明显。 在飞机设计中存在许多空气动力学问题，黏性效应占主导地位，只有透过解决RANS方程才能正确模拟它们。 内部流动问题（入口、扩散器、喷嘴等）和高升系统（多元件

438 RAJ和解决方案的可靠性继续受制于湍流型号的不足。 图6显示了一些流行模型在机翼上准确模拟冲击/边界层相互作用方面的局限性。 正如Marvin和Huang [39]所指出的，只要压力梯度小且激波弱，涡黏性模型就可以自信地用于预测压力分布；预测皮肤摩擦和热传递到工程精度需要仔细关注网格加密和自由流边界条件的影响。 在作者看来，一个简单的通用湍流模型的前景相当暗淡；将湍流的复杂性质捕捉到一个具有少数自由参数的模型确实是一个长途。 继续使用IV级代码的最佳理由--尽管它们有局限性--可能来自Bradshaw [40]：“......我们无法计算工程准确性的所有工程兴趣流。 然而，最好的现代方法允许几乎所有的流量的计算精度都高于最明智的猜测，这意味着即使这些方法不能取代实验，它们也是真正有用的。” 如果不提及Digital Physics™技术[41]，IV级代码的简要概述就不完整，该技术在不诉诸传统的湍流建模的情况下模拟黏性流！ 据说该技术比RANS代码具有根本优势，因为它没有离散化的人工。 根据成功将技术扩展到可压缩流和附加验证/演示，Digital Physics™可以为RANS方法提供有吸引力的替代方案。 摘要。-上述讨论清楚地指出了过去三十年取得的巨大进展。 可以引用更多示例，以演示如何应用 CFD 模拟复杂配置的流程。 迄今为止，众多的应用帮助我们学会了利用CFD和风洞的互补性。 当智能地与风洞结合使用时，CFD可以更好地理解流场和空气动力学特性，从而为

CFD AT A CROSSROADS 439设计降低风险和提高质量。 然而，风洞将继续成为为整个设计包络生成空气动力学数据和最终设计验证的必要条件。 未来几年，CFD能力的进一步增强将扩大其在飞机设计中的作用。 然而，要应对未来飞机设计的挑战，需要的不仅仅是下一节所讨论的渐进式进展。 23.5 CFD在哪里？ 未来二十年 在本节中，我们研究了CFD面临的一些挑战，以及为了更有效地应对挑战而需要追求的途径。 不断变化的设计环境。-自20世纪80年代中后期以来，美国航空航天业一直严重关注在竞争日益激烈的全球市场中保持其领导地位。 对于一般的航空航天业，特别是军用飞机制造商来说，飞机设计的挑战已经从任何成本更高的性能急剧转变为以负担得起的成本获得高质量。 挑战不是质量或可负担性，而是两者同时。 工业界和政府联合进行的许多研究得出了一个被广泛接受的结论，即工业界必须过渡到综合产品和流程开发（IPPD）环境。 IPPD的特点是整合了产品开发的各个方面，包括设计、营销、制造和产品支持。 IPPD方法依赖于从一开始就考虑所有要求和约束，而不是在后期更改设计以促进制造或满足产品支持需求。 因此，可以尽早做出适当的权衡，以后对设计更改的需求会大大减少。 结果是提高了质量，提高了生产力。 在IPPD的背景下，设计被视为一个综合的多学科过程。 它不应该被解释为“自动化流程”。 人类的创造力和独特的合成能力是无可替代的。 其目的是通过迅速为设计团队提供在早期阶段做出更明智的决策所需的数据来缩短设计周期时间，从而减轻传统设计流程的严重缺点。 历史数据库历来是早期设计研究的核心。 然而，当新的配置和相关流程（制造、支持等）与旧的配置和相关流程大相去甚远时，此类数据库的价值不大--这种情况在军用飞机设计中越来越常见。 新设计通常由越来越严格的任务要求驱动。 正在对实施真正集成IPPD设计环境的具有成本效益的方法进行大量研究。 例如，DeLaurentis等人[42]提出的方法使用统计技术，通过代表来自高级的大量知识，在搜索设计空间时更加灵活

440个RAJ计算代码（或物理实验）通过响应曲面方程。 重点是提供更多有关设计的知识，以影响早期设计阶段的决策。 IPPD设计流程成功的关键要求之一是快速、准确和具有成本效益的方式为每个贡献学科以及多学科互动生成数据。 CFD的关键作用。-在作者看来，CFD将在IPPD设计过程中发挥关键作用。 CFD和其他学科的高级计算方法（例如结构、控制、推进、声学）的结合提供了模拟多学科交互的唯一实用手段，这是IPPD设计过程的基石。 风洞只是不适合以及时和具有成本效益的方式产生所需的信息。 此外，CFD提供了一种强大的计算定义和/或改进几何形状的手段，以产生特定的流动特性，同时满足规定的约束；这在风洞中是不切实际的。 但为了满足IPPD设计过程的期望，CFD的有效性必须大幅提高。 尽管人们非常重视提高CFD的有效性，但在周转时间、准确性和成本方面仍面临许多挑战。 周转时间。-周转时间受到CFD代码级别的强烈影响，即I级、II级、III级或IV级。 较低级别的代码提供快速周转，而更高级别的代码需要更长的时间。 由于低级代码在建模许多感兴趣的流程方面的局限性（见第23.4节），绝对必须使高级代码的周转时间与低级代码的周转时间相当。 预计到本世纪之交，使用非结构化网格RANS代码将在不到24小时内对完整飞机配置进行稳定的空气动力学分析。 然而，即使是这种速度也不够。 当在单个模型上进行大量运行以产生满足设计需求所需的空气动力学数据时，总时间跨度达到了不可接受的比例。 当配置几何形状也发生变化，并且必须对每个变体进行多重分析时，就会出现更具挑战性的情况。 这正是IPPD设计环境的要求！ 还必须指出的是，迄今为止，更高层次的方法在及时和具有成本效益的方式提供不稳定的空气动力学数据方面相当无效；大多数设计项目目前都依赖于线性势能方法。 CFD社区的挑战是显而易见的：开发适当的技术，并以使每次分析的周转时间只有几分钟的方式整合它们。 潜在的赋能技术列表包括计算机辅助设计（CAD）系统的简化接口；标准数据交换协议；“自动化”网格生成；流动求解器软件的并行处理；以及用于数据分析和管理的智能系统，仅举几例。 正在进行的研究和开发给CFD从业人员带来了相当大的希望和鼓励，即在未来十到十五年内将实现这一目标。

CFD AT A CROSSROADS 441 准确性。-计算解决方案的准确性是最大的担忧之一。 正如Bangert等人[9]所指出的，由于当前CFD代码在模拟黏性效应方面的局限性，F-22设计团队严重依赖风洞数据，特别是当应用于完整的飞机几何形状和全速、高度和机动飞行包络时。 为了支持F-22的开发，已经进行了4万多小时的风洞测试。 解决方案的准确性有两个组成部分：数值和物理。 如果一个解决方案对网格的变化以及与算法相关的数值参数几乎没有灵敏度，则该解决方案可以被认为是数值准确的。 （假设有关代码已经过验证，其数值公式在解决控制方程时是否充分。） 目前，系统参数研究是估计网格分辨率和耗散和色散等数值参数影响的唯一手段。 典型设计工作的时间表和成本限制不允许进行广泛的调查来确定“最佳”网格和数值参数。 CFD分析师最终依赖以前的经验和专业知识。 随着我们进入间隔更远的飞机计划减少的时代，严重依赖过去的经验可能变得太危险了。 我们需要的是量化准确性水平的内置手段。 当CFD方法与其他学科的方法相结合以模拟多学科互动时，对解决方案的准确性水平和相关误差范围的评估更加重要。 诚然，误差估计是一个难题，但如果要在IPPD环境中有效利用CFD，迫切需要解决方案。 采用求解自适应技术是最小化数值误差的一种方法。 图7显示了解决方案对网格自适应的灵敏度的示例。 翼体中心线尾部模型的升力系数值，使用两个非结构化网格欧拉代码计算，相互比较，并与0.4 马赫数的实验数据进行比较。 该模型有一个鳍状的前身、锋利的翅膀和垂直鳍。 与目前没有自适应网格功能的四面体网格USM3D代码相比，使用解决方案自适应的笛卡尔网格SPLITFLOW代码预测了30度攻角的升力值要大得多。 差异很可能是由于与USM3D代码相比，SPLITFLOW代码提供的涡流分辨率更好。 在更高的攻角下，前缘涡流从配置移动到USM3D网格太稀疏，无法充分解析它们的区域。 这些分析的更多细节可以在参考资料43和44中找到。

442 RAJ 即使代码产生了一个数值上准确的解决方案，确定解决方案与实际流量的叠加程度并非易事--这是衡量其物理准确性的指标。 迄今为止，CFD社区倡导并进行了广泛的“验证”练习，以生成计算数据与实验数据相关性，并用它们来证实声称的物理准确性水平。 然而，这些练习对数据库的扩散的贡献更大，而不是提高为飞机设计提供高精度空气动力学数据所需的知识库和信心水平。 与设计团队考虑的几何形状和流动条件有很大不同的几何形状和流动条件的广泛相关性没有什么价值。 这种情况与军用飞机项目特别相关，因为新设计通常与旧设计有很大不同（CFD可能存在相关性）。 传统的代码验证方法充满了困难。 在代码被声明为完全验证之前，我们必须考虑多少个测试用例，每个测试用例的流程条件组合，以及哪些湍流模型？ 使用任何合理的测试案例集和流程条件的运行矩阵迅速成长为一项不朽的任务。 即使我们假设有足够的资源和高质量的测量数据来执行这样的任务，我们也会逆着技术动态的潮流。 硬件、数值算法以及湍流建模和过渡的快速进步助长了一个代码永远不会完全“完成”的环境。 有时变化是名义上的，很多时候不是。 对计划的任何合理成本/效益评估，如果计划分配大量资源来验证第二天可能被“新的和改进的”方法取代的代码，都不支持传统的验证方法。 一个更有价值的方法可能是选择一套标准的测试用例，即几何形状和空气动力学数据库。 数据库可能包括理论（或数据数值）解决方案和/或已知精度的测量数据。 代码开发人员可以使用标准测试案例对其代码进行全面的参数化研究，以制定电网分布、数值参数等指南。 标准边界层和剪切流测试用例可用于确定网格点、数值耗散参数、湍流模型等的数量和分布，从而产生可接受的精度解决方案。 例如，Dacles-Mariani等人[45]最近对近场机翼-翼尖涡流动的数值/实验研究发现，其结构网格代码需要图7。 结构化网格欧拉解对网格自适应的灵敏度

CFD AT A CROSSROADS 443 涡流核心至少15点，网格单元长宽比为1和五阶方案，以充分解决流场。 然后，用户将有适当的指南来“正确”设置他们的计算问题。 因此，在设计环境中生成有意义的数据的机会将会提高。 拟议的方法还将为比较不同代码的整体有效性提供一个健全的基础。 随着越来越多的第三方供应商提供CFD软件，采用拟议方法的需求更大。 大涡流模拟（LES）和直接数值模拟（DNS）的持续进步为开发CFD能力提供了非常有吸引力的选择，这种能力可以真正预测，不像RANS方法将继续受到湍流和过渡模型的不确定性的影响。 然而，如图8所示，即使是对雷诺数航班完整飞机分析所需的计算资源的粗略估计也令人震惊。 由于Re3计算工作障碍，DNS似乎对几个图8来说不切实际。 CFD计算挑战数十年。 然而，如果计算机内存和速度增加约四个数量级，LES可以在大约20年内变得实用。 考虑到自1975年以来，速度和内存每7年增加十倍，可以想象到2015年至2020年所需的计算资源将可用。 当务之急是将更多精力放在几个起搏项目上，如子网格比例建模；数值算法；边界条件；网格生成；以及分析、可视化和管理大量数据的工具。 成本。-与CFD应用程序相关的成本包括人工和计算。 目前，人工费用主要与前期和后期处理步骤有关。 对于更高级别的CFD代码，人工费用仍然超出可接受的范围。 例如，使用结构化网格方法需要几个人周来产生第一个解决方案，而理想值则接近几个人小时。 一般来说，非结构化网格方法，特别是笛卡尔网格，在减少努力水平方面似乎很有希望。 在网格生成方法和CAD系统之间开发简化接口方面取得的进展对于减少几何获取/建模时间至关重要。 这些改进也将有助于以低廉的方式评估设计变化。 事实上，减少工作时间所需的技术与减少周转时间所需的技术基本相同。

444 RAJ计算成本主要与运行流动求解器有关，在某些情况下可能包括网格生成。 通常需要具有高处理速度和大内存的计算机来按计划生成所需数量的数据。 IPPD设计过程很可能需要在几天内而不是几个月内生成全部数据。 计算成本以适应这种时间框架，不得太大，使产品开发总成本增加而不是减少。 因此，整个硬件和软件系统的成本和计算效率是未来有效使用CFD的非常重要的考虑因素。 提高计算效率和降低成本的策略必须是未来所有CFD开发和应用计划的组成部分。 市场上有各种竞争的计算机架构，从单个工作站到共享内存处理器再到分布式内存处理器；未来几年无疑将推出新的架构。 匹配软件和硬件特性将是未来有效开发和应用CFD更具挑战性的方面之一。 23.6 结论 在本文中，从行业角度研究了CFD的现状，该行业将CFD视为一种产生空气动力学数据的工程工具，以支持飞机设计需求。 CFD功能已从六十年代末和七十年代初的简化几何、简化物理模型发展到今天常用的复杂几何、现实物理模型。 这些进步的动机是需要为风洞测试提供具有成本效益的替代方案。 在此期间，测试成本一直在稳步增长，而计算成本（以每个浮点操作的美元计算）正在迅速下降。 CFD现在已经进入了一个阶段，即没有飞机项目完全依赖风洞的所有空气动力学数据。 然而，CFD不是风洞的替代品，只是一个完美的补充。 目前，人们非常重视提高更先进的欧拉/N-S方法的整体有效性。 正在采用策略来缩短周转时间、降低成本并提高准确性。 这些策略包括简化与CAD系统的接口，提高数值算法的效率，开发用于数据处理的智能工具，以及改进湍流模型，以更准确地表示复杂流的物理性。 然而，渐进式变化不足以应对未来的挑战。 航空航天业务的最新变化引发了飞机设计流程的巨大变化。 预计不断发展的过程将更加一体化，重点是尽早向设计团队提供必要的数据，以便在早期阶段做出更明智的决定。 它将要求

CFD AT A CROSSROADS 445，如果我们继续走过去三十年所走的道路，对CFD的期望比我们能够交付的要多得多。 我们需要真正大幅减少周转时间和成本。 生产可靠和可接受的准确性的解决方案的问题需要正面解决。 在选择构建未来CFD能力的“正确”技术时，评估其对周期时间、准确性和成本的影响至关重要。 CFD社区今天面临的挑战是以使CFD充分响应未来飞机设计需求的方式引导他们的努力和资源。 将先进的CFD方法纳入设计流程将带来许多好处。 使用提供快速周转能力的先进方法将缩短设计周期时间。 然后，设计团队可以在典型产品开发工作的时间表和成本限制范围内探索比目前可行的更广泛的替代方案。 快速、准确和经济实惠的方法将提高生产力，并减少支援设计资料需求所需的昂贵测试数量。 使用高级方法也可能减少设计闭合所需的周期数。 设计团队将能够进行必要的广泛权衡，以指导配置的演变朝着最小化获取和生命周期成本的方向发展。 提高对组件交互的了解将允许尽早进行设计更改，从而降低风险并增加满足所有客户要求的概率。 参考资料1。 Landahl，M，“CFD和湍流”，ICAS-90-0.1，瑞典斯德哥尔摩，1990年9月。 2. Rai，M.，Gatski，T.和Erlbacher，G.，“空间进化可压缩湍流边界层的直接模拟”，AIAA论文95-0583，1995年1月。 3. Moin，P.和Mahesh，K.，“直接数值模拟：湍流研究中的工具”，CTR手稿166，斯坦福大学湍流研究中心，斯坦福，加利福尼亚州，1997年4月。 4. 米兰达，L.R.，“计算空气动力学在飞机设计中的应用”，飞机杂志21，1984年6月，pp. 355-370。 5. 汉考克，G.J.，“空气动力学：计算机的作用”，航空杂志，1985年8月/9月，pp. 269-279。 6. Miranda，L.R.，“超音速和战斗机：CFD的挑战和机遇”，NASA CP 3020，卷。 I，第1部分，1988年4月，pp. 153-173。 7. Cosner，R.R.，“CFD分析的航空航天应用问题”，AIAA论文94-0464，1994年1月。 8. Rubbert，P.E.，“CFD和不断变化的飞机设计世界”，ICAS-94-0.2，1994年9月。

446 RAJ 9。 Bangert，L.H.，Johnston，C.E.和Schoop，M.J.，“F-22设计中的CFD应用”，AIAA论文93-3055，1993年7月。 10. Raj，P.，“在航空航天设计中有效使用CFD的要求”，美国宇航局CP 3291，1995年5月，第15-28页。 11. Miranda，L.R.，Elliott，R.D.和Baker，W.M.，“亚音速和超声速流动应用的广义涡旋晶格方法”，NASA-CR 2865，1977年12月。 12. Youngren，H.H.，Bouchard，E.E.，Coopersmith，R.M.和Miranda，L.R.，「面板方法配方的比较及其对QUADPAN（一种先进的低阶面板方法）发展的影响」，AIAA论文83-1827，1983年7月。 13. Magnus，A.E.，Ehlers，F.E.和Epton，M.A.，“PANAIR 使用高阶面板方法预测亚音速或超音速线性势流的任意配置的计算机程序”，美国宇航局CR-3251，1980年4月。 14. Johnston，C.E.，Youngren，H.H.和Sikora，J.S.，“高级低阶面板方法的工程应用”，SAE文件851793，1985年10月。 15. Giesing，J.P.，Kalman，T.P.和Rodden，W.P.，“多管干扰机翼和机体的亚音速稳定和振荡空气动力学”，飞机杂志9，10月。 1972年，pp. 693-702。 16. Chen, P.C., Lee, H.W., and Liu, D.D.,“包括唤醒效应在内的外部存储的身体和机翼的非稳态亚音速空气动力学》，《飞机杂志》30，1993年5月，pp. 618-628。 17. Chen，P.C.和Liu，D.D.，“包括外部存储在内的任意机翼体配置的非稳超音速计算”，《飞机杂志》27日，2月。 1990年，pp. 108-116。 18. 《超音速空气动力学》，航天学和航空学进展，卷。 81，美国航空航天研究所，华盛顿特区，1982年，尼克松，D。 （Ed.） 19. Murman，E.M.和Cole，J.D.，“平面稳定跨声速流动的计算”，AIAA杂志9，1971年1月，pp. 114-121。 20. Boppe，C，“机翼机身配置的跨声速流场分析”，美国宇航局CR 3243，1980年。 21. Caughey，D.A.和Jameson，A.，“机翼-机身组合有限体积计算的进展”，AIAA杂志18日，11月。 1980年，第1281-1288页。 22. Hicks，R.M.和Henne，P.A.，“通过数值优化的翼设计”，AIAA Paper 77-1247，1977年。 23. Borland，C.J.和Rizetta，D.P.，“用于空气弹性应用的超音速非稳态空气动力学，卷。 I，XTRANS3S的技术开发摘要，”AFWAL-TR-80-3107，1982年。

CFD在十字路口447 24。 Samant, S.S., Bussoletti, J.E., Johnson, F.T., Burkhart, R.H., Everson, B.L., Melvin, R.G., Young, D.P., Erickson, L.L., Madson, M.D., and Woo, A.C., “TRANAIR：任意配置跨声速分析的计算机代码”，AIAA论文87-0034，1987年1月。 25. Jameson，A.，Schmidt，W.和Turkel，E.，“使用Runge-Kutta时间步进方案的有限体积方法对Euler 方程的数值解”，AIAA论文81-1259，1981年6月。 26. Rizzi，A.，“计算跨声速流动关于翼体组合的减振欧拉方程法”，AIAA杂志20日，10月。 1982年，pp. 1321-1328。 27. “应用计算空气动力学”，航天学和航空学进步，卷。 125，美国航空航天研究所，华盛顿特区，1990年，Henne，P.A. （Ed.） 28. “基于Euler 方程的计算空气动力学”，AGARD-AG-325，1994年9月，Sloof，J.W.和Schmidt，W. （编辑。） 29. Raj，P.和Harris，B.W.，“使用欧拉方法的表面蒸腾进行具有成本效益的空气动力学分析”，AIAA论文93-3506，1993年8月。 30. Goble，B.D.，King，S.，Terry，J.和Schoop，M.，“使用3D非稳欧拉/Navier-Stokes代码的入口锤击分析”，AIAA论文96-2547，1996年7月。 31. Tidd，D.M.，Strash，D.J.，Epstein，B.，Luntz，A.，Nachshon，A.和Rubin，T.，“高效3D多网格尤拉方法（MGAERO）在完成飞机配置中的应用”，AIAA论文91-3236-CP，1991年。 32. Karman，S.，“SPLITFLOW：复杂几何的3D非结构笛卡尔/棱镜网格CFD代码”，AIAA论文95-0343，1995年1月。 33. Guruswamy，G.P.，“使用Euler 方程的机翼的非稳态空气动力学和空气弹性计算”，AIAA杂志28，1990年，pp. 461-469。 34. 坎贝尔，R.L.和史密斯，洛杉矶，“跨声速翼型和机翼设计的混合算法”，AIAA论文87-2552，1987年。 35. Reuther，J.J.，Jameson，A.，Alonso，J.J.，Rimlinger，M.J.和Saunders，D.，“使用辅助配方和并行计算机的受限多点气动外形优化”，AIAA论文97-0103，1997年1月。 36. Ghaffari，F.，“F/A-18飞机的Navier-Stokes、飞行和风洞流量分析”，美国宇航局技术论文3478，1994年12月。 37. Gomez，R.J.和Ma，E.C.，“航天飞机运载火箭大型嵌套网格系统的验证”，AIAA文件94-1859-CP，1994年。 38. Venkatakrishnan，V.，“非结构化网格流求解器的观点”，AIAA论文95-0667，1995年。 39. Marvin，J.G.和Huang，G.P.，“湍流建模-进展和未来展望”，Proc。 第15届流体力学数值方法国际会议，1996年6月，加利福尼亚州蒙特雷。

448 RAJ 40。 Bradshaw，P.，“湍流次级流动”，流体力学年度评论，19，1987年，pp. 53-74。 41. 《数字物理学》，《机器设计》，第66卷，12月。 1994年，pp. 96-100。 42. DeLaurentis，D.，Mavris，D.N.和Schrage，D.P.，「使用统计方法初步飞机设计中的系统合成」，ICAS 96-3.4.4，1996年9月8日至13日。 43. Kinard，T.A.，Finley，D.B.和Karman，S.L.，“使用非结构化尤拉分析预测涡流主导流场的可压缩性效应”，AIAA论文96-2499，1996年6月。 44. Raj, P., Kinard, T.A. and Vermeersch, S.A., “使用非结构网格欧拉方法的自来模拟”，ICAS-96-1.4.5，1996年9月。 45. Dacles-Mariani，J.，Zilliac，G.G.，Chow，J.S.和Bradshaw，P.，“近场翼尖涡的数值/实验研究”，AIAA杂志33，9月。 1995年，pp. 1561-1568。

\clearpage
\chapter{航空航天工程2000：综合实践课程}
航空航天工程2000：综合实践课程A。 Richard seebass的和Lee D. peterson的24.1“有史以来最大的礼物......” “纽约市的一个基金会对工程教育状况感到沮丧，承诺投入2亿美元从头开始建造一所学院”[41]。 这似乎代表了高度的挫败感。 劳伦斯·W。 米拉斯，F。 W。 《纽约时报》的这篇文章援引奥林基金会主席的话说：“我们完全致力于让这个机构启动并运行，并且拥有3亿至4亿美元的强大捐赠，...... 我们将强调跨学科方法、低学生与教师的比例、指导和实践体验。” 太棒了！ 当然，这种方法并不新鲜，但这种礼物反映了人们对工程教育没有满足国家需求的担忧。 这所新的工程学院将是一项值得注意的努力，也许是哈维·穆德的马萨诸塞州类似物，也许是独一无二的。 这篇文章所暗示的批评是否合理？ 麻省理工学院的航空航天课程只是一个例子，长期以来在其大二时提供了一种综合方法[34]，当Earll Murman担任主席时，其航空航天课程在许多方面得到了进一步加强，包括采用这种系统方法来学习fdl课程，并强调重要的“隐式课程”[21]。 另一个例子是，在科罗拉多大学，我们计划、筹集资金、建造，并于1997年4月24日奉献了我们的综合教学实验室。 '航空航天工程科学，科罗拉多大学，科罗拉多州博尔德，80309-0429。 计算流体动力学的前沿--1998年。 编辑：David A. Caughey和Mohamed M. 哈菲兹。 01998 世界科学

450 SEEBASS \& PETERSON工程教育已经研究了一个多世纪，包括《威肯登报告》[52]、《哈蒙德报告》[53-54]、《格林特报告》[8]、《哈达德和皮斯特报告》[44-45]以及ASEE的《不断变化的世界工程教育》[9]。 正如Monteith提醒我们的那样，Burr在1893年所说的话至今仍然是真实的：“青年工程师理想教育的首要和基本要求，即哲学和艺术方面的广泛、自由教育，[是]纯专业培训的先例。... 主要目的应该是培养人类品质，使工程师能够满足人类和物质。” [16，38]。 关于这个主题，另请参阅Florman和Moulakis的书籍[32，39]。 工程教育不是我们最严重的教育问题；K-12教育才是。 在这个主题上，有不同的观点；一个观点在几个方面与我们计划的课程平行，是阿德勒的观点[3]。 一般来说，大学教育有批评者和评估者。 也许对大多数声誉的评估是由于Boyer [15]。 这个国家在20世纪70年代和80年代面临的经济困难（由于经济强劲、低失业率和通胀率以及日元疲软，现在有些处于后台）在很大程度上仍然存在。 我们现在是最大的债务国，我们每月进口的商品继续比出口多7100亿美元（12个月移动平均）。 航空航天公司感到自豪的是，在制成品中，航空航天产品提供了最大的年度贸易顺差，1992年超过310亿美元[5]。 我们的经济困难似乎更多的源于我们的管理实践，而不是工程教育的失败，尽管后者不能被免除贡献[26-28，33，35]。 日本人的成功至少部分归功于他们从美国工程师和统计物理学家W那里学到的东西。 Edwards Deming。 在1980年6月NBC新闻白皮书之后，「如果日本能...... 为什么我们不能？”，这个国家的企业领导人开始学习德明必须教给他们的东西。 其他许多人，其中第一作者，也从这个人那里学到了很多。 24.2“尽最大的努力是不行的！” 在20世纪80年代，工程系主任和系主任以及美国工程教育学会（ASEE）的工程系长理事会开始担心，虽然我们可能培养出非常优秀的工程师，但他们可能不够好，无法维持我们的经济福祉。 我们开始问，正如其他人后来问的那样[38]，“我们应该如何培训我们的工程师，使他们成为世界上最好的？” 我们当时不知道，现在也不知道。 然而，我们担心的是，工程和技术认证委员会（ABET）工程认证委员会的认证学位标准并没有鼓励必要的改变。 只要我们满足ABET的详细标准，认证通常就得到了保证。 一些学校表现出了偏离ABET标准的想象力，使他们未来的认证面临风险。

航空航天课程2000 451 ABET标准代表了确保这个国家获得令人满意的工程教育的“最大努力”。 但令人满意是不足以满足的。 作为W。 Edwards Deming一向说：“尽最大的努力是行不了的！” 而且他们不会。 有些人认为他们知道我们需要做什么；其他人只是知道我们需要改进。 在各种工程院长团体（Big Ten Plus、工程院长小组、ASEE工程院长委员会）的压力下，工程认证委员会重新思考了其标准，并制定了新的标准，这值得称赞。 这些将于2001年生效，将于1998年在自愿基础上实施，并通过1996-97年的试点认证进行研究和改进[1-2]。 新标准为项目提供了实施实验和改革所需的灵活性，但它们要求我们更多。 应用得当，它们提供了持续改进的可能性，以实现最好的工程教育。 24.3 我们的客户和其他人怎么说？ 行业经常认为自己是我们的客户。 他们雇佣了我们的学生。 但试着把他们雇用的人的教育账单发给他们。 您很快就会知道他们不是我们的客户。 我们的学生和他们的父母是我们的客户。 对于公共教育来说，社会也是我们的客户。 因此，我们应该询问我们的学生关于他们的教育。 在这里，我们有很多选择需要考虑，包括那些对美国不满意的客户终止学习的学生，可以说[51]。 但是我们的客户知道我们应该做什么吗？ 不，他们没有。 但他们肯定希望我们知道。 为此，也许我们应该听取那些领导主要行业但考虑社会和学生需求以及他们自己需求的人的意见；我们应该听取那些对这一主题有深入思考的教育工作者的意见[7、10-11、13、25、38、40、56]，以及我们的毕业生。 在学术界，我们早就认识到，企业招聘人员和企业领导人对他们在毕业生身上寻求的东西有不同的看法。 因此，令人鼓舞的是，行业现在渴望全面看待我们毕业生的理想属性[14，36-37]。 除了我们长期且在某种程度上成功解决的问题（例如少数民族和妇女更多地参与工程教育，需要控制我们的成本，以及我们的毕业生需要了解他们必须开始终身学习），对于我们需要做什么，大家普遍同意。 最重要的是，我们必须继续做我们现在做得很好的事情--强调基本面。 此外，我们的学生需要培养更好的团队和沟通技能，有能力整合不同学科，并在人文、全球和政治背景下理解工程。 尤其是，我们的学生需要更多的实践经验，更强的整合理论能力，

452 SEEBASS \& PETERSON分析、实验以及参与真实设备的设计、制造和测试。 有些人还觉得是时候攻读硕士学位，作为工程学的第一个专业学位了[见，例如，10]。 这是一个古老而合理的想法，其时代可能已经到来。 这种形式上的改变的后果是深远的，长期以来一直受到工业界和学术界的抵制。 工业抵制，因为它使工程教育更容易获得，工程师更昂贵；学术界抵制，因为它意味着更长的学位时间，因此可能更少的学生。 如果没有为此目的的正式改变，无论如何，这在很大程度上是通过更加强调学生通过硕士学位继续他们的大学学习，以及通过他们在职业生涯早期追求硕士学位而发生的。 24.4 综合教学实验室 在20世纪90年代初，我们学院的教师制定了第二个五年战略计划，并独立规划了学院的百年纪念活动。 从我们的战略规划中可以看出，我们需要做很多事情来改善我们的本科教育，为了资助这种改进，我们必须让本科课程改革成为我们百年筹款活动的基石[22-24]。 我们特别关注本科生的实践、实验室和设计经验有限。 当然，我们也关心提高我们吸引的学生的素质，尽管我们收取了非常可观的非居民学费，但超过三分之一的一年级学生来自州外。 这导致1992年，在机械工程的Larry Carlson和土木工程的Jacquelyn Sullivan的指导下任命了一个教师团队：“......开创一个多学科学习环境，将理论与实践相结合，并促进创造性、以团队为导向的解决问题的技能。” 为了实现这一变化，我们计划、建造并于1997年4月24日投入了价值1100万美元的综合教学实验室（ITLL）[55]。 该设施由34,400平方英尺组成，配备了我们自己设计的30个实验室站和约600万美元的计算机、仪器和网络（主要来自惠普，值得一提的是ITLL众多企业支持者中最大的），该建筑及其系统作为活生生的实验室。

航空航天课程 2000 453 24.4.1 ITLL课程组成部分 ITLL的课程组成部分如上面的“ITLL轮”所示。 它们是第一年的项目课程、讲座演示、实践作业、实验模块、计算机模拟、基于学科的实验室课程，以及构成这个轮子辐条的跨学科设计项目。 这些都由广泛的计算机网络提供服务。 ITLL的主要组成部分是：两个大型（4000平方英尺）实验室广场，每个有15个实验室站，总共可容纳60到90名学生在两到三人小组中工作，相邻的分组讨论区；一个带25个UNIX工作站的模拟实验室；四个ITLL高级项目设计工作室；两个一年级设计工作室；一个制造中心，一个电子中心，一个主动学习中心，两个学术“客厅”（采用华盛顿大学新物理大楼的想法）；十个带个人电脑和互联网连接的小组学习室；一个动能雕塑 展示许多科学和工程原理的画廊；以及一个学生休息室。 第一年的项目

454 SEEBASS \& PETERSON教室、主动学习中心和两个分组区域是配备网络访问和视频投影仪的“智能空间”。 这座建筑本身的仪器仪表很好。 所有组件共享通用的数据采集和分析软件。 LabStations LabStations由第二作者领导的团队设计，可容纳独立的实验模块，这些模块是独立于序列的实验，需要最少的监督，并且适合开放式探索。 每个LabStation都包含两台HP Pentium PC，配有Windows NT、LabView、MatLab、Microsoft Office、数据采集和信号调节硬件，连接到100 MB/秒的网络。 它们为惠普台式仪表（由示波器、功能发生器和万用表组成）提供接口，以及各种级别的交流/直流电源和压缩空气供应。 推车上的实验模块很容易连接到LabStations，这允许每天对LabPlazas进行多次重新编程[17]。 模块模块是移动实验，可以与LabStations一起使用，用于课堂演示和动手作业。 迄今为止，已经开发了35个实验。 它们涵盖了测量和仪器、电子和微处理器、动态系统和控制、流体力学、传热以及结构和材料。 动手作业ITLL为教师提供了向学生进行“动手”作业的机会，他们可以使用简单的设备和材料来展示他们在课堂上学到的概念。 作为实验室的建筑，建筑本身是一个实验室，有部分暴露和仪器系统，包括各种结构元素、供暖和通风系统以及用于太阳能增益和热量损失的玻璃仪表[17]。 第一年工程项目课程 作为全校系课程改革的一部分，我们推出了许多新课程。 充分利用ITLL的一个是第一年的项目课程。 它于1994年以原型为基础推出，并不断改进。 早期结果表明，进入的保留率提高了50\%

航空航天课程2000 455名学生参加本课程的工程研究。 我们的航空航天课程希望航空航天专业学生学习这门课程。 然而，我们全校的无过错第一年意味着我们不需要这门课程。 在这里，我们以马里兰州第一年设计经验为指导[25，48-49]。 该课程以破冰“神秘人工制品”挑战开始，继续进行设计项目，并以逆向工程项目结束。 它向学生介绍团队合作、口头和书面沟通方法、项目管理以及低级计算技能，这些技能是隐性工程课程的一部分。 它包括关于团队动态、学习风格和小组沟通、AutoCad、基本电子产品和手工工具使用的研讨会。 我们渴望利用这个项目课程来鼓励我们的学生利用他们不断增长的微积分、化学和物理学知识，并帮助他们理解这种材料是工程学不可或缺的，而不是无关的。 马里兰大学的类似课程证明了这是可能的。 大一的其他值得注意的课程变化包括Drexel的E4课程（将基础、设计和实验室与工程研究所需的数学和科学相结合[20，46]），以及Rose-Hulman的综合第一年[50]。 关于新生设计的主题，另见Ercolano [30]。 菲茨霍恩审查了第一年工程课程的趋势[31]。 24.5 航空航天工程科学课程 2000 航空航天业经历了充满挑战的时期。 以固定美元和行业就业的国防支出是十年前的一半。 然而，今天工作机会比比皆是。 我们正在从军事和美国宇航局太空计划过渡到商业太空计划的丰富性，从我们在太空中的反复存在到我们（事实上，世界）在那里的永久存在，从大型、昂贵和长时间延迟的空间科学实验到在更多大学工程参与下频繁进入太空。 商用飞机和燃气轮机发动机的产量持续增长。 新的大型交通工具和超音速交通工具都在研究中，通用航空正在复兴。 我们的行业每年为国家提供超过300亿美元的航空航天贸易顺差。 我们必须尽自己的一份力量来确保它的持续健康。 结合ITLL的规划，航空航天工程科学课程和教学委员会在第二作者的指导下，彻底审查了我们的课程。 在1995-1996学年，预计1997年秋季ITLL将满员，并打算充分利用这一设施，课程和教学委员会提出了一个全新的大二和不同的上级课程。 今年秋天，我们将为我们的毕业班实施新的大二课程

456 SEEBASS \& PETERSON 2000。 我们现在正在为这个课程和连续的课程定义上级课程。 24.5.1 目标 本课程改革的目标是提高我们为本科生提供的教育质量，以确保他们在工业、政府和学术界拥有最成功的职业生涯。 为了实现这一目标，我们得出结论，我们必须：建立一个核心课程，将材料整合到这个核心中，使课程与应用程序相关，使其具有体验性，即“动手”，将沟通和团队合作技能的发展纳入所有课程中，在上级部门提供更多的课程选择，实施持续改进程序，我们还认识到，随着作为工程师有效贡献所需的知识不断增加，MS/ME学位正在成为事实上的专业学位。 为了鼓励我们前四分之一的学生留在我们这里攻读一学年的硕士或ME学位，我们正在实施一个为五年的BS/MS学位课程。 我们之前对研究生课程的改革，加上我们的本科课程改革，使那些在大四时有足够的能力学习两门研究生课程的学生有可能。 24.5.2 为了完成我们的教育改革，我们必须：为在校学生和新生划定过渡，从1997-98年到1999-2000年制定教学时间表，确定并启动所需的教师发展和培训，将我们的课程计划纳入行业设计审查，与学生密切合作，以确保他们重视我们代表他们做出的改变，我们现在必须完成对上级部门的定义，以实现我们的教育目标，并满足AIAA的航空航天项目特定标准[6]。 我们还必须为我们所有的教育努力开发评估方法，以确保在德明实施舍哈特循环后，我们可以不断改进它们[26]。

航空航天课程 2000 457 24.5.3 环境 为了在这项努力中取得成功，我们必须确保我们的教师有一个环境：为课程和课程开发提供时间，为学生互动和咨询提供额外的时间，为教学支持增加额外的资源，为教师发展提供时间（例如，学习LabView、Mathematica、MatLab等） 确保公平的教学责任，确保进一步改进我们的研究和研究生课程的时间，奖励他们对教育改进的参与 24.5.4 新大二 改革 我们大二在很大程度上来自麻省理工学院的大二核心课程“统一工程”[34]。 我们的本科生虽然对于州立大学来说非常有能力，但并不是麻省理工学院的统一高质量，所以我们的经验可能与麻省理工学院不同。 他们在综合大二的成功已经持续了很长时间。 虽然它是从小额注册开始的，但它也通过了大额注册的测试。 当我们承诺进入新的大二时，我们的招生人数正在下降。 现在我们发现，更多的学院新生选择了航空航天而不是所有其他学位课程来学习。 这就是航空航天工程的世界。 我们的大二核心课程由四、五个学分课程组成：静力学、结构和材料介绍、热力学和空气动力学介绍、动力学和系统介绍、航空航天飞行器设计和性能介绍。这些课程都教授航空航天应用背景下的基础知识；它们使学生能够在学习的早期进行航空航天工程，既能为整个领域提供背景，又激励他们在上级进行深入研究。 这些课程是团队授课的，每周有两节课（不是讲座）让学生参与ITLL主动学习教室和ITLL LabPlazas的两个实验室。 这被数量有限的家庭作业所增强。 家庭作业将结合分析、设计和计算。 实验室工作将强调实验室和团队技能，以及书面和口头沟通（例如，见Ercolano [29]）。 每门课程由两个教师教授，包括我们教师的八个（或三分之一）。 我们的目的是每门课程都会

458 SEEBASS和PETERSON借鉴了其他人。 为了确保我们的上级课程借鉴这些大二课程的内容，所有教师将在六年内参与这些课程的团队教学。 24.5.5 上级 在编写本文时，我们仍在为指定上级课程的程度而纠结。 这部分将由AIAA制定的航空航天特定标准决定。 该标准允许沿着航空和航天线进行划分，例如马里兰州[4]采用的，麻省理工学院开发的学科“支柱”方法[21]，或辛辛那提采用的综合课程方法[57]。 高级项目课程 我们的新课程还包括一个新的综合高级项目，在为期一年的10个学分序列中实施。 旧课程包括几门高级设计课程：涵盖航天器或飞机的初步设计的3学分论文研究，以及长达6学分的高级设计实验室。 新的10学分课程结合了这些内容，并又迈出了一步。 在新的高级项目课程中，在综合产品开发团队中工作的学生将执行航空航天相关产品或设备的完整设计/构建/测试。 学生将遵循国际标准分阶段采购周期，其中包括航空航天专业人士熟悉的A、B、C、D和E阶段。 项目定义和设计（A和B阶段）将跨越第一学期，C/D/E阶段将跨越第二学期。 本课程中学到的最重要的概念之一是“要求驱动设计”，这是航空航天业广泛使用的高效产品开发的有纪律的方法，但在本科设计项目中经常被忽视。 24.5.6 资源 学院的百年纪念活动筹集了约5500万美元，其中近三分之一用于支持ITLL。 我们与学院的其他本科学位课程共享这个独特的设施。 它为我们的课程改革提供了基础，如果没有这个设施，我们不可能考虑。 我们对课程重新制定的成本估计，包括高级项目设计设施的开发，约为100万美元，加上目前提供的上级选修课的丰富性减少，随之而来的教育成本。 其他地方的经验使我们得出结论，新课程在全面发展后，将使教师的教学负荷增加25-30\%。 这只能通过减少我们提供的上级选修课的数量来满足。

航空航天课程 2000 459 幸运的是，科罗拉多州高等教育委员会于1997年将航空航天工程科学选为全州优秀课程，以表彰我们的卓越研究，以及我们与大气和空间物理实验室在三家USRA学生教育发现探测器之一，学生一氧化氮探测器上的合作。 这项五年赠款由洛克希德·马丁公司的类似赠款增强，为我们的课程修订和洛克希德·马丁/K提供了所需的资源。 D。 木材高级项目实验室。 24.6 两个警告，我们在课程改革中执行的许多工作将使工程学学习更有趣，更难教授。 计算机和互联网也适合这种乐趣[47]。 除非这种乐趣激发了我们学生更努力、更聪明的工作，从而让他们取得更大的成就，否则我们面临用乐趣取代基本原理的风险。 而且，除非我们为教师保留他们成为其领域领导者所需的时间，否则我们面临他们成为差教师的风险：“...... 我见过一位老师迷迷地抱着一百五十名学生，教什么是错的。 他的学生评价他是一位伟大的老师。 相比之下，我自己在大学里最伟大的两位老师会被评为差老师...... 那为什么人们从世界各地来和他们一起学习，包括我？ 原因很简单，这些人有东西要教。 他们激励学生进行进一步的研究。 他们是思想的领导者...... 他们的作品将在几个世纪里保持经典。” W。 Edwards Deming [26] 24.7 结论 学院新的综合教学实验室将于秋季开始全时运营。 有了这个全校的设施来帮助我们，在科罗拉多州高等教育委员会和洛克希德·马丁公司的大力支持下，航空航天工程科学建立了一门第一年项目课程，一个新的大二课程由四门课程组成，整合了理论、分析和实验（以及学科），以提供系统视角，以及顶点，两个学期的高级设计实验室，由为该系计划的新设施提供支持。 这项本科改革同时增加了我们的研究生学位要求的具体性，并采用了五年制BS/MS联合毕业学位。 我们现在正在完成上级课程的计划。 我们已经实施或即将实施的东西，响应了我们从企业和教育领导人那里听到的，包括两个国家科学院

460 SEEBASS \& PETERSON工程主席-Norman Augustine和Richard Morrow。 然而，他们都会问，我们正在做什么来鼓励我们的毕业生更多地参与政治进程。 当然，他们是对的，我们需要改革我们的政治领导，工程师对我们的治理没有做出足够的贡献，更不用说他们有能力的贡献了。 我们学院通过题为“现实世界中的创业、领导力和道德”的高级研讨会，朝着这个方向取得了一些进步。 公民身份是麻省理工学院在Earll Murman领导下采用的隐性课程的一部分。 这些是朝着正确方向迈出的一步，但我们必须做得更多，否则我们的国家将面临风险。 致谢 作者感谢他们的同事Brian Argrow、Penny Axelrad、Larry Carlson和Jacquelyn Sullivan对本文最有价值的评论；科罗拉多州高等教育委员会、洛克希德·马丁公司和John R. Woodhull，以支援这项努力。 参考资料1。 ABET，工程展望：新十年，新世纪，新千年，ABET全国会议记录，1990年。 2. ABET，2000年工程标准（草案），工程和技术认证委员会工程认证委员会，1997年。 3. 阿德勒，M。 J.，改革教育，纽约：麦克米伦，1988年。 4. 阿金，D。 L.，巴洛，J. B.和施密特，D. K.，为下个世纪设计航空航天工程课程：马里兰大学的经验，ASEE年度会议记录，1994年，pp. 27-35。 5. AIA，航空航天事实和数据，1993-1994年，美国航空航天工业协会，1993年。 6. AIAA，航空航天和类似名称项目的项目标准，AIAA网页，1997年。 7. 阿姆斯特朗，J. A.，行业应该对学术工程研究有什么期望？，塑造美国的力量。 学术工程研究企业，华盛顿特区：国家科学院出版社，1995年，pp. 59-66。

航空航天课程2000 461 8。 ASEE（Grinter报告），工程教育评估委员会的报告，工程教育杂志，46，1995年，pp. 26-60。 9. ASEE，不断变化的世界的工程教育，1994年10月。 10. 奥古斯丁，N。 R.，社会工程（以及奥古斯丁的第二定律），桥梁，1994年秋季，pp. 3-14。 11. 奥古斯丁，N。 R.，重建工程教育，高等教育纪事，1996年5月24日，pp. B1-B2。 12. 贝达德，A。 J.，Jr.和Meyer，D. G.，实践工程作业：课外学习的新方法，ASEE年度会议记录，1996年，第1626节。 13. Bordogna，J.，Fromm，E.，Ernst，E.，工程教育：通过整合进行创新，工程教育杂志，82（1），1993年，pp. 3-8。 14. Bowman，D.，Lang，J.和McMasters，J.，加强工程教育圆桌会议-更新，AIAA Paper 97-0844，1997年。 15. 博耶，E。 L.，大学：美国的本科生经历，纽约：Harper \& Row，1987年。 16. 伯尔，W。 M。 H.，理想的工程教育，工程教育，卷。 1，1893，pp. 17-49。 17. Carlson，L.E.和Brandemuehl，M. J.，一个活生生的实验室，ASEE年度会议记录，1997年，第3226次会议。 18. 卡尔森，L。 E.，彼得森，L. D.，隆德，W. S.和Schwartz，T. L.，使用LabStation促进跨学科实践学习，ASEE年度会议记录，1997年，第2659次会议。 19. 卡尔森，L。 E.等人，第一年工程项目：跨学科、实践工程介绍，ASEE年度会议记录，1995年，pp. 2039-2043。 20. Carr R.等人，综合工程课程的数学和科学基础，工程教育杂志，84（2），1995年，pp. 137-150。

462 SEEBASS \& PETERSON 21。 Crawley等人，麻省理工学院航空航天课程改革，工程教育杂志，83（1），1994年，pp. 47-56。 22. CUE，为我们的学生准备21世纪，CUEngineering，1991年8日，pp. 2-5。 23. CUE，工程教育，本科经验，本科教育，CU工程，7，1990年，pp. 6-33. 24. CUE，工程卓越战略计划，CUEngineering，9，1992/1993。 25. 达利，J。 W.和张，G. M.，新生工程设计课程，工程教育杂志，82（2）1993，pp. 83-91。 26. 德明，W。 E.，摆脱危机，马萨诸塞州剑桥：麻省理工学院高级工程研究中心，1982年。 27. 德明，W。 E.，新经济学，马萨诸塞州剑桥：麻省理工学院高级工程研究中心，1993年。 28. 德图佐斯，M。 L.，莱斯特，R. K.，索洛，R. M.和麻省理工学院工业生产力委员会，美国制造：重新获得生产力优势，马萨诸塞州剑桥：麻省理工学院出版社，1989年。 29. Ercolano，V.，通过合作学习，ASEE PRISM，1994年11月，pp. 26-29。 30. Ercolano，V.，设计新生，ASEE PRISM，1996年4月，pp. 20-25。 31. Fitzhorn，P。 A.和Johnson，G. R.，第一年工程课程有收敛吗？，ASEE年度会议记录，1994年，pp. 813-819。 32. 弗洛曼，S。 C，土木工程师，纽约：圣。 马丁出版社，1987年。 33. Halberstam，D.，《清算》，纽约：Morrow，1986年。 34. 霍利斯特，W。 M.，克劳利，E. F.和Amir，A. R“统一工程：麻省理工学院航空航天教育二十年实验，工程教育杂志，85（1），1995年，pp. 13-19。

航空航天课程2000 463 35。 拉德西克，J。 G.和Hazen，D. C，工程教育课程更正，美国航空航天，1995年5月，pp. 22-27。 36. 麦克马斯特，J。 H.和Lang，J. D.，加强工程和制造教育：行业需求，行业角色，ASEE年会记录，1995年，第2502次会议。 37. McMasters，J.和Matsch，L.，工程毕业生的期望属性-行业视角，AIAA论文96-2241，1996年。 38. 蒙特斯，L。 K.，工程教育-一个世纪的机会，工程教育杂志，83（1），1994年，pp. 22-25。 39. Moulakis，A.，超越效用：技术时代的自由教育，密苏里州哥伦比亚：密苏里大学出版社，1994年。 40. 莫罗，R。 M.，工程教育面临的问题，工程教育杂志，83（1），1994年，pp. 15-18。 41. 《纽约时报》，2亿美元，有史以来最大的礼物，捐赠新工程学院，1997年6月6日，p. A18。 42. NRC，美国的竞争地位。 工业，华盛顿特区：国家科学院出版社，1985年。 43. NRC，美国的竞争地位。 民航制造业-概述，华盛顿特区：国家科学院出版社，1985年。 44. NRC（Haddad报告），美国工程教育和实践：我们技术经济未来的基础，国家研究委员会，华盛顿特区：国家科学院出版社，1985年。 45. NRC（Pister报告），工程教育：设计适应系统，国家研究委员会，华盛顿特区：国家科学院出版社，1995年。 46. 奎因，R. G.，Drexel的E4Program：工程学生和教师的不同专业经验，工程教育杂志，82（4），1993年10月，pp. 196-202。 47. Panitz，B.，网络怀疑论者的观点，ASEE PRISM，1997年5月至6月，pp. 18-20。

464 SEEBASS \& PETERSON 48. 里根，T。 M.和Minderman，Jr.，P. A.，重新设计新生工程，化学工程进步，90（4）1994年，pp. S1-S4。 49. 里根，T。 M.和Minderman，Jr.，P. A.，600名新生的工程设计-规模化成功，会议记录，教育前沿，1993年，pp. 56-60。 50. 罗杰斯，G.和温克尔，B. J.，科学、工程和数学、自然、进化和评估的综合第一年课程，ASEE会议记录，1993年，pp. 186-191。 51. Seymour，E.和Hewitt，N. M.，谈论离开：导致科学、数学和工程本科专业高自然减退率的因素，给Alfred P.的最终报告。 斯隆基金会，科罗拉多州博尔德：科罗拉多大学社会学研究局，1997年。 52. SPEE（威肯登报告），工程教育调查报告，伊利诺伊州厄巴纳：工程教育促进协会，1930年。 53. SPEE（哈蒙德报告I），工程课程的目的和范围，工程教育杂志，30（7），1940年，pp. 555-566。 54. SPEE（哈蒙德报告II），战后工程教育委员会的报告，工程教育杂志，34（9），1944年，p. 589-614。 55. Sullivan，J.和Etter，D. M.，“综合教学计划：迎接21世纪工程教育的挑战”，The Interface，IEEE，1，1996年，pp. 1-4。 56. 范·瓦尔肯伯格，M。 E.，准备21世纪的工程师，工程教育，78（10），1988年，p. 103. 57. 沃克，B。 K.，Wade，J. E.，Orkwis，P. D.，Jeng，S.-M.，Khosla，P. K.和Slater，G. L.，21世纪拟议的航空航天工程课程的开发，AIAA论文97-0737，1997年。 58. 白色，R。 M.，科学、工程和巫师的学徒，The Bridge，21（1），1991年，pp. 13-20。

\clearpage
\chapter{流体力学入门型计算机教科书}
流体力学入门计算机教科书David A. Caughey*和James A. Liggett*摘要作者开发了一本利用个人电脑的力量教授流体力学入门的互动教科书。 这种演示模式集成了超文本导航和搜索功能，展示了视频和动画来说明现象和概念，以及计算，以允许展示各种参数值的结果，并解决非线性问题，而无需学生对表格查找或迭代感到繁重。 作者使用该书的早期版本教授机械工程、土木和环境工程的初级学生的经验表明，学生欣赏动态数字、易于访问数据、快速定位定义和特定材料的能力，以及最重要的是计算设施。 25.1 导言 流体力学是一门对大多数工程分支具有根本重要性的工程科学，包括航空航天、化学、土木、环境和机械工程，以及电气工程和材料工程的某些方面。 流体力学通常在上述领域的课程中向工程学生教授，从大三的一门或多门课程开始。 尽管我们一生都被流体所包围并沉浸在流体中，但大多数学生认为这门课程很困难--主要是因为流体力学中许多问题的表述具有抽象性质，典型的学生没有形成直觉，以及“康奈尔大学，伊萨卡，纽约14853-7501的频率。 计算流体动力学的前沿-1998年。 编辑：David A. Caughey A Mohajned J4，Hjifez ©1998 世界科学

466 CAUGHEY \& LIGGETT 非线性是制定最常见的工程问题的一个因素。 作者开发了一本教科书3，该教科书利用个人电脑的力量来尝试解决现象可视化、流体现象及其数学描述之间的联系以及有意义计算的固有困难等问题。 这本书是作为独立文本设计的，而不是作为现有文本的补充。 虽然纸质版本将与电子形式一起提供，但该书旨在在计算机上阅读，它集成了超文本功能、说明运动学和动态现象的动画和视频序列、动态呈现数据的图形，以及解决复杂问题的计算设施，而没有与经典（图形或表格迭代）方法相关的无味。 导航功能包括广泛的关键词词汇表，其定义和相关关键词可以在上下文中找到，以前访问过的页面的历史列表，以及几种类型的目录，允许学生快速跳转到文本的任何特定部分。 说明关键概念或复杂推导的动画包含在“QuickTime”电影中，以及说明各种流体流动现象的视频剪辑。 所有数值计算都在MATLAB*中完成，但数据是输入的，结果是通过图形用户界面呈现的，因此学生不需要编程。 这种基础的计算能力允许纳入动态图形，学生可以在其中更改一个或多个参数的值，并立即看到绘图结果中的效果，以及“主动方程”实用程序，允许学生将任何有意义的因变量绘制为许多方程的任何有意义的自变量的函数。 十几个计算实用程序与文本集成，包括单位转换程序、求解非线性或超越方程组的通用实用程序、用于摩擦管道流动问题的一维能量方程求解器、二维势流问题的叠加和边界积分方程方法、通常以表格形式呈现的可压缩流动函数的电子表格实现，以及许多其他。 这些计算实用程序使学生能够做足够的练习，以便他/她能够对流体系统的行为发展一些直觉。 本文概述了这种新形式的教科书的特点，以及作者使用文本的初步版本教授机械和土木工程初级课程的一些经验。 导航*MATLAB是由马萨诸塞州纳蒂克的MathWorks公司开发和分发的多功能数值分析和图形软件包。 MATLAB的可执行版本将随本书一起分发，但仅可用于运行作者开发的实用程序。

CAUGHEY \& LIGGETT 467 图1。 交互式教科书的典型页面（或屏幕）。下一节描述了功能，然后是动画和视频功能的描述。 计算功能用几个例子来描述。 最后，总结了关于新教科书如何可能改变教师/学生互动的性质的观察，以及使用文本的初步版本向土木和环境工程以及机械工程学生教授初级课程的经验。 25.2 导航和集成 教科书在电脑屏幕上的视觉呈现旨在模仿传统（纸质）教科书的外观。 屏幕是面向页面的（而不是滚动的），因此保留了页面上对象出现的视觉记忆，并提供了许多视觉提示来提醒学生在书中的当前位置。 典型页面的外观如图所示。 1. 基本导航控件位于页面的左下角；“下一步”和“上一个”按钮将学生带到当前级别文本中的下一个或上一个线性页面。 *“返回”按钮让学生返回*文本中的部分分为三个级别：第1级包含概述材料，第2级包含所有第1级材料以及所需的材料

468 CAUGHEY \& LIGGETT最近访问的页面--即通过“历史”列表--无论它们相对于当前页面的位置如何。 最近访问的页面的历史列表可以从“历史”菜单访问，学生可以跳转到历史列表中的任何选定的页面。 章节中当前页面的位置由靠近页面底部的滑块中红色矩形的位置指示；图中显示的页面。 1大约是第11章的20\%。 滑块左侧和右侧的数字表示书中的其他章节（由于空间限制，仅显示奇数；偶数章节由项目符号标记表示）；单击章节编号（或标记）将学生带到相应章节的第一页。 可以通过单击页面左上角的“章节”随时查看当前章节的详细目录，任何章节的目录都可以从列出所有章节的图书目录中查看。 学生还可以使用目录菜单下的“图搜索”项目在整本书（或任何选定的章节子集）中搜索任何图（包括动画、视频等），如果需要，该项还能够仅查找字幕中包含特定文本字符串的图。 教科书提供了与超文本演示相关的大部分功能。 这些包括能够直接从文本中出现的最重要技术术语的简短定义调用它们，从词汇表中调用相关术语的定义，以及跳转到文本中首次引入该术语或相关术语的页面。 以这种方式被视为“热文本”的词汇表术语使用粗体字体表示：图中页面上出现的术语过渡、分离点、层压和湍流。 1是例子。 当前页面上引用的前几页的方程和数字也可以以类似的方式查看。 学生还可以在任何页面上写“笔记”，或在任何页面上放置书签，以便使用“书签”菜单轻松找到（并返回）。 25.3 插图和数据 文本中提供了三种类型的插图：（1）静态图形，（2）动态动画/视频，以及（3）活动图形。 静态数字类似于了解如何进行工程应用的数字，第3级包含所有2级材料以及复杂衍生或更专业性质材料的详细背景。 当前级别由页面右上角的单选按钮指示。 虽然第1级和第2级对读者“隐藏”了一些材料，但通过更改级别或使用目录，这些材料很容易获得。

CAUGHEY \& LIGGETT 469 图。 2 视频动画的单帧说明了路径线、条纹线和流线之间的区别，对于理想的不可压缩流动，经过一个半径随时间变化的圆形圆柱体。传统的纸质文本。 简单的数字全尺寸显示在页面的左页边空白处。 更复杂的数字，如图中页面上出现的数字。 1，由空白处带有“放大镜”图标的缩略图表示。 点击此类图的缩略图草图，将其全尺寸显示在屏幕上，以便进行详细研究。 动态动画或视频用于说明动态特征，提供实际流动现象的视频显示，并提供几种类型的介绍性信息。 本介绍性信息包括一些章节的概述（出现在章节的第一页）和教程会议，说明了作为文本一部分的许多计算实用程序。 这些动画和视频由页面左页边空白处的缩略图表示，并由“胶片”图标标识。 通过单击其识别缩略图草图，动画或视频将作为“QuickTime”电影启动。 “QuickTime”电影可以向前或向后播放，可以在任何帧处停止，或者可以一次向前或向后移动一帧。 一个动画序列的例子，说明了非定常流动经过一个圆形时间圆柱的路径线、条纹线和流线之间的区别-

图中显示了470 CAUGHEY和LIGGETT的变化半径。 2. 不幸的是，在本文的格式中，只能显示视频的单帧。 活动图表用于显示定量结果，这些结果取决于学生可能更改的一个或多个参数。 它们由活动“计算器”图标标识，如图所示。 3. 单击图表缩略图中的任意位置可启动图形的活动版本。 在教科书中呈现的大多数定量数据图中，任何曲线上特定点附近的区域都可以放大，以便更准确地读取值。 在图中的任何点上单击鼠标按钮，即可将该点附近的比例翻倍，或者可以通过在相反的角落之间单击和拖动来定义“缩放框”，以定义要放大到整个图区域的区域。 双击鼠标按钮在图区域的任意位置，即可将图缩放到原始比例。 图中所示页面上的活动图形。 3展示了通过给定角度转动超声速流动的压力比和冲击角。 结果取决于上游马赫数和气体的比热比（\& = c p I cv）的比率。 当学生更改马赫数或特定热量的比率时，图会自动重新绘制，以对应于新值。 图3 具有活动图形的页面。 “计算器图标”表示可以单击缩略图以启动活动图形。

CAUGHEY \& LIGGETT 471 图。 4 显示压力比作为超声速流动转弯角度的函数的有源图。 该图说明了单个上游马赫数的曲线和比热比，但学生可以为这些参数输入任何值来查看相应的图。 图4显示了该图的主动版本（对于上游马赫数 M = 2.0和比热比k = 1.667）。 当压力比被选为因变量（而不是冲击角）时，也绘制了等熵流动的相应曲线。 学生可以通过单击相应的单选按钮，选择以弧度或度显示角度。 有源图形还用于显示流体属性，如黏度系数或导热系数，可以通过从“数据”菜单中选择相关的流体属性来启动。 图5显示了几种常见液体的导热系数对温度的依赖性。 作为最后一个例子，图。 6显示了美国宇航局的前身国家航空咨询委员会（NACA）测试的几架机翼的升力、阻力和俯仰矩特性。 启动活动图形后，学生可以显示几个不同翼型形状的升力、阻力和俯仰矩特征。

472 CAUGHEY \& LIGGETT图。 5 显示几种常见液体的导热系数对温度的依赖性的有源图。 学生可以通过单击所需的曲线或在相应的数据框中输入温度值，并在选择特定曲线后读取结果的导热率来读取值。 曲线是使用Reid、Prausnitz和Poling5的半经验公式和实验值计算的。 25.4 计算和实用程序 计算实用程序用于解决可能需要求解非线性（或非线性）方程的各种工程问题，并集成到教科书中。 从历史上看，这些问题是通过图形方法、需要重复插值的表格函数的迭代或要求学生实施数值方法来解决的。 互动教科书为MATLAB中编写的许多例程提供了图形用户界面（GUI），允许学生在相同的时间内（或需要类似的努力）解决更多种类和数量的此类问题，以解决过去的单个问题。 提供此类基于GUI的实用程序来解决：•具有等熵流和正态激波流的面积变化的可压缩流；•任何诺伊曼和狄利克雷边界条件组合的任意二维几何中的势流问题；•维分析问题，为任何一组维度输入变量找到无维变量组；

CAUGHEY \& LIGGETT 473 图。 6 各种NACA 翼型形状的升力、阻力和力矩系数，（a）升力和力矩系数绘制为攻角的函数；（b）阻力系数绘制为升力系数的函数。 在Abbott \& von Doenhoff 1之后。 •稳态管道网络流动的解决方案；•具有摩擦的恒面积管道中的可压缩流动；•函数的数值积分；•不可压缩管道流量问题，用于确定管道尺寸、流量、头部损失等，用于湍流摩擦损失的各种公式；•函数和数据的一维和二维绘图，以确定表面图下的体积；•通过基本解（源、汇、涡流、双元和均匀）的叠加来叠加平面和轴对称不可压缩势流问题 溪流）；•开放通道剖面问题的解决方案；•恒定面积管道中的可压缩流动与热量加成；•任何非线性或超越方程系统；•标准大气中的属性；•从任何系统到几乎任何其他系统的单位转换问题；•弹性管道中的水锤问题。 在这里，我们将提供其中三个实用程序的简要描述，以说明其功能。

474 CAUGHEY \& LIGGETT图。 7 全开发湍流管道流量中摩擦系数的穆迪图表。 摩擦系数是非维管道粗糙度和雷诺数的函数。 25.4.1 管道流 在直径为d的圆管的长度L中，完全发育的湍流摩擦引起的头损失h f由关系hf=t LVL dig'（25-4.1）给出，其中V是平均流速，g是重力加速度，f是（Darcy-Weisbach）摩擦系数。 摩擦系数仅是管道表面非维粗糙度和雷诺数 Red = Vd I v的函数，其中V是流体的运动黏度。 Colebrook2的公式提供了光滑和粗糙管道1 •2.0 log（eld 3.7 2.51 Re，f 1/2（25.4.2）摩擦系数作为雷诺数函数的图，通常称为穆迪图4）之间的插值如图所示。 7. 对于某些问题，使用穆迪图表相对简单。 例如，通过指定长度和直径的管道的给定流速的头部损失可以以封闭的形式找到。 在这种情况下，雷诺数可以直接从给定的数量中计算，摩擦系数确定为

CAUGHEY \& LIGGETT 475（通过迭代）来自Eq。 （25.4.2）或直接来自穆迪图表。 然后，摩擦系数可用于使用Eq.确定黏性头损失。 （25.4.1）。 另一方面，有些问题要困难得多。 例如，通过给定长度L的管道的指定头部损失hf和体积流量Q的直径需要求解同步的非线性方程。 在这种情况下，雷诺数必须确定为解决方案的一部分，因为平均流速V和管道直径d都不是先验的。 方程（25.4.1）和（25.4.2）在体积流速Q方面重新铸造，必须同时求解，以确定与给定hf相对应的速度V和直径d。 由于没有通用技术来求解此类非线性方程组，人们通常采用简单的迭代方案，对管道直径进行初步估计，然后根据R\textasciicircum{}=-\textasciicircum{}计算雷诺数。 （25.4.3）穆迪图表（或等同于，Eq. （25.4.2））然后可用于计算相应的摩擦系数f，并且可以通过求解直径d = 8f L\_ <\$\_ K1 h的方程\% d g，（25.4.4）\textbackslash\{\}" "f 6，如果流量包含入口或出口损失、轻微损失或泵或涡轮机，从头部损失中计算出管道直径的新近似值，方程变得更加复杂。 如果穆迪图表或等式，这个过程非常乏味。 （25.4.2）必须用于确定每次迭代时雷诺数的摩擦系数，但它可以在计算机上轻松实现自动化。 可以使用基于GUI的实用程序PipeFlow找到此问题以及许多其他问题的解决方案。 PipeFlow解决一维机械能方程、Darcy-Weisbach方程（25.4.1）和Colebrook方程（25.4.2），单独或组合，对于摩擦系数f加上另一个变量，允许各种边界条件和轻微损失。 它可以从管道类型（例如混凝土）中选择粗糙度，查找水属性（或允许指定其他流体的属性），并从商业管道尺寸表中找到实际（而不是标称）的管道尺寸。 用于指定这些问题的PipeFlow的主屏幕如图所示。 8. 图中显示的屏幕。 8说明了刚才描述的问题类型的一个例子的解决方案。 在长度为L = 100米的管道中，由Ap = 10,000 Pa的压力差驱动的完全开发的湍流所需的管道直径为2= 1.0 m3/s的溶液。

476 CAUGHEY \& LIGGETT IPi 机械能方程：管道流解决方案 I 图。 8 PipeFlow的主面板，这是一个基于GUI的实用程序，用于解决涉及管道中湍流损失的各种流体问题。 25.4.2 带面积变化的可压缩流动 另一个问题示例，使用廉价的个人电脑可以作为优势，以解决其他相当繁琐的问题，涉及不同横截面面积的管道中的可压缩流动。 作为喷嘴中等熵流动的横截面面积A的函数，马赫数 M可以表示为1（2 + (fc-l）M 2 (25.4.5) 2(*-l) M V fc + 1，其中k是比热的比率，A*是喷嘴的有效声波面积（根据操作条件，它可能对应于考虑的喷嘴的最小面积，也可能不对应）；例如，见Shapiro6。 一旦知道马赫数，其他流动变量可以直接从等熵关系中确定；例如，压力p由P1= (> + \textasciicircum{} M ')" C5.4.6给出，其中P0是等熵停滞压力。 对于完全等熵流，将流特性确定为喷嘴横截面面积的函数的问题只是求解方程之一。 （25.4.5）或者，等价地，插值此函数的表格值，以确定马赫

CAUGHEY \& LIGGETT 477 编号 M，然后使用公式，例如 Eq. （25.4.6）从马赫数中确定流动特性。 如果喷嘴中存在激波，问题就会变得更加复杂。 一个典型的问题是确定给定出口压力的给定几何形状（即给定的A（x））的喷嘴中的激波位置，通常太乏味了，甚至不能作为家庭作业练习分配。 *这个问题和上一节中描述的许多湍流管道流问题一样，需要迭代。 对于给定的激波位置，压力的跳跃可以从正常的冲击关系t= 1+ F>, -« (25-4.7，其中下标（X和( ) 2分别指冲击的上游和下游的状态。 冲击紧靠下游的马赫数由2 \_\_ k + 1 + (k -1) (M, 2 - 1) fc+1 + 2/t(M, 2 -1) ' ( 25A8 )给出，这可以与Eq一起使用。 （25.4.5）确定新的有效A'，而Po2 Po，（fc + l）M 2 "\textbackslash\{\}iZ7 k 1 v 2 + Ot-l)M，2 y f i., i A k + l IkM* - (k-l) <:-i（25.4.9）可用于确定新的停滞压力。 方程（25.4.9）以及从冲击到出口的面积比，可以与上述公式一起使用来计算喷嘴出口处的压力。 将这种计算压力与所需的出口压力进行比较，表明冲击是否需要向上游或下游移动，以更接近所需的出口压力。 基于GUI的实用程序AreaFlow可以轻松访问不同横截面面积和正态冲击关系的管道中等熟流量的完整公式集，使上述计算序列变得不那么乏味。 数据以电子表格的形式输入和呈现，每次输入或更改值时都会更新，以便始终保持内部一致。 AreaFlow显示面板如图所示。 9，显示计算的值，以确定当出口面积是喷嘴喉部面积的两倍时，产生出口压力与上游停滞压力比为0.5所需的冲击强度。 这个解决方案是用大约半打迭代计算的，从开始到结束需要不到3分钟的使用者互动时间。 *开放通道流动中存在类似的问题，其中液压跳跃的位置需要为指定的边界条件和通道几何形状确定。 通道剖面求解器与超越方程求解器相结合，为该问题提供了一个简单的解法。

47 8 CAUGHEY \& LIGGETT图。 9 AreaFlow的数据面板，这是一个基于GUI的实用程序，提供对可压缩流动方程的访问。 25.4.3 Inpressable 势流 我们集成到教科书中的计算实用程序的最后一个例子是PotFlow，它绘制了选定的流线、恒定速度势线和二维势流（平面或轴对称）的恒压线，由任意数量的源、汇、涡流、双流和均匀流的叠加生成。 PotFlow的主控制面板如图所示。 10. 使用面板右上角的控件定义元素并添加到流程中，可以使用面板右下角的控件更改现有元素的属性，当前元素的属性显示在面板的左上角。 面板左下角的控件用于清除数据空间，从文件中读取元素数据，并选择要绘制的结果形式。 图11显示了由PotFlow生成的图，该图显示了常数压力系数的流线和选定的轮廓，用于在20度攻角下通过平板翼型的离散涡流表示。 沿着平板均匀间隔的32个涡流的强度，使用平板翼型（教科书中的“3级”主题）的薄翼型理论计算。

CAUGHEY \& LIGGETT 479 图。 10 PotFlow的主控制面板，PotFlow是一个基于GUI的实用程序，用于研究由源、涡流、双元和均匀流叠加生成的二维势流。 图。 11 PotFlow图形显示了恒定压力的流线型和轮廓，用于离散涡流近似，以20度攻角通过平板翼型的流动。

480 CAUGHEY \& LIGGETT 25.5 教学经验 两位作者都使用互动教科书的初步版本向土木和环境工程以及机械工程的初级学生教授部分或全部流体力学入门课程。 学生们通常对新媒体提供的图形和计算实用工具充满热情，但对从屏幕上阅读大量文本和方程则不太热情。 这并不奇怪，因为计算和视频是计算机擅长的事情。 这些动态互动具有吸引力，并鼓励主动学习。 如果有替代方案，我们都不愿意花很多时间在电脑屏幕上阅读文本。 这一经验导致我们提供了文本的纸质版本来补充电子形式，但我们仍然认为将图形演示和计算与纹理材料相结合很重要。 我们尝试了几种在讲座和朗诵部分使用教科书的方法。 在讲座中，我们发现最好有可用的计算机投影设备，并以有限的方式使用教科书的功能来说明主题，而不是将整个讲座建立在投影页面、插图和实用工具上。 演讲厅越来越多地提供计算机投影设备，在大多数情况下，在传统的黑板模式、视频和动画序列以及计算之间来回切换的便利性预示着将此类工具更多地纳入讲座演示。 讲师可以使用书签轻松找到动画、电影、图表和活动方程。 讲师可以在家里在自己的机器上使用这些相同的功能来进一步学习和解决练习。 我们发现，至少在每个学生都可以使用电脑的设施中举行一些朗诵部分是有效的，特别是在课程的早期。 通过这种方式，学生可以在有使用新媒介经验的教师或助教面前获得早期新媒介的体验。 （随着教科书变得更加强大，学生对新媒介越来越熟悉，使用文本的指导可能会变得不那么重要。） 在「工作室」环境中，鼓励学生在巡回教练的注视下使用各种计算工具探索练习，也被证明是有效的。

CAUGHEY \& LIGGETT 481 致谢 该教科书是使用MacroMedia AuthorWare编写的，并将在Macintosh、Windows 3.1、Windows 95和Windows NT操作系统下运行的CD-ROM上分发。 视频片段和动画作为QuickTime电影运行，所有数值计算都在MATLAB中完成。 作者想借此机会感谢美国土木工程师学会出版物董事总经理David Dresia对该项目的持续支持。 我们还对康奈尔大学多媒体课件工作室的人员的贡献表示深深的感谢。 特别是，如果没有Rob Levine、Kate Mink、Melanie Swain、Mike Tolomeo、Kenny Unice、Dave Wickstrom、John Wolf以及康奈尔大学的众多本科生的奉献精神和专业知识，该项目不可能完成，他们为该项目的成功做出了重大贡献。 特别感谢凯特·明克和约翰·沃尔夫，他们自该项目启动以来就一直与该项目有关，他们的非凡努力在很大程度上促成了该项目的及时完成。 参考资料1。 我。 H。 阿博特和A。 E. von Doenhoff，机翼截面理论，纽约多佛，1959年。 2. C. F。 科尔布鲁克，湍流在管道中，特别提到平滑和粗糙管道定律之间的过渡，J。 Inst。 Civ。 英语。 伦敦，卷。 11，pp. 133-156，1938-39。 3. J。 A. 利格特和D。 A. Caughey，流体力学：互动文本，美国土木工程师学会，1998年。 4. L。 F。 穆迪，管道流动的摩擦系数，ASME Trans.，卷。 66，pp. 671-684，1944年。 5. R。 C. 里德，J。 M。 普拉斯尼茨和B。 E。 波林，气体和液体的特性，麦格劳-希尔，纽约，1987年。 6. A. H。 Shapiro，可压缩流体流动的动力学和热力学，卷。 1，罗纳德，纽约，1953年。

\chapter*{正文后所附补充}
按字母顺序排列的作者列表P。 J。 布鲁克斯第15章 威尔士大学斯旺西土木工程系 SA2 8PP，英国David A. Caughey第1章，25 Sibley机械与航空航天工程学院康奈尔大学伊萨卡，纽约14853-7501 Jean-Jacques Chattot第11章机械与航空工程系加州大学戴维斯分校戴维斯，加利福尼亚州95616 Connie K. 陈第六章 科兰特数学科学研究所 纽约大学 纽约，NY 10012 H。 K。 Cheng 第5章 南加州大学航空航天工程系 洛杉矶，加利福尼亚州 90089-1191 Julian D. 科尔第2章 数学科学系 伦塞勒理工学院 特洛伊，纽约 12180-3590 L. 帕梅拉·库克第4章 特拉华大学数学科学系 纽瓦克，DE 19716 J. Dacles-Mariani第18章 加州大学戴维斯分校，加利福尼亚州David L. 达莫法尔第13章航空航天工程系德克萨斯A\&M大学学院站，德克萨斯州77843-3141 计算流体动力学的前沿-1998编辑：David A. Caughey和Mohamed M. 哈菲兹 ©1998 世界科学

484指数Mark Drela第19章航空航天系M.I.T. 马萨诸塞州剑桥 02139 Nathalie Duquesne 第17章 航空系 KTH 皇家理工学院 S-10044 瑞典斯德哥尔摩 D. Scott Eberhardt 第12章 华盛顿大学 Box 352400 华盛顿州西雅图 98195-2400 Kozo Fujii 第21章 相模原空间和航天研究所，神奈川 229 日本 Cyril Gacherieu 第17章 欧洲研究和形成中心 Avancee en Calcul Scientinque F-31057 法国图卢兹 Paul R. Garabedian 第6章 科兰特数学科学研究所 纽约大学 纽约，纽约 10012 Michael B. Giles第9章 劳斯莱斯阅读器在CFD牛津大学计算实验室，英国牛津。 穆罕默德·M。 哈菲兹第1章，第5章，加州大学戴维斯分校机械和航空工程系，加利福尼亚州95616 O。 Hassan 第15章 土木工程系 威尔士大学 斯旺西 SA2 8PP, UK Wen-Huei Jou 第20章 波音商用飞机集团 P. 哦。 Box 3707 MS 67-LL 西雅图，华盛顿州 98124-2207 C。 Kiris第18章 MCAT公司 山景城，加利福尼亚州D。 Kwak 第18章 高级计算方法分支 美国宇航局艾姆斯研究中心 Moffett Field, California 94035

按字母顺序排列的作者列表 485 James A. Liggett 第25章 土木与环境工程学院 康奈尔大学伊萨卡，纽约 14853-7501 刘平 第7章 工程发展中心 成都飞机工业公司。 中国成都，罗世军第7章 西北理工大学飞机工程系，中国西安，Robert W. MacCormack第10章 航空和航天系 斯坦福大学 斯坦福，加利福尼亚州 94305-4035 M。 T。 曼扎里第15章 威尔士大学斯旺西土木工程系 SA2 8PP，英国K。 摩根第15章 威尔士大学斯旺西土木工程系 SA2 8PP，英国 艾尔莎·纽曼 第4章 数学系 马里蒙特大学阿灵顿，弗吉尼亚州 22207 小川隆信 第21章 清水公司 越岛 3-4-17 日本Koto-Ku Lee D. 彼得森第24章 航空航天工程科学 科罗拉多大学博尔德，科罗拉多州 80309-0429 Pradeep Raj 第23章 洛克希德·马丁航空系统 佐治亚州玛丽埃塔 30063-0685 Arthur Rizzi 第17章 航空系 KTH 皇家理工学院 S-10044 斯德哥尔摩 瑞典 Thomas W. 罗伯茨第14章空气动力学和声学方法分部美国宇航局兰利研究中心汉普顿，弗吉尼亚州23681-0001

S。 罗杰斯高级计算方法分部美国宇航局艾姆斯研究中心莫菲特菲尔德，加利福尼亚州94035 Paul E. Rubbert 波音公司 P. 0. Box 3707 西雅图，华盛顿州 98124-2207 Oleg S. Ryzhov数学科学系Rensselaer理工学院特洛伊，纽约12180-3590 R。 威尔士斯旺西大学土木工程系SA2 8PP，英国Donald W. 施温德曼数学科学系伦斯勒理工学院特洛伊，纽约12180-3590 A。 Richard Seebass 航空航天工程科学 科罗拉多大学博尔德，科罗拉多州 80309-0429 Huili Shen 航空发动机工程系 西北理工大学 西安，中国 David Sidilkover ICASE 美国宇航局兰利研究中心 汉普顿，弗吉尼亚州 23681-0001 Helmut Sobieczky DLR 德国航空航天研究机构 Bunsenstr. 10 D-37073 Gfittingen R。 C. 斯旺森空气动力学和声学方法分部 美国宇航局兰利研究中心 汉普顿，弗吉尼亚州 23681-0001 E。 D。 Terent'ev计算中心 俄罗斯科学院 40 Vavilov Street 117333 莫斯科，俄罗斯联邦 Loic Tourrette Aerospatiale Avions F-31060 Toulouse France INDEX 第18章 第22章 第8章 第15章 第2章 第24章 第7章 第14章 第3章 第14章 第8章 第17章

作者字母列表487 Susan A. Triantaflllou 第2章 数学科学系 Rensselaer理工学院 纽约特洛伊 12180-3590 Bram van Leer 第13章 密歇根大学航空航天工程系 安阿伯，密歇根州 48109-2140 Jan Vos 第17章 液压机械和流体力学研究所（IMHEF） 瑞士联邦理工学院（EPFL） CH-1015 洛桑 N. P。 Weatherill第15章 威尔士大学斯旺西土木工程系 SA2 8PP，英国Carlos Weber 第17章 欧洲研究与形成中心 Avancee en Calcul Scientifique F-31057 法国图卢兹 Pratomo Wibowo 第12章 华盛顿大学 Box 352400 华盛顿州西雅图 98195-2400 Laurence B. 威格顿第16章波音商用飞机集团P。 0. 华盛顿州西雅图3707号信箱98124-2207 S。 Yoon第18章高级计算方法分部美国宇航局艾姆斯研究中心莫菲特菲尔德，加利福尼亚州94035 Anders Ytterstrom第17章航空系KTH皇家理工学院S-10044瑞典斯德哥尔摩

\backmatter
\chapter*{翻译与版式说明}
全文按原书章节顺序生成，并采用连续的纯中文书籍版式。
\end{document}
